% Paper deep reading note zone.
% Add topic sub-zones below.

\part*{论文精读专区}
\phantomsection
\addcontentsline{toc}{part}{论文精读专区}

\begin{conceptbox}
\textbf{用途：} 这里沉淀单篇论文的精读笔记。日报负责发现论文，Survey 负责组织专题，精读专区负责把一篇论文拆到问题、方法、实验、局限和可复用点。\\
\textbf{组织方式：} 每个方向一个小专区，单篇精读放进对应子目录；跨方向论文优先按“最主要贡献”归档，必要时在其它专区索引里加交叉引用。\\
\textbf{建议节奏：} 每篇精读控制在一个稳定模板内，避免只摘摘要；重点写清楚它解决什么问题、相对前作改了哪里、关键假设是什么、哪些结论可以迁移到自己的研究。
\end{conceptbox}

\section*{专区索引}
\phantomsection
\addcontentsline{toc}{section}{专区索引}

\begin{itemize}[leftmargin=1.5em]
  \item \textbf{后训练与 RL 对齐}：RLHF/RLVR/GRPO/DPO/奖励模型/信用分配。
  \item \textbf{图像生成与编辑模型}：T2I、image editing、视觉文本生成、可控生成；子专题：Pixel-Space Diffusion \& Tokenizer-Generation Tradeoff（DiT / REPA / JiT / Latent Forcing）。
  \item \textbf{视频生成与编辑模型}：T2V/I2V、视频编辑、长视频 DiT、物理一致性。
  \item \textbf{统一理解与生成模型}：BAGEL/Janus/Tuna/BLIP3o/any-to-any。
  \item \textbf{世界模型与 3D}：world model、3D/4D 生成、空间智能、仿真。
  \item \textbf{具身智能与 VLA}：VLA、机器人后训练、仿真环境、操作策略。
  \item \textbf{Style / Identity / Design}：风格、身份保持、设计工作流、审美奖励。
  \item \textbf{位置编码专题}：RoPE/MRoPE/3D-RoPE、segment offset、多模态/图像/视频位置组织。
  \item \textbf{Idea 专区}：记录读论文、做实验和讨论时产生的研究想法、假设和最小验证实验。
  \item \textbf{评测、安全与红队}：benchmark、safety、provenance、failure discovery。
  \item \textbf{系统与推理效率}：训练/推理系统、attention、采样加速、工程优化。
\end{itemize}

% Post-training and RL alignment deep-reading notes.

\section*{后训练与 RL 对齐}
\phantomsection
\addcontentsline{toc}{section}{后训练与 RL 对齐}

\begin{conceptbox}
\textbf{收录范围：} RLHF / RLVR / GRPO / DPO / 奖励模型 / 偏好优化 / 视觉生成后训练 / credit assignment / rollout side（正负样本利用）。\\
\textbf{当前主线：} Flow-GRPO 在图像/视频生成里的时间步信用分配，包括 advantage side、ratio side、chunk-level 和 multi-objective credit；以及 LLM RLVR 上 rollout 侧的正负样本利用问题（如 POPO 的 positive-only 路线）。
\end{conceptbox}

% ============================================================
% Paper deep reading: DenseGRPO
% ============================================================

\section{论文精读：DenseGRPO}

\begin{conceptbox}
\textbf{论文：} \href{https://arxiv.org/abs/2601.20218}{DenseGRPO: From Sparse to Dense Reward for Flow Matching Model Alignment}\\
\textbf{版本：} arXiv 2601.20218v2，ICLR 2026。\\
\textbf{方向标签：} Flow Matching、Text-to-Image、GRPO、Dense Reward、Credit Assignment。\\
\textbf{阅读日期：} 2026-04-27。\\
\textbf{前置：} 需要已经理解 Flow-GRPO 的 SDE sampling、group relative advantage、policy ratio 和 KL regularization。
\end{conceptbox}

\subsection{0. 一句话结论}

DenseGRPO 解决的是 Flow-GRPO 里的核心信用分配问题：\textbf{最终图像 reward 不应该被原封不动分给所有 denoising steps}。它把终局 reward 改造成每个 timestep 的 dense reward，并进一步用 dense reward 来调每个 timestep 的探索噪声。

\begin{summarybox}
\textbf{我的判断：} 这是 Flow-GRPO 后最值得优先精读的一篇。因为它不是单纯堆一个技巧，而是把“reward 应该分给哪一步”这个问题具体化为可计算的 \key{step-wise reward gain}。
\end{summarybox}

\begin{figure}[htbp]
  \centering
  \includegraphics[width=\linewidth]{densegrpo/fig1_sparse_vs_dense_reward.png}
  \caption{Sparse reward 与 DenseGRPO dense reward 的核心差异。左侧是把 terminal reward 直接分给所有中间步骤；右侧是对每个 denoising step 估计独立贡献。图源：arXiv HTML。}
  \label{fig:densegrpo-sparse-vs-dense}
\end{figure}

\subsection{1. 这篇论文为什么重要}

Flow-GRPO 已经把 GRPO 接到了 flow matching / diffusion-like 视觉生成上。它的基本思路是：

\[
\text{prompt } c
\rightarrow
\text{sample a group of trajectories}
\rightarrow
\text{score final images}
\rightarrow
\text{GRPO update}
\]

问题在于，Flow-GRPO 的 reward 是 terminal reward：

\[
R(x_0, c)
\]

它只评价最终生成图像。但训练时每条轨迹有很多 denoising steps：

\[
x_T \rightarrow x_{T-1} \rightarrow \cdots \rightarrow x_1 \rightarrow x_0
\]

如果把同一个最终 reward 分给每一个 timestep，就隐含了一个很强的假设：

\[
\text{每一步对最终结果的贡献相同}
\]

这个假设通常是错的。某些步骤决定构图、颜色、对象关系；某些步骤只是细节 refinement。DenseGRPO 的核心价值，就是把这件事从直觉变成训练信号。

\subsection{2. Flow-GRPO 的问题在哪里}

Flow-GRPO 中，同一个 prompt 会采样一个 group。对第 \(i\) 个样本，最终图像是 \(x^i_0\)，reward model 给出：

\[
R(x^i_0, c)
\]

然后 group relative advantage 大致是：

\[
\hat{A}^i
=
\frac{
R(x^i_0,c)-\mathrm{mean}\{R(x^j_0,c)\}_{j=1}^{G}
}{
\mathrm{std}\{R(x^j_0,c)\}_{j=1}^{G}
}
\]

关键问题是：这个 \(\hat{A}^i\) 不依赖 timestep。也就是说，在第 \(t\) 步更新 policy 时，用的还是同一个整图 reward。

\begin{warnbox}
\textbf{问题本质：} Flow-GRPO 优化的是每一步 transition 的 logprob，但 advantage 来自整条 trajectory 的最终分数。这会导致 feedback-contribution mismatch。
\end{warnbox}

更具体地说，图像最后得分高，可能是因为早期 layout 选对了；也可能是因为后期文字细节修好了。把同一个 reward 分给所有步骤，会让模型不知道到底该强化哪一步。

\subsection{3. DenseGRPO 的核心想法}

DenseGRPO 要做两件事：

\begin{enumerate}[leftmargin=1.5em]
  \item \textbf{把 sparse terminal reward 变成 step-wise dense reward。}
  \item \textbf{根据 dense reward 分布，给不同 timestep 设置不同探索噪声。}
\end{enumerate}

这两件事对应论文的两个关键模块：

\[
\text{Step-wise Dense Reward}
\quad + \quad
\text{Exploration Space Calibration}
\]

\subsection{4. 模块一：Step-wise Dense Reward}

\subsubsection{4.1 从“最终分数”变成“每步增益”}

DenseGRPO 定义每一步的 reward 为相邻 latent 对应 reward 的提升量：

\[
\Delta R^i_t
=
R^i_{t-1} - R^i_t
\]

直觉是：

\[
x^i_t \rightarrow x^i_{t-1}
\]

这一小步 denoising 如果真的让图像更符合 reward model，那么从 \(x^i_t\) 对应的 clean image 到 \(x^i_{t-1}\) 对应的 clean image，reward 应该上升。这个上升量就是该步骤的贡献。

\begin{intubox}
\textbf{直觉版：} 不再问“最终图像好不好”，而是问“从这一步到下一步，图像有没有变得更好”。这就是 dense reward。
\end{intubox}

\subsubsection{4.2 中间 latent 怎么打分}

难点是：reward model 通常只能给 clean image 打分，不能直接给 noisy latent \(x_t\) 打分。

DenseGRPO 的办法是利用 flow matching ODE sampler 的确定性。给定中间 latent \(x^i_t\)，如果从它继续走 deterministic ODE denoising，可以得到一个 clean latent：

\[
\hat{x}^i_{t,0}
=
\mathrm{ODE}_n(x^i_t, c)
\]

再把 \(\hat{x}^i_{t,0}\) decode 成 clean image，并用普通 reward model 打分：

\[
R^i_t
:=
R(\hat{x}^i_{t,0}, c)
\]

这样就把“中间 latent 的 reward”转化成了“从这个 latent 确定性还原出来的 clean image 的 reward”。

\begin{summarybox}
\textbf{关键点：} DenseGRPO 没有额外训练 process reward model，也没有训练 critic。它复用已有 reward model，只是通过 ODE denoising 把中间 latent 映射回 clean image 后再打分。
\end{summarybox}

\subsubsection{4.2.1 疑问澄清：ODE denoising 不是探索}

这里容易产生一个误解：如果对某个 timestep 的 latent \(x_t\) 采取 ODE denoising，那 ODE 是 deterministic 的，怎么探索？

答案是：\textbf{ODE denoising 这一步本身不是用来探索的}。DenseGRPO 里有两条路径，需要分开看。

\textbf{第一条是采样/探索路径：SDE trajectory。}

训练时真正产生 rollout 的是 SDE sampler：

\[
x_T \rightarrow x_{T-1} \rightarrow \cdots \rightarrow x_0
\]

这里每一步有随机噪声，所以可以探索。Group 里不同样本、不同 trajectory 的差异，来自 SDE 注入噪声。

\textbf{第二条是打分/归因路径：ODE readout。}

对于已经由 SDE 采样得到的某个中间 latent \(x_t\)，DenseGRPO 额外做一次 deterministic ODE denoising：

\[
x_t
\xrightarrow{\mathrm{ODE}}
\hat{x}_{t,0}
\]

这一步不是为了探索，而是为了回答一个评估问题：

\begin{intubox}
如果当前已经走到这个 \(x_t\)，并且之后不再随机，只按确定性路径生成，它大概会通向什么 clean image？
\end{intubox}

然后对这个 clean prediction 打 reward：

\[
R_t = R(\hat{x}_{t,0}, c)
\]

所以 ODE 的角色是 \key{probe / readout / value estimation}，不是 policy rollout。

\[
\text{SDE = 产生探索样本}
\qquad
\text{ODE = 把中间 latent 映射成可打分的 clean prediction}
\]

DenseGRPO 的 dense reward：

\[
\Delta R_t = R_{t-1} - R_t
\]

可以理解为：沿着真实 SDE trajectory，从 \(x_t\) 走到 \(x_{t-1}\) 后，如果分别用 ODE 读出它们未来会变成什么 clean image，那么 clean-image reward 提升了多少。

\begin{warnbox}
\textbf{不要把两件事混在一起：} SDE 负责探索，ODE 负责评估。DenseGRPO 不是用 ODE 探索，而是先用 SDE 得到中间 latent，再用 ODE 把这些 latent 变成 reward model 可以评价的 clean-image proxy。
\end{warnbox}

\begin{figure}[htbp]
  \centering
  \includegraphics[width=\linewidth]{densegrpo/fig2_method_overview.png}
  \caption{DenseGRPO 总体流程：对中间 latent 做 ODE denoising，得到对应 clean image 的 reward，再用相邻 reward gain 构造 dense reward。图源：arXiv HTML。}
  \label{fig:densegrpo-method-overview}
\end{figure}

\subsubsection{4.3 新 advantage 怎么算}

拿到每个样本、每个 timestep 的 dense reward gain \(\Delta R^i_t\) 后，DenseGRPO 不再用整图 reward 做 group normalization，而是在每个 timestep 上单独做 group normalization：

\[
\hat{A}^i_t
=
\frac{
\Delta R^i_t - \mathrm{mean}\{\Delta R^j_t\}_{j=1}^{G}
}{
\mathrm{std}\{\Delta R^j_t\}_{j=1}^{G}
}
\]

这一步非常重要。它意味着：

\[
\text{第 } t \text{ 步的 policy update}
\quad
\text{只看第 } t \text{ 步的相对贡献}
\]

于是，不同 timestep 可以得到不同方向、不同强度的学习信号。

\subsection{5. 模块二：Exploration Space Calibration}

Flow-GRPO 为了让 deterministic flow model 可以做 RL 探索，把 ODE sampler 改成 SDE sampler，在每一步注入噪声：

\[
\sigma_t \epsilon
\]

原来的做法通常用一个统一噪声强度 \(a\)。DenseGRPO 认为这也不合理：不同 timestep 的作用不同，所需探索强度也不同。

\subsubsection{5.1 为什么统一噪声不对}

如果某个 timestep 噪声太大，生成轨迹会偏离合理分布，dense reward 可能几乎全是负的；如果噪声太小，group 内样本差异不够，RL 没有有效探索。

论文用 dense reward 的正负分布来判断探索是否合适：

\[
\Delta R^i_t > 0
\quad \text{vs.} \quad
\Delta R^i_t < 0
\]

如果一个 timestep 上正负样本比较均衡，说明该步既有好探索也有坏探索，训练信号更有区分度；如果几乎全正或全负，说明探索空间不合适。

\subsubsection{5.2 校准规则}

DenseGRPO 为每个 timestep 学一个噪声强度函数 \(\psi(t)\)。校准逻辑很简单：

\begin{itemize}[leftmargin=1.5em]
  \item 如果某个 timestep 的 positive / negative dense rewards 比较均衡，就轻微增加 \(\psi(t)\)，鼓励更多探索。
  \item 如果不均衡，就轻微降低 \(\psi(t)\)，避免过度或无效探索。
\end{itemize}

最终采样时使用：

\[
\sigma_t = \psi(t)
\]

而不是所有 timestep 共享一个标量 \(a\)。

\begin{intubox}
\textbf{直觉版：} Dense reward 不只是训练 signal，它还变成了诊断工具：用来判断每个时间步的探索噪声是否合适。
\end{intubox}

\begin{figure}[htbp]
  \centering
  \includegraphics[width=\linewidth]{densegrpo/fig3_dense_reward_noise_calibration.png}
  \caption{不同噪声强度下的 dense reward 分布。固定噪声会在部分 timestep 造成奖励极度不平衡；DenseGRPO 用 timestep-specific \(\psi(t)\) 校准探索空间。图源：arXiv HTML。}
  \label{fig:densegrpo-noise-calibration}
\end{figure}

\subsection{6. 整体训练流程}

可以把 DenseGRPO 的训练理解成下面几步：

\begin{enumerate}[leftmargin=1.5em]
  \item 对同一个 prompt 采样 \(G\) 条 SDE trajectories。
  \item 对每条 trajectory 的每个 intermediate latent \(x^i_t\)，用 ODE denoising 得到 \(\hat{x}^i_{t,0}\)。
  \item 对 \(\hat{x}^i_{t,0}\) decode 后用 reward model 打分，得到 \(R^i_t\)。
  \item 用相邻 reward 差值计算 step-wise dense reward：
  \[
  \Delta R^i_t = R^i_{t-1} - R^i_t
  \]
  \item 对每个 timestep，在 group 内 normalize \(\Delta R^i_t\)，得到 \(\hat{A}^i_t\)。
  \item 用原来的 GRPO clipped objective 更新 policy，只是 advantage 从 trajectory-level 换成 step-level。
  \item 训练前用 dense reward 分布校准 \(\psi(t)\)，让每个 timestep 的 SDE exploration 更合理。
\end{enumerate}

\subsection{7. 实验看什么}

论文沿用 Flow-GRPO 的三类 text-to-image 任务：

\begin{itemize}[leftmargin=1.5em]
  \item \textbf{Compositional Image Generation：} reward model 用 GenEval。
  \item \textbf{Human Preference Alignment：} reward model 用 PickScore。
  \item \textbf{Visual Text Rendering：} reward 来自 OCR accuracy。
\end{itemize}

主要对比：

\begin{itemize}[leftmargin=1.5em]
  \item Flow-GRPO：terminal reward 直接分配给所有 timestep。
  \item Flow-GRPO+CoCA：用 latent similarity 给 step-wise feedback。
  \item DenseGRPO：用 ODE-based intermediate clean image 估计 dense reward gain。
\end{itemize}

论文报告 DenseGRPO 在三类任务上都优于这些 baseline。比较关键的是，DenseGRPO 相对 Flow-GRPO+CoCA 也有明显收益，这说明它的优势不只是“用了 step-wise reward”，而是 \key{ODE-based latent reward 比 latent similarity 更贴近真实贡献}。

\begin{figure}[htbp]
  \centering
  \includegraphics[width=\linewidth]{densegrpo/fig4_main_results.png}
  \caption{主实验结果柱状图。DenseGRPO 在组合图像生成、人类偏好对齐和视觉文字渲染任务上都优于主要 baseline。图源：arXiv HTML。}
  \label{fig:densegrpo-main-results}
\end{figure}

\begin{figure}[htbp]
  \centering
  \includegraphics[width=0.92\linewidth]{densegrpo/fig5_qualitative_results.png}
  \caption{质性结果对比。论文用这些样例说明 DenseGRPO 在组合关系、语义一致性和文本渲染上有更好的对齐效果。图源：arXiv HTML。}
  \label{fig:densegrpo-qualitative}
\end{figure}

\subsection{8. Ablation 该怎么读}

\begin{center}
  \includegraphics[width=\linewidth]{densegrpo/fig6_ablation_analysis.png}

  {\small \textbf{图 6：} 消融实验。Dense reward、timestep-specific noise 和 ODE denoising steps 分别对应这里的三个核心问题：step-wise reward 是否有用、\(\psi(t)\) 是否比固定噪声更好、ODE readout 步数是否影响 reward 估计质量。图源：arXiv HTML。}
\end{center}

\subsubsection{8.1 Dense reward 是否真的有用}

论文做了一个 Dense Reward Baseline：不用 reward gain，只直接用 \(R^i_{t-1}\) 来优化第 \(t\) 步。它比 Flow-GRPO 好，说明 dense reward 本身有价值；但 DenseGRPO 更好，说明“相邻 reward gain”比“下一步 latent reward 绝对值”更合理。

\[
R^i_{t-1}
\quad \text{不如} \quad
R^i_{t-1} - R^i_t
\]

原因是绝对 reward 仍然混入了未来步骤的累计贡献，而 reward gain 更接近当前一步的边际贡献。

\subsubsection{8.2 Timestep-specific noise 是否有用}

论文比较了固定噪声 \(a=0.7\) 和校准后的 \(\psi(t)\)。结论是 \(\psi(t)\) 更好。

这说明 DenseGRPO 的第二个模块不是附属技巧：一旦你有了 step-wise dense rewards，就能看到哪些时间步探索过强或过弱，并据此调 sampling noise。

\subsubsection{8.3 ODE denoising steps 的成本与收益}

DenseGRPO 需要对每个 intermediate latent 做 ODE denoising。论文比较 \(n=1,2,t\)，发现 \(n=t\) 最好，但计算也更贵。

\[
\text{full readout cost}
\approx T + (T-1) + \cdots + 1
= \frac{T(T+1)}{2}
\]

也就是说，如果 rollout 本身有 \(T\) 个 denoising steps，而每个 \(x_t\) 都用完整 ODE 读出到 clean image，那么 dense reward 的额外模型调用会从线性的 \(O(T)\) 接近变成平方级 \(O(T^2)\)。这还没有算 VAE decode 和 reward model scoring 的额外开销。

论文在 Fig. 6(c) 里确实报告了训练成本：在相同实验设置下，\(n=1\)、\(n=2\)、\(n=t\) 分别需要 11、13、19 GPU hours 来训练 20 steps。也就是说，最准确的 \(n=t\) readout 明显更贵；论文的说法是更大的 \(n\) 虽然有更高计算开销，但同等 GPU training time 下收益更好。

\begin{warnbox}
\textbf{工程代价：} DenseGRPO 的核心收益来自更准的 dense reward，但代价是多次 ODE denoising、VAE decode 和 reward model scoring。它不是免费的加速方法，而是更精细、更贵的训练信号。对于视频模型，这个问题会被进一步放大，因为单次 denoising forward 和 reward evaluation 都远贵于图像模型。
\end{warnbox}

\textbf{我的判断：} full \(n=t\) 更像是一个证明 dense reward gain 上限有效的版本。实际做大图像模型或视频模型后训练时，更现实的实现可能是：只对少量关键 timesteps 做 readout，或者固定 \(n=1,2,k\) 做 truncated readout，再用插值、chunk 或 learned proxy 来近似 step-wise advantage。

\subsection{9. 和 TempFlow-GRPO / Chunk-GRPO 的关系}

这三篇都在回答 Flow-GRPO 的 credit assignment 问题，但切入点不同：

\begin{itemize}[leftmargin=1.5em]
  \item \textbf{TempFlow-GRPO：} 问“哪个 timestep 更关键”，用 trajectory branching 看 reward 差异。
  \item \textbf{Chunk-GRPO：} 问“step 太细、trajectory 太粗怎么办”，把连续 steps 合成 chunk。
  \item \textbf{DenseGRPO：} 问“每一步真实贡献怎么直接估计”，用 ODE-based reward gain 做 dense reward。
\end{itemize}

如果只读一篇来理解这条线，DenseGRPO 最适合作为核心，因为它最直接地把 terminal reward 拆成了 step-wise signal。

\subsection{10. 局限与风险}

\begin{itemize}[leftmargin=1.5em]
  \item \textbf{计算成本高：} 每个 timestep 的 intermediate latent 都可能需要 ODE denoising、VAE decode 和 reward scoring。若使用完整 \(n=t\) readout，额外 denoising 调用接近 \(T(T+1)/2\)，比普通 \(O(T)\) rollout 贵很多。
  \item \textbf{论文只给了局部成本：} 实验使用 16 张 NVIDIA A100。论文报告 \(n=1,2,t\) 训练 20 steps 分别需要 11、13、19 GPU hours，但没有给出完整 4500/1500/4500 iteration 训练的 wall-clock time。
  \item \textbf{依赖 reward model 质量：} 如果 reward model 有偏，dense reward 会把偏差更精细地传播到每一步。
  \item \textbf{可能放大奖励黑客：} 论文附录也提到，更准确的 step-wise reward 可能让模型更容易过拟合某些 reward model。
  \item \textbf{适用前提：} 需要能从 intermediate latent 通过 deterministic sampler 映射到 clean image。对没有稳定 deterministic path 的模型，需要额外改造。
\end{itemize}

\subsection{11. 对图像编辑后训练的启发}

DenseGRPO 对图像编辑很有启发，因为图像编辑的 reward 更不应该只看最终图：

\[
\text{edit region 应该变}
\quad
\text{preserve region 不应该变}
\]

如果未来做 image editing RL，可以借鉴 DenseGRPO 的想法：

\begin{itemize}[leftmargin=1.5em]
  \item 对每个 timestep 还原 intermediate clean image。
  \item 分别计算 edit reward、preservation reward、identity reward、aesthetic reward。
  \item 用相邻 timestep 的 reward gain 判断哪一步真正造成了目标区域变化或非目标区域破坏。
  \item 再把不同区域、不同 reward component 的 gain 分给对应 timestep。
\end{itemize}

\begin{summarybox}
\textbf{可复用点：} DenseGRPO 的真正可迁移贡献不是某个公式，而是一个范式：
\[
\begin{aligned}
\text{intermediate latent}
&\rightarrow \text{deterministic clean prediction}\\
&\rightarrow \text{reward model scoring}\\
&\rightarrow \text{step-wise reward gain}
\end{aligned}
\]
这个范式可以迁移到图像编辑、视频生成、甚至 discrete diffusion 的后训练。
\end{summarybox}

\subsection{12. 读完要记住的三个问题}

\begin{enumerate}[leftmargin=1.5em]
  \item \textbf{为什么 terminal reward 不能直接分给所有 timestep？}\\
  因为最终结果是所有 steps 的累计贡献，单步更新需要的是当前 step 的边际贡献。
  \item \textbf{DenseGRPO 怎么给 intermediate latent 打分？}\\
  用 deterministic ODE denoising 把 \(x_t\) 映射到 clean image，再用已有 reward model 打分。
  \item \textbf{为什么 dense reward 还能调探索噪声？}\\
  因为 dense reward 的正负分布可以反映某个 timestep 的探索是否过强或过弱。
\end{enumerate}

% ============================================================
% Paper deep reading: Chunk-GRPO
% ============================================================

\section{论文精读：Chunk-GRPO}

\begin{conceptbox}
\textbf{论文：} \href{https://arxiv.org/abs/2510.21583}{Sample By Step, Optimize By Chunk: Chunk-Level GRPO For Text-to-Image Generation}\\
\textbf{简称：} Chunk-GRPO。\\
\textbf{版本：} arXiv 2510.21583v1，submitted on 2025-10-24。\\
\textbf{方向标签：} Flow Matching、Text-to-Image、GRPO、Chunk-level Optimization、Credit Assignment。\\
\textbf{阅读日期：} 2026-04-27。\\
\textbf{前置：} 已理解 Flow-GRPO 的 step-level GRPO objective，以及 DenseGRPO 为什么要用 step-wise reward gain。
\end{conceptbox}

\subsection{0. 一句话结论}

Chunk-GRPO 的核心不是“重新设计 reward”，而是把 GRPO 的优化粒度从 \textbf{step-level} 改成 \textbf{chunk-level}：

\[
\text{sample by step}
\quad \rightarrow \quad
\text{optimize by chunk}
\]

它要解决的问题是：单个 denoising step 太细，整条 trajectory 又太粗；更合理的优化单位可能是一段连续 timesteps，也就是 chunk。

\begin{summarybox}
\textbf{我的判断：} Chunk-GRPO 是 DenseGRPO 后最应该读的一篇，因为它正好回应 DenseGRPO 的工程痛点：full step-wise ODE readout 太贵。Chunk-GRPO 提供了一个低成本方向：不逐步估计 reward gain，而是把生成过程切成若干阶段，在 chunk 级别平滑 policy update。
\end{summarybox}

\begin{warnbox}
\textbf{先把边界说清楚：} Chunk-GRPO 并没有像 DenseGRPO 那样真的计算每个 chunk 的 reward improvement。它主要仍然使用整条 trajectory 的 terminal advantage \(A_i\)，改的是 policy ratio 和优化粒度。因此它缓解 credit assignment，但没有彻底解决“哪个 chunk 真正带来 reward improvement”的问题。
\end{warnbox}

\begin{center}
  \includegraphics[width=\linewidth]{chunkgrpo/fig1_qualitative_results.jpg}

  {\small \textbf{图 1：} Chunk-GRPO 的质性结果。论文用这些样例说明 chunk-level optimization 可以改善结构、光照和细节质量。图源：arXiv HTML。}
\end{center}

\subsection{1. 它接在 DenseGRPO 后面读的原因}

DenseGRPO 的关键贡献是：

\[
\Delta R_t = R_{t-1} - R_t
\]

这比把 terminal reward 平均分给所有 steps 更合理。但问题是，如果每个 \(x_t\) 都完整 ODE readout 到 clean image，成本接近：

\[
T + (T-1) + \cdots + 1
= \frac{T(T+1)}{2}
\]

Chunk-GRPO 的思路不是继续追求每一步都准，而是换一个问题：

\begin{quote}
如果 step 太细、trajectory 太粗，那么能不能找一个中间时间尺度？
\end{quote}

这个中间尺度就是 chunk：

\[
x_T \rightarrow \cdots \rightarrow x_0
\quad \Rightarrow \quad
\mathrm{ch}_1,\mathrm{ch}_2,\ldots,\mathrm{ch}_K
\]

所以它和 DenseGRPO 的关系可以这样理解：

\[
\text{DenseGRPO：每一步的 reward gain}
\]

\[
\text{Chunk-GRPO：每一段的 policy update unit}
\]

\subsection{2. Flow-GRPO 的问题：step-level 更新太细}

Flow-GRPO / Dance-GRPO 这类方法通常对同一个 prompt 采样一组 trajectories，然后用最终 reward 得到 group relative advantage：

\[
A_i
=
\frac{
R(x_0^i,c)-\mathrm{mean}\{R(x_0^j,c)\}_{j=1}^{G}
}{
\mathrm{std}\{R(x_0^j,c)\}_{j=1}^{G}
}
\]

然后每个 step 都拿同一个 \(A_i\) 做更新：

\[
\sum_{t=1}^{T}
\min
\left(
r_t^i(\theta)A_i,\;
\mathrm{clip}(r_t^i(\theta),1-\epsilon,1+\epsilon)A_i
\right)
\]

问题是：最终图像 reward 高，不代表这条 trajectory 的每个中间 step 都比别的 trajectory 好。

\begin{center}
  \includegraphics[width=0.72\linewidth]{chunkgrpo/fig2_motivation.png}

  {\small \textbf{图 2：} Chunk-GRPO 的动机图。Trajectory 1 最终 reward 更高，但其中某个中间 timestep 可能比 Trajectory 2 更差；标准 GRPO 仍会把正 advantage 平等分给所有 steps。图源：arXiv HTML。}
\end{center}

\textbf{这和 DenseGRPO 的问题意识一致：} terminal reward 的全局平均分配是不合理的。区别在于，DenseGRPO 尝试重新估计每步 reward；Chunk-GRPO 则尝试改变优化单位。

\subsection{3. 方法一：Chunk-level policy ratio}

给定一条生成 trajectory：

\[
(x_T,x_{T-1},\ldots,x_1,x_0)^i
\]

Chunk-GRPO 把它切成 \(K\) 个 chunk：

\[
\{\mathrm{ch}_1,\mathrm{ch}_2,\ldots,\mathrm{ch}_K\}^i
\]

其中第 \(j\) 个 chunk 的长度是 \(cs_j\)，所有 chunk size 加起来覆盖 denoising trajectory：

\[
\sum_{j=1}^{K} cs_j = T
\]

在实验里，论文使用 \(T=17\) 个 sampling steps，并采用：

\[
\{cs_j\}_{j=1}^{4}=[2,3,4,7]
\]

这里实际只覆盖 16 个 transitions，因为论文跟随 Dance-GRPO 忽略最后一个不引入 stochasticity 的 timestep。

\subsubsection{3.1 从 step ratio 到 chunk ratio}

标准 step-level GRPO 的 ratio 是：

\[
r_t^i(\theta)
=
\frac{
p_\theta(x_{t-1}^i|x_t^i,c)
}{
p_{\mathrm{old}}(x_{t-1}^i|x_t^i,c)
}
\]

Chunk-GRPO 把一个 chunk 内所有 step 的 likelihood ratio 合起来，定义 chunk-level ratio：

\[
r_j^i(\theta)
=
\left(
\prod_{t\in \mathrm{ch}_j}
\frac{
p_\theta(x_{t-1}^i|x_t^i,c)
}{
p_{\mathrm{old}}(x_{t-1}^i|x_t^i,c)
}
\right)^{1/cs_j}
\]

这个 \(1/cs_j\) 很关键。它相当于几何平均，而不是简单乘积。否则 chunk 越长，ratio 的尺度越容易爆。

\textbf{几何平均的定义：} 对 \(n\) 个正数 \(a_1,a_2,\ldots,a_n\)，它们的几何平均定义为：

\[
\mathrm{GM}(a_1,\ldots,a_n)
=
\left(
\prod_{k=1}^{n} a_k
\right)^{1/n}
\]

等价地，也可以写成 log-space 的平均：

\[
\mathrm{GM}(a_1,\ldots,a_n)
=
\exp
\left(
\frac{1}{n}
\sum_{k=1}^{n}\log a_k
\right)
\]

所以几何平均不是把所有数简单乘起来，而是先得到总的乘法变化，再除以长度 \(n\)，还原成“平均每一个元素贡献了多少倍的乘法变化”。如果所有 \(a_k=r\)，那么 \(\mathrm{GM}=r\)，而简单乘积会变成 \(r^n\)。

更具体地说，chunk 内每一步都有一个 ratio：

\[
r_t^i(\theta)
=
\frac{
p_\theta(x_{t-1}^i|x_t^i,c)
}{
p_{\mathrm{old}}(x_{t-1}^i|x_t^i,c)
}
\]

如果直接把它们相乘：

\[
\prod_{t\in \mathrm{ch}_j} r_t^i(\theta)
\]

得到的是整个 chunk trajectory likelihood 的 ratio。这在数学上没错，但有一个工程问题：\textbf{chunk 越长，乘积越容易偏离 1}。例如如果每一步 ratio 都只是轻微变大，\(r_t=1.05\)，那么：

\[
1.05^2=1.10,\quad
1.05^{10}=1.63,\quad
1.05^{20}=2.65
\]

同样的单步变化，因为 chunk 长度不同，会被放大成完全不同的 ratio 尺度。这样 PPO/GRPO 的 clipping 会更频繁触发，长 chunk 和短 chunk 也不再可比。

因此论文用：

\[
\left(
\prod_{t\in \mathrm{ch}_j}r_t^i(\theta)
\right)^{1/cs_j}
=
\exp
\left(
\frac{1}{cs_j}
\sum_{t\in \mathrm{ch}_j}
\log r_t^i(\theta)
\right)
\]

这就是几何平均。它等价于先对 chunk 内每一步的 \(\log\)-ratio 求平均，再 exponentiate 回 ratio 空间。

\begin{summarybox}
\textbf{直觉：} 简单乘积衡量的是“整个 chunk 的联合概率变了多少”；几何平均衡量的是“chunk 内平均每一步的概率变了多少”。Chunk-GRPO 需要后者，因为它要把不同长度的 chunks 放进同一个 PPO/GRPO clipping 框架里比较。
\end{summarybox}

\textbf{一个更直观的例子：} 如果两个 chunks 的每一步平均变化都差不多，那么它们应该有差不多的 policy ratio，而不应该因为一个 chunk 长度是 2、另一个长度是 10，就让后者天然更容易被 clip。\(1/cs_j\) 的作用就是做长度归一化。

\subsubsection{3.2 Chunk-level objective}

Chunk-GRPO 的目标变成：

\[
J(\theta)
=
\mathbb{E}
\left[
\frac{1}{GK}
\sum_{i=1}^{G}
\sum_{j=1}^{K}
\min
\left(
r_j^i(\theta)A_i,\;
\mathrm{clip}(r_j^i(\theta),1-\epsilon,1+\epsilon)A_i
\right)
-\beta D_{\mathrm{KL}}(\pi_\theta||\pi_{\mathrm{ref}})
\right]
\]

注意这里仍然是 \(A_i\)，不是 \(A_{i,j}\)。

\begin{warnbox}
\textbf{关键判断：} Chunk-GRPO 的 chunk-level objective 没有给每个 chunk 计算独立 reward 或 independent advantage。它把多个 step 的 ratio 合成一个 chunk ratio，从而让一个 chunk 内的 transition 作为整体被优化。这个改动主要带来 smoother gradient 和更合理的时间粒度，但不是 DenseGRPO 那种显式 reward gain。
\end{warnbox}

\subsection{4. 为什么同样用 terminal advantage 还能有用}

这点容易误解：如果 Chunk-GRPO 仍然用 \(A_i\)，为什么它还能改善 credit assignment？

我的理解是，它缓解的不是“reward 从哪里来”，而是“policy update 以什么单位发生”。

\[
\text{step-level：每个 transition 单独吃同一个 } A_i
\]

\[
\text{chunk-level：一段 transitions 共同吃同一个 } A_i
\]

当单个 step 的局部质量和最终 reward 不一致时，step-level 会把噪声直接打到每个 transition 上；chunk-level 会把相邻 steps 合成一个单位，等价于在时间维度上做 smoothing。

所以 Chunk-GRPO 的收益更像是：

\[
\text{减少 step-wise noisy credit}
\quad
\neq
\quad
\text{得到真正 chunk-wise reward}
\]

这也是它和 DenseGRPO 最大的区别。

\subsection{5. 方法二：用 temporal dynamics 决定 chunk}

如果只是随便切 chunk，效果不一定好。论文先做了固定 chunk size 的实验，发现不同 chunk size 的性能差异明显。

\begin{center}
  \includegraphics[width=0.58\linewidth]{chunkgrpo/fig5_chunk_size_ablation.png}

  {\small \textbf{图 5：} 不同 chunk size 的性能不同，说明 chunk 不是随便切就行。TD 表示 temporal dynamics guided chunking。图源：arXiv HTML。}
\end{center}

论文认为 flow matching 的 denoising trajectory 有明显的 temporal dynamics：不同时间步的 latent change 不一样。它用相邻 latent 的相对 \(L1\) 距离来刻画：

\[
L1_{\mathrm{rel}}(x,t)
=
\frac{\|x_t-x_{t-1}\|_1}{\|x_t\|_1}
\]

直觉是：

\[
L1_{\mathrm{rel}} \text{ 大}
\Rightarrow
\text{latent 变化快，通常是高噪声 / 结构阶段}
\]

\[
L1_{\mathrm{rel}} \text{ 小}
\Rightarrow
\text{latent 变化慢，通常是后期稳定 refinement}
\]

\begin{center}
  \includegraphics[width=\linewidth]{chunkgrpo/fig3_temporal_dynamics.png}

  {\small \textbf{图 3：} flow matching 的 relative \(L1\) dynamics。论文观察到这种动态模式对 prompt 相对稳定，但随 timestep 明显变化，因此可以用它指导 chunk 边界。图源：arXiv HTML。}
\end{center}

论文的原则是：

\[
\text{dynamics 相似的 timesteps 放进同一个 chunk}
\]

\[
\text{dynamics 差异大的 timesteps 分到不同 chunk}
\]

最终默认设置是 4 个 chunk：

\[
[2,3,4,7]
\]

它不是严格证明出的最优切分，而是基于 relative \(L1\) 曲线、二阶变化和训练前后 dynamics 观察得到的工程选择。

\begin{center}
  \includegraphics[width=0.66\linewidth]{chunkgrpo/fig10_l1_before_after_dancegrpo.png}

  {\small \textbf{图 10：} Dance-GRPO 训练前后 relative \(L1\) dynamics 的变化。论文据此说明 RL 主要改变后半段 timesteps，并把最后一个 chunk 设置得更长。图源：arXiv HTML。}
\end{center}

\subsection{6. 方法三：Weighted Sampling}

论文还提出一个 optional weighted sampling strategy。它不是每次更新所有 chunks，而是像 Dance-GRPO 一样只选一部分 chunks，例如 fraction \(0.5\)。区别是 Chunk-GRPO 不做 uniform sampling，而是按 chunk 的平均 relative \(L1\) 给采样权重：

\[
w(\mathrm{ch}_j)
=
\frac{
\frac{1}{cs_j}\sum_{t\in \mathrm{ch}_j}L1_{\mathrm{rel}}(x,t)
}{
\frac{1}{T}\sum_{t=1}^{T}L1_{\mathrm{rel}}(x,t)
}
\]

这个公式会偏向高噪声 / 高变化 chunk。

\begin{summarybox}
\textbf{直觉：} 高噪声阶段通常决定结构、构图和大尺度语义，所以更值得优化；低噪声阶段更多是 refinement，优化信号可能小一些。
\end{summarybox}

但是它也有明显风险：高噪声阶段一旦过度优化，容易改变整体结构，甚至导致语义崩塌。

\begin{center}
  \includegraphics[width=\linewidth]{chunkgrpo/fig9_weighted_sampling_failure.jpg}

  {\small \textbf{图 9：} weighted sampling 的失败例子。论文指出它可能在高噪声区域错误改变图像结构，导致语义元素缺失。图源：arXiv HTML。}
\end{center}

\subsection{7. 整体框架}

Chunk-GRPO 的流程可以概括为：

\begin{enumerate}[leftmargin=1.5em]
  \item 用 SDE 采样生成 trajectory。
  \item 计算 relative \(L1\) dynamics。
  \item 根据 dynamics 预设 chunk boundaries。
  \item 把 step-level ratio 合成 chunk-level ratio。
  \item 用 terminal group relative advantage \(A_i\) 做 chunk-level GRPO update。
  \item 可选：按 \(w(\mathrm{ch}_j)\) 采样更重要的 chunks。
\end{enumerate}

\begin{center}
  \includegraphics[width=\linewidth]{chunkgrpo/fig4_framework.png}

  {\small \textbf{图 4：} Chunk-GRPO 框架。核心是 temporal-dynamic-guided chunking、chunk-level importance ratio，以及可选 weighted sampling。图源：arXiv HTML。}
\end{center}

\subsection{8. 实验结果怎么读}

\subsubsection{8.1 主实验}

实验设置：

\begin{itemize}[leftmargin=1.5em]
  \item Base model：FLUX.1 Dev。
  \item Baseline：Dance-GRPO。
  \item 训练 reward model：HPSv3。
  \item 数据：HPDv2.1。
  \item 评估：HPSv3、ImageReward、WISE、GenEval。
  \item 默认 chunk：\(T=17\)，\([2,3,4,7]\)。
\end{itemize}

偏好对齐结果：

\begin{center}
\begin{tabular}{lcc}
\toprule
\textbf{Model} & \textbf{HPSv3} & \textbf{ImageReward} \\
\midrule
FLUX & 13.804 & 1.086 \\
Dance-GRPO & 15.080 & 1.141 \\
Chunk-GRPO w/o ws & 15.236 & 1.147 \\
Chunk-GRPO w/ ws & 15.373 & 1.149 \\
\bottomrule
\end{tabular}
\end{center}

这里可以看到：weighted sampling 对 preference metrics 有正收益。

但在 WISE 上，Chunk-GRPO w/o ws overall 是 0.76，Dance-GRPO 是 0.75，Chunk-GRPO w/ ws 反而降到 0.73。这说明 weighted sampling 的偏好提升可能以结构或语义稳定性为代价。

\subsubsection{8.2 Chunk setting 消融}

论文对不同 chunk size 做了消融。关键结果是：

\begin{itemize}[leftmargin=1.5em]
  \item 固定 chunk size 通常能略优于 step-level Dance-GRPO。
  \item 不同 chunk size 表现差异明显。
  \item temporal-dynamics-guided chunking \([2,3,4,7]\) 最好，HPSv3 达到 15.236。
\end{itemize}

\[
\text{chunk size matters}
\quad \Rightarrow \quad
\text{不能随便把 steps 平均切块}
\]

\subsubsection{8.3 只训练特定 chunk}

论文还只训练某一个 chunk，观察哪个 chunk 贡献大。

\begin{center}
  \includegraphics[width=0.66\linewidth]{chunkgrpo/fig6_specific_chunk.png}

  {\small \textbf{图 6：} 单独训练不同 chunks。高噪声 chunk 收益更大，但后期也更不稳定。图源：arXiv HTML。}
\end{center}

结论是：早期高噪声 chunks 带来更大提升，但训练后期容易不稳定。这正是 weighted sampling 的动机，也是它失败风险的来源。

\subsubsection{8.4 不同 reward model 和 GenEval}

论文还换了 reward model，例如 PickScore、CLIP。Chunk-GRPO 仍然优于 Dance-GRPO。

GenEval overall：

\begin{center}
\begin{tabular}{lc}
\toprule
\textbf{Model} & \textbf{GenEval Overall} \\
\midrule
FLUX & 0.66 \\
Dance-GRPO & 0.67 \\
Chunk-GRPO w/o ws & 0.69 \\
Chunk-GRPO w/ ws & 0.67 \\
\bottomrule
\end{tabular}
\end{center}

这里再次说明：\textbf{w/o weighted sampling 更稳}。Weighted sampling 会追求偏好分数，但可能损害 instruction following。

\subsection{9. 和 DenseGRPO 的本质区别}

\begin{center}
{\small
\begin{tabular}{p{0.20\linewidth}p{0.36\linewidth}p{0.36\linewidth}}
\toprule
\textbf{维度} & \textbf{DenseGRPO} & \textbf{Chunk-GRPO} \\
\midrule
核心动作 & 计算中间 latent 的 reward gain & 改变 policy optimization 粒度 \\
信用粒度 & step-wise & chunk-wise update unit \\
advantage & \(\Delta R_t=R_{t-1}-R_t\) & 主要仍用 terminal \(A_i\) \\
成本 & ODE readout + decode + reward scoring，贵 & 几乎不额外打 reward，便宜 \\
解决程度 & 更接近真实 step contribution & 平滑 step-level noisy attribution \\
主要风险 & 成本高、reward model bias 被细粒度传播 & chunk boundary 固定、chunk 内仍可能被 terminal advantage 误导 \\
\bottomrule
\end{tabular}
}
\end{center}

所以我会把两者关系总结成：

\[
\text{DenseGRPO = 更准但更贵}
\]

\[
\text{Chunk-GRPO = 更便宜但更粗}
\]

它们不是互相替代，而是可以组合。

\subsection{10. 对我们后续 idea 的启发}

读完 DenseGRPO 和 Chunk-GRPO 后，最自然的改进方向是：

\[
\text{Chunk-GRPO 的 chunk unit}
\quad + \quad
\text{DenseGRPO 的 reward gain}
\]

也就是不要每个 step 都 ODE readout，而是在 chunk boundary 上做 sparse readout：

\[
R_{\tau_K},R_{\tau_{K-1}},\ldots,R_{\tau_0}
\]

然后定义 chunk-level improvement：

\[
\Delta R_{\mathrm{ch}_j}
=
R_{\tau_{j-1}} - R_{\tau_j}
\]

这样成本从 DenseGRPO 的：

\[
O(T^2)
\]

变成更接近：

\[
O(K\cdot \text{readout cost})
\]

其中 \(K \ll T\)。

\begin{summarybox}
\textbf{最值得做的组合：} 用 Chunk-GRPO 的 temporal-dynamic-guided chunks 决定 readout 边界，再用 DenseGRPO 的 ODE-based reward gain 计算 chunk-level improvement。这样既保留“reward improvement 不该全局平分”的核心收益，又避免每个 step 都 full ODE readout。
\end{summarybox}

\subsection{11. 局限与风险}

\begin{itemize}[leftmargin=1.5em]
  \item \textbf{没有真正 chunk-specific reward：} Chunk-GRPO 主要仍用 terminal \(A_i\)，所以不同 chunk 没有独立 reward gain。
  \item \textbf{chunk boundary 固定：} 论文默认 \([2,3,4,7]\)，训练过程中不动态变化。论文自己也把 adaptive chunking 列为未来方向。
  \item \textbf{weighted sampling 有偏好-语义 tradeoff：} 它提升 HPSv3 / ImageReward，但可能降低 WISE / GenEval，甚至造成结构崩塌。
  \item \textbf{对视频还不够：} 视频里的 temporal consistency、motion、identity 可能需要不同 chunk reward，而不是只按 latent \(L1\) 切分。
  \item \textbf{对图像编辑还缺 region-awareness：} 编辑任务不能只按 denoising time 切，还要区分目标区域和非目标区域。
\end{itemize}

\subsection{12. 读完要记住的三句话}

\begin{enumerate}[leftmargin=1.5em]
  \item Chunk-GRPO 的核心是 \textbf{优化粒度}：step 太细、trajectory 太粗，chunk 是折中。
  \item 它的 chunk-level ratio 是 step ratios 的几何平均，能平滑 step-level noisy updates。
  \item 它还没有真正解决 chunk-wise reward assignment；最有价值的下一步是把 DenseGRPO 的 reward gain 搬到 chunk boundary 上。
\end{enumerate}

% ============================================================
% Paper deep reading: TempFlow-GRPO
% ============================================================

\section{论文精读：TempFlow-GRPO}

\begin{conceptbox}
\textbf{论文：} \href{https://arxiv.org/abs/2508.04324}{TempFlow-GRPO: When Timing Matters for GRPO in Flow Models}\\
\textbf{简称：} TempFlow-GRPO。\\
\textbf{版本：} arXiv 2508.04324v4，submitted on 2025-08-06，last revised on 2025-10-15；OpenReview 为 ICLR 2026 Poster。\\
\textbf{方向标签：} Flow Matching、GRPO、Temporal Credit Assignment、Trajectory Branching、Noise-aware Weighting。\\
\textbf{阅读日期：} 2026-04-27。\\
\textbf{前置：} 已读 DenseGRPO 和 Chunk-GRPO，当前关心“step-level reward 是否必要、chunk 是否只是平滑 credit、如何避免 A+B 式组合”。
\end{conceptbox}

\subsection{0. 一句话结论}

TempFlow-GRPO 最值得关注的点不是“又一个 Flow-GRPO variant”，而是它提出了一种构造 step-level process reward 的新方式：

\[
\begin{aligned}
&\text{不是直接给中间 latent 打分}\\
&\rightarrow\quad
\text{在某个 timestep 做 controlled branching，}\\
&\hspace{2.2em}\text{再用最终 reward 反推该 step 的贡献}
\end{aligned}
\]

它的核心动作是：

\[
\begin{aligned}
\text{ODE 到 } x_k
&\rightarrow
\text{只在 } k \text{ 做一次 SDE 探索}\\
&\rightarrow
\text{ODE 到 } x_0
\rightarrow
\text{用最终 reward 归因给 } k
\end{aligned}
\]

\begin{summarybox}
\textbf{我的判断：} 这篇比 Chunk-GRPO 更直接地回答了我们刚才的问题：\key{每个 step 的 reward 可以不靠中间 reward model，而靠 controlled intervention + final reward 构造出来。} DenseGRPO 是“从中间 latent 垂直 readout”；TempFlow-GRPO 是“在指定时刻水平分叉探索，再把最终差异归因给分叉点”。
\end{summarybox}

\subsection{1. 它要解决什么}

Flow-GRPO / Dance-GRPO 的基本问题是：

\[
\hat{A}_t^i
=
\frac{
R(x_0^i,c)-\mathrm{mean}(\{R(x_0^g,c)\}_{g=1}^G)
}{
\mathrm{std}(\{R(x_0^g,c)\}_{g=1}^G)
}
\]

这里 \(\hat{A}_t^i\) 实际上没有真正的 \(t\)-specific 信息。它虽然写成 \(\hat{A}_t^i\)，但 reward 来自最终 \(x_0\)，所以整条 trajectory 的所有 timesteps 基本共用同一个 terminal advantage。

TempFlow-GRPO 指出这有两个问题：

\begin{enumerate}[leftmargin=1.5em]
  \item \textbf{sparse terminal reward：} 只知道最终结果好坏，不知道哪个 timestep 的探索带来提升。
  \item \textbf{temporal uniformity：} 所有 timesteps 被等权优化，但生成过程里不同时间点的探索能力不同。
\end{enumerate}

作者先做了一个关键观察：只在某一个 denoising step 上引入 SDE stochasticity，其余地方都用 ODE，最后看 PickScore 的 reward std。结果显示 early steps 的 reward variance 最大，late steps 接近 0。

\begin{center}
  \includegraphics[width=0.46\linewidth]{tempflowgrpo/fig2_left_reward_variance.png}
  \hfill
  \includegraphics[width=0.46\linewidth]{tempflowgrpo/fig2_right_branching.png}

  {\small \textbf{图 1：} TempFlow-GRPO 的动机。左图：只在单个 timestep 注入 SDE 噪声后，reward variance 在 early steps 最大。右图：从确定性 ODE 轨迹上分叉出 stochastic branch，并把最终 reward 差异归因给分叉点。图源：arXiv HTML。}
\end{center}

\begin{summarybox}
\textbf{这张图很重要：} 它把“不同 step 贡献不同”从直觉变成了一个可实验测量的问题。早期 step 的随机扰动会显著改变最终 reward，说明早期确实承担结构性探索；后期 step 的扰动影响小，更像细节 refinement。
\end{summarybox}

\subsection{2. 方法一：Trajectory Branching}

\subsubsection{2.1 为什么不直接训练 process reward model}

传统 process reward 的想法是直接评价中间状态 \(x_t\)。但在 diffusion / flow 模型里，中间 latent 或 partially-denoised image 很难被 reward model 稳定理解：

\[
x_t \quad \text{semantic ambiguous}
\]

所以 TempFlow-GRPO 换了一个角度：不评价 \(x_t\) 本身，而是做一个 intervention：

\[
\text{如果只在 } t=k \text{ 这个点探索，最终 reward 会怎么变？}
\]

这就把问题从“中间状态评分”变成“受控分叉后的最终结果评分”。

\subsubsection{2.2 Branching 的定义}

先从初始噪声 \(x_T\) 出发，用 ODE 确定性走到指定 timestep \(k\)：

\[
x_T \xrightarrow{\mathrm{ODE}} x_k
\]

然后只在 \(k\) 这个点引入一次 SDE stochasticity：

\[
x_{k-1}
=
\mathrm{SDE}(x_k,\epsilon)
\]

最后从 \(x_{k-1}\) 到 \(x_0\) 全部用 ODE 确定性完成：

\[
x_0
=
\mathrm{ODE}^{k-1}(x_{k-1})
\]

所以第 \(i\) 个 branch 的最终可评分样本是：

\[
\tilde{x}_{0,k}^i
=
\mathrm{ODE}^{k-1}
\left(
\mathrm{SDE}(x_k^i,\epsilon^i)
\right)
\]

对应 reward 是：

\[
R_k^i
=
R(\tilde{x}_{0,k}^i,c)
=
R\left(
\mathrm{ODE}^{k-1}
\left(
\mathrm{SDE}(x_k^i,\epsilon^i)
\right),
c
\right)
\]

然后用同一个 \(k\) 下多个 branches 的 reward 做 group normalization：

\[
\hat{A}_k^i
=
\frac{
R_k^i-\mathrm{mean}(\{R_k^g\}_{g=1}^{G})
}{
\mathrm{std}(\{R_k^g\}_{g=1}^{G})+\epsilon
}
\]

\begin{summarybox}
\textbf{关键理解：} 这个 \(\hat{A}_k^i\) 不是 DenseGRPO 的 \(\Delta R_k=R_{k-1}-R_k\)。它衡量的是：在相同 \(x_k\) 附近，不同 stochastic branch 最终会产生怎样的 reward 差异。因此它更像 \key{branch-level counterfactual advantage}。
\end{summarybox}

\subsubsection{2.3 为什么它能做 credit localization}

因为除了 \(k\) 这个点之外，其余路径都被 ODE 固定住了。也就是说，在同一个 initial noise 和 prompt 下，如果多个 branches 最后 reward 不同，差异主要来自：

\[
\epsilon_k
\quad \text{at branching timestep } k
\]

这就是它所谓的 credit localization：

\[
\text{reward variation}
\quad \Rightarrow \quad
\text{branching action at } k
\]

这和 DenseGRPO 有本质区别：

\[
\text{DenseGRPO：}
\quad
x_k \xrightarrow{\mathrm{ODE\ readout}} \hat{x}_0
\quad
\text{给 } x_k \text{ 打分}
\]

\[
\text{TempFlow-GRPO：}
\quad
x_k \xrightarrow{\mathrm{SDE\ branch}} x_{k-1}
\xrightarrow{\mathrm{ODE}} \hat{x}_0
\quad
\text{给 branch action 打分}
\]

\begin{warnbox}
\textbf{需要注意：} 它不是证明“第 \(k\) 步独立决定最终 reward”。它只是通过控制变量，让 stochastic difference 只发生在 \(k\)。后续 ODE 仍然会放大、修正或过滤这个差异，所以它更准确的说法是：\key{把探索差异定位到 branching point}，而不是完全因果地分解整条 denoising 链。
\end{warnbox}

\begin{center}
  \includegraphics[width=0.90\linewidth]{tempflowgrpo/fig5_framework.png}

  {\small \textbf{图 2：} TempFlow-GRPO 框架图。核心是 ODE-SDE-ODE 的分叉采样：在指定 timestep 注入 stochasticity，之后用最终 reward 形成该 timestep 的 process signal。图源：arXiv HTML。}
\end{center}

\subsection{3. 方法二：Noise-Aware Policy Weighting}

Trajectory branching 解决的是：

\[
\text{每个 timestep 的 advantage 怎么来？}
\]

Noise-aware weighting 解决的是：

\[
\text{每个 timestep 的更新强度应该一样吗？}
\]

答案是不应该一样。SDE sampler 在不同 timestep 的噪声注入幅度不同：

\[
\sigma_t\sqrt{\Delta t}
\]

早期噪声大，探索空间大；后期噪声小，主要是细节 refinement。作者观察到 reward std 和 noise level 强相关，所以用 noise level 作为 exploration capacity 的 proxy。

修改后的 policy loss 是：

\[
\mathcal{J}_{\text{policy}}(\theta)
=
\frac{1}{G}
\sum_{i=1}^{G}
\frac{1}{T}
\sum_{t=0}^{T-1}
\mathrm{Norm}(\sigma_t\sqrt{\Delta t})
\min
\left(
r_t^i(\theta)\hat{A}_t^i,
\mathrm{clip}(r_t^i(\theta),1-\epsilon,1+\epsilon)\hat{A}_t^i
\right)
\]

这里的 \(\mathrm{Norm}(\sigma_t\sqrt{\Delta t})\) 会被归一化，使所有 timestep 权重均值为 1。这样总 loss scale 不变，但不同时间点的相对优化强度改变。

\begin{center}
  \includegraphics[width=0.46\linewidth]{tempflowgrpo/fig6_noise_reward_correlation.png}
  \hfill
  \includegraphics[width=0.46\linewidth]{tempflowgrpo/fig6_scale_term.png}

  {\small \textbf{图 3：} 左：noise level 和 reward std 的相关性。右：标准 GRPO 的 scale term 与 noise-aware reweighting 后的 scale term 对比。图源：arXiv HTML。}
\end{center}

\subsubsection{3.1 理论推导在说明什么}

论文从 policy gradient 写起：

\[
\nabla_\theta \mathcal{J}(\theta)
=
\sum_{k=0}^{T-1}
\mathbb{E}
\left[
\nabla_\theta
\log p_\theta(x_{k-1}|x_k)
\hat{A}_k
\right]
\]

标准形式下，推导出的 timestep scale term 近似为：

\[
\sqrt{\frac{\Delta k(1-k)}{k}}
\]

加入 noise-aware reweighting 后，scale term 变成：

\[
\Delta k
\]

作者的解释是：标准 GRPO 的梯度贡献和真实 exploration capacity 不匹配，会让低噪声 refinement steps 在优化里占不合理的比重；noise-aware weighting 后，各 timestep 的梯度贡献更接近均衡。

\begin{summarybox}
\textbf{这里的重点不是“早期一定更大越好”。} 更准确地说，TempFlow-GRPO 想让优化强度匹配每个 timestep 的探索空间。早期有结构性探索，不能被低估；后期是细节修复，不能被过度扰动。
\end{summarybox}

\subsection{4. 方法三：Seed Group}

Trajectory branching 要比较不同 branches 的 reward。如果这些 branches 的初始噪声也不同，那 reward 差异会混入 initial noise 的影响，不能干净地归因到 branching point。

所以作者提出 seed group：

\[
\text{same prompt}
\quad + \quad
\text{same initial noise}
\quad + \quad
\text{different branch noise at } k
\]

这样 group 内 reward variance 更接近来自指定 timestep 的 exploration，而不是来自不同 \(x_T\) 的初始化差异。

实验配置里总 group size 保持为 24，默认是：

\[
4\ \text{initial seeds}
\times
6\ \text{branches}
=
24
\]

\begin{summarybox}
\textbf{这点对我们很有启发：} 如果要研究 credit assignment，group 的构造本身就是方法的一部分。不是随便采 \(G\) 条 trajectory，而是要控制住 prompt、initial noise、branching point，才能让 reward 差异有可解释性。
\end{summarybox}

\subsection{5. 实验结果}

\subsubsection{5.1 GenEval}

在 GenEval 上，论文报告：

\begin{center}
\begin{tabular}{lccc}
\toprule
方法 & step & overall & 备注 \\
\midrule
SD3.5-M base & - & 0.63 & base flow model \\
Flow-GRPO & 3800 & 0.88 & 同步数对比 \\
Flow-GRPO & 5600 & 0.95 & 更长训练 \\
TempFlow-GRPO & 3800 & 0.97 & 最优 \\
\bottomrule
\end{tabular}
\end{center}

论文还说 TempFlow-GRPO 约 2000 steps 就达到 0.95，而 Flow-GRPO 需要约 5600 steps。

\subsubsection{5.2 PickScore 和 FLUX}

在 PickScore benchmark 上：

\begin{itemize}[leftmargin=1.5em]
  \item 相比原始 Flow-GRPO，TempFlow-GRPO 大约提升 1.7\%。
  \item 相比改进后的 Flow-GRPO Prompt baseline，大约提升 1.0\%。
  \item 只需 100--200 training steps 就能达到 Flow-GRPO 的水平，300--400 steps 达到 Flow-GRPO Prompt 的水平。
\end{itemize}

在 FLUX.1-dev + HPSv3 上：

\begin{itemize}[leftmargin=1.5em]
  \item 约 80 steps 达到 Flow-GRPO 300 steps 的性能。
  \item Appendix 报告 300 steps 后相比 Flow-GRPO 约有 0.3 的性能提升。
  \item KL loss 更低、更稳定。
\end{itemize}

\subsubsection{5.3 和 DanceGRPO 的比较}

Appendix 里还和 DanceGRPO 做了比较。论文报告：

\begin{itemize}[leftmargin=1.5em]
  \item DanceGRPO 300 iterations 后 HPSv2 约 37.2。
  \item TempFlow-GRPO 280 iterations 达到 38.5。
  \item TempFlow-GRPO 约 150 iterations 就匹配 DanceGRPO 300 iterations 的最终效果，相当于约 \(2\times\) convergence speedup。
\end{itemize}

\subsubsection{5.4 Ablation}

\begin{center}
  \includegraphics[width=0.31\linewidth]{tempflowgrpo/fig9_branch_ablation.png}
  \hfill
  \includegraphics[width=0.31\linewidth]{tempflowgrpo/fig9_pickscore_ablation.png}
  \hfill
  \includegraphics[width=0.31\linewidth]{tempflowgrpo/fig9_geneval_ablation.png}

  {\small \textbf{图 4：} TempFlow-GRPO 的 ablation。左：不同 seed/branch 组合。中/右：trajectory branch、noise-aware weighting、seed group 的贡献。图源：arXiv HTML。}
\end{center}

几个关键结论：

\begin{itemize}[leftmargin=1.5em]
  \item 在总 group size 固定为 24 时，比较了 \(2\times 12\)、\(4\times 6\)、\(6\times 4\)。
  \item 更多 initial noises 有利于 early training 的快速提升。
  \item 更多 branches 在训练后期更好，因为同一个 branching point 下的探索比较更充分。
  \item 最终默认选择 \(4\times 6\)，平衡 initial diversity 和 branch comparison。
  \item 在 GenEval 上，noise-aware policy reweighting 从 0.82 提升到 0.92；trajectory branch 约带来 5\% gain；seed group 约带来 2\% gain。
\end{itemize}

\subsection{6. 成本问题}

TempFlow-GRPO 并不是免费午餐。由于每个 timestep 要做 \(K\)-branch exploration，采样成本高于 Flow-GRPO。论文在 Appendix A.6 里给了一个估计：

\[
K=10
\quad \Rightarrow \quad
\text{average branch factor}
\approx
\frac{9+\cdots+1}{10}
=
4.5
\]

也就是说，单个样本的采样过程可能比 Flow-GRPO 更贵。但论文认为它整体训练效率更高，因为 convergence 更快；在高分辨率任务上达到 Flow-GRPO 最终性能只需要约 33\%--50\% 的训练时间。

\begin{center}
  \includegraphics[width=0.46\linewidth]{tempflowgrpo/fig14_training_time_geneval.png}
  \hfill
  \includegraphics[width=0.46\linewidth]{tempflowgrpo/fig14_training_time_pickscore.png}

  {\small \textbf{图 5：} 训练时间对比。TempFlow-GRPO 单次采样更贵，但由于收敛更快，达到同等性能的 wall-clock time 更低。图源：arXiv HTML。}
\end{center}

\begin{warnbox}
\textbf{我的看法：} 这里的成本叙事要谨慎。TempFlow-GRPO 的优势是 \key{sample efficiency / convergence speed}，不是单步计算便宜。对图像任务这个 tradeoff 可能成立；但如果迁移到视频模型，多 branch 采样、decode、reward evaluation 的成本会被放大，需要重新评估。
\end{warnbox}

\subsection{7. 和 DenseGRPO / Chunk-GRPO 的关系}

\begin{center}
{\small
\setlength{\tabcolsep}{3pt}
\begin{tabular}{p{0.20\linewidth}p{0.25\linewidth}p{0.22\linewidth}p{0.22\linewidth}}
\toprule
方法 & reward / advantage & 更新粒度 & 核心代价 \\
\midrule
DenseGRPO & 对每个 \(x_t\) ODE readout，计算 \(\Delta R_t\) & step-level & full ODE readout 接近 \(O(T^2)\) \\
Chunk-GRPO & 主要仍用 terminal \(A_i\) & chunk-level ratio & credit 仍不是真正 chunk-specific \\
TempFlow-GRPO & ODE-SDE-ODE branch 后用最终 reward 得到 \(\hat{A}_k\) & step-level with temporal weighting & 多 branch exploration 成本高 \\
\bottomrule
\end{tabular}
}
\end{center}

三篇论文其实给了三种不同的 credit assignment 路径：

\[
\text{DenseGRPO：}
\quad
\text{readout-based local improvement}
\]

\[
\text{Chunk-GRPO：}
\quad
\text{chunk-level optimization granularity}
\]

\[
\text{TempFlow-GRPO：}
\quad
\text{branching-based counterfactual advantage}
\]

\begin{summarybox}
\textbf{新的认识：} step-level reward 不一定只能来自 \(\Delta R_t=R_{t-1}-R_t\)。还可以通过 controlled branching 构造：固定其他路径，只在一个 timestep 注入探索，再用最终 reward 差异反推该 timestep 的 advantage。
\end{summarybox}

\subsection{8. 局限与疑问}

\begin{enumerate}[leftmargin=1.5em]
  \item \textbf{它不是 DenseGRPO 式 improvement。} TempFlow-GRPO 的 \(R_k^i\) 是 branch outcome reward，不是相邻状态差分 \(R_{k-1}-R_k\)。因此它回答的是“在 \(k\) 探索会导致什么最终结果”，不是“从 \(x_k\) 到 \(x_{k-1}\) 本身让 reward 提升了多少”。
  \item \textbf{credit localization 依赖 controlled setup。} 如果初始 noise、prompt、ODE remainder 没控制好，reward 差异就不再只来自 branch noise。
  \item \textbf{多 branch 成本可能很高。} 对视频生成尤其明显，因为每个 branch 都可能意味着完整后续 ODE、VAE decode 和 reward scoring。
  \item \textbf{noise-aware weighting 是 proxy。} noise level 和 reward variance 相关，但不等于真实 semantic importance。某些任务里 late-stage detail 或 identity preservation 也可能很关键。
  \item \textbf{仍然依赖 outcome reward model。} 如果最终 reward model 本身偏了，branch advantage 也会跟着偏。
\end{enumerate}

\subsection{9. 对我们后续思考的启发}

这篇读完之后，我们可以把当前问题改得更深一点：

\[
\text{credit assignment 的单位是什么？}
\]

现在不只是 step vs chunk，还多了一个维度：

\[
\text{credit signal 是 readout difference，还是 controlled branching difference？}
\]

也就是说，未来方法不一定只是：

\[
\text{Chunk-GRPO} + \text{DenseGRPO}
\]

更有意思的方向可能是：

\[
\text{temporal segment}
\quad + \quad
\text{controlled branching}
\quad + \quad
\text{low-cost readout}
\]

例如，对一个 segment \([\tau_j,\tau_{j-1}]\)，不是机械做 DenseGRPO 式 boundary readout，也不是简单用 terminal reward，而是做 segment-level intervention：

\[
x_{\tau_j}
\xrightarrow{\mathrm{SDE\ segment}}
x_{\tau_{j-1}}
\xrightarrow{\mathrm{ODE}}
\hat{x}_0
\xrightarrow{R}
R_{\mathrm{seg}_j}
\]

然后同一个 segment 下比较多个 branches：

\[
\hat{A}_{\mathrm{seg}_j}^i
=
\frac{
R_{\mathrm{seg}_j}^i-\mathrm{mean}(\{R_{\mathrm{seg}_j}^g\})
}{
\mathrm{std}(\{R_{\mathrm{seg}_j}^g\})+\epsilon
}
\]

这会比“Chunk-GRPO + DenseGRPO”更自然，因为它把新问题定义成：

\begin{quote}
如何在合适的 temporal resolution 上构造可归因的 counterfactual reward signal？
\end{quote}

\begin{summarybox}
\textbf{读完要记住的三句话：}

1. TempFlow-GRPO 证明 step-level credit 不一定要靠中间 reward model，可以靠 controlled branching + final reward 构造。\\
2. 它把探索和评价分开：指定 timestep 用 SDE 探索，其余路径用 ODE 固定，再用 outcome reward 评价。\\
3. 它进一步说明 temporal credit assignment 的关键不是“越细越好”，而是要同时考虑 \key{归因可解释性、探索能力、更新强度和计算成本}。
\end{summarybox}

% Paper deep reading: PCPO
% Source: https://arxiv.org/html/2509.25774v3

\section{论文精读：PCPO}

\textbf{论文：} \href{https://arxiv.org/abs/2509.25774}{PCPO: Proportionate Credit Policy Optimization for Aligning Image Generation Models}\\
\textbf{简称：} PCPO。\\
\textbf{版本：} arXiv v3，2025-09-30 提交，2026-02-24 更新，ICLR 2026。\\
\textbf{代码：} \href{https://github.com/jaylee2000/pcpo/}{github.com/jaylee2000/pcpo}\\
\textbf{阅读口径：} 本笔记只从 \textbf{Flow Matching / Flow-GRPO SDE} 视角讲 PCPO，不走论文前半部分的其他 sampler 入口。

\begin{conceptbox}
\textbf{一句话总结：} PCPO 在 flow matching 里的关键发现是：Flow-GRPO / DanceGRPO 的 SDE policy ratio 展开后，会自带一个 timestep weight \(w(t_i)\)。这个 weight 不是我们设计出来的 credit，而是由 SDE variance、timestep interval 和 sampler 结构隐式产生的。PCPO 要做的是把这个 native weight 校准成和积分区间 \(\Delta t_i\) 成比例，避免某些 flow steps 被 objective 侧过度放大。
\end{conceptbox}

\subsection{0. 它为什么还在这条主线里}

我们现在看的主线是：

\[
\textbf{连续 flow 的信用分配}
\]

核心问题是：

\begin{align*}
&\text{terminal reward 太粗，}\\
&\text{不能说明哪个 timestep / chunk / branch / objective 真正导致好结果。}
\end{align*}

前面三篇分别处理了不同层面：

\begin{itemize}[leftmargin=2em]
  \item \textbf{DenseGRPO：} 从 reward 侧做 step-wise reward gain，问每一步到底让最终 reward 提升了多少。
  \item \textbf{Chunk-GRPO：} 从时间粒度侧做 chunk-level ratio，问 step 太细、trajectory 太粗时怎么更新更稳。
  \item \textbf{TempFlow-GRPO：} 从 branch 侧做 ODE-SDE-ODE，对某个 timestep 做反事实探索。
\end{itemize}

PCPO 接在这里，问的是另一个更底层的问题：

\begin{summarybox}
即使我们暂时还用同一个 terminal advantage \(A\)，Flow-GRPO 的 \(\log \rho_{t_i}\) 本身是否已经在不公平地放大某些 timestep？\\
PCPO 的回答是：是的。因此 credit assignment 不只发生在 reward 侧，也发生在 objective / ratio 侧。
\end{summarybox}

\subsection{1. Flow Matching 里为什么要先变成 SDE}

Flow Matching 模型原本通常按 ODE 采样：

\[
\frac{d x_t}{dt}=u_\theta(x_t,t,c).
\]

这个 ODE 是 deterministic 的。如果完全 deterministic，就没有常规 policy gradient 需要的 step transition probability：

\[
p_\theta(x_{t_{i-1}}\mid x_{t_i},c).
\]

所以 Flow-GRPO / DanceGRPO 的基本操作是：把 ODE 改写成一个有相同边缘分布的 reverse SDE：

\[
d x_t
=
\left(
u_t
-
\frac{1}{2}\sigma_t^2
\nabla \log p_t(x_t)
\right)dt
+
\sigma_t d w .
\]

离散化后，每一步可以理解成一个 Gaussian policy：

\[
x_{t_{i-1}}
\sim
p_\theta^{(i)}
\left(
x_{t_{i-1}}\mid x_{t_i},c
\right).
\]

这样才可以对每一步写 policy ratio：

\[
\rho_{t_i}(\theta)
=
\frac{
p_\theta^{(i)}(x_{t_{i-1}}\mid x_{t_i},c)
}{
p_{\theta_{\mathrm{old}}}^{(i)}(x_{t_{i-1}}\mid x_{t_i},c)
}.
\]

\begin{intubox}
\textbf{直觉：} Flow-GRPO 的第一步不是“怎么分 reward”，而是“怎么让 flow transition 变成一个可采样、可算 logprob 的随机策略”。PCPO 研究的就是这个 step-level ratio 里有没有隐藏的不公平权重。
\end{intubox}

\subsection{2. 标准 Flow-GRPO 的问题不只在 advantage}

简化的 per-step clipped objective 可以写成：

\[
\sum_i
\min
\left(
\rho_{t_i} A,\,
\mathrm{clip}(\rho_{t_i},1-\epsilon,1+\epsilon) A
\right).
\]

这里最容易被我们盯上的问题是：

\[
A \text{ 是 terminal reward 得到的 trajectory-level advantage。}
\]

DenseGRPO、Chunk-GRPO、TempFlow-GRPO 都是在不同方式上修这个问题。PCPO 进一步指出，还有一个更隐蔽的问题：

\[
\rho_{t_i} \text{ 自己也不一定公平。}
\]

也就是说，即使两个 timestep 拿到同一个 \(A\)，它们的 \(\log\rho_{t_i}\) 也可能因为 SDE schedule 的不同而产生完全不同的梯度尺度。

\subsection{3. Flow 版本 log-ratio 里藏着什么}

PCPO 对 Flow Matching SDE 推导出的 step log-ratio 可以写成：

\[
\log\rho_{t_i}
=
-
\left[
w_{\mathrm{native}}(t_i)
\left(
u_\theta - u_{\mathrm{old}}
\right)
\cdot
\epsilon_{\mathrm{old}}^{(t_i)}
+
\frac{1}{2}
\left\|
w_{\mathrm{native}}(t_i)
\left(
u_\theta - u_{\mathrm{old}}
\right)
\right\|^2
\right].
\]

其中 flow SDE 自带的 native weight 是：

\[
w_{\mathrm{native}}(t_i)
=
\frac{\sqrt{\Delta t_i}}{\sigma_{t_i}}
\left(
1+
\frac{(1-t_i)\sigma_{t_i}^2}{2t_i}
\right).
\]

这个式子是 PCPO 在 flow 视角下最重要的地方。它说明：

\[
\text{step } i \text{ 的更新强度}
\quad
\propto
\quad
w_{\mathrm{native}}(t_i).
\]

而 \(w_{\mathrm{native}}(t_i)\) 由三个东西共同决定：

\begin{itemize}[leftmargin=2em]
  \item \(\Delta t_i\)：这个 flow step 覆盖的积分区间长度。
  \item \(\sigma_{t_i}\)：SDE 注入随机性的强度。
  \item \(t_i\)：flow trajectory 上的位置，以及对应的 schedule 修正项。
\end{itemize}

\begin{warnbox}
这就解释了为什么 PCPO 不能只被理解成“又一个 timestep weighting trick”。它指出的是：Flow-GRPO 的 policy ratio 数学形式本身已经内置了一个 timestep credit 分配器。问题在于，这个分配器不一定和我们想要的 temporal credit 一致。
\end{warnbox}

\subsection{4. 什么叫 proportionate credit}

PCPO 的原则是：

\[
\text{credit should be proportional to integration interval.}
\]

在 flow trajectory 上，第 \(i\) 步覆盖的区间是 \(\Delta t_i\)。如果某一步代表更长的积分区间，它自然应该承担更多 trajectory update；如果某一步只是一个很短的区间，却被 SDE schedule 放大很多，那就是 disproportionate credit。

所以 PCPO 希望目标权重变成：

\[
w_{\mathrm{target}}(t_i)
=
\zeta \Delta t_i,
\]

其中

\[
\zeta
=
\sum_{i=1}^{N}
\frac{\sqrt{\Delta t_i}}{\sigma_{t_i}}
\left(
1+
\frac{(1-t_i)\sigma_{t_i}^2}{2t_i}
\right).
\]

一种直观理解是：

\[
\text{native: }
w_{\mathrm{native}}(t_i)
\quad
\Longrightarrow
\quad
\text{PCPO: }
w_{\mathrm{target}}(t_i)\propto \Delta t_i .
\]

\textbf{关键点：} PCPO 在 flow model 里不想大幅改 SDE variance schedule 或 timestep shifting，因为那会改变原本调好的 sampler。它采用的是 objective reweighting。可以把每个 step 的 loss 乘上：

\[
\omega_i
=
\frac{
w_{\mathrm{target}}(t_i)
}{
w_{\mathrm{native}}(t_i)
}.
\]

于是 flow 版本的 PCPO objective 可以理解成：

\[
\mathcal{L}_{\mathrm{PCPO\mbox{-}flow}}
=
\mathbb{E}
\left[
\sum_i
\omega_i
\max
\left\{
0,\,
\xi|A|-A\log\rho_{t_i}
\right\}
\right].
\]

\begin{figure}[h]
  \centering
  \includegraphics[width=0.82\linewidth]{pcpo/fig2_flow_sde_weights.png}
  \caption{PCPO 的 Flow SDE 权重校准。图中只保留 DanceGRPO SDE 和 Flow-GRPO SDE 的 native weights 与 proportional weights 对比。图源：arXiv HTML，经本地裁剪。}
\end{figure}

\subsection{5. 它和 TempFlow-GRPO 的关键区别}

TempFlow-GRPO 的直觉是：

\[
\text{某些 timestep 的 branch reward variance 大}
\Rightarrow
\text{这些 timestep 可能更该被优化。}
\]

PCPO 的直觉更底层：

\[
\text{当前 objective 到底已经放大了哪些 timestep？}
\]

所以二者关注点不同：

\begin{center}
\begin{tabular}{p{0.28\linewidth}p{0.62\linewidth}}
\toprule
\textbf{方法} & \textbf{它问的问题} \\
\midrule
TempFlow-GRPO & 如果只在 \(t_i\) 注入随机性，最终 reward 会怎样变？ \\
PCPO & Flow SDE 的 \(\log\rho_{t_i}\) 是否已经把某些 \(t_i\) 的梯度放大过头？ \\
\bottomrule
\end{tabular}
\end{center}

\begin{summarybox}
一个 timestep 有影响，不代表 objective 还应该继续放大它。\\
PCPO 的价值在于提醒我们：temporal credit signal 和 temporal optimization weight 是两件事，需要分别校准。
\end{summarybox}

\subsection{6. 实验结果按 flow 口径怎么读}

\subsubsection{6.1 DanceGRPO / FLUX 上的加速}

论文在 DanceGRPO + FLUX 的 setting 上报告了明显的训练加速：

\[
\text{DanceGRPO FLUX + HPSv2.1 reward: } 41.2\% \text{ speedup.}
\]

这里的收益不是因为 PCPO 少采样，也不是因为少评估 reward，而是因为每次 policy update 的 timestep credit 更成比例，clipping fraction 更低，训练更稳。

\begin{figure}[h]
  \centering
  \includegraphics[width=0.96\linewidth]{pcpo/fig3_training_dynamics.png}
  \caption{PCPO 的 training dynamics。注意这里只按 flow 主线阅读：重点看 DanceGRPO / FLUX 相关曲线中 PCPO 的 clipping 更低、更稳定。图源：arXiv HTML。}
\end{figure}

\subsubsection{6.2 Flow-GRPO / SD3.5-M 上的泛化}

论文还把 PCPO 接到 SD3.5-M 的 Flow-GRPO setting 上，任务包括 OCR 和 PickScore，并带 KL penalty。结果显示 PCPO 相比原始 Flow-GRPO 有：

\begin{itemize}[leftmargin=2em]
  \item 更高 reward。
  \item 更低 KL。
  \item 更低 clipping fraction。
\end{itemize}

这说明 PCPO 不是只服务于某一个 SDE 实现，而是可以迁移到 Flow-GRPO 框架。

\begin{figure}[h]
  \centering
  \includegraphics[width=0.96\linewidth]{pcpo/fig5_flow_grpo_generalization.png}
  \caption{PCPO 在 SD3.5-M + Flow-GRPO 设置上的泛化结果。PCPO 提高 reward，同时降低 KL 和 clipping。图源：arXiv HTML。}
\end{figure}

\subsubsection{6.3 不要把它误解成 high-noise heuristic}

一个容易误解的点是：TempFlow-GRPO 发现高噪声阶段有更大 reward variance，那是不是直接放大 high-noise timesteps 就行？

PCPO 的实验给出反例。对于 DanceGRPO SDE，native weights 本来就已经偏向某些 short high-noise intervals。如果继续按 high-noise heuristic 放大，最终反而不如 proportional weighting。

\begin{figure}[h]
  \centering
  \includegraphics[width=0.96\linewidth]{pcpo/fig7_heuristic_ablation.png}
  \caption{PCPO 对简单 heuristic 的反驳：减少 timestep 或盲目强调 high-noise steps 都可能牺牲质量；proportional credit 更稳。图源：arXiv HTML。}
\end{figure}

\subsection{7. 和 DenseGRPO / Chunk-GRPO / TempFlow-GRPO 的关系}

现在四篇可以放进一个更清楚的分层：

\begin{center}
\begin{tabular}{p{0.22\linewidth}p{0.26\linewidth}p{0.42\linewidth}}
\toprule
\textbf{层面} & \textbf{代表论文} & \textbf{解决什么} \\
\midrule
reward signal & DenseGRPO & 用 ODE readout 估计每步 reward improvement \\
temporal unit & Chunk-GRPO & 把 step 合成 chunk，改变 update 粒度 \\
branch credit & TempFlow-GRPO & 在某个 timestep branch，看最终 reward 差异 \\
objective weight & PCPO & 校准 Flow SDE ratio 里隐藏的 timestep weight \\
\bottomrule
\end{tabular}
\end{center}

所以 PCPO 不替代前三篇。它补的是另一个缺口：

\[
\text{即使你已经有了更好的 } A_i(t),
\quad
\text{objective 也可能用 } w_{\mathrm{native}}(t_i)
\text{ 把它重新扭曲。}
\]

\subsection{8. 局限}

\begin{warnbox}
PCPO 解决的是 objective-side disproportionate credit，不是 reward-side sparse credit。它不会告诉你某个 flow step 真的让最终图像变好还是变差；它只是在防止 Flow SDE 的 native ratio weight 把某些 steps 过度放大。
\end{warnbox}

几个边界：

\begin{itemize}[leftmargin=2em]
  \item \textbf{没有 step reward gain。} PCPO 不构造 \(\Delta R_t\)，所以不能替代 DenseGRPO。
  \item \textbf{没有 branch intervention。} PCPO 不做 ODE-SDE-ODE counterfactual branch，所以不能替代 TempFlow-GRPO。
  \item \textbf{没有改 temporal unit。} PCPO 仍然可以逐 step update，不像 Chunk-GRPO 那样把多个 step 合成 chunk ratio。
  \item \textbf{依赖 proportionality 假设。} \(w_{\mathrm{target}}(t_i)\propto \Delta t_i\) 是一个强原则，但如果某个任务本身强烈 stage-specific，仅靠 \(\Delta t_i\) 可能不够。
\end{itemize}

\subsection{9. 对我们后续 idea 的启发}

昨天我们觉得 ``Chunk-GRPO + DenseGRPO'' 有点 A+B。PCPO 给了更好的 framing：

\[
\text{temporal credit signal}
\quad
\neq
\quad
\text{temporal optimization weight}.
\]

一个更完整的问题应该是：

\begin{summarybox}
如何估计某个 flow segment 的 reward improvement，同时保证 Flow SDE objective 里的 native timestep weight 不会把这个 credit 再次扭曲？
\end{summarybox}

这就比简单拼模块更自然。真正的研究问题变成：

\[
\text{reward-side credit}
\quad +
\quad
\text{objective-side credit calibration}.
\]

\subsection{10. 今天读 PCPO 时只抓三句话}

\begin{enumerate}[leftmargin=2em]
  \item PCPO 在 Flow Matching 里的核心观察是：SDE policy ratio 里有隐藏的 \(w_{\mathrm{native}}(t_i)\)。
  \item PCPO 的核心原则是：\textbf{这个 weight 应该和 integration interval \(\Delta t_i\) 成比例}。
  \item PCPO 对我们的启发是：\textbf{以后设计 step-wise / chunk-wise reward 时，还必须校准 Flow-GRPO objective 侧的 temporal weight}。
\end{enumerate}

% ============================================================
% Mainline summary: Flow-GRPO credit assignment
% ============================================================

\section{主线总结：Flow-GRPO 信用分配四篇阶段总结}

\begin{conceptbox}
\textbf{覆盖论文：} DenseGRPO、Chunk-GRPO、TempFlow-GRPO、PCPO。\\
\textbf{总结日期：} 2026-04-28。\\
\textbf{当前判断：} 这四篇构成一个比较清晰的主线：\textbf{Flow-GRPO 后训练的 credit assignment 不能只看 terminal reward，也不能只看 advantage；policy ratio / logprob scale 本身也会改变每个 step 的实际更新强度。}
\end{conceptbox}

\subsection{0. 最核心的一句话}

标准 Flow-GRPO 的问题可以压缩成一个式子：

\[
\text{实际每个 step 的更新强度}
\approx
\underbrace{\text{advantage / credit}}_{\text{reward side}}
\times
\underbrace{\text{policy ratio / logprob scale}}_{\text{policy side}}.
\]

更具体地，PPO / GRPO 的 step-level objective 里有：

\[
\min
\left(
\rho_t^i A_i,\,
\mathrm{clip}(\rho_t^i,1-\epsilon,1+\epsilon)A_i
\right).
\]

所以不能只问：

\[
A_i \text{ 应该怎么分给每个 timestep？}
\]

还要问：

\[
\rho_t^i \text{ 本身是不是已经在放大或压小某些 timestep？}
\]

\begin{summarybox}
\textbf{阶段性结论：} 好的 Flow-GRPO 后训练，需要同时校准 \textbf{advantage side} 和 \textbf{ratio side}。前者决定“奖励应该分给谁”，后者决定“同样的奖励乘到 objective 里以后，实际更新强度会不会偏”。
\end{summarybox}

\subsection{1. 第一类问题：advantage 分布不平衡}

标准 Flow-GRPO / DanceGRPO 的做法是：同一个 prompt 采样 \(G\) 条 trajectories，最终每条 trajectory 得到一个 terminal reward，再归一化成 trajectory-level advantage：

\[
A_i
=
\frac{
R(x_0^i,c)-\mathrm{mean}_{j}R(x_0^j,c)
}{
\mathrm{std}_{j}R(x_0^j,c)+\epsilon
}.
\]

然后每个 denoising step 都拿同一个 \(A_i\)：

\[
\sum_{t=1}^{T}
\min
\left(
\rho_t^i A_i,\,
\mathrm{clip}(\rho_t^i,1-\epsilon,1+\epsilon)A_i
\right).
\]

这等价于默认：

\[
\text{一条 trajectory 最终好}
\Rightarrow
\text{这条 trajectory 的每个 step 都好。}
\]

这个假设太强。因为最终 reward 高，不代表所有中间 steps 都贡献了正收益；最终 reward 低，也不代表所有中间 steps 都应该被惩罚。

\begin{warnbox}
\textbf{这就是 advantage 分布不平衡：} terminal advantage 被全局广播到所有 timesteps，导致真正有贡献的 step 和没贡献的 step 拿到同一个训练信号。它不是 reward model 不准的问题，而是 reward credit 的时间分配太粗。
\end{warnbox}

\subsection{2. DenseGRPO：最直接地修 advantage side}

DenseGRPO 的核心想法是把 terminal reward 拆成 step-wise reward gain：

\[
\Delta R_t
=
R_{t-1}-R_t.
\]

这里的 \(R_t\) 不是直接给 noisy latent 打分，而是从中间 latent \(x_t\) 出发，用 ODE denoising readout 到 clean image，再 decode、再用 reward model 打分。

所以 DenseGRPO 本质上在问：

\[
\text{经过第 \(t\) 个 denoising transition 后，最终可读出的 reward 提升了多少？}
\]

这比 terminal advantage 更合理，因为它更接近：

\[
A_i
\quad \rightarrow \quad
A_t^i.
\]

\begin{summarybox}
\textbf{DenseGRPO 的价值：} 它是四篇里对 advantage side 修得最直接的一篇。它试图估计每个 step 的 actual improvement，而不是把最终 reward 平均分配给所有 steps。
\end{summarybox}

但 DenseGRPO 的问题也很清楚：

\[
T+(T-1)+\cdots+1
=
\frac{T(T+1)}{2}.
\]

如果每个中间 latent 都要完整 ODE readout 到 clean image，成本接近 \(O(T^2)\)。所以 DenseGRPO 的核心收益很硬，工程成本也很硬。

\subsection{3. Chunk-GRPO：不直接修 reward，而是修优化粒度}

Chunk-GRPO 的判断是：单个 step 太细，整条 trajectory 太粗，所以可以用 chunk 作为中间粒度。

它把 denoising trajectory 切成：

\[
\mathrm{ch}_1,\mathrm{ch}_2,\ldots,\mathrm{ch}_K.
\]

然后把 step-level ratio 合成 chunk-level ratio：

\[
r_j^i(\theta)
=
\left(
\prod_{t\in \mathrm{ch}_j}
\frac{
p_\theta(x_{t-1}^i|x_t^i,c)
}{
p_{\mathrm{old}}(x_{t-1}^i|x_t^i,c)
}
\right)^{1/cs_j}.
\]

其中 \(1/cs_j\) 是几何平均，防止 chunk 越长 ratio 尺度越容易爆。

Chunk-GRPO 没有像 DenseGRPO 那样真正计算 chunk reward improvement。它更多是在说：

\[
\text{更新单位不一定要是单 step，也不应该是整条 trajectory。}
\]

\begin{summarybox}
\textbf{Chunk-GRPO 的价值：} 它让 credit 分配更平滑，减少 step-level 噪声，也提供了一个很实用的工程方向：如果 DenseGRPO 的 step-wise readout 太贵，可以尝试 chunk-level improvement。
\end{summarybox}

它的局限是：如果仍然用同一个 terminal \(A_i\)，那它只是缓解了 credit assignment，没有彻底解决“哪个 chunk 真正带来 reward improvement”的问题。

\subsection{4. TempFlow-GRPO：用 branch 做局部反事实}

TempFlow-GRPO 试图回答的是另一个问题：

\[
\text{某个 timestep 的选择，是否真的导致了最终结果差异？}
\]

它用 trajectory branching 的方式，在某个时刻 \(t\) 产生分支，然后观察不同 branch 的后续结果。直觉上，这比单纯看一条 trajectory 更接近反事实：

\[
\text{如果在同一个前缀状态下，换一个 branch，最终 reward 会怎样？}
\]

所以 TempFlow-GRPO 的 credit localization 不是 DenseGRPO 那种 step-wise reward gain，而是 branch-level counterfactual。

\begin{summarybox}
\textbf{TempFlow-GRPO 的价值：} 它提醒我们，credit 不只是 timestep 的问题，也和局部探索有关。一个 step 的重要性，最好通过“同一位置不同 branch 的后果差异”来判断。
\end{summarybox}

但它也有两个问题：

\begin{enumerate}[leftmargin=1.5em]
  \item branch 怎么设计会影响 credit 结论。
  \item 分支越多，成本越高；分支太少，反事实又不够充分。
\end{enumerate}

\subsection{5. PCPO：ratio side 也会影响 credit}

前面三篇主要都在处理 advantage side：

\[
A_i
\rightarrow
A_t^i
\quad \text{or} \quad
A_{\mathrm{chunk}}^i
\quad \text{or} \quad
A_{\mathrm{branch}}^i.
\]

PCPO 的关键提醒是：即使 advantage 先不变，policy ratio 本身也不是中性的。

Flow-GRPO 里每一步都有：

\[
\rho_t^i(\theta)
=
\frac{
p_\theta(x_{t-1}^i|x_t^i,c)
}{
p_{\theta_{\mathrm{old}}}(x_{t-1}^i|x_t^i,c)
}.
\]

PCPO 展开 Flow SDE 的 log-ratio 后指出，不同 timestep 的 log-ratio scale 会受到 step size、SDE variance、timestep position 等因素影响。也就是说，即使用同一个 \(A_i\)，不同 timestep 的实际更新强度也可能不一样。

可以把它理解成：

\[
\rho_t^i A_i
\quad
\text{里面的}
\quad
\rho_t^i
\quad
\text{已经带了一个隐式 timestep weight。}
\]

\begin{warnbox}
\textbf{PCPO 的关键意义：} credit assignment 不只发生在 \(A\) 里，也发生在 \(\rho\) 里。只修 advantage side，不看 ratio side，可能仍然会让某些 steps 被 objective 数学形式过度放大或压小。
\end{warnbox}

所以 PCPO 把我们的主线从：

\[
\text{如何分配 reward？}
\]

推进到：

\[
\text{如何校准每个 step 的实际 policy update 强度？}
\]

\subsection{6. 四篇放在一起看}

\begin{center}
\resizebox{\linewidth}{!}{
\begin{tabular}{llll}
\toprule
\textbf{论文} & \textbf{修哪一侧} & \textbf{核心机制} & \textbf{主要局限}\\
\midrule
DenseGRPO & advantage side & step-wise reward gain \(\Delta R_t\) & ODE readout 成本高，接近 \(O(T^2)\)\\
Chunk-GRPO & advantage side / 粒度 & chunk-level ratio 和 chunk-level update & 多数情况下仍依赖 terminal advantage\\
TempFlow-GRPO & advantage side / 探索 & branch-level counterfactual credit & branch 设计和额外采样成本复杂\\
PCPO & ratio side & 校准 Flow SDE ratio 的 hidden timestep weight & 不直接提供 step-wise reward gain\\
\bottomrule
\end{tabular}
}
\end{center}

这张表的重点不是“谁最好”，而是它们在拆同一个目标的不同部分：

\[
\textbf{让每个 timestep 的实际更新强度更接近它对最终质量的真实贡献。}
\]

\subsection{7. 当前形成的统一视角}

我们现在可以把 Flow-GRPO 的每个 step 更新写成更完整的概念式：

\[
\Delta \theta_t
\propto
\underbrace{
\nabla_\theta \log p_\theta(x_{t-1}|x_t,c)
}_{\text{policy gradient direction}}
\cdot
\underbrace{
A_t
}_{\text{advantage credit}}
\cdot
\underbrace{
w_{\rho}(t)
}_{\text{ratio / SDE scale}}
.
\]

标准 Flow-GRPO 的问题是：

\[
A_t
\approx
A_{\mathrm{terminal}},
\qquad
w_{\rho}(t)
\text{ 没有显式校准。}
\]

这就导致两类失真：

\begin{enumerate}[leftmargin=1.5em]
  \item \textbf{Advantage 侧失真：} 该 step 不一定真的贡献了 reward，但它拿到了同一个 \(A_i\)。
  \item \textbf{Ratio 侧失真：} 该 step 就算拿到了合理 \(A_t\)，它的 logprob / ratio scale 也可能让实际更新偏大或偏小。
\end{enumerate}

\begin{summarybox}
\textbf{最重要的阶段结论：}
\[
\boxed{
\text{Flow-GRPO 的 credit assignment = advantage side + ratio side。}
}
\]
DenseGRPO、Chunk-GRPO、TempFlow-GRPO 主要让我们看到 advantage 分布不平衡；PCPO 让我们看到 policy ratio 本身也会影响 credit。
\end{summarybox}

\subsection{8. 对 OTCA 的定位}

OTCA 可以暂时降级为旁支启发，而不是这条主线里的关键节点。

它有一个有用提醒：

\[
\text{多 reward objective 也会影响 advantage 怎么组织。}
\]

也就是除了问：

\[
\text{which timestep?}
\]

还可以问：

\[
\text{which objective?}
\]

但 OTCA 自己的实现没有前四篇那么硬：

\begin{itemize}[leftmargin=1.5em]
  \item TCD 的 latent similarity gain 只是 proxy，不是真正 causal reward gain。
  \item MOCA 主要是在 advantage scalar 上做融合，不是真正的 gradient-space MOO。
  \item 主公式里的 \(c_k^i\) 没有 \(t\)，所以“timestep-aware objective mixture”这个说法需要保留。
\end{itemize}

所以后续可以借 OTCA 的 when + what 框架，但不要直接照搬它的 TCD / MOCA。

\subsection{9. 后续读论文时的抓手}

之后继续读这条主线的论文，可以先问五个问题：

\begin{enumerate}[leftmargin=1.5em]
  \item 它是在修 \textbf{advantage side}，还是在修 \textbf{ratio side}？
  \item 如果修 advantage，它是 step、chunk、branch，还是 objective-level credit？
  \item 它的 credit 是 causal / counterfactual，还是 proxy / heuristic？
  \item 它有没有解决计算成本，还是只是把 credit 算得更密？
  \item 它是否考虑了 Flow SDE ratio 的 timestep scale，还是默认 ratio 是中性的？
\end{enumerate}

\begin{intubox}
\textbf{一个更高标准的目标：} 最理想的后续方法，应该同时做到：
\[
\begin{aligned}
&\text{低成本估计 temporal reward contribution}\\
&\quad +\ \text{校准 ratio-side timestep scale}\\
&\quad +\ \text{在多 reward 场景下避免 objective 冲突。}
\end{aligned}
\]
现在四篇论文各自解决了一部分，但还没有一篇把三件事都做干净。
\end{intubox}

% ============================================================
% Paper deep reading: OTCA
% Source: https://arxiv.org/html/2604.19234v2
% ============================================================

\section{论文精读：OTCA}

\begin{conceptbox}
\textbf{论文：} \href{https://arxiv.org/abs/2604.19234}{Learning to Credit the Right Steps: Objective-aware Process Optimization for Visual Generation}\\
\textbf{简称：} OTCA。\\
\textbf{版本：} arXiv:2604.19234v2，2026-04-27。\\
\textbf{方向标签：} Flow Matching、Visual GRPO、Credit Assignment、Multi-objective Reward、Image / Video Generation。\\
\textbf{阅读日期：} 2026-04-28。\\
\textbf{阅读口径：} 只放在我们当前的主线里读：\textbf{连续 flow / visual GRPO 的信用分配}。论文覆盖 image 和 video，术语里会说 denoising trajectory；这里统一理解成 Flow-SDE 采样轨迹上的 policy optimization。
\end{conceptbox}

\begin{summarybox}
\textbf{一句话总结：} OTCA 把原本一个 trajectory-level advantage 拆成两个问题：\textbf{什么时候更新}和\textbf{更新哪个 reward objective}。核心公式是
\[
\tilde{A}_{t}^{i}
=
w_t^i
\sum_{k=1}^{K} c_k^i A_k^i .
\]
其中 \(w_t^i\) 来自 TCD，负责 temporal credit；\(c_k^i\) 来自 MOCA，负责多目标 reward credit。它和我们前面几篇的关系是：DenseGRPO / TempFlow 更偏 reward gain 或 branch 反事实，PCPO 更偏 ratio 侧权重，OTCA 则进一步问 multi-reward setting 下 advantage 本身应该怎么组织。
\end{summarybox}

\begin{center}
  \includegraphics[width=0.95\linewidth]{otca/fig1_framework_comparison.png}

  {\small \textbf{图 1：} OTCA 和普通 GRPO 的区别。普通做法把多个 reward 压成一个 scalar，再广播给所有 timesteps；OTCA 试图同时做 timestep credit 和 objective credit。图源：arXiv HTML。}
\end{center}

\subsection{0. 它接在 PCPO 后面读的原因}

我们这条主线一直在拆同一个问题：

\[
\text{terminal reward 太粗，不能说明哪个 timestep 真正导致好结果。}
\]

前四篇可以这样放：

\begin{center}
\resizebox{\linewidth}{!}{
\begin{tabular}{lll}
\toprule
\textbf{论文} & \textbf{主要改动的位置} & \textbf{它问的问题}\\
\midrule
DenseGRPO & reward gain & 哪个 step 真的提升了最终 reward？\\
Chunk-GRPO & optimization unit & step 太细、trajectory 太粗时，能不能按 chunk 更新？\\
TempFlow-GRPO & branch / counterfactual & 某个时间点的分叉到底带来什么影响？\\
PCPO & policy ratio / objective scale & ratio 数学形式是否隐式偏向某些 steps？\\
OTCA & temporal credit + reward objective & 哪个 step 重要，以及这个 step 应该优化哪个目标？\\
\bottomrule
\end{tabular}
}
\end{center}

所以 OTCA 不是在重复 DenseGRPO。它的切入点是 multi-reward visual generation：

\[
\begin{aligned}
R
&=
\text{CLIP / text alignment}
+
\text{aesthetic / visual quality}\\
&\quad
+
\text{motion / temporal consistency}
+
\cdots
\end{aligned}
\]

如果简单把它们加成一个 scalar，再把这个 scalar advantage 平均给所有 timesteps，就会同时犯两个错误：

\begin{enumerate}[leftmargin=1.5em]
  \item 不知道 \textbf{哪个 timestep} 应该承担更大 credit。
  \item 不知道 \textbf{哪个 reward objective} 在当前生成阶段更应该被强调。
\end{enumerate}

\begin{intubox}
\textbf{直觉：} 早期 denoising 可能更影响 layout / semantic structure，中后期可能更影响 texture / aesthetic / motion details。把所有目标都压成一个 \(A_i\)，再平等广播给所有 steps，其实是在假设“每个阶段都应该优化同一种东西”，这个假设很粗。
\end{intubox}

\subsection{1. 标准 GRPO 的问题：一个 advantage 管所有 steps}

普通 visual GRPO 对同一个 prompt 采样 \(G\) 个样本，得到每个样本的 reward：

\[
r_1,r_2,\ldots,r_G .
\]

然后做 group-wise normalized advantage：

\[
A_i
=
\frac{
r_i-\mathrm{mean}(r_1,r_2,\ldots,r_G)
}{
\mathrm{std}(r_1,r_2,\ldots,r_G)
}.
\]

训练目标是：

\[
\mathcal{J}(\theta)
=
\mathbb{E}
\left[
\frac{1}{G}\sum_{i=1}^{G}
\frac{1}{T}\sum_{t=1}^{T}
\min
\left(
\rho_t^i A_i,\,
\mathrm{clip}(\rho_t^i,1-\epsilon,1+\epsilon)A_i
\right)
\right].
\]

这里的问题非常直接：

\[
A_i
\quad
\text{和}
\quad
t
\quad
\text{无关。}
\]

也就是说，一张图最后 reward 高，整条 trajectory 所有 steps 都被正向更新；一张图最后 reward 低，整条 trajectory 所有 steps 都被压低。OTCA 要把这个单一的 \(A_i\) 换成 timestep-specific effective advantage：

\[
A_i
\quad \Longrightarrow \quad
\tilde{A}_t^i .
\]

\subsection{2. OTCA 总公式：when 乘 what}

OTCA 的核心公式是：

\[
\tilde{A}_{t}^{i}
=
w_t^i
\cdot
\sum_{k=1}^{K} c_k^i A_k^i .
\]

这个式子可以直接读成：

\[
\text{effective advantage}
=
\text{temporal importance}
\times
\text{objective mixture advantage}.
\]

各个符号含义是：

\begin{itemize}[leftmargin=1.5em]
  \item \(i\)：同一个 prompt 下的第 \(i\) 条 sampled trajectory。
  \item \(t\)：flow / denoising trajectory 上的某个 timestep。
  \item \(k\)：第 \(k\) 个 reward objective，例如 CLIP-T、Aesthetic、HPS，或者 video 里的 visual / motion / text alignment。
  \item \(A_k^i\)：第 \(i\) 个样本在第 \(k\) 个 reward objective 上的 group-normalized advantage。
  \item \(w_t^i\)：第 \(t\) 个 step 对最终结果的相对贡献。
  \item \(c_k^i\)：第 \(i\) 条 trajectory 上，不同 reward objectives 的融合权重。
\end{itemize}

\begin{center}
  \includegraphics[width=0.92\linewidth]{otca/fig2_method_overview.png}

  {\small \textbf{图 2：} OTCA 方法概览。TCD 产生 timestep 权重 \(w_t^i\)，MOCA 产生多目标权重 \(c_k^i\)，最后合成 \(\tilde{A}_t^i\)。图源：arXiv HTML。}
\end{center}

\begin{warnbox}
\textbf{一个细节要注意：} 论文文字里常说 timestep-aware objective mixture，但主公式里的 MOCA 权重是 \(c_k^i\)，不是 \(c_{t,k}^i\)。也就是说，严格从公式看，OTCA 的 objective mixture 是 sample / trajectory conditioned；真正随 timestep 变化的是 \(w_t^i\)。它不是完全为每个 timestep 单独求一组 objective weights。
\end{warnbox}

\subsection{3. Flow-SDE 背景：为什么有 trajectory 可以做 credit}

Flow matching 的确定性采样通常写成 ODE：

\[
\mathrm{d}\mathbf{z}_t
=
\mathbf{u}_t\,\mathrm{d}t .
\]

但 GRPO 需要 stochastic trajectory，才能把每一步看成 policy transition 并计算 ratio。因此论文沿用 prior work，把 ODE 换成 reverse-time SDE：

\[
\mathrm{d}\mathbf{z}_t
=
\left(
\mathbf{u}_t
-
\frac{1}{2}\varepsilon_t^2
\nabla_{\mathbf{z}_t}\log p_t(\mathbf{z}_t)
\right)\mathrm{d}t
+
\varepsilon_t\,\mathrm{d}\mathbf{w}.
\]

这样每个 step 都可以被理解为：

\[
\mathbf{a}_{t,i}
\sim
\pi_\theta(\cdot|\mathbf{s}_{t,i}),
\]

于是可以写：

\[
\rho_t^i
=
\frac{
\pi_\theta(\mathbf{a}_{t,i}|\mathbf{s}_{t,i})
}{
\pi_{\theta_{\mathrm{old}}}(\mathbf{a}_{t,i}|\mathbf{s}_{t,i})
}.
\]

\textbf{这一点和 PCPO 是同一个入口：} 都先把 flow ODE 放到 SDE / stochastic policy 的框架里，再谈 GRPO objective。区别是 PCPO 研究 \(\rho_t^i\) 本身的 hidden weight；OTCA 研究 \(\tilde{A}_t^i\) 应该怎么构造。

\subsection{4. TCD：用 latent similarity 估计哪个 step 重要}

TCD 全称是 Trajectory-Level Credit Decomposition。它负责估计：

\[
w_t^i
=
\text{第 \(i\) 条 trajectory 的第 \(t\) 个 step 应该拿多少 temporal credit。}
\]

它没有像 DenseGRPO 那样对每个 \(x_t\) 做 ODE readout 再打 reward，而是用一个更便宜的 latent proxy。

\subsubsection{4.1 先看中间 latent 和 final latent 的相似度}

对第 \(i\) 条 trajectory，记中间 latent 为 \(z_t^i\)，最终 latent 为 \(z_{\mathrm{final}}^i\)。TCD 计算它们的 cosine similarity：

\[
S_t^i
=
\frac{
z_t^i \cdot z_{\mathrm{final}}^i
}{
\|z_t^i\|\,\|z_{\mathrm{final}}^i\|
}.
\]

这个 \(S_t^i\) 的直觉是：

\[
\text{当前 latent 和最终结果方向越接近，说明它越像最终要生成的东西。}
\]

但论文不直接用 \(S_t^i\)，而是用相邻两步的变化量：

\[
\Delta S_t^i
=
S_t^i-S_{t+1}^i .
\]

注意这里的符号和时间方向有关。因为 denoising / flow trajectory 通常从 noisy 走向 clean，论文按它自己的 \(t\) 编号写成 \(S_t^i-S_{t+1}^i\)。我们读的时候只要抓住意思：

\[
\Delta S_t^i
\approx
\text{这个 transition 让 latent 朝 final latent 靠近了多少。}
\]

如果某一步让 latent 更明显地对齐最终状态，就认为它对最终结果更重要。

\subsubsection{4.2 把 similarity gain 变成 temporal weight}

OTCA 把负的贡献截掉，然后归一化：

\[
w_t^i
=
\frac{
\max(0,\Delta S_t^i)+\epsilon
}{
\sum_{\tau}
\left(
\max(0,\Delta S_\tau^i)+\epsilon
\right)
}.
\]

这保证：

\[
w_t^i \ge 0,
\qquad
\sum_t w_t^i = 1.
\]

于是每条 trajectory 的 temporal credit 被重新分配到不同 timesteps：

\[
\text{重要 steps}
\Rightarrow
\text{更大的 }w_t^i,
\qquad
\text{不重要 steps}
\Rightarrow
\text{更小的 }w_t^i.
\]

\begin{intubox}
\textbf{和 DenseGRPO 的区别：} DenseGRPO 的 \(\Delta R_t\) 更像真实 reward gain，但要 ODE readout 和 reward scoring，成本高。OTCA 的 \(\Delta S_t\) 是 latent proxy，不需要额外 reward model 调每个中间点，便宜很多，但因果性弱一些。
\end{intubox}

\begin{center}
  \includegraphics[width=0.9\linewidth]{otca/fig6_tcd_proxy_validation.png}

  {\small \textbf{图 3：} 论文在 appendix 里用 \(\Delta S_t^i\) 和真实 reward gain \(\Delta R_t^i\) 的趋势对比来支持 TCD。它们大体相关，但不是严格相等。图源：arXiv HTML。}
\end{center}

\subsection{5. MOCA：多目标 reward 不要粗暴相加}

MOCA 全称是 Multi-Objective Credit Allocation。它负责估计：

\[
c_k^i
=
\text{第 \(i\) 条 trajectory 上第 \(k\) 个 reward objective 的权重。}
\]

假设我们有 \(K\) 个 reward objectives。对第 \(i\) 个样本，每个 objective 都可以得到一个 group-normalized advantage：

\[
\mathbf{A}^i
=
\left[
A_1^i,A_2^i,\ldots,A_K^i
\right].
\]

如果直接加权平均：

\[
A_{\mathrm{sum}}^i
=
\sum_{k=1}^{K}\alpha_k A_k^i,
\]

就有两个问题：

\begin{enumerate}[leftmargin=1.5em]
  \item 不同 reward 可能冲突，比如 aesthetic 上升但 text alignment 下降。
  \item 固定 \(\alpha_k\) 不知道当前样本和当前 trajectory 更需要哪个 objective。
\end{enumerate}

\subsubsection{5.1 从 gradient-space MOO 到 advantage-space MOO}

标准 multi-objective optimization 可以在 gradient space 里做：

\[
\min_{\{c_k\}_{k=1}^{K}}
\left\|
\sum_{k=1}^{K}
c_k
\nabla_\theta \mathcal{L}_k(\theta)
\right\|^2,
\qquad
\sum_{k=1}^{K}c_k=1,\quad c_k\ge 0.
\]

但对大视觉生成模型来说，为每个 objective 单独算 gradient 再解 MOO 太贵。OTCA 的观察是：对同一个 sample，在同一次 policy update 里，不同 objective 的梯度主要差在 scalar advantage \(A_k^i\)。于是把问题投影到 advantage space：

\[
\min_{\{c_k^i\}_{k=1}^{K}}
\left\|
\sum_{k=1}^{K}
c_k^i A_k^i
\right\|^2,
\qquad
\sum_{k=1}^{K}c_k^i=1,\quad c_k^i\ge 0.
\]

这就是 MOCA 的主体：不用真正存每个 objective 的梯度，而是在低维 advantage vector 上决定 reward mixture。

\subsubsection{5.2 为什么还要 exploration bias}

单纯 minimum-norm 有一个问题：它容易太保守。因为它倾向于找到一个让综合 advantage 接近 0 的组合，避免冲突，但也可能把真正高收益的方向抵消掉。

所以 OTCA 构造了一个 trajectory-level exploration signal：

\[
q_i
=
\sum_{t=1}^{T} w_t^i \Delta S_t^i,
\qquad
e_i
=
\frac{
|q_i|
}{
\mathrm{std}(\{q_j\}_{j=1}^{G})+\epsilon
}.
\]

其中 \(q_i\) 表示这条 trajectory 的结构变化强度，\(e_i\) 表示它相对组内其他 samples 有多 distinctive。然后定义 exploration strength：

\[
\lambda_i
=
\frac{
\max\left(0,\sum_{k=1}^{K} A_k^i\right)
}{
\left|\sum_{k=1}^{K} A_k^i\right|+\epsilon
}
\cdot
\tanh(e_i).
\]

这个设计有两个门控：

\begin{itemize}[leftmargin=1.5em]
  \item 只有综合 advantage 为正时才鼓励探索，避免放大坏方向。
  \item \(\tanh(e_i)\) 把探索强度压到有界范围，避免异常 trajectory 造成不稳定。
\end{itemize}

最终的 MOCA solver 变成：

\[
\min_{\{c_k^i\}}
\left\|
\sum_{k=1}^{K} c_k^i A_k^i
\right\|^2
-
\lambda_i
\sum_{k=1}^{K} c_k^i A_k^i,
\qquad
\sum_{k=1}^{K}c_k^i=1,\quad c_k^i\ge 0.
\]

前半项负责 conflict-aware fusion，后半项负责不要太保守。

\subsection{6. 最终训练目标}

OTCA 最后只是把原本的 \(A_i\) 换成 \(\tilde{A}_t^i\)：

\[
\mathcal{J}_{\mathrm{OTCA}}(\theta)
=
\mathbb{E}
\left[
\frac{1}{G}\sum_{i=1}^{G}
\frac{1}{T}\sum_{t=1}^{T}
\min
\left(
\rho_t^i \tilde{A}_t^i,\,
\mathrm{clip}(\rho_t^i,1-\epsilon,1+\epsilon)\tilde{A}_t^i
\right)
\right].
\]

把 \(\tilde{A}_t^i\) 展开就是：

\[
\rho_t^i
\left(
w_t^i
\sum_{k=1}^{K}c_k^i A_k^i
\right).
\]

所以每个 step 的实际更新强度由三部分共同决定：

\[
\text{policy ratio}
\times
\text{temporal credit}
\times
\text{objective credit}.
\]

\begin{summarybox}
\textbf{和 PCPO 的连接：} PCPO 提醒我们 \(\rho_t^i\) 里也有 timestep bias；OTCA 处理的是 \(\tilde{A}_t^i\) 的结构。更完整的 credit assignment 其实需要同时看 ratio 侧和 advantage 侧。
\end{summarybox}

\subsection{7. 实验怎么证明它有效}

\subsubsection{7.1 Image generation}

图像实验用 FLUX.1-dev。训练 reward 包括 CLIP-T、LAION-Aesthetic、HPSv2.1；同时用 PickScore 和 ImageReward 看泛化。

\begin{center}
\resizebox{\linewidth}{!}{
\begin{tabular}{lccccc}
\toprule
\textbf{Method} & \textbf{CLIP-T} & \textbf{Aesthetic} & \textbf{HPS} & \textbf{PickScore} & \textbf{ImageReward}\\
\midrule
Flux & 0.2682 & 6.2508 & 0.2986 & 22.66 & 1.1001\\
DanceGRPO & 0.2619 & \textbf{6.8917} & 0.3018 & 22.70 & 1.0172\\
VIPO & 0.2722 & 6.7813 & 0.3026 & 22.69 & 1.1128\\
OTCA & \textbf{0.3071} & 6.6028 & \textbf{0.3225} & \textbf{22.97} & \textbf{1.1998}\\
\bottomrule
\end{tabular}
}

{\small \textbf{表 1：} 图像指标。OTCA 在 CLIP-T、HPS、PickScore、ImageReward 上最好；Aesthetic 是 DanceGRPO 更高。}
\end{center}

这里最值得注意的不是“每个指标都赢”，而是 \textbf{trade-off 更均衡}。DanceGRPO 的 aesthetic 很高，但 semantic / preference 指标不一定一起变好；OTCA 通过 MOCA 把多个 objective 更平衡地组织起来。

\begin{center}
  \includegraphics[width=0.95\linewidth]{otca/fig3_image_qualitative.png}

  {\small \textbf{图 4：} 图像生成质性结果。OTCA 主要展示语义一致性、结构清晰度和局部细节上的提升。图源：arXiv HTML。}
\end{center}

\subsubsection{7.2 Video generation}

视频实验用 Wan2.2-T2V-14B-480P，reward 来自 VideoAlign，包含 visual quality、motion quality、text alignment。

\begin{center}
\resizebox{\linewidth}{!}{
\begin{tabular}{lccccccc}
\toprule
\textbf{Method} & \textbf{Color} & \textbf{Dynamic} & \textbf{Imaging} & \textbf{Multiple} & \textbf{Scene} & \textbf{Spatial} & \textbf{Subject}\\
\midrule
Wan2.2 & 89.79 & 70.83 & 68.34 & 63.79 & 35.90 & 80.54 & 93.78\\
DanceGRPO & 90.35 & 71.45 & 67.93 & 64.85 & 35.18 & 83.25 & 93.76\\
OTCA & \textbf{92.88} & \textbf{75.00} & \textbf{68.58} & \textbf{68.45} & \textbf{38.45} & \textbf{88.57} & \textbf{93.95}\\
\bottomrule
\end{tabular}
}

{\small \textbf{表 2：} VBench 细分指标。OTCA 在列出的 video 维度上都优于 Wan2.2 和 DanceGRPO。}
\end{center}

\begin{center}
\resizebox{0.7\linewidth}{!}{
\begin{tabular}{lccc}
\toprule
\textbf{Method} & \textbf{Quality} & \textbf{Semantic} & \textbf{Total}\\
\midrule
Wan2.2 & 83.30 & 73.10 & 81.26\\
DanceGRPO & 83.31 & 73.31 & 81.31\\
OTCA & \textbf{83.72} & \textbf{75.19} & \textbf{82.01}\\
\bottomrule
\end{tabular}
}

{\small \textbf{表 3：} VBench 总分。提升主要体现在 semantic score。}
\end{center}

\begin{center}
  \includegraphics[width=0.95\linewidth]{otca/fig4_video_qualitative.png}

  {\small \textbf{图 5：} 视频生成质性结果。论文强调 OTCA 改善物理合理性、运动演化和语义一致性。图源：arXiv HTML。}
\end{center}

\subsubsection{7.3 Ablation：TCD 和 MOCA 都有用}

\begin{center}
\resizebox{\linewidth}{!}{
\begin{tabular}{lccccc}
\toprule
\textbf{Ablation} & \textbf{CLIP-T} & \textbf{Aesthetic} & \textbf{HPS} & \textbf{PickScore} & \textbf{ImageReward}\\
\midrule
Flux & 0.2682 & 6.2508 & 0.2986 & 22.66 & 1.1001\\
MOCA-only & 0.2787 & 6.2120 & 0.3049 & 22.64 & 1.1250\\
TCD-only & 0.2708 & 6.3270 & 0.3104 & 22.87 & \textbf{1.2618}\\
OTCA Full & \textbf{0.3071} & \textbf{6.6028} & \textbf{0.3225} & \textbf{22.97} & 1.1998\\
\bottomrule
\end{tabular}
}

{\small \textbf{表 4：} Ablation。TCD-only 和 MOCA-only 都有提升，完整 OTCA 在多数指标上最好。}
\end{center}

\begin{center}
  \includegraphics[width=0.95\linewidth]{otca/fig5_reward_curves.png}

  {\small \textbf{图 6：} Reward curves。完整 OTCA 的收敛速度和最终 reward 普遍更好。图源：arXiv HTML。}
\end{center}

\subsection{8. 训练成本和实现细节}

论文给出的主要实现设置是：

\begin{itemize}[leftmargin=1.5em]
  \item \textbf{Image：} FLUX.1-dev，group size \(G=12\)，训练分辨率 \(512\times512\)，16 sampling steps。
  \item \textbf{Video：} Wan2.2-T2V-14B-480P，group size \(G=32\)，训练分辨率 \(240\times416\times53\)，16 sampling steps。
  \item \textbf{Inference：} Flux 用 \(512\times512\)、50 steps；Wan2.2 用 \(480\times832\times53\)、50 steps。
  \item \textbf{Hardware：} image experiments 用 8 张 NVIDIA H100；video experiments 用 32 张 NVIDIA H100。
  \item \textbf{SDE stochasticity：} \(\eta=0.3\) for Flux，\(\eta=0.8\) for Wan2.2。
  \item \textbf{Learning rate：} image \(1\times10^{-5}\)，video \(5\times10^{-6}\)。
\end{itemize}

\begin{warnbox}
\textbf{成本判断：} OTCA 比 DenseGRPO 那种每个 timestep ODE readout 便宜，因为 TCD 用 latent similarity proxy；但它并不是低成本小实验，尤其视频部分仍然是 32 张 H100 级别。它的主要节省来自“不对每个中间 latent 都 decode / reward scoring”，不是来自 GRPO 本身变便宜。
\end{warnbox}

\subsection{9. 我对这篇的理解和保留意见}

\subsubsection{9.1 最值得学的点}

\textbf{OTCA 的最大价值是把 credit assignment 拆成了两个正交轴：}

\[
\text{when}
\quad
\text{and}
\quad
\text{what}.
\]

这比只说“每个 step 都有不同 reward”更进一步。因为在真实 visual post-training 里，reward 往往不是一个，而是一组：

\[
\text{text alignment},\quad
\text{aesthetic},\quad
\text{human preference},\quad
\text{motion},\quad
\text{physical plausibility}.
\]

如果只做 temporal credit，不处理 reward objective 冲突，仍然可能把某些 steps 往错误目标上推。

\subsubsection{9.2 最大的局限}

\begin{enumerate}[leftmargin=1.5em]
  \item \textbf{TCD 是 proxy，不是真正 causal reward gain。} \(\Delta S_t^i\) 衡量的是 latent 朝 final latent 的对齐变化，不等价于 reward improvement。它可能和最终 reward 有相关性，但不是 DenseGRPO 式的 readout counterfactual。
  \item \textbf{MOCA 的 \(c_k^i\) 不是逐 timestep 的。} 主公式里 \(c_k^i\) 没有 \(t\)，所以它更多是在 trajectory / sample 层面做 objective mixture，再通过 \(w_t^i\) 分配到各 steps。
  \item \textbf{advantage-space MOO 是近似。} 它避免了 per-objective gradient 的计算，但也意味着它没有真正解决 gradient conflict，只是在 scalar advantage 层面做了一个 proxy。
  \item \textbf{实验成本仍然很高。} 特别是视频部分用了 32 张 H100，这说明方法虽然避免了 Dense readout，但 visual GRPO 本身仍是重计算路线。
\end{enumerate}

\subsection{10. 和我们当前主线的关系}

读完 OTCA 后，这条主线可以更新成：

\begin{center}
\resizebox{\linewidth}{!}{
\begin{tabular}{llll}
\toprule
\textbf{层面} & \textbf{代表论文} & \textbf{解决什么} & \textbf{还没解决什么}\\
\midrule
step reward gain & DenseGRPO & 每步真实 reward improvement & 成本高，readout 太密\\
chunk unit & Chunk-GRPO & 更新粒度更平滑 & 仍多用 terminal advantage\\
branch counterfactual & TempFlow-GRPO & 单步分叉的局部影响 & branch 设计和成本复杂\\
ratio / objective scale & PCPO & ratio 自带 timestep bias & 不直接重构 reward credit\\
temporal + objective credit & OTCA & when + what & TCD / MOCA 都是 proxy\\
\bottomrule
\end{tabular}
}
\end{center}

\begin{summarybox}
\textbf{当前阶段我的判断：} OTCA 值得学，但不要把它看成 credit assignment 的最终答案。它更像一个框架提醒：真实 visual GRPO 的有效 advantage 应该同时包含 temporal importance 和 objective importance。对于我们后续想做的改进，更关键的问题仍然是：能否用低成本方式估计更接近 causal reward gain 的 \(w_t\) 或 chunk-level \(w_{\mathrm{ch}}\)，同时保留多目标 reward 的动态分配。
\end{summarybox}

% ============================================================
% Paper deep reading: OP-GRPO
% Source: https://arxiv.org/html/2604.04142v1
% ============================================================

\section{论文精读：OP-GRPO}

\begin{conceptbox}
\textbf{论文：} \href{https://arxiv.org/abs/2604.04142}{OP-GRPO: Efficient Off-Policy GRPO for Flow-Matching Models}\\
\textbf{简称：} OP-GRPO。\\
\textbf{版本：} arXiv 2604.04142v1。\\
\textbf{方向标签：} Flow Matching、GRPO、Off-policy RL、Replay Buffer、Sample Efficiency、Image / Video Generation。\\
\textbf{阅读日期：} 2026-05-04。\\
\textbf{阅读口径：} 这篇不是前面 DenseGRPO / Chunk-GRPO / PCPO 那条 ``credit assignment'' 主线，而是 \textbf{sample efficiency / off-policy reuse} 主线：核心问题不是 reward 该分给哪一步，而是 \textbf{已经采出来的好轨迹能不能安全地重复利用}。
\end{conceptbox}

\subsection{0. 一句话结论}

OP-GRPO 的核心不是单独加一个 replay buffer，而是把 Flow-GRPO 里 ``采一次、用一次、立刻丢掉'' 的流程，改成了下面这个三段式：

\[
\begin{aligned}
\text{OP-GRPO}
&= \text{高质量轨迹复用}\\
&\quad + \text{sequence-level 分布纠偏}\\
&\quad + \text{late-step truncation}
\end{aligned}
\]

最后效果是：在 image 和 video 上，它都能用明显更少的训练步数达到 Flow-GRPO 的最终水平，而且最好结果还略高。

\begin{summarybox}
\textbf{我的判断：} 这是视觉生成里 ``GRPO 样本效率'' 方向非常值得读的第一篇。它不是只说 ``把旧样本拿回来再训''，而是把真正会导致失败的三个点都正面处理了：\key{怎么存}、\key{怎么纠偏}、\key{哪些 off-policy steps 根本不该用}。
\end{summarybox}

\begin{figure}[htbp]
  \centering
  \includegraphics[width=\linewidth]{opgrpo/fig1_framework.png}
  \caption{OP-GRPO 总体框架，包括 rollout 和 training 两部分。蓝色区域表示来自 replay buffer 的样本，绿色区域表示直接来自数据集的新样本。图源：arXiv HTML。}
  \label{fig:opgrpo-framework}
\end{figure}

\subsection{1. 它为什么重要}

Flow-GRPO 已经证明了：GRPO 可以有效地接到 flow-matching 的图像/视频生成后训练上。但它有一个非常直接的工程问题：\textbf{完全 on-policy}。

这意味着每一轮训练都要重新生成一整组 denoising trajectories，而上一轮里那些已经算过 reward、甚至已经发现很高质量结果的 trajectories，更新完以后就被丢掉。

对 LLM 来说，丢掉历史 token rollout 已经很可惜；对图像和视频生成来说，这个代价更高，因为一次样本不只是一个 response，而是一整条 denoising trajectory：

\[
z_T \rightarrow z_{T-1} \rightarrow \cdots \rightarrow z_1 \rightarrow z_0 .
\]

每条轨迹都带着多步 transition log-probability、一次 reward model 打分、以及相当昂贵的采样成本。所以视觉生成里的 on-policy 浪费，比语言模型里更刺眼。

\begin{warnbox}
\textbf{但这件事不能直接照搬 LLM 的 replay buffer。} Flow-matching 模型的 off-policy 难点在于：不同 denoising steps 的 transition log-probability 数值尺度并不均匀，尤其在低噪声区域会变得非常尖锐。这使得 ``旧轨迹重放'' 在这里不只是一个数据管道问题，而是一个 objective 稳定性问题。
\end{warnbox}

所以 OP-GRPO 真正回答的是：

\begin{quote}
如果我想把历史高质量 trajectory 重新用起来，那么怎样做才能既提高 sample efficiency，又不把 GRPO 的 clipping 语义和训练稳定性搞坏？
\end{quote}

\subsection{2. 标准 Flow-GRPO 的样本效率问题}

论文先沿用标准 flow-matching / Flow-GRPO 的基本设定。

\subsubsection{2.1 生成过程本质上是一条 trajectory}

Flow model 学习一个速度场：

\[
\frac{d x_t}{dt} = v_{\theta}(x,t), \qquad t \in [0,1].
\]

推理时通常用 Euler 形式离散化：

\[
x_{k-1} = x_k + \Delta t \, v_{\theta}(x_k, t_k).
\]

一旦把 deterministic flow sampling 改写成可计算 transition probability 的 rollout diffusion policy，每条 latent trajectory 的对数概率就可以按时间步分解：

\[
\log p(z_{0:T} \mid c)
=
\sum_{t=1}^{T}
\log p(z_{t-1} \mid z_t, c).
\]

这一步很关键，因为后面所有 GRPO / PPO 风格的比值和 off-policy 校正都建立在这条分解上。

\subsubsection{2.2 标准 GRPO 的 advantage 仍然是 terminal reward}

给定同一个 prompt \(c\)，模型采样一个 group，最终只对终点样本 \(x_0^i\) 打 reward：

\[
\hat{A}^i
=
\frac{
R(x_0^i,c)-\mathrm{mean}\{R(x_0^j,c)\}_{j=1}^{G}
}{
\mathrm{std}\{R(x_0^j,c)\}_{j=1}^{G}
}.
\]

然后每个 denoising step 都拿这个同一个 trajectory-level advantage 去做 clipped GRPO update。

这一点在 OP-GRPO 里并没有被改掉。也就是说，这篇不是 DenseGRPO 那样去改 reward density；它默认还是 Flow-GRPO 的 reward 设定，只改 \textbf{样本复用机制}。

\subsubsection{2.3 真正低效的地方：现采现用}

标准 Flow-GRPO 的训练可以概括成：

\[
\text{采一批新 trajectories}
\rightarrow
\text{算 reward / advantage}
\rightarrow
\text{更新一次}
\rightarrow
\text{全部丢弃}.
\]

这里的问题不是 ``没有探索''，而是 ``探索出来的好东西留不住''。尤其当任务较难时，模型偶尔碰到高质量样本，但如果只用一次，那条 beneficial trajectory 对后续训练几乎没有持续影响。

\begin{intubox}
\textbf{直觉：} Flow-GRPO 更像 ``今天找到一条好路，走一次就忘了''；OP-GRPO 想做的是 ``把这条好路先记下来，但只在安全的时候拿出来重复参考''。
\end{intubox}

\subsubsection{2.4 为什么 naive off-policy 会失败}

如果只是把旧 trajectory 硬塞回训练，会遇到两个结构性问题：

\begin{enumerate}[leftmargin=1.5em]
  \item 这些样本不是当前 rollout policy \(\pi_{\mathrm{old}}\) 采出来的，直接优化会有偏。
  \item 在 flow-matching 的 late denoising steps，off-policy ratio 会因为低噪声数值病态而变得极不稳定。
\end{enumerate}

所以 OP-GRPO 后面三个模块，正好一一对应这三个问题：

\[
\begin{aligned}
\text{replay buffer} &\rightarrow \text{如何存},\\
\text{sequence correction} &\rightarrow \text{如何纠偏},\\
\text{truncation} &\rightarrow \text{哪些步不要复用}.
\end{aligned}
\]

\subsection{3. 方法一：High-quality Replay Buffer}

\subsubsection{3.1 它存什么}

对每条被保留的 trajectory，OP-GRPO 在 buffer 里存的是：

\[
x = \big(c, z_{0:T}, \log p_{\mathrm{off}}(z_{0:T}\mid c), R\big).
\]

也就是：

\begin{itemize}[leftmargin=1.5em]
  \item 条件输入 \(c\)；
  \item 整条 latent trajectory \(z_{0:T}\)；
  \item 这条轨迹在旧策略下的 trajectory log-probability；
  \item 最终 reward \(R\)。
\end{itemize}

这说明它不是只存最后一张图，而是把后面做 sequence-level correction 所需的信息一起存下来。

\subsubsection{3.2 它怎么选样本}

论文的 replay 策略很克制，不是无脑把所有历史数据都放进去，而是只保留 \textbf{高质量且相对新鲜} 的轨迹。

具体规则有三条：

\begin{enumerate}[leftmargin=1.5em]
  \item 每个 sampled group 里，只挑 reward 最高的 trajectory 作为候选进入 buffer。
  \item 新候选会和 buffer 里当前最低 reward 的项比较，只有更好才替换。
  \item 对同一个 conditioning input \(c\)，buffer 最多保留一条 trajectory。
\end{enumerate}

第三条很重要。它意味着 replay buffer 的目标不是把某个 prompt 的高分样本堆很多份，而是保证跨 prompt 的覆盖和多样性，避免 buffer collapse 成几个反复出现的 trivial prompt-template。

\begin{summarybox}
\textbf{这一步的设计思路很清楚：} OP-GRPO 不是想把训练慢慢变成离线 RL，而是想拿到一小批 ``最值得再看一眼'' 的历史轨迹。因此它主动牺牲了 buffer 容量上的贪心，换来更小的 off-policy 偏移和更好的多样性。
\end{summarybox}

\subsubsection{3.3 为什么还要做 stepwise decay}

旧轨迹越老，就越 off-policy。论文因此对 buffer 中轨迹的 reward 做 stepwise decay，让旧样本自然更容易被新样本替换。

这相当于在 replay buffer 里手动加入一个 ``时间衰减''：

\[
\text{同样 reward 的旧轨迹，比新轨迹更不值得保留。}
\]

这一步没有复杂理论，但非常符合这里的问题结构，因为 OP-GRPO 并不追求长期长期的离线数据复用，反而追求 ``历史样本足够近，能帮当前 policy 提速，但又别离得太远''。

\subsubsection{3.4 训练时怎么混用 on-policy 和 off-policy}

OP-GRPO 不会让一个 group 全部由旧样本组成。它采用 hybrid rollout：

\begin{itemize}[leftmargin=1.5em]
  \item 对直接来自原始数据集的 prompt，像普通 Flow-GRPO 一样生成 \(G\) 条新 trajectories。
  \item 对来自 replay buffer 的 prompt，只生成 \((G-1)\) 条新的 on-policy trajectories，再把 buffer 里的那一条旧 trajectory 补进来。
\end{itemize}

也就是：

\[
\mathcal{G}(c)
=
\Big\{z_i^{\mathrm{on}} \sim p_{\mathrm{old}}(\cdot\mid c)\Big\}_{i=1}^{G-1}
\cup
\Big\{z^{\mathrm{off}} \sim \mathcal{B}_{\mathrm{off}}\Big\}.
\]

这个设计很巧，因为它没有破坏 GRPO 的 group-relative 结构。每个 buffer prompt 仍然在和 \textbf{当前策略 freshly sampled 的同 prompt 轨迹} 一起做相对比较，只是组里多了一条 ``历史参考''。

\begin{intubox}
\textbf{可以把它理解成：} OP-GRPO 不是拿旧数据单独训一个 batch，而是把旧轨迹插回它熟悉的 prompt 语境里，让它继续参加 group competition。这样既复用了旧经验，又保留了 GRPO 最核心的组内相对排序。
\end{intubox}

\subsection{4. 方法二：Sequence-level Importance Sampling Correction}

这是整篇论文最关键的数学部分。

\subsubsection{4.1 为什么直接把 off-policy correction 塞进 per-step ratio 是错的}

假设我们已经拿到一条旧策略 \(\pi_{\mathrm{off}}\) 采出的轨迹。一个看似自然的做法，是把分布偏移直接并到每个 timestep 的 importance ratio 里。

但这样一来，原本 PPO / GRPO 中非常重要的参考策略 \(p_{\mathrm{old}}\) 会被约掉：

\[
r_t^i(\theta)
=
\frac{p_{\theta_{\mathrm{old}}}(z_{t-1}^i\mid z_t^i,c)}{p_{\theta_{\mathrm{off}}}(z_{t-1}^i\mid z_t^i,c)}
\times
\frac{p_{\theta}(z_{t-1}^i\mid z_t^i,c)}{p_{\theta_{\mathrm{old}}}(z_{t-1}^i\mid z_t^i,c)}
=
\frac{p_{\theta}(z_{t-1}^i\mid z_t^i,c)}{p_{\theta_{\mathrm{off}}}(z_{t-1}^i\mid z_t^i,c)}.
\]

问题在于，clipping 的本意是约束：

\[
\text{当前更新后的策略 } p_{\theta}
\quad \text{相对} \quad
\text{当前更新起点 } p_{\mathrm{old}}
\quad \text{不要走太远。}
\]

如果 ratio 变成 \(p_{\theta}/p_{\mathrm{off}}\)，那 clipping 监控的就不再是 ``本轮更新幅度''，而变成了 ``当前策略离某个更老的历史策略有多远''。这会直接破坏 PPO / GRPO 的近端更新语义。

论文实证上也观察到：这种 naive substitution 会让 \textbf{超过 40\%} 的 off-policy samples 被 clip 掉。

\begin{warnbox}
\textbf{这里一定要抓住：} OP-GRPO 的重点不是 ``要不要 importance sampling''，而是 ``importance sampling 应该加在什么层级''。如果加错层级，虽然形式上也在纠偏，但已经把 GRPO 的 clipping reference 搞坏了。
\end{warnbox}

\subsubsection{4.2 正确做法：保留 step ratio，额外乘一个 sequence-level correction}

OP-GRPO 的解决方案是：

\begin{enumerate}[leftmargin=1.5em]
  \item 每个 timestep 内部，仍然保留原来那个 \(p_{\theta}/p_{\mathrm{old}}\) 的 step ratio，让 clipping 继续对准当前 update 的起点。
  \item 在整个 loss 外层，再乘一个 trajectory-level 的分布修正因子，补偿 \(\pi_{\mathrm{off}}\) 和 \(\pi_{\mathrm{old}}\) 的采样分布差异。
\end{enumerate}

它的目标写成：

\[
\mathcal{J}_{\mathrm{OP\mbox{-}GRPO}}(\theta)
=
\mathbb{E}
\left[
\frac{P_{\pi_{\mathrm{old}}}(\tau)}{P_{\pi_{\mathrm{off}}}(\tau)}
\cdot
f(r, \hat{A}, \theta, \epsilon, \beta)
\right],
\]

其中 sequence-level correction 是：

\[
\frac{P_{\pi_{\mathrm{old}}}(\tau)}{P_{\pi_{\mathrm{off}}}(\tau)}
=
\prod_{t=1}^{T}
\frac{p_{\theta_{\mathrm{old}}}(z_{t-1}^i \mid z_t^i, c)}{p_{\theta_{\mathrm{off}}}(z_{t-1}^i \mid z_t^i, c)}.
\]

这里的 \(f(r,\hat{A},\theta,\epsilon,\beta)\) 就是普通 GRPO 那个 clipped objective。

\subsubsection{4.3 这个设计为什么成立}

这个式子有两个非常直观的好处。

第一，\textbf{clipping 的语义被保住了}。因为 step-level ratio 还是比较 \(p_{\theta}\) 和 \(p_{\mathrm{old}}\)，所以 clipped / unclipped 仍然准确表示 ``本轮 update 是否太激进''。

第二，\textbf{off-policy correction 只有在样本没有被 clip 掉时才真正发挥作用}。这就像一个 ``先过安全检查，再纠偏'' 的机制：

\[
\text{先确认更新幅度不过激}
\rightarrow
\text{再让历史轨迹参与无偏校正}.
\]

结果是，论文报告被 clip 的 off-policy 样本比例从 over 40\% 降到 \textbf{11.8\%}。

\begin{summarybox}
\textbf{我认为这是本文最聪明的一点：} 它不是把 off-policy 问题粗暴地塞进每个 timestep，而是承认 ``当前 update 稳不稳定'' 和 ``历史分布偏不偏'' 是两个层级不同的问题，所以分别在 step-level 和 sequence-level 处理。
\end{summarybox}

\subsection{5. 方法三：Late-step Trajectory Truncation}

即便有了 sequence-level correction，论文仍然观察到一个严重现象：late denoising steps 的 off-policy ratio 依旧会出问题。

\begin{figure}[htbp]
  \centering
  \includegraphics[width=0.96\linewidth]{opgrpo/fig2_logprob_truncation.png}
  \caption{On-policy 和 off-policy 样本在不同 denoising steps 上的 log-probability 变化。虚线表示 truncation 的起始位置。图源：arXiv HTML。}
  \label{fig:opgrpo-logprob-truncation}
\end{figure}

\subsubsection{5.1 经验现象：最后几步会出现 cliff-like drop}

论文把 on-policy 和 off-policy 样本沿着 denoising trajectory 的 log-probability 画出来，发现 off-policy 曲线在大多数 steps 都还相对平稳，但到了最后几个低噪声 steps，会出现非常陡的 cliff-like drop。

这意味着问题不是 ``整条轨迹 uniformly off-policy''，而是 \textbf{late steps 特别病态}。

\subsubsection{5.2 理论原因：当噪声趋近于 0 时 transition 近乎 deterministic}

论文把 transition log-probability 写成显式 Gaussian 形式：

\[
\log p_{\theta}(z_{t-1}^i \mid z_t^i, c)
=
\log \mathcal{N}(z_{t-1}^i; \mu_t, \sigma_t)
=
-\frac{(z_{t-1}^i-\mu_t)^2}{2\sigma_t^2}
-\log \sigma_t
-\frac{1}{2}\log(2\pi).
\]

当 \(\sigma \to 0\) 时，transition variance \(\sigma_t^2\) 迅速变小，过程越来越接近 deterministic。此时上式里那个除以 \(\sigma_t^2\) 的二次项会变得极其敏感，任何一点离散化误差、policy mismatch、数值波动，都会被放大成巨大的 log-probability 差异。

所以这里的问题不是普通意义上的 ``history sample 有点旧''，而是：

\[
\text{flow model 在低噪声尾部的 transition 本来就极尖锐，}
\text{off-policy mismatch 在这里会被数值上放大。}
\]

\subsubsection{5.3 解法：只复用前半段 off-policy 轨迹，后半段重新采样}

OP-GRPO 直接采用 truncation：只使用 off-policy trajectory 到某个截断步 \(t_{\mathrm{off}}\)，之后剩余的 late steps 由当前策略重新生成：

\[
\tau_m
=
\Big\{z_i^{\mathrm{off}}\Big\}_{i=T}^{t_{\mathrm{off}}}
\cup
\Big\{z_j\Big\}_{j=t_{\mathrm{off}}}^{0}.
\]

也就是说，这条 mixed-policy trajectory 的前半段来自历史轨迹，后半段换成当前 policy 接着走。

这个设计有两个含义：

\begin{enumerate}[leftmargin=1.5em]
  \item 早中期的高价值语义结构仍然被复用了。
  \item 最数值病态的低噪声尾部，不再强行做 off-policy correction。
\end{enumerate}

\begin{intubox}
\textbf{直觉版：} OP-GRPO 的态度不是 ``整条旧轨迹都可信''，而是 ``旧轨迹在高噪声到中噪声阶段大体还值得借鉴，但越到最后越应该交还给当前 policy 自己完成''。
\end{intubox}

论文进一步主张：late denoising steps 对最终感知质量的边际贡献相对更小，所以把 off-policy 使用截在这里，对质量影响很有限，但对稳定性帮助很大。

\subsection{6. 实验设置}

\subsubsection{6.1 模型}

论文在两个 backbone 上评估：

\begin{itemize}[leftmargin=1.5em]
  \item \textbf{Stable-Diffusion-3.5-medium}（SD3.5-M），用于图像生成实验。
  \item \textbf{Wan2.1-1.4B}，用于视频生成实验。
\end{itemize}

\subsubsection{6.2 任务}

图像部分有三类任务：

\begin{enumerate}[leftmargin=1.5em]
  \item \textbf{Compositional Image Generation：} 关注对象数量、颜色、空间关系等组合约束，主要指标是 EvalGen / GenEval。
  \item \textbf{Visual Text Rendering：} 关注图中显式文字渲染能力，主要指标是 OCR 规则识别准确率。
  \item \textbf{Human Preference Alignment：} 关注和人类审美偏好的对齐，训练 reward 用 PickScore。
\end{enumerate}

视频部分则在 \textbf{visual text generation / OCR} 任务上测试可扩展性。

\subsubsection{6.3 泛化评测}

为了验证它不是只会刷单一 reward，论文还汇报了多个 unseen metrics：

\begin{itemize}[leftmargin=1.5em]
  \item 图像质量：Aesthetic Predictor、DeQA。
  \item 偏好打分：ImageReward、PickScore、UnifiedReward。
\end{itemize}

这点很重要，因为 off-policy replay 很容易被担心会过拟合当前 reward model；作者专门用 cross-task / cross-metric 评测来回应这个问题。

\subsection{7. 主要结果：效率、图像、视频、泛化}

\subsubsection{7.1 训练效率是最强结论}

论文最核心的 headline result 是：

\begin{itemize}[leftmargin=1.5em]
  \item OP-GRPO 总训练步数是 1150，而 Flow-GRPO 是 1600。
  \item OP-GRPO 平均只需要 \textbf{34.2\%} 的训练步，就能达到 Flow-GRPO 的最终性能。
  \item 最终性能平均还比 Flow-GRPO 高 \textbf{1.17\%}。
\end{itemize}

\begin{figure}[htbp]
  \centering
  \includegraphics[width=\linewidth]{opgrpo/fig3_training_curves.png}
  \caption{OP-GRPO 与 Flow-GRPO 的训练曲线。论文用这张图展示在三个图像任务和一个视频任务上，OP-GRPO 都能更快达到基线最终性能。图源：arXiv HTML。}
  \label{fig:opgrpo-training-curves}
\end{figure}

这说明 replay buffer 在这里不是只带来一点点方差下降，而是真的让 \textbf{达到同一目标所需的探索-优化回合数} 明显减少了。

\subsubsection{7.2 图像任务：中期收益尤其明显}

我把最重要的 task metric 单独摘出来看，会更清楚：

\begin{center}
\resizebox{\linewidth}{!}{
\begin{tabular}{lcccc}
\toprule
\textbf{任务} & \textbf{Flow 中间} & \textbf{OP 中间} & \textbf{Flow 最好} & \textbf{OP 最好}\\
\midrule
Compositional Image Generation & 0.68 @ 1020 & 0.88 @ 1020 & 0.95 & 0.96\\
Visual Text Rendering & 0.70 @ 720 & 0.81 @ 720 & 0.92 & 0.93\\
Human Preference Alignment & 22.19 @ 720 & 23.27 @ 720 & 23.32 & 23.64\\
\bottomrule
\end{tabular}
}
\end{center}

从这张表能直接看出两个模式：

\begin{enumerate}[leftmargin=1.5em]
  \item \textbf{中间 checkpoint 提前拉开差距。} 这最符合 sample efficiency 的主张，因为 OP-GRPO 在训练还没跑完时就已经更强。
  \item \textbf{最好结果基本不掉，甚至略涨。} 这说明 off-policy reuse 没有用效率换上限。
\end{enumerate}

更细一点地看，三项任务的中期收益都相当大：

\begin{itemize}[leftmargin=1.5em]
  \item 组合生成从 0.68 提到 0.88，提升非常明显，说明 buffer 对 ``稀有但高质量的 compositional success case'' 特别有帮助。
  \item 文本渲染从 0.70 提到 0.81，说明 replay 对视觉文字这种难任务也有增益。
  \item 偏好对齐从 22.19 提到 23.27，连 Aesthetic、DeQA、ImageReward 等 cross-metric 也同步上升，说明不是只把 PickScore 刷高了。
\end{itemize}

\begin{figure}[htbp]
  \centering
  \includegraphics[width=\linewidth]{opgrpo/fig4_image_results.png}
  \caption{SD3.5-M 上三类图像任务的质性结果。论文用这些案例说明 OP-GRPO 在组合约束、可读文本和偏好质量上都更早表现出优势。图源：arXiv HTML。}
  \label{fig:opgrpo-image-results}
\end{figure}

质性图也和定量结论一致：组合任务里对象数量和属性更准，文字渲染更清楚，人类偏好任务里细节和整体语义一致性更好。

\subsubsection{7.3 泛化：不像 reward hacking}

如果 off-policy replay 只是帮助模型更快记住当前 reward 的漏洞，我们通常会看到 task metric 上升，但 unseen metrics 下降。

而这篇的结果不是这样。以组合生成任务为例，在 1020 step 时，虽然 OP-GRPO 的主指标比 Flow-GRPO 高很多（0.88 vs. 0.68），但 Aesthetic、DeQA、PickScore、UnifiedReward 也整体保持竞争力，最好 checkpoint 下 UnifiedReward 还从 3.51 提到 3.57。

作者还在更难的 \textbf{T2I-CompBench++} 上做了扩展验证：

\begin{itemize}[leftmargin=1.5em]
  \item 中间 checkpoint average：Flow-GRPO 0.520，OP-GRPO \textbf{0.567}。
  \item 最好 checkpoint average：Flow-GRPO 0.573，OP-GRPO \textbf{0.579}。
\end{itemize}

这个结果我会更看重中间值，因为它更直接体现 ``更快学会组合约束''，而不是只在训练末期略微抬高 ceiling。

\subsubsection{7.4 视频任务：难任务上收益更明显}

视频实验更能说明问题，因为它更难、成本更高、也更依赖 bootstrap。

论文在 Wan2.1-1.4B 的 OCR 任务上报告：OP-GRPO 只用 \textbf{30.1\%} 的训练时间，就达到 Flow-GRPO 的最终性能，而且最终质量还更好。

作者给出的解释也很合理：难任务下，纯 on-policy 方法经常是 ``偶尔撞到一次高质量样本，但马上丢掉''，所以很难持续强化正确模式；而 off-policy replay 恰好能把这些稀有但关键的 trajectories 留下来反复利用。

\begin{figure}[htbp]
  \centering
  \includegraphics[width=\linewidth]{opgrpo/fig5_video_results.png}
  \caption{Wan2.1-1.4B 上视频 OCR 任务的质性结果。论文用这些例子展示 OP-GRPO 在更困难的视频文字生成上也能明显优于 Flow-GRPO。图源：arXiv HTML。}
  \label{fig:opgrpo-video-results}
\end{figure}

\subsection{8. Ablation 与我的判断}

\begin{figure}[htbp]
  \centering
  \includegraphics[width=\linewidth]{opgrpo/fig6_ablation.png}
  \caption{OP-GRPO 的消融实验。左图比较去掉 sequence correction 或 truncation 后的稳定性，右图比较不同 off-policy 比例下的收敛速度与稳定性。图源：arXiv HTML。}
  \label{fig:opgrpo-ablation}
\end{figure}

\subsubsection{8.1 三个模块不是独立加分项，而是相互依赖}

消融最值得记住的结论是：

\begin{itemize}[leftmargin=1.5em]
  \item 去掉 \textbf{sequence correction}，训练会明显不稳定，甚至发散。
  \item 去掉 \textbf{truncation}，同样会出现显著 instability。
\end{itemize}

这说明 replay buffer 本身并不神奇。真正让它可用的是：

\[
\text{replay}
\quad + \quad
\text{distribution correction}
\quad + \quad
\text{late-step stabilization}
\]

也就是说，这篇论文不是 ``buffer matters'' 那种单点工程 tweak，而是一个完整的 off-policy recipe。

\subsubsection{8.2 off-policy 比例不是越高越好}

作者还测了 off-policy sample proportion 的影响：

\begin{itemize}[leftmargin=1.5em]
  \item 5\% 的 off-policy 比例也能带来 \textbf{23.1\%} 的收敛加速。
  \item 更高比例会更快，但也更不稳定，甚至可能发散。
\end{itemize}

这和方法设计完全一致。因为 off-policy 的收益来自复用高价值样本，但风险也来自分布偏移；所以最好的区间一定是一个折中，而不是 ``越多 replay 越好''。

\begin{summarybox}
\textbf{我的判断：} OP-GRPO 真正揭示的是一个 ``可用的 off-policy 窗口''。buffer 太弱，收益不明显；buffer 太强，distribution shift 和 late-step numeric issue 会反噬训练。它通过 decay、sequence correction 和 truncation，把这个窗口撑开到一个工程上可用的范围。
\end{summarybox}

\subsubsection{8.3 我认为它的局限}

这篇工作很强，但也有边界：

\begin{enumerate}[leftmargin=1.5em]
  \item 它的 buffer prioritization 基本只看 reward，高 reward 不一定等于最有训练价值。
  \item 同一个 prompt 最多保留一条 trajectory，这保证了多样性，但也可能损失 ``同 prompt 下多条不同成功模式'' 的复用机会。
  \item truncation 在这篇的 image/video 任务上成立，但如果未来任务更依赖最终精细局部控制，\(t_{\mathrm{off}}\) 可能需要重新调。
  \item 论文目前只验证到 image / video；扩展到 audio、3D、长程世界模型还没有证据。
\end{enumerate}

\subsection{9. 它和 Flow-GRPO / DenseGRPO / PCPO 的关系}

可以把几篇论文放在一张图里理解：

\begin{center}
\resizebox{\linewidth}{!}{
\begin{tabular}{lll}
\toprule
\textbf{论文} & \textbf{主要改哪里} & \textbf{核心问题}\\
\midrule
Flow-GRPO & on-policy GRPO baseline & 让 flow model 能做 GRPO 后训练\\
DenseGRPO & advantage / reward side & 最终 reward 应该怎样分给不同 steps\\
PCPO & ratio / objective scale side & step ratio 是否隐式偏向某些 timesteps\\
OP-GRPO & data reuse / off-policy correction side & 历史高质量轨迹能否安全复用\\
\bottomrule
\end{tabular}
}
\end{center}

它们之间的关系我会这样记：

\begin{itemize}[leftmargin=1.5em]
  \item \textbf{和 Flow-GRPO：} OP-GRPO 直接建立在 Flow-GRPO 上，reward 形式、group-relative advantage 和 clipped GRPO 骨架都没变，变的是数据来源和分布校正。
  \item \textbf{和 DenseGRPO：} DenseGRPO 回答 ``同一条 trajectory 里，哪个 step 贡献大''；OP-GRPO 回答 ``旧 trajectory 能不能再用''。两者是明显正交的，理论上完全可以组合。
  \item \textbf{和 PCPO：} PCPO 研究的是 on-policy 情况下，flow SDE ratio 自带的 timestep weight 是否失衡；OP-GRPO 研究的是 off-policy 情况下，ratio / logprob 在 late steps 为什么会数值病态，以及如何校正。两篇都碰 ratio side，但关心的问题不同。
\end{itemize}

\begin{summarybox}
\textbf{我最后会这样概括：} DenseGRPO / PCPO 更像是在问 ``同样一批样本，怎么学得更准''；OP-GRPO 问的是 ``已经花大代价采出来的样本，能不能别只学一次''。如果后面要继续看 E-GRPO、MaskFocus、LeapAlign，这篇就是整条 ``采样效率'' 支线的第一个支点。
\end{summarybox}

% ============================================================
% Paper deep reading: E-GRPO
% Source: https://arxiv.org/html/2601.00423v1
% ============================================================

\section{论文精读：E-GRPO}

\begin{conceptbox}
\textbf{论文：} \href{https://arxiv.org/abs/2601.00423}{E-GRPO: High Entropy Steps Drive Effective Reinforcement Learning for Flow Models}\\
\textbf{作者：} Shengjun Zhang, Zhang Zhang, Chensheng Dai, Yueqi Duan。\\
\textbf{简称：} E-GRPO。\\
\textbf{版本：} arXiv 2601.00423v1。\\
\textbf{代码：} \href{https://github.com/shengjun-zhang/VisualGRPO}{VisualGRPO}。\\
\textbf{方向标签：} Flow Matching、GRPO、Entropy-aware Exploration、Step Merging、Visual Preference Alignment。\\
\textbf{阅读日期：} 2026-05-06。\\
\textbf{阅读口径：} 这篇和 OP-GRPO 都属于 ``sample efficiency / exploration efficiency'' 支线，但切入点不同：OP-GRPO 问 \textbf{旧轨迹能不能复用}，E-GRPO 问 \textbf{哪些 denoising steps 真正值得随机探索和做 RL 更新}。
\end{conceptbox}

\subsection{0. 一句话结论}

E-GRPO 的核心判断很直接：Flow-GRPO 不应该在所有 denoising steps 上均匀做 SDE 探索和 GRPO 更新，因为不同时间步的探索熵不一样。

\[
\begin{aligned}
\text{E-GRPO}
&= \text{entropy-aware step selection}\\
&\quad + \text{low-entropy step merging}\\
&\quad + \text{multi-step group-normalized advantage}
\end{aligned}
\]

它认为 high-entropy steps 会产生更有区分度的 rollouts，更适合用 reward 做相对比较；low-entropy steps 的随机扰动太小，生成结果差异不明显，继续把这些 steps 当成独立 RL action 只会增加计算和 reward attribution 噪声。

\begin{summarybox}
\textbf{我的判断：} E-GRPO 不是在改 reward，也不是在做 replay buffer，而是在改 \key{探索发生的位置}。它把 Flow-GRPO 里的 ``每一步都随机、每一步都更新'' 改成 ``只在高信息量的 consolidated SDE step 上随机和更新''。这使它和 OP-GRPO 正好互补：一个提高轨迹复用效率，一个提高单次 rollout 内部的步骤效率。
\end{summarybox}

\subsection{1. 它为什么重要}

Flow-GRPO 把 deterministic flow ODE 改写成 stochastic SDE policy，从而能计算 transition probability、做 GRPO-style policy update。这个改写解决了 ``flow model 没有 logprob 怎么做 RL'' 的问题，但也引入了一个新的默认假设：

\[
\text{所有 denoising steps 都值得注入随机性，也都值得接受同样形式的 RL 更新。}
\]

E-GRPO 认为这个假设是有问题的。因为在 flow / diffusion 采样里，不同时间步的随机扰动影响完全不同：

\begin{itemize}[leftmargin=1.5em]
  \item 早期 high-noise / high-entropy steps 会改变整体构图、语义布局和主体形态，生成结果之间差异大。
  \item 后期 low-noise / low-entropy steps 更多是在局部细节附近微调，rollouts 之间区分度小。
\end{itemize}

如果后期低熵步骤产生的几条样本几乎看不出差异，但我们还强行对它们做 group-relative reward normalization，那么 reward 里的随机波动、scorer 噪声、下游步骤误差就会被误当成该 step 的学习信号。

\begin{warnbox}
\textbf{关键问题：} Flow-GRPO 的稀疏 reward 不只是不知道 ``哪一步贡献大''，还不知道 ``哪一步本来就没有足够探索空间''。如果某一步的 rollout 差异太小，组内相对优势本身就很难可靠。
\end{warnbox}

所以 E-GRPO 的问题可以概括成：

\begin{quote}
在 flow-model RL 中，怎样根据 step entropy 选择真正有探索价值的时间步，并避免低熵多步随机性带来的 reward attribution ambiguity？
\end{quote}

\subsection{2. 熵视角：为什么不是所有 SDE steps 都一样}

论文从 Flow-GRPO 的 SDE 采样形式出发。为了让 flow model 可以探索，原本的 ODE update 会被改写成带随机性的 SDE transition：

\[
x_{t+\Delta t}
=
x_t
+
\mu_{\theta}(x_t,t)\Delta t
+
\sigma_t\sqrt{\Delta t}\epsilon,
\qquad
\epsilon\sim\mathcal{N}(0,I).
\]

其中论文使用的噪声系数形式是：

\[
\sigma_t
=
a\sqrt{\frac{t}{1-t}}.
\]

对于反向 SDE step，它的 transition distribution 可以看成 Gaussian，因此熵主要由 covariance 控制。论文给出的 step entropy 可写成：

\[
h(t)
=
\frac{d}{2}\log\left(2\pi e\sigma_t^2\Delta t\right).
\]

代入 \(\sigma_t=a\sqrt{t/(1-t)}\) 后，有：

\[
h(t)
=
\frac{d}{2}
\log\left(
2\pi e\cdot a^2\cdot \frac{t}{1-t}\cdot \Delta t
\right).
\]

也就是说，在忽略常数后：

\[
e^{h(t)}
\propto
\frac{t}{1-t}\Delta t.
\]

这给了 E-GRPO 一个很明确的判断标准：哪个 step 的 SDE entropy 更大，哪个 step 的随机探索范围就更大，也更可能产生 reward 上可区分的 rollouts。

\begin{intubox}
\textbf{直觉版：} high-entropy step 像是在问 ``我要往哪个语义方向走''；low-entropy step 更像是在问 ``边缘再抹平一点还是锐一点''。前者更适合 RL 探索，后者如果也用同样的 GRPO 强度去训，收益小、噪声大。
\end{intubox}

\begin{figure}[htbp]
\centering
\includegraphics[width=0.95\linewidth]{egrpo/01_entropy_influence.png}
\caption{E-GRPO 对 step entropy 的经验观察。左侧比较不同 SDE sampling 策略下的生成差异，中间展示合并步数如何提升 entropy，右侧展示只训练 high-entropy early steps 的 reward 曲线更优。图源：\href{https://arxiv.org/html/2601.00423v1}{E-GRPO arXiv HTML Figure 1}。}
\label{fig:egrpo-entropy-influence}
\end{figure}

\subsection{3. 方法一：Entropy-driven Step Merging}

\subsubsection{3.1 用阈值区分 high-entropy 和 low-entropy steps}

E-GRPO 首先为每个 denoising timestep 计算 entropy proxy，然后用阈值 \(\tau\) 将时间步分成两类：

\[
\begin{aligned}
\text{high-entropy steps} &: e^{h(t_k)}\geq \tau,\\
\text{low-entropy steps} &: e^{h(t_k)}<\tau.
\end{aligned}
\]

实际 flow sampling 里通常还会做 time shift：

\[
t_k
=
\frac{s\hat{t}_k}{1+(s-1)\hat{t}_k},
\qquad
\hat{t}_k=\frac{k}{T}.
\]

论文将这个时间变换代入 entropy proxy 后，用它来判断哪些 steps 具有足够探索熵。

\subsubsection{3.2 低熵步骤不单独做 SDE，而是合并成一个更大的 SDE step}

对连续低熵步骤 \(\{t_m,\ldots,t_{m-l}\}\)，E-GRPO 不再逐步注入 SDE noise，而是把它们合并成一个 consolidated SDE step：

\[
x_{t_{m-l}}
=
x_{t_m}
+
\mu_{\theta}(x_{t_m},t_m)(t_m-t_{m-l})
+
\sigma_t\sqrt{t_m-t_{m-l}}\epsilon.
\]

合并后的 entropy 近似满足：

\[
e^{h(t_k)}
\propto
\frac{t_m}{1-t_m}(t_m-t_{m-l}).
\]

因为 \((t_m-t_{m-l})\) 变大，所以多个低熵小 step 可以合成一个达到阈值附近的高熵 step。

\begin{summarybox}
\textbf{这一步的设计很关键：} E-GRPO 不是简单跳过 low-entropy steps，而是把连续 low-entropy steps 合并成一个 ``刚好有足够探索空间'' 的 SDE event。这样既减少了独立随机步骤数量，又避免完全失去这段轨迹上的探索能力。
\end{summarybox}

\begin{figure}[htbp]
\centering
\includegraphics[width=0.82\linewidth]{egrpo/02_pipeline.png}
\caption{E-GRPO 的整体 sampling strategy：先为不同 timesteps 构造 anchor noise latents，再根据 entropy analysis 为每个 active SDE timestep 选择 merged steps，最后在对应 group 内生成结果并计算 advantage。图源：\href{https://arxiv.org/html/2601.00423v1}{E-GRPO arXiv HTML Figure 2}。}
\label{fig:egrpo-pipeline}
\end{figure}

\subsubsection{3.3 为什么不用固定 2-step / 4-step / 6-step merge}

一个简单做法是固定每 \(l\) 步合并一次。但 E-GRPO 的消融显示，固定合并不如 adaptive merging。

原因也很自然：不同 timestep 的 entropy 本来就不一样。如果固定 4-step merge，有些位置合并后还是太低熵，有些位置可能熵过高，导致探索方向太散。论文因此为每个 low-entropy timestep 自适应选择 \(l\)，让合并后的 entropy 刚好超过阈值 \(\tau\)。

\[
\text{choose }l
\quad\text{s.t.}\quad
 e^{h(t_m, l)} \gtrsim \tau.
\]

这相当于把每个 consolidated SDE step 的探索尺度校准到一个相近范围。

\subsection{4. 方法二：Multi-step Group Normalized Advantage}

合并步骤之后，advantage 也不能再按原来的单步方式理解。

对一个 active SDE timestep \(t_n\)，它对应一个合并后的时间段：

\[
\mathcal{T}_n
=
\{t_n,t_{n-1},\ldots,t_{n-l_n}\}.
\]

给定同一个 prompt \(c\)，论文为每个 active timestep \(t_n\) 生成 \(G^{(n)}\) 条 trajectories，并只在共享同一个 consolidated SDE step 的样本内部计算 group-normalized advantage：

\[
A_{t_n}^{(i)}
=
\frac{
R(x_{0,t_n}^{(i)},c)
-
\mathrm{mean}\{R(x_{0,t_n}^{(j)},c)\}_{j=1}^{G^{(n)}}
}{
\mathrm{std}\{R(x_{0,t_n}^{(j)},c)\}_{j=1}^{G^{(n)}}
}.
\]

对应的 ratio 是合并 step 的 transition ratio：

\[
r_{t_n}^{(i)}(\theta)
=
\frac{
p_{\theta}(x_{t_{n-l_n}}^{(i)}\mid x_{t_n}^{(i)},c)
}{
p_{\theta_{\mathrm{old}}}(x_{t_{n-l_n}}^{(i)}\mid x_{t_n}^{(i)},c)
}.
\]

于是 clipped surrogate 仍然保持 GRPO / PPO 风格：

\[
\frac{1}{\hat{T}}
\sum_{n=N}^{T}
\frac{1}{G^{(n)}}
\sum_{i=1}^{G^{(n)}}
\min\left(
 r_{t_n}^{(i)}A_{t_n}^{(i)},
 \mathrm{clip}(r_{t_n}^{(i)},1-\epsilon,1+\epsilon)A_{t_n}^{(i)}
\right).
\]

\begin{figure}[htbp]
\centering
\includegraphics[width=0.72\linewidth]{egrpo/03_ambiguous_reward.png}
\caption{连续多步 SDE sampling 会让同一个 final advantage 同时对应多个 timesteps，从而产生 ambiguous reward signal；E-GRPO 通过 merged-step SDE sampling 把最终 advantage 对齐到一个 merged SDE step，降低单步 credit assignment 的歧义。图源：\href{https://arxiv.org/html/2601.00423v1}{E-GRPO arXiv HTML Figure 3}。}
\label{fig:egrpo-ambiguous-reward}
\end{figure}

\begin{warnbox}
\textbf{这里和 DenseGRPO 不一样：} DenseGRPO 是想给每个 step 估计更密的 reward gain；E-GRPO 是先把低信息量 steps 合成一个有意义的 exploration event，再在这个 event 内做 group-relative advantage。它不是 ``更细''，而是 ``该细的地方细，该合的地方合''。
\end{warnbox}

\subsection{5. 算法流程：训练时到底怎么做}

可以把 E-GRPO 的训练流程压缩成五步：

\begin{enumerate}[leftmargin=1.5em]
  \item 计算所有 timesteps 的 entropy proxy \(h(t_k)\)。
  \item 根据阈值 \(\tau\) 识别 high-entropy steps 和 low-entropy steps。
  \item 对连续 low-entropy steps 自适应选择 \(l_n\)，形成 consolidated SDE steps。
  \item 对非 active steps 使用 ODE sampling，对 consolidated active steps 使用 SDE sampling。
  \item 在每个 merged-step group 内计算 advantage 和 clipped GRPO objective。
\end{enumerate}

因此，E-GRPO 实际上把 rollout policy 改成了一个 ODE/SDE hybrid：

\[
\text{ODE on uninformative steps}
\quad + \quad
\text{SDE on entropy-calibrated merged steps}.
\]

\begin{intubox}
\textbf{直觉：} 它不是减少采样步数本身，而是减少 ``无效随机性'' 和 ``无效 RL 更新''。ODE steps 负责稳定推进生成，SDE merged steps 负责真正探索。
\end{intubox}

\subsection{6. 实验设置}

论文的实验设置如下：

\begin{itemize}[leftmargin=1.5em]
  \item \textbf{Backbone：} FLUX.1-dev。
  \item \textbf{数据集：} HPD，约 103K training prompts 和 400 testing prompts。
  \item \textbf{采样步数：} 总采样步数 \(T=16\)。
  \item \textbf{噪声参数：} \(a=0.7\)。
  \item \textbf{默认 entropy threshold：} \(\tau=2.2\)。
  \item \textbf{训练步数：} 300 iterations。
  \item \textbf{训练 reward：} 先使用 HPS-v2.1；为了缓解 reward hacking，又使用 HPS-v2.1 + CLIP 的联合 reward。
  \item \textbf{评测指标：} HPS、CLIP、PickScore、ImageReward。
\end{itemize}

这套设置有一个值得注意的点：作者明确指出只用 HPS 训练可能导致 reward hacking，所以主结果里同时报告了单 reward 和 HPS+CLIP 两种训练方式。

\subsection{7. 主要结果：E-GRPO 的收益在哪里}

\subsubsection{7.1 单 HPS reward 下，HPS 指标最好}

在只用 HPS-v2.1 作为训练 reward 时，E-GRPO 在 HPS 上达到 \textbf{0.391}，高于几个主要 baseline：

\begin{center}
\resizebox{\linewidth}{!}{
\begin{tabular}{lcccc}
\toprule
\textbf{方法} & \textbf{HPS} & \textbf{CLIP} & \textbf{PickScore} & \textbf{ImageReward}\\
\midrule
FLUX.1-dev & 0.311 & 0.388 & 0.231 & 1.089\\
DanceGRPO & 0.353 & 0.375 & 0.228 & 1.233\\
MixGRPO & 0.378 & 0.358 & 0.225 & 1.266\\
GranularGRPO & 0.385 & 0.355 & 0.229 & 1.313\\
TempFlowGRPO & 0.382 & 0.357 & 0.231 & 1.264\\
E-GRPO & \textbf{0.391} & 0.355 & \textbf{0.232} & \textbf{1.324}\\
\bottomrule
\end{tabular}
}
\end{center}

论文强调，E-GRPO 相比 DanceGRPO 在 HPS 上提升约 \textbf{10.8\%}。但我更关注另一个现象：HPS 提高的同时，CLIP 并没有同步提高，说明单 reward 训练确实存在向 HPS scorer 偏移的风险。

\subsubsection{7.2 HPS+CLIP 联合 reward 下，泛化指标更稳}

在 HPS+CLIP 联合训练下，E-GRPO 的表现更有说服力：

\begin{center}
\resizebox{\linewidth}{!}{
\begin{tabular}{lcccc}
\toprule
\textbf{方法} & \textbf{HPS} & \textbf{CLIP} & \textbf{PickScore} & \textbf{ImageReward}\\
\midrule
FLUX.1-dev & 0.311 & 0.388 & 0.231 & 1.089\\
DanceGRPO & 0.331 & 0.389 & 0.227 & 1.128\\
MixGRPO & 0.363 & 0.399 & 0.230 & 1.436\\
GranularGRPO & 0.377 & 0.400 & 0.236 & 1.490\\
BranchGRPO & 0.342 & 0.384 & 0.230 & 1.243\\
TempFlowGRPO & 0.310 & 0.388 & 0.230 & 1.106\\
E-GRPO & \textbf{0.382} & \textbf{0.401} & \textbf{0.237} & \textbf{1.494}\\
\bottomrule
\end{tabular}
}
\end{center}

论文报告，相比 DanceGRPO，E-GRPO 在 ImageReward 上提升 \textbf{32.4\%}，在 PickScore 上提升 \textbf{4.4\%}。这说明 entropy-guided step selection 不只是刷某一个训练 reward，而是能带来更广泛的人类偏好指标改善。

\subsubsection{7.3 训练曲线：早期增长更快，收敛更平滑}

论文还展示了 reward trajectory：E-GRPO 早期 reward 增长更快，最终 reward 更高，曲线也更平滑。

这和方法假设是一致的：如果训练集中在 high-entropy / consolidated informative steps 上，那么每次更新都更可能改变真正影响 reward 的生成方向，而不是在低熵后期步骤上浪费梯度。

\begin{figure}[htbp]
\centering
\includegraphics[width=0.95\linewidth]{egrpo/04_reward_curve.png}
\caption{E-GRPO 与 baseline 的训练 reward 曲线对比。曲线说明 entropy-guided exploration 能更快提升训练 reward，并且相比全 timestep 更新更稳定。图源：\href{https://arxiv.org/html/2601.00423v1}{E-GRPO arXiv HTML Figure 4}。}
\label{fig:egrpo-reward-curve}
\end{figure}

\subsubsection{7.4 可视化结果：语义和细节更均衡}

论文的可视化对比显示，E-GRPO 相比 baseline 更容易同时保留 prompt 语义和细粒度视觉细节。这个现象和前面的机制解释是一致的：high-entropy steps 负责更大的语义探索，merged-step SDE 又避免后期低熵步骤把 reward 噪声放大成错误更新。

\begin{figure}[htbp]
\centering
\includegraphics[width=0.95\linewidth]{egrpo/05_visual_comparison.png}
\caption{E-GRPO 与其他 baseline 的生成结果可视化对比。论文认为 E-GRPO 能更好地整合语义对齐和细节质量。图源：\href{https://arxiv.org/html/2601.00423v1}{E-GRPO arXiv HTML Figure 5}。}
\label{fig:egrpo-visual-comparison}
\end{figure}

\subsection{8. 消融：哪些设计真的重要}

\subsubsection{8.1 Adaptive merging 优于固定步数 merging}

论文比较了固定 2-step、4-step、6-step merging 和 adaptive merging：

\begin{center}
\resizebox{\linewidth}{!}{
\begin{tabular}{lcccc}
\toprule
\textbf{Merging strategy} & \textbf{HPS} & \textbf{CLIP} & \textbf{PickScore} & \textbf{ImageReward}\\
\midrule
2-step & 0.382 & 0.290 & 0.232 & 1.223\\
4-step & 0.374 & 0.302 & 0.230 & 1.216\\
6-step & 0.337 & \textbf{0.372} & 0.226 & 1.298\\
Adaptive & \textbf{0.391} & 0.355 & \textbf{0.232} & \textbf{1.324}\\
\bottomrule
\end{tabular}
}
\end{center}

这说明固定 merge 的问题在于不看局部 entropy：合并太少，探索仍然不足；合并太多，单个 SDE event 过于随机，reward 方向又会变模糊。Adaptive merging 的优势就是把每个 active step 的探索强度拉回到一个合适范围。

\subsubsection{8.2 只训前 8 个 high-entropy steps 最好}

论文还比较了训练不同 timestep subset：

\begin{center}
\resizebox{\linewidth}{!}{
\begin{tabular}{lcccc}
\toprule
\textbf{训练 steps} & \textbf{HPS} & \textbf{CLIP} & \textbf{PickScore} & \textbf{ImageReward}\\
\midrule
First 4 Steps & 0.370 & 0.348 & 0.231 & 1.252\\
First 8 Steps & \textbf{0.391} & 0.355 & \textbf{0.232} & \textbf{1.324}\\
Second 8 Steps & 0.357 & \textbf{0.381} & 0.231 & 1.250\\
Full Steps & 0.366 & 0.359 & 0.231 & 1.169\\
\bottomrule
\end{tabular}
}
\end{center}

结论非常清楚：前 8 个 high-entropy steps 已经足以获得最强性能；训练所有 steps 不但成本更高，效果还不如只训前半段；只训后 8 个 low-entropy steps 则明显更弱。

\begin{summarybox}
\textbf{我会这样记这组消融：} E-GRPO 不是 ``少训一些 steps 也能省钱'' 这么简单，而是证明了 \key{后半段低熵 steps 的 RL 信号质量确实差}。全步骤训练把有用信号和低质量信号混在一起，反而拖累优化。
\end{summarybox}

\subsection{9. 和 OP-GRPO / DenseGRPO / TempFlow-GRPO 的关系}

把 E-GRPO 放进这条 Flow-GRPO 后训练主线里，它的位置大概是：

\begin{center}
\resizebox{\linewidth}{!}{
\begin{tabular}{lll}
\toprule
\textbf{论文} & \textbf{主要改哪里} & \textbf{核心问题}\\
\midrule
TempFlow-GRPO & temporal weighting / branching & 哪个 timestep 更该分到 credit\\
DenseGRPO & dense reward gain & terminal reward 怎样分给每个 step\\
OP-GRPO & off-policy replay & 已经采出来的高质量 trajectory 能否复用\\
E-GRPO & entropy-aware exploration & 哪些 steps 值得随机探索和 RL 更新\\
\bottomrule
\end{tabular}
}
\end{center}

\begin{itemize}[leftmargin=1.5em]
  \item \textbf{和 OP-GRPO：} 两者都关心效率，但 OP-GRPO 是跨 iteration 的样本复用，E-GRPO 是单条 trajectory 内的 step selection。前者解决 ``好轨迹别浪费''，后者解决 ``低熵步骤别乱训''。
  \item \textbf{和 DenseGRPO：} DenseGRPO 把 reward 变密，E-GRPO 把探索步骤变少但更有信息量。一个偏 reward attribution，一个偏 exploration attribution。
  \item \textbf{和 TempFlow-GRPO：} TempFlow-GRPO 也注意到 timestep 不应等权；E-GRPO 则从 SDE entropy 出发，给出更直接的 step merging 和 ODE/SDE hybrid sampling 机制。
\end{itemize}

\begin{intubox}
\textbf{如果组合使用：} 一个自然方向是 OP-GRPO 复用高 reward trajectories，但只复用 / 更新其中 E-GRPO 认为 high-entropy 或 consolidated informative 的 steps。这样可能同时减少采样成本和无效更新成本。
\end{intubox}

\subsection{10. 我认为它的局限}

E-GRPO 的思想很清楚，但也有几个边界：

\begin{enumerate}[leftmargin=1.5em]
  \item \textbf{Entropy proxy 不等于真实任务重要性。} 高熵 step 更可能产生差异，但差异大不一定等于对某个 reward 最关键；任务相关 credit 仍然需要 reward 侧方法补充。
  \item \textbf{阈值 \(\tau\) 仍是超参数。} 论文默认设为 2.2，并在附录做了阈值消融，但换模型、换 scheduler、换采样步数后可能需要重新校准。
  \item \textbf{主要验证在 FLUX.1-dev + HPD。} 它是否同样适用于视频、图像编辑、3D 或 unified multimodal generation，还需要更多证据。
  \item \textbf{reward hacking 没有被根治。} 论文自己也承认 reward design 仍是瓶颈；E-GRPO 能减少无效步骤优化，但不能保证 reward 本身完全符合人类偏好。
\end{enumerate}

\begin{figure}[htbp]
\centering
\includegraphics[width=0.95\linewidth]{egrpo/08_failure_cases.png}
\caption{E-GRPO 的 failure cases。这个图提醒我们：entropy-aware step selection 能提升 RL 更新效率，但不能替代更可靠的 reward model、prompt grounding 和视觉一致性约束。图源：\href{https://arxiv.org/html/2601.00423v1}{E-GRPO arXiv HTML Figure 8}。}
\label{fig:egrpo-failure-cases}
\end{figure}

\subsection{11. 读完后的记忆钩子}

\begin{summarybox}
\textbf{我最后会这样概括：} OP-GRPO 说 ``昂贵采样出来的好轨迹别只用一次''；E-GRPO 说 ``昂贵轨迹里的每一步也别都当成有效探索''。前者是 trajectory-level efficiency，后者是 timestep-level exploration efficiency。两篇连起来，正好构成 Flow-GRPO 采样效率支线的两个层级。
\end{summarybox}

\subsection{参考来源}

\begin{itemize}[leftmargin=1.5em]
  \item \textbf{论文主页：} \href{https://arxiv.org/abs/2601.00423}{E-GRPO: High Entropy Steps Drive Effective Reinforcement Learning for Flow Models}。
  \item \textbf{arXiv HTML：} \href{https://arxiv.org/html/2601.00423v1}{arXiv HTML version}。
  \item \textbf{论文配图：} \href{https://arxiv.org/html/2601.00423v1}{arXiv HTML figures}，已用于本文的 entropy analysis、pipeline、reward ambiguity、training curve、visual comparison 和 failure cases 图示。
  \item \textbf{代码仓库：} \href{https://github.com/shengjun-zhang/VisualGRPO}{shengjun-zhang/VisualGRPO}。
\end{itemize}

% ============================================================
% Paper deep reading: MaskFocus
% Source: https://arxiv.org/abs/2512.18766
% ============================================================

\section{论文精读：MaskFocus}

\begin{conceptbox}
\textbf{论文：} \href{https://arxiv.org/abs/2512.18766}{MaskFocus: Focusing Policy Optimization on Critical Steps for Masked Image Generation}\\
\textbf{作者：} Guohui Zhang, Hu Yu, Xiaoxiao Ma, Yaning Pan, Hang Xu, Feng Zhao。\\
\textbf{简称：} MaskFocus。\\
\textbf{版本：} arXiv 2512.18766v1，submitted on 21 Dec 2025。\\
\textbf{代码：} 论文页标注为 \href{https://github.com/zghhui/MaskFocus}{zghhui/MaskFocus}，但我阅读时仓库接口未返回有效公开信息，因此这里只把它作为论文给出的代码入口。\\
\textbf{方向标签：} Masked Generative Models、Meissonic、GRPO、Critical Step Selection、Information Gain、Dynamic Routing Sampling。\\
\textbf{阅读日期：} 2026-05-06。\\
\textbf{阅读口径：} 这篇和 E-GRPO 都在问 ``哪些 step 值得优化''，但对象不同：E-GRPO 面向 flow / diffusion 的 denoising entropy，MaskFocus 面向 masked image generation 的 iterative unmasking trajectory。
\end{conceptbox}

\subsection{0. 一句话结论}

MaskFocus 的核心判断是：masked generative model 的每个 sampling step 对最终图像的贡献并不相同，所以 RL 不应该对整条 unmasking trajectory 均匀做 policy optimization，也不应该随机挑 step 训练；更合理的做法是先用 intermediate image 和 final image 的 embedding similarity 估计 step-level information gain，再只在最关键的 steps 上做 GRPO。

\[
\begin{aligned}
\text{MaskFocus}
&= \text{Critical Step Selection (CSS)}\\
&\quad + \text{all-masked-token probability estimation}\\
&\quad + \text{Dynamic Routing Sampling (DR-Sampling)}.
\end{aligned}
\]

\begin{summarybox}
\textbf{我的判断：} MaskFocus 是 masked generative model 版的 ``focused credit assignment''。它不是像 DenseGRPO 那样把 credit 分到所有 steps，而是先判断哪些 unmasking steps 真正改变了最终图像，再把 RL 更新预算集中到这些 critical steps 上。它和 E-GRPO 的共同点是减少无效 step 更新；差异是 E-GRPO 用 entropy 作为探索价值指标，MaskFocus 用 image embedding similarity change 作为信息增益指标。
\end{summarybox}

\subsection{1. 为什么 masked generative model 的 GRPO 不一样}

Masked Generative Models，简称 MGM，例如 MaskGIT / Meissonic，并不是 diffusion 那种从噪声连续去噪，也不是 AR 模型那种从左到右逐 token 生成。它的基本机制是：

\begin{enumerate}[leftmargin=1.5em]
  \item 初始图像 token 全部是 \([\mathrm{MASK}]\)。
  \item 每一步并行预测所有 masked tokens 的分布。
  \item 按 confidence / mask schedule 只保留一部分 token。
  \item 剩余 token 继续保持 masked，留到后续步骤 refinement。
\end{enumerate}

预训练目标可以写成：

\[
p(z_M\mid z_V)
=
\prod_{i\in M}p(z_i\mid z_V),
\]

其中 \(z_M\) 是 masked tokens，\(z_V\) 是 visible tokens。

这给 RL 后训练带来一个特殊问题：MGM 的 sampling 是多步 iterative refinement，每一步都有 mask、visible tokens、masked tokens、token probability；如果像 Mask-GRPO 那样把整个 unmasking trajectory 都纳入 policy optimization，就要沿整条轨迹估计 step likelihood，成本很高。

\begin{warnbox}
\textbf{关键矛盾：} MGM 的每个 step 都有 probability，可以做 policy gradient；但不是每个 step 都同等重要。全轨迹优化会很贵，随机挑 step 又可能挑到低价值 refinement step，导致训练不稳定或收益不足。
\end{warnbox}

MaskFocus 因此提出一个更直接的问题：

\begin{quote}
在 MGM 的 iterative unmasking 过程中，哪些 steps 真正决定最终图像的结构、主体和质量？能不能只在这些 steps 上做 RL？
\end{quote}

\subsection{2. 动机：sampling trajectory 的贡献是不均匀的}

论文观察到，MGM 的 early steps 往往快速建立图像的整体 appearance 和 structure，而 later steps 更多是在补局部细节。这个现象和 diffusion / flow 中 ``前期决定语义、后期修细节'' 的直觉很像，但 MGM 的载体不是连续 latent，而是 masked visual tokens。

为了量化这个不均匀性，论文记录每个 step 的 intermediate image embedding \(E_t\)，并和最终图像 embedding \(E_T\) 比较相似度。若某一步之后 similarity 变化很大，就说明这个 step 对最终图像的形成贡献更大。

\begin{figure}[htbp]
\centering
\includegraphics[width=0.95\linewidth]{maskfocus/01_motivation.png}
\caption{MaskFocus 的动机图：early steps 的 masked tokens 会更强地决定图像外观和结构；不同 step 到最终图像 embedding 的 cosine similarity 变化不均匀；不同样本的 entropy trajectory 也不同，低熵样本更确定、探索不足。图源：\href{https://arxiv.org/html/2512.18766v1}{MaskFocus arXiv HTML Figure 1}。}
\label{fig:maskfocus-motivation}
\end{figure}

\begin{intubox}
\textbf{直觉版：} 如果某一步之后 intermediate image suddenly becomes much closer to the final image，那么这一步很可能就是一个 ``结构决定点''。这种 step 比单纯的局部补纹理 step 更适合承载 reward credit。
\end{intubox}

\subsection{3. 方法一：Critical Step Selection (CSS)}

\subsubsection{3.1 用 final-image similarity change 定义 information gain}

MaskFocus 首先计算每个 sampling step 的 intermediate image embedding \(E_t\) 与 final image embedding \(E_T\) 的 cosine similarity：

\[
S_t
=
\mathrm{CosSim}(E_t,E_T).
\]

然后用相邻 step 的 similarity 变化来定义 step-level information gain：

\[
V_t
=
\left|\Delta S_t\right|
=
\left|S_{t+1}-S_t\right|.
\]

直觉上，\(V_t\) 越大，说明从 \(t\) 到 \(t+1\) 的 unmasking / refinement 让图像在语义 embedding 空间里发生了更明显的变化，因此这个 step 对最终图像更关键。

接着，MaskFocus 从所有 steps 中选出 information gain 最大的 \(K\) 个 critical steps：

\[
\mathcal{C}
=
\operatorname{TopK}_{t\in\{1,\ldots,T-1\}} V_t.
\]

论文实验里使用的 backbone 是 Meissonic，sampling steps 为 \(T=64\)，critical steps 数量为 \(K=6\)。也就是说，它不是全 64 步都训练，而是把 RL 更新集中到少数最有信息增益的位置。

\begin{summarybox}
\textbf{这一步和 E-GRPO 的差别：} E-GRPO 的 step selection 主要来自 SDE entropy，回答 ``哪里有探索空间''；MaskFocus 的 CSS 来自 embedding similarity change，回答 ``哪里真正改变了最终图像''。前者偏 exploration informativeness，后者偏 generation trajectory informativeness。
\end{summarybox}

\subsubsection{3.2 为什么不是随机挑 early steps}

一个看似合理的 baseline 是：既然 early steps 重要，那随机从前 40\% steps 里挑一些做训练即可。论文的讨论显示这种做法仍然不稳，因为 early stage 内部也不是均匀重要；有些 early step 只是 minor refinement，有些才是真正的结构转折点。

更严重的是，随机 early-step selection 可能在后期训练中诱发 reward hacking：模型会过拟合某些容易提升 reward 的微小细节，而不是稳定改善图像结构和指令遵循。

\begin{figure}[htbp]
\centering
\includegraphics[width=0.95\linewidth]{maskfocus/04_more_discussions.png}
\caption{论文对 step selection、sampling strategy、CFG 和 mask strategy 的补充比较。重点是：固定区间或随机选择 early steps 不如基于 information gain 的 CSS 稳定。图源：\href{https://arxiv.org/html/2512.18766v1}{MaskFocus arXiv HTML Figure 4}。}
\label{fig:maskfocus-more-discussions}
\end{figure}

\subsection{4. 方法二：critical step 上如何估计 probability}

只选出 critical steps 还不够，还要解决一个 MGM 特有的问题：某个 step 实际上只会 unmask 一部分高 confidence tokens，但模型在该 step 对所有 masked tokens 都有预测分布。

MaskFocus 的做法是：对 critical step \(k\) 记录生成结果 \(o_k\)、step index \(k\) 和该 step 的 mask \(M_k\)。训练时随机 shuffle mask 得到新的 \(M'_k\)，再重新构造 masked completion \(\hat{o}_i\)，让模型预测这些 masked tokens 的 probabilities，并计算 \(\pi_\theta\)、\(\pi_{\mathrm{old}}\)、\(\pi_{\mathrm{ref}}\) 的 log-likelihood。

它的关键点是：

\begin{quote}
对每个 critical step，不只使用该步最终采出来的 token，而是把所有 masked tokens 的 probability 都纳入优化目标。
\end{quote}

这样做的理由是 MGM 的预训练目标本来就是并行预测所有 masked tokens；即使某些 token 在当前 step 没有被选中，它们的 probability distribution 仍然包含关于最终图像的有效估计。

\begin{intubox}
\textbf{直觉：} MGM 在每一步其实都 ``想象'' 了整张图，只是采样策略只兑现其中一部分 token。MaskFocus 利用这个特性，把 critical step 的全 masked-token prediction 变成更直接的 probability estimation。
\end{intubox}

\subsection{5. 方法三：Dynamic Routing Sampling (DR-Sampling)}

\subsubsection{5.1 为什么需要 dynamic routing}

MGM 常用 confidence-based sampling：优先保留高 confidence tokens。这个策略稳定，但容易导致早期大量生成简单区域，例如背景，而主体、复杂属性和细节的探索不足。

论文认为这和 sample-level entropy 有关。低熵 trajectory 表示模型过于确定，采样路径更 deterministic，探索新 mask / token 的机会少；高熵 trajectory 已经不确定，如果再强行提高 temperature，反而容易引入不稳定甚至 reward hacking。

所以 MaskFocus 使用 group 内动态路由：

\begin{itemize}[leftmargin=1.5em]
  \item \textbf{高熵样本：} 已经有足够不确定性，走 exploitation branch，使用标准 confidence-based MaskGIT sampling。
  \item \textbf{低熵样本：} 过于确定，走 exploration branch，使用 entropy-based dynamic temperature modulation，温和增加探索。
\end{itemize}

\subsubsection{5.2 entropy 与 temperature 公式}

论文先计算 sample / token 预测分布的 entropy：

\[
H_i
=
-
\sum_{v\in\mathcal{V}}p(v)\log p(v),
\]

其中 \(\mathcal{V}\) 是 codebook，\(p(v)\) 是 token 的预测概率。

对低熵样本，论文使用 entropy-based temperature：

\[
T_i
=
T\exp\left(-\frac{H_{i,j}}{\alpha}\right)+\theta,
\]

其中 \(T\) 是最大 temperature，\(\theta\) 是 lower bound，\(\alpha\) 控制 entropy 增大时 temperature 衰减的速度。这个公式的效果是：entropy 越低，temperature 越高；entropy 越高，temperature 越接近 lower bound。

\begin{figure}[htbp]
\centering
\includegraphics[width=0.95\linewidth]{maskfocus/02_overview.png}
\caption{MaskFocus 方法总览：先用 DR-Sampling 在组内平衡 exploration / exploitation，再通过 CSS 找出 information gain 最高的 critical steps，最后对这些 steps 的 masks 做 shuffle 与 re-mask，计算 masked-token likelihood 并更新 GRPO 目标。图源：\href{https://arxiv.org/html/2512.18766v1}{MaskFocus arXiv HTML Figure 2}。}
\label{fig:maskfocus-overview}
\end{figure}

\begin{warnbox}
\textbf{这里的 routing 很容易误解：} 它不是对高熵样本继续加噪，而是只对低熵样本增强探索。原因是高熵样本已经不稳定，继续升温会破坏训练；低熵样本才是真的 ``太确定、探索不足''。
\end{warnbox}

\subsection{6. MaskFocus 的 GRPO 目标怎么理解}

论文仍然使用 GRPO / PPO 风格的 clipped objective。对同一个 prompt 采样 \(G\) 个 completions，reward 为 \(R^i\)，group-normalized advantage 为：

\[
A^i
=
\frac{
R^i-\mathrm{mean}(\{R^i\}_{i=1}^{G})
}{
\mathrm{std}(\{R^i\}_{i=1}^{G})
}.
\]

importance ratio 写作：

\[
r_t^i(\theta)
=
\frac{
\pi_\theta(o_t^i\mid q,o_{<t}^i)
}{
\pi_{\theta_{\mathrm{old}}}(o_t^i\mid q,o_{<t}^i)
}.
\]

然后 objective 仍然是 clipped surrogate 加 KL penalty：

\[
\begin{aligned}
\mathcal{J}
&=
\mathbb{E}_{\{o^i\}\sim\pi_{\mathrm{old}}}
\Bigg[
\frac{1}{\sum_i |o^i|}
\sum_i\sum_t
\Bigg(
\min\Big(
 r_t^i(\theta)A^i,
\mathrm{clip}(r_t^i(\theta),1-\epsilon,1+\epsilon)A^i
\Big)\\
&\qquad\qquad\qquad\qquad
-
\beta D_{\mathrm{KL}}(\pi_\theta\Vert\pi_{\mathrm{ref}})
\Bigg)
\Bigg].
\end{aligned}
\]

但和普通全轨迹 Mask-GRPO 不同的是，MaskFocus 并不把所有 \(t\) 都当成同等训练对象，而是在数据 buffer 中只保留 CSS 选出的 critical steps：

\[
D
\leftarrow
\{(o_k,k,M_k): k\in\mathcal{C}\}.
\]

\begin{summarybox}
\textbf{所以它的本质不是换 objective，而是换训练样本分布：} GRPO 的数学外壳仍然熟悉，但被优化的 timestep / mask event 变成了 information-gain-filtered critical steps。
\end{summarybox}

\subsection{7. 算法流程：训练时到底怎么跑}

可以把 Algorithm 1 压缩成下面几步：

\begin{enumerate}[leftmargin=1.5em]
  \item 初始化 \(\pi_\theta\leftarrow\pi_{\mathrm{ref}}\)，准备 critical-step buffer \(C\) 和 data buffer \(D\)。
  \item 对每个 prompt，用 \(\pi_{\mathrm{old}}\) 采样 \(G\) 个 completions；采样过程中使用 DR-Sampling，并记录每步 image embeddings \(\{E_t\}_{1:T}\) 与 masks \(\{M_t\}_{1:T}\)。
  \item 对每个 completion 计算 reward、group-normalized advantage，以及每个 step 的 information gain \(V_t\)。
  \item 用 CSS 选出 \(K\) 个 critical steps，并把 \((o_k,k,M_k)\) 放入训练 buffer。
  \item 对 mini-batch 中的 critical-step samples，shuffle mask、re-mask completion，重新计算 \(\pi_\theta\)、\(\pi_{\mathrm{old}}\)、\(\pi_{\mathrm{ref}}\) 的 likelihood。
  \item 用 clipped GRPO objective + KL penalty 更新 \(\pi_\theta\)。
\end{enumerate}

\begin{intubox}
\textbf{我会这样记：} DR-Sampling 负责在 rollout 时让低熵样本多探索；CSS 负责在 update 时只保留最有信息增益的 steps；all-masked-token likelihood 负责让 critical step 的 probability estimation 更贴合 MGM 的并行预测机制。
\end{intubox}

\subsection{8. 实验设置}

论文实验主要围绕 Meissonic 这个 masked generative model：

\begin{itemize}[leftmargin=1.5em]
  \item \textbf{Backbone：} Meissonic，开源 MGM，输出分辨率为 \(1024\times1024\)。
  \item \textbf{Sampling steps：} \(T=64\)。
  \item \textbf{Critical steps：} 默认 \(K=6\)。
  \item \textbf{CFG：} training 和 inference 都设为 5。
  \item \textbf{Compositional generation 数据：} 50K GenEval-style prompts，来自 Flow-GRPO，且不与官方 GenEval prompts 重叠。
  \item \textbf{Preference alignment 数据：} HPS v2 的 10K training samples。
  \item \textbf{Reward：} compositional task 使用 CLIP + GenEval reward；preference alignment task 使用 HPS。
  \item \textbf{评测：} GenEval、T2I-CompBench、DrawBench；指标包括 DEQA、PickScore、HPS、ImageReward。
\end{itemize}

\subsection{9. 主要结果：MaskFocus 的收益在哪里}

\subsubsection{9.1 GenEval：比 MaskGRPO 更强的组合生成能力}

在 GenEval 上，MaskFocus 把 Meissonic 的 overall score 从 \textbf{0.54} 提升到 \textbf{0.76}，也高于 MaskGRPO 的 \textbf{0.73}。

\begin{center}
\resizebox{\linewidth}{!}{
\begin{tabular}{lccccccc}
\toprule
\textbf{方法} & \textbf{Overall} & \textbf{Single Obj.} & \textbf{Two Obj.} & \textbf{Counting} & \textbf{Color} & \textbf{Position} & \textbf{Color Attr.}\\
\midrule
FLUX.1-dev & 0.66 & 0.98 & 0.81 & 0.74 & 0.79 & 0.22 & 0.45\\
Janus-Pro-1B & 0.73 & 0.98 & 0.82 & 0.51 & 0.89 & 0.65 & 0.56\\
Mask-GRPO & 0.73 & 0.99 & 0.90 & 0.69 & 0.85 & 0.35 & 0.59\\
Meissonic & 0.54 & 0.99 & 0.66 & 0.42 & 0.86 & 0.10 & 0.22\\
Meissonic + MaskGRPO & 0.73 & 1.00 & 0.87 & 0.83 & 0.87 & 0.39 & 0.48\\
Meissonic + MaskFocus & \textbf{0.76} & \textbf{1.00} & \textbf{0.91} & \textbf{0.85} & \textbf{0.87} & \textbf{0.42} & \textbf{0.54}\\
\bottomrule
\end{tabular}
}
\end{center}

这张表最值得关注的是 position 和 color attribute。它们通常更依赖结构化关系和属性绑定，而不是单纯的图像美观。MaskFocus 在这些子项上相对 Meissonic 和 MaskGRPO 都有提升，说明 critical-step optimization 确实更有利于指令遵循和组合关系。

\subsubsection{9.2 Human preference：接近或超过强 diffusion baseline}

在人类偏好相关指标上，MaskFocus 也提升了 Meissonic：

\begin{center}
\resizebox{\linewidth}{!}{
\begin{tabular}{lcccc}
\toprule
\textbf{方法} & \textbf{DEQA} & \textbf{PickScore} & \textbf{HPS} & \textbf{ImageReward}\\
\midrule
SDXL & 4.17 & 22.27 & 30.06 & 0.56\\
SD3.5-M & 4.24 & 22.50 & 30.17 & 0.98\\
FLUX.1-dev & 4.37 & \textbf{22.97} & 31.13 & 1.06\\
Meissonic & 4.00 & 21.63 & 28.89 & 0.39\\
Meissonic + MaskGRPO & 4.35 & 22.34 & 35.48 & 1.06\\
Meissonic + MaskFocus & \textbf{4.39} & 22.39 & \textbf{35.52} & \textbf{1.09}\\
\bottomrule
\end{tabular}
}
\end{center}

需要注意的是，MaskFocus 并不是所有指标都超过 FLUX.1-dev；例如 PickScore 仍低于 FLUX.1-dev。但它把 Meissonic 的 HPS 和 ImageReward 明显推高，说明 MGM 经过 focused RL 后训练之后，偏好指标可以接近强 diffusion model。

\begin{figure}[htbp]
\centering
\includegraphics[width=0.95\linewidth]{maskfocus/03_qualitative.png}
\caption{MaskFocus 的 qualitative comparison。论文展示了其在图像质量、人类偏好，以及 counting、color、attribute binding、position 等指令遵循任务上的可视化提升。图源：\href{https://arxiv.org/html/2512.18766v1}{MaskFocus arXiv HTML Figure 3}。}
\label{fig:maskfocus-qualitative}
\end{figure}

\subsubsection{9.3 Ablation：CSS 和 DR-Sampling 都有贡献}

论文消融如下：

\begin{center}
\resizebox{0.78\linewidth}{!}{
\begin{tabular}{lccc}
\toprule
\textbf{方法} & \textbf{GenEval} & \textbf{DEQA} & \textbf{PickScore}\\
\midrule
MaskFocus & \textbf{0.76} & \textbf{4.39} & \textbf{22.39}\\
w/o Critical Step Selection & 0.72 & 4.34 & 22.34\\
w/o Dynamic Routing Sampling & 0.74 & 4.35 & 22.31\\
\bottomrule
\end{tabular}
}
\end{center}

去掉 CSS 后，GenEval 从 0.76 降到 0.72，说明 information-gain-based critical steps 不是装饰项；去掉 DR-Sampling 后，GenEval 和偏好指标也下降，说明低熵样本的探索路由对训练同样有帮助。

\subsection{10. Pre-/Post-RL trajectory：它到底改变了什么}

论文还可视化了 RL 前后的 sampling trajectories。这个图的意义不是单看最终图，而是看 MGM 在每个 unmasking step 中如何逐渐形成图像。

\begin{figure}[htbp]
\centering
\includegraphics[width=0.95\linewidth]{maskfocus/05_sampling_trajectories.png}
\caption{Pre-/Post-RL sampling trajectories。MaskFocus 不只是提高最终 reward，而是改变了 masked generation 的中间演化路径，使关键 unmasking steps 更有利于最终图像质量和指令对齐。图源：\href{https://arxiv.org/html/2512.18766v1}{MaskFocus arXiv HTML Figure 5}。}
\label{fig:maskfocus-trajectories}
\end{figure}

\begin{summarybox}
\textbf{我的理解：} MaskFocus 的真正目标不是 ``让 64-step MGM 采样更少步''，而是 ``让 64-step 采样里被用于 RL 更新的 steps 更少、更准''。它主要省的是 policy optimization 的轨迹 likelihood / update 成本，而不是 necessarily 把 inference steps 从 64 降到 6。
\end{summarybox}

\subsection{11. 和 E-GRPO / MaskGRPO / DenseGRPO 的关系}

\begin{center}
\resizebox{\linewidth}{!}{
\begin{tabular}{llll}
\toprule
\textbf{方法} & \textbf{生成范式} & \textbf{核心选择信号} & \textbf{解决的问题}\\
\midrule
DenseGRPO & diffusion / flow & dense reward gain & terminal reward 太粗，需要细粒度 credit\\
E-GRPO & flow model & step entropy & 低熵 denoising steps 探索价值低\\
MaskGRPO & masked generation & full / mask-ratio-related trajectory & MGM 如何做 stepwise policy optimization\\
MaskFocus & masked generation & embedding similarity information gain & MGM 全轨迹优化太贵，随机 step 不稳定\\
\bottomrule
\end{tabular}
}
\end{center}

\begin{itemize}[leftmargin=1.5em]
  \item \textbf{和 E-GRPO：} 两者都不是简单追求更细 credit，而是减少无效 step。E-GRPO 通过 entropy-aware step merging 选择值得探索的 denoising steps；MaskFocus 通过 embedding similarity change 选择真正改变最终图像的 unmasking steps。
  \item \textbf{和 DenseGRPO：} DenseGRPO 倾向于把 reward 更细地分给每个 step；MaskFocus 则认为并不是所有 steps 都值得分 credit，应该先筛 critical steps。
  \item \textbf{和 MaskGRPO：} MaskGRPO 证明 MGM 可以做 GRPO，但全轨迹优化成本高；MaskFocus 在此基础上加入 CSS 和 DR-Sampling，把 update 集中到更有价值的 steps。
\end{itemize}

\begin{intubox}
\textbf{和我们前面讨论的 ``10 个关键点'' idea 的关系：} MaskFocus 不是先完整跑 64 步再为了 inference 省步数，而是在 rollout 过程中记录 trajectory，再只把 high-information-gain steps 放进 RL update。它更像 ``update-efficient critical step selection''，不是 ``pre-scheduled fast sampler''。
\end{intubox}

\subsection{12. 我认为它的局限}

\begin{enumerate}[leftmargin=1.5em]
  \item \textbf{它主要提升训练 / update 效率，不等于减少采样步数。} CSS 需要 intermediate embeddings 和 final image embedding，因此默认仍然要完成 rollout 后才能判断 critical steps。它不是直接把 64-step inference 变成 6-step inference。
  \item \textbf{information gain 是后验指标。} \(E_t\) 和 \(E_T\) 的 similarity change 很适合解释哪一步重要，但 \(E_T\) 只有采样结束后才知道。若要进一步做采样加速，需要把 CSS 蒸馏成前验 schedule 或低成本 selector。
  \item \textbf{embedding similarity 不等于 reward causality。} 某步让 image embedding 变化大，不一定说明它对 HPS / GenEval reward 的因果贡献最大；这仍然是一个 proxy。
  \item \textbf{方法强依赖 MGM 的并行 masked-token prediction 结构。} all-masked-token probability estimation 很适合 MGM，但不一定能直接迁移到 flow / diffusion。
  \item \textbf{代码可访问性需要再确认。} arXiv 页面给出 GitHub 链接，但我阅读时仓库元信息未能通过公开接口获取；复现实验时需要以论文发布后的实际仓库状态为准。
\end{enumerate}

\subsection{13. 读完后的记忆钩子}

\begin{summarybox}
\textbf{我最后会这样概括：} DenseGRPO 说 ``细粒度 credit 有用''；E-GRPO 说 ``低熵 denoising step 不值得都随机探索''；MaskFocus 说 ``MGM 的 unmasking steps 也不是都值得做 RL，应该只优化 information gain 最大的 critical steps''。它是 masked generation 方向上把 RL update 从 full trajectory 推向 critical-step trajectory 的代表工作。
\end{summarybox}

\subsection{参考来源}

\begin{itemize}[leftmargin=1.5em]
  \item \textbf{论文主页：} \href{https://arxiv.org/abs/2512.18766}{MaskFocus: Focusing Policy Optimization on Critical Steps for Masked Image Generation}。
  \item \textbf{arXiv HTML：} \href{https://arxiv.org/html/2512.18766v1}{MaskFocus HTML version}，用于核对章节、图表、公式和实验结果。
  \item \textbf{PDF：} \href{https://arxiv.org/pdf/2512.18766}{arXiv PDF}，用于核对方法公式、Algorithm 1 和实验设置。
  \item \textbf{代码入口：} \href{https://github.com/zghhui/MaskFocus}{zghhui/MaskFocus}，来源于 arXiv 页面 comments；当前笔记未将其视为已验证可复现代码。
\end{itemize}

% ============================================================
% Paper deep reading: LeapAlign
% Source: https://arxiv.org/abs/2604.15311
% ============================================================

\section{论文精读：LeapAlign}

\begin{conceptbox}
\textbf{论文：} \href{https://arxiv.org/abs/2604.15311}{LeapAlign: Post-Training Flow Matching Models at Any Generation Step by Building Two-Step Trajectories}\\
\textbf{作者：} Zhanhao Liang, Tao Yang, Jie Wu, Chengjian Feng, Liang Zheng。\\
\textbf{简称：} LeapAlign。\\
\textbf{版本：} arXiv 2604.15311v2，CVPR 2026。\\
\textbf{项目主页：} \href{https://rockeycoss.github.io/leapalign/}{LeapAlign project page}。\\
\textbf{代码：} \href{https://github.com/RockeyCoss/LeapAlign_Code}{RockeyCoss/LeapAlign\_Code}。\\
\textbf{方向标签：} Flow Matching、Rectified Flow、Direct Reward Backpropagation、Preference Alignment、Two-step Leap Trajectory、Gradient Discounting。\\
\textbf{阅读日期：} 2026-05-06。\\
\textbf{阅读口径：} 这篇不是 GRPO / policy-gradient 路线，而是 flow matching 的 \textbf{direct-gradient post-training} 路线：它问的是 ``能不能让 reward gradient 稳定传到任意生成步骤，尤其是早期构图步骤''。
\end{conceptbox}

\subsection{0. 一句话结论}

LeapAlign 的核心判断是：Flow Matching / Rectified Flow 的采样轨迹本身是连续可微的，所以偏好对齐不一定非要把生成过程改造成 stochastic policy 再做 GRPO；如果 reward model 可微，可以直接把 reward gradient 反传回 flow trajectory。

问题是，完整多步采样轨迹太长，直接反传显存开销大、梯度也容易爆炸；而只训最后几步又会错过最决定全局布局的早期步骤。LeapAlign 的解法是：

\[
\begin{aligned}
\text{LeapAlign}
&= \text{two-step leap trajectory}\\
&\quad + \text{latent connector}\\
&\quad + \text{gradient discounting}\\
&\quad + \text{trajectory-similarity weighting}.
\end{aligned}
\]

\begin{summarybox}
\textbf{我的判断：} LeapAlign 是 Flow 后训练里和 Flow-GRPO 很不一样的一条路线。Flow-GRPO / OP-GRPO / E-GRPO 关心的是如何构造 policy、logprob、advantage、ratio 和更高效的 rollout；LeapAlign 则绕开 logprob，直接利用 flow trajectory 可微性，把长反传路径压缩成两跳，让 reward gradient 能够到达任意 timestep。
\end{summarybox}

\begin{figure}[htbp]
\centering
\includegraphics[width=0.95\linewidth]{leapalign/01_performance_overview.png}
\caption{LeapAlign 的性能概览：在 compositional alignment 训练曲线、多种偏好评测指标和 GenEval 上，LeapAlign 都相对 Flux base 和 direct-gradient / GRPO baseline 有提升。图源：\href{https://arxiv.org/html/2604.15311v2}{LeapAlign arXiv HTML Figure 1}。}
\label{fig:leapalign-performance-overview}
\end{figure}

\subsection{1. 它为什么重要：Flow 后训练不只有 GRPO}

在前面的 Flow-GRPO 主线里，我们一直在处理一个问题：flow matching 默认是 deterministic ODE，没有天然的 step-level logprob，所以要先把 ODE 改写成 SDE，再把每一步 transition 看成 Gaussian policy，最后才能做 GRPO / PPO-style ratio update。

LeapAlign 站在另一个角度：既然 flow matching 的采样过程是连续可微的，那么如果 reward model 也是可微的，是否可以直接优化最终图像 reward？

\[
\max_{\theta} r(x_0),
\]

其中 \(x_0\) 是最终生成图像 latent，\(r(\cdot)\) 是偏好 / 对齐 reward。

这条路线有明显吸引力：

\begin{itemize}[leftmargin=1.5em]
  \item \textbf{不需要 logprob：} 不必把 flow ODE 强行 SDE 化，也不需要计算 step transition ratio。
  \item \textbf{不需要 group sampling：} 理论上不必为同一 prompt 采样一组候选来估计 advantage。
  \item \textbf{信号更直接：} reward gradient 直接告诉模型 ``怎样改图像能提高 reward''。
\end{itemize}

但直接梯度路线的瓶颈也很明确：完整采样链可能有几十步。如果让 reward gradient 穿过完整链条：

\[
x_1 \rightarrow x_{t_{T-1}} \rightarrow \cdots \rightarrow x_{t_1} \rightarrow x_0,
\]

就要保存大量中间激活，显存开销很高；更麻烦的是长链式反传会让梯度数值不稳定。

\begin{warnbox}
\textbf{核心矛盾：} 早期 timestep 往往决定整体布局、主体关系和空间结构，但直接反传到早期最贵、也最不稳定；后期 timestep 容易训练，却更多是在修局部纹理和细节。
\end{warnbox}

所以 LeapAlign 真正要解决的问题是：

\begin{quote}
如何在不反传完整长轨迹的情况下，让最终 reward 的梯度稳定影响任意 generation step？
\end{quote}

\subsection{2. Flow Matching 基础：为什么可以 ``跳步''}

论文采用的是 rectified flow / flow matching 视角。设 \(x_1\sim\mathcal{N}(0,I)\) 是噪声端，\(x_0\sim p_{\mathrm{data}}\) 是数据端，时间 \(t\in[0,1]\)，其中 \(t=1\) 对应噪声，\(t=0\) 对应干净图像。

Rectified flow 的线性插值可以写成：

\[
x_t = (1-t)x_0 + t x_1.
\]

模型学习速度场 \(v_{\theta}(x_t,t,c)\)，采样时从 \(x_1\) 沿着 ODE 走到 \(x_0\)：

\[
\frac{dx_t}{dt}=v_{\theta}(x_t,t,c).
\]

离散采样时，普通 Euler 步是一步一步走：

\[
x_{t-\Delta t}
=
x_t - \Delta t\, v_{\theta}(x_t,t,c).
\]

LeapAlign 利用的关键性质是：如果速度场足够好，那么可以从较早的 \(k\) 时刻直接一步预测较晚的 \(j\) 时刻：

\[
\hat{x}_{j\mid k}
=
x_k - (k-j)v_{\theta}(x_k,k,c),
\qquad k>j.
\]

同理，也可以从中间点 \(x_j\) 直接预测最终 clean latent：

\[
\hat{x}_{0\mid j}
=
x_j - j\,v_{\theta}(x_j,j,c).
\]

\begin{intubox}
\textbf{直觉：} 普通采样是 ``小步走完整条路''；LeapAlign 训练时构造的是 ``从早期点直接跳到中间点，再从中间点直接跳到终点''。它不声称推理一定少走这些步，而是用这种跳跃路径来缩短 reward gradient 的反传链条。
\end{intubox}

\subsection{3. 方法一：Two-step Leap Trajectory}

\subsubsection{3.1 从完整轨迹中随机挑两个点}

先照常在线采样一条完整 flow trajectory：

\[
x_1 \rightarrow \cdots \rightarrow x_k \rightarrow \cdots \rightarrow x_j \rightarrow \cdots \rightarrow x_0,
\qquad k>j>0.
\]

LeapAlign 不让 reward gradient 沿完整链条回传，而是在这条真实轨迹里随机选两个时间点 \(k,j\)，构造两跳：

\[
x_k
\Rightarrow
\hat{x}_{j\mid k}
\Rightarrow
x_j
\Rightarrow
\hat{x}_{0\mid j}
\Rightarrow
x_0.
\]

第一跳负责把早期点 \(x_k\) 连接到中间点 \(x_j\) 附近；第二跳负责把 \(x_j\) 连接到最终输出 \(x_0\) 附近。

\begin{figure}[htbp]
\centering
\includegraphics[width=0.95\linewidth]{leapalign/02_method_overview.png}
\caption{LeapAlign 方法总览：从完整生成轨迹中随机选择 \(k>j\)，构造 \(x_k\to\hat{x}_{j\mid k}\) 和 \(x_j\to\hat{x}_{0\mid j}\) 两个 leap，并用 latent connector 把跳跃估计接回真实轨迹点，使前向数值等于真实采样结果，反向梯度沿短路径传播。图源：\href{https://arxiv.org/html/2604.15311v2}{LeapAlign arXiv HTML Figure 2}。}
\label{fig:leapalign-method-overview}
\end{figure}

\subsubsection{3.2 为什么是 two-step，而不是 one-step / full-step}

One-step 只从某个中间点直接预测 \(x_0\)，反传路径更短，但只能覆盖一个训练点，对 step interaction 的建模能力较弱。Full trajectory 能保留所有依赖，但显存和梯度稳定性都很差。

Two-step 是中间折中：

\begin{itemize}[leftmargin=1.5em]
  \item 比 full trajectory 短得多，反传成本基本固定。
  \item 比 one-step 多保留一个中间交互点，可以让早期 step 和后续 step 的耦合关系进入梯度。
  \item 随机选择 \(k,j\) 后，长期训练中所有 timestep 都有机会被更新。
\end{itemize}

论文消融也显示，two-step 在性能和显存成本之间最好；three-step 更贵，但收益不明显。

\begin{figure}[htbp]
\centering
\includegraphics[width=0.70\linewidth]{leapalign/05_design_leap_steps.png}
\caption{不同 leap step 数量的消融：two-step 在效果和内存之间取得最好折中。图源：\href{https://arxiv.org/html/2604.15311v2}{LeapAlign arXiv HTML Figure 4(b)}。}
\label{fig:leapalign-leap-steps}
\end{figure}

\subsection{4. 方法二：Latent Connector}

如果只用 \(\hat{x}_{j\mid k}\) 和 \(\hat{x}_{0\mid j}\)，跳跃预测可能偏离真实长轨迹，导致 reward 其实对应的是一个近似图像，而不是模型真实采样出来的最终图像。

LeapAlign 因此引入 latent connector：

\[
x_j
=
\hat{x}_{j\mid k}
+
\operatorname{sg}\left(x_j-\hat{x}_{j\mid k}\right),
\]

\[
x_0
=
\hat{x}_{0\mid j}
+
\operatorname{sg}\left(x_0-\hat{x}_{0\mid j}\right),
\]

其中 \(\operatorname{sg}(\cdot)\) 表示 stop-gradient。

这个技巧的作用非常关键：

\begin{itemize}[leftmargin=1.5em]
  \item \textbf{前向数值：} 等式右边精确等于真实轨迹点 \(x_j\) 和 \(x_0\)。所以 reward 看的是完整采样得到的真实最终图像。
  \item \textbf{反向梯度：} stop-gradient 项不传梯度，梯度会把 \(x_j\) 和 \(x_0\) 当作由 leap prediction 产生。
\end{itemize}

\begin{summarybox}
\textbf{这一步是 LeapAlign 的工程核心：} 前向用真实长轨迹保证 reward input 可靠；反向用短 leap trajectory 保证梯度路径短且能覆盖早期步骤。
\end{summarybox}

\subsection{5. 方法三：Nested Gradient 与 Gradient Discounting}

\subsubsection{5.1 两步 leap 会产生跨步骤依赖梯度}

沿两步 leap trajectory 反传时，最终 \(x_0\) 对参数 \(\theta\) 的梯度不只是两个单步梯度，还会出现一个跨步骤交互项。论文将其展开为：

\[
\frac{\partial x_0}{\partial \theta}
=
- j\frac{\partial v_{\theta}(x_j)}{\partial \theta}
-
(k-j)\frac{\partial v_{\theta}(x_k)}{\partial \theta}
+
j(k-j)
\frac{\partial v_{\theta}(x_j)}{\partial x_j}
\frac{\partial v_{\theta}(x_k)}{\partial \theta}.
\]

前两项是比较直观的 single-step gradients：一个更新 \(j\) 点，一个更新 \(k\) 点。最后一项是 nested gradient，它表示 \(x_k\) 处的变化会先影响 \(x_j\)，再通过 \(x_j\) 处的速度场影响最终 \(x_0\)。

\begin{intubox}
\textbf{直觉：} nested gradient 捕捉的是 ``早期一步怎么改变后续一步的输入，从而影响最终图像''。如果把它完全删掉，就等于假设不同 generation steps 之间相互独立，这显然不符合 flow 生成过程。
\end{intubox}

\subsubsection{5.2 为什么不能直接完整保留 nested gradient}

Nested gradient 有用，但数值上也危险。它包含 \(\partial v_{\theta}(x_j)/\partial x_j\) 这样的 Jacobian 项，乘上 \(j(k-j)\) 后可能带来很大的梯度范数，导致训练不稳定。

过去一些 direct-gradient 方法倾向于直接截断这类跨步梯度。LeapAlign 的选择更温和：不删除，而是打折。

\subsubsection{5.3 Gradient Discounting 的写法}

LeapAlign 在第二跳速度场输入处使用：

\[
\hat{x}_{0\mid j}
=
x_j
-
j\,v_{\theta}\left(
\alpha x_j + (1-\alpha)\operatorname{sg}(x_j), j, c
\right),
\qquad \alpha\in[0,1].
\]

前向时，因为 \(\alpha x_j+(1-\alpha)x_j=x_j\)，所以数值不变；反向时，经过 \(x_j\) 的梯度会被乘上 \(\alpha\)。于是 nested gradient 变成：

\[
\alpha\,j(k-j)
\frac{\partial v_{\theta}(x_j)}{\partial x_j}
\frac{\partial v_{\theta}(x_k)}{\partial \theta}.
\]

三个极端很好理解：

\begin{itemize}[leftmargin=1.5em]
  \item \(\alpha=1\)：完整保留 nested gradient，信息最多，但最不稳定。
  \item \(\alpha=0\)：完全截断 nested gradient，稳定但丢掉跨步依赖。
  \item \(0<\alpha<1\)：保留跨步信号，但降低幅度。
\end{itemize}

论文主实验中，HPSv2.1 reward 下 \(\alpha=0.3\) 表现最好；附录中对 nested gradient 的单独分析也支持 ``适度折扣比完整保留或完全删除更好''。

\begin{figure}[htbp]
\centering
\includegraphics[width=0.78\linewidth]{leapalign/04_design_gradient_discount.png}
\caption{Gradient discounting 的主文消融：\(\alpha=0.3\) 在保留跨步骤信号和控制梯度幅度之间取得较好平衡。图源：\href{https://arxiv.org/html/2604.15311v2}{LeapAlign arXiv HTML Figure 4(a)}。}
\label{fig:leapalign-gradient-discount}
\end{figure}

\begin{figure}[htbp]
\centering
\includegraphics[width=0.78\linewidth]{leapalign/11_nested_gradient_analysis.png}
\caption{Nested gradient 的附录分析：直接使用完整 nested gradient 会显著增大梯度范数并损害性能；适度 discount 能降低梯度范数并改善结果。图源：\href{https://arxiv.org/html/2604.15311v2}{LeapAlign arXiv HTML Figure 6}。}
\label{fig:leapalign-nested-gradient}
\end{figure}

\subsection{6. 方法四：Reward Objective 与 Trajectory-Similarity Weighting}

\subsubsection{6.1 Reward 打在真实最终图像 \(x_0\) 上}

LeapAlign 不是把 reward 打在近似的一步预测 \(\hat{x}_{0\mid j}\) 上，而是打在完整 rollout 得到的真实最终图像 \(x_0\) 上。

这是一个重要选择。因为 \(\hat{x}_{0\mid j}\) 只是 one-step approximation，可能存在模糊、伪影或轨迹偏差；如果 reward model 在这种近似图像上打分，信号可能不可靠。真实 \(x_0\) 才是模型实际会输出的图像。

论文使用 hinge-style reward objective：

\[
\mathcal{L}_{\mathrm{raw}}
=
\max(0,\lambda-r(x_0)).
\]

如果 reward 已经超过阈值 \(\lambda\)，loss 为 0；如果 reward 不够高，就继续推动模型提升 reward。这个阈值式写法比无上限地最大化 reward 更克制，可以缓解过度优化和 reward hacking 风险。

论文默认设置中，HPSv2.1 的阈值为 \(\lambda=0.55\)。附录阈值消融显示 \(\lambda=0.55\) 在 HPSv2.1、HPSv3、PickScore 和 ImageReward 上整体最好。

\subsubsection{6.2 不是每条 leap trajectory 都同样可信}

如果两步 leap 预测与真实长轨迹差距很大，那么沿这条短路径反传的梯度可能会误导训练。LeapAlign 因此计算 leap approximation 和真实轨迹点之间的偏差：

\[
d_j
=
\operatorname{mean}\left(|x_j-\hat{x}_{j\mid k}|\right),
\]

\[
d_0
=
\operatorname{mean}\left(|x_0-\hat{x}_{0\mid j}|\right).
\]

然后定义 trajectory-similarity weight：

\[
w_{\mathrm{sim}}
=
\frac{1}{\max(d_j,\tau)+\max(d_0,\tau)},
\]

最终损失为：

\[
\mathcal{L}
=
\operatorname{sg}(w_{\mathrm{sim}})\,\mathcal{L}_{\mathrm{raw}}.
\]

其中 \(\tau\) 是下限，论文经验值为 \(0.1\)，用于防止权重无限变大。

\begin{summarybox}
\textbf{直觉：} 如果 two-step leap 跳得很像真实长轨迹，就更信任这条短路径传回来的梯度；如果跳得很偏，就降低它对训练的影响。
\end{summarybox}

\begin{figure}[htbp]
\centering
\includegraphics[width=0.78\linewidth]{leapalign/06_design_reward_input.png}
\caption{Reward input 的消融：使用真实最终图像 \(x_0\) 作为 reward 输入优于使用 one-step approximation。图源：\href{https://arxiv.org/html/2604.15311v2}{LeapAlign arXiv HTML Figure 4(c)}。}
\label{fig:leapalign-reward-input}
\end{figure}

\begin{figure}[htbp]
\centering
\includegraphics[width=0.78\linewidth]{leapalign/07_design_similarity_weight.png}
\caption{Trajectory-similarity weighting 的消融：同时使用 \(d_j\) 和 \(d_0\) 衡量 leap 与真实轨迹的偏差，比只看单个偏差或不加权更好。图源：\href{https://arxiv.org/html/2604.15311v2}{LeapAlign arXiv HTML Figure 4(d)}。}
\label{fig:leapalign-similarity-weight}
\end{figure}

\subsection{7. 算法流程：训练时到底怎么跑}

把 LeapAlign 压缩成训练流程，可以写成下面几步：

\begin{enumerate}[leftmargin=1.5em]
  \item 用当前 flow model 对 prompt \(c\) 正常在线采样，得到完整 trajectory \(x_1,\ldots,x_0\)。
  \item 在完整 timestep 范围内随机选择 \(k>j>0\)。
  \item 在 \(x_k\) 处计算速度，构造 \(\hat{x}_{j\mid k}=x_k-(k-j)v_{\theta}(x_k,k,c)\)。
  \item 用 latent connector 将 \(\hat{x}_{j\mid k}\) 接回真实 \(x_j\)。
  \item 在 \(x_j\) 处用 gradient-discounted 输入计算速度，构造 \(\hat{x}_{0\mid j}=x_j-jv_{\theta}(\cdot,j,c)\)。
  \item 用 latent connector 将 \(\hat{x}_{0\mid j}\) 接回真实 \(x_0\)。
  \item 用真实 \(x_0\) 计算 reward hinge loss，并乘上 trajectory-similarity weight。
  \item 沿 two-step leap trajectory 反传更新模型。
\end{enumerate}

\begin{figure}[htbp]
\centering
\includegraphics[width=0.78\linewidth]{leapalign/08_design_timestep_range.png}
\caption{训练 timestep 范围的消融：在完整 \([0,1]\) 范围内随机选择 timestep 更好，说明早期步骤确实需要被覆盖。图源：\href{https://arxiv.org/html/2604.15311v2}{LeapAlign arXiv HTML Figure 4(e)}。}
\label{fig:leapalign-timestep-range}
\end{figure}

\begin{figure}[htbp]
\centering
\includegraphics[width=0.78\linewidth]{leapalign/09_design_kj_selection.png}
\caption{\(k,j\) 选择策略的消融：随机选择优于固定间距，说明多样化的 leap pair 对训练有帮助。图源：\href{https://arxiv.org/html/2604.15311v2}{LeapAlign arXiv HTML Figure 4(f)}。}
\label{fig:leapalign-kj-selection}
\end{figure}

\subsection{8. 实验设置}

论文主要在 FLUX.1-dev 上做实验，并补充验证到 SD3.5-M。

\begin{itemize}[leftmargin=1.5em]
  \item \textbf{主模型：} FLUX.1-dev。
  \item \textbf{补充模型：} Stable Diffusion 3.5 Medium。
  \item \textbf{默认 reward：} HPSv2.1。
  \item \textbf{额外 reward：} PickScore、HPSv3。
  \item \textbf{训练数据：} HPDv2 的 50K prompts；组合对齐使用 GenEval-style training prompts；额外也在 MJHQ-30K 上验证。
  \item \textbf{训练采样：} 720\(\times\)720，25 steps，CFG=3.5。
  \item \textbf{评测采样：} 720\(\times\)720，50 steps，CFG=3.5。
  \item \textbf{优化器：} AdamW，学习率 \(1\times10^{-5}\)，batch size 64，weight decay \(1\times10^{-4}\)，EMA 0.995。
  \item \textbf{训练轮数：} 300 iterations。
  \item \textbf{关键超参：} trajectory similarity 的 \(\tau=0.1\)，HPSv2.1 下 gradient discounting 默认 \(\alpha=0.3\)，hinge threshold \(\lambda=0.55\)。
\end{itemize}

\subsection{9. 主结果：LeapAlign 强在哪里}

\subsubsection{9.1 一般偏好对齐：HPSv2.1 训练也能提升 OOD 指标}

在 FLUX 上用 HPSv2.1 训练时，LeapAlign 在多个指标上都超过 direct-gradient baseline 和 policy-gradient baseline。核心数字如下：

\begin{center}
\resizebox{\linewidth}{!}{
\begin{tabular}{lcccccc}
\toprule
\textbf{方法} & \textbf{HPSv2.1} & \textbf{HPSv3} & \textbf{PickScore} & \textbf{UR-Align} & \textbf{UR-IQ} & \textbf{ImageReward}\\
\midrule
Flux & 0.3078 & 13.5020 & 22.7902 & 3.4514 & 3.5708 & 1.0455\\
DanceGRPO & 0.3451 & 14.8336 & 23.1186 & 3.4660 & 3.6199 & 1.2347\\
MixGRPO & 0.3692 & 14.7530 & 23.5184 & 3.4393 & 3.6241 & \textbf{1.6155}\\
DRTune & 0.3882 & 15.5606 & 23.5185 & 3.4793 & 3.6679 & 1.3562\\
LeapAlign & \textbf{0.4092} & \textbf{15.7678} & \textbf{23.7137} & \textbf{3.4984} & \textbf{3.7244} & 1.5104\\
\bottomrule
\end{tabular}
}
\end{center}

这里最值得注意的是：LeapAlign 只用 HPSv2.1 训练，但 HPSv3、PickScore、UnifiedReward 等 out-of-domain metrics 也同步提升。这说明它不只是简单刷训练 reward，至少在论文设置下有一定泛化收益。

\subsubsection{9.2 GenEval：组合生成受益于早期步骤更新}

GenEval 上，LeapAlign 的 overall score 达到 \textbf{0.7420}，高于 Flux base 的 0.6535，也高于 MixGRPO 的 0.7232 和 DRTune 的 0.7101。

\begin{center}
\resizebox{\linewidth}{!}{
\begin{tabular}{lccccccc}
\toprule
\textbf{方法} & \textbf{Overall} & \textbf{Single Obj.} & \textbf{Two Obj.} & \textbf{Count} & \textbf{Color} & \textbf{Position} & \textbf{AttrB}\\
\midrule
Flux & 0.6535 & 99.38 & 86.62 & 66.88 & 74.47 & 19.50 & 45.25\\
DanceGRPO & 0.6775 & 99.38 & 90.15 & 69.38 & 76.33 & 22.25 & 49.00\\
MixGRPO & 0.7232 & \textbf{99.69} & 93.69 & \textbf{80.00} & 80.05 & 24.25 & 56.25\\
DRTune & 0.7101 & 99.38 & 93.69 & 73.12 & 76.86 & 27.50 & 55.50\\
LeapAlign & \textbf{0.7420} & 99.38 & \textbf{96.46} & 72.50 & \textbf{80.59} & \textbf{30.25} & \textbf{66.00}\\
\bottomrule
\end{tabular}
}
\end{center}

这个结果和论文动机是对齐的：two objects、position、attribute binding 这些任务更依赖全局结构和对象关系；如果方法真的能更新早期构图步骤，就应该在这些子任务上更明显。

\begin{figure}[htbp]
\centering
\includegraphics[width=0.95\linewidth]{leapalign/03_qualitative_geneval.png}
\caption{LeapAlign 在 GenEval 上的定性比较。论文强调，ReFL / DRaFT-LV 这类偏后期步骤的方法更难改变整体布局，而 LeapAlign 可以更明显地修正全局 composition。图源：\href{https://arxiv.org/html/2604.15311v2}{LeapAlign arXiv HTML Figure 3}。}
\label{fig:leapalign-qualitative-geneval}
\end{figure}

\begin{figure}[htbp]
\centering
\includegraphics[width=0.82\linewidth]{leapalign/10_geneval_improvement_curve.png}
\caption{GenEval 分数提升曲线：LeapAlign 相比 ReFL、DRaFT-LV、DRTune 有更强的 compositional alignment 提升。图源：\href{https://arxiv.org/html/2604.15311v2}{LeapAlign arXiv HTML Figure 5}。}
\label{fig:leapalign-geneval-curve}
\end{figure}

\subsubsection{9.3 换 reward / prompt set 仍然有效}

论文还验证了不同 reward 和 prompt set：在 HPDv2 上用 PickScore 训练、在 MJHQ-30K 上用 HPSv3 训练时，LeapAlign 仍然优于 ReFL、DRaFT-LV 和 DRTune。

\begin{center}
\resizebox{0.78\linewidth}{!}{
\begin{tabular}{lccc}
\toprule
\textbf{方法} & \textbf{HPSv2.1} & \textbf{PickScore} & \textbf{HPSv3}\\
\midrule
Flux & 0.3078 & 22.7902 & 10.7624\\
ReFL & 0.3852 & 25.2373 & 11.7642\\
DRaFT-LV & 0.3859 & 24.9596 & 11.2701\\
DRTune & 0.3882 & 25.1021 & 12.0023\\
LeapAlign & \textbf{0.4092} & \textbf{25.7589} & \textbf{12.5855}\\
\bottomrule
\end{tabular}
}
\end{center}

\subsubsection{9.4 换到 SD3.5-M 也有提升}

在 Stable Diffusion 3.5 Medium 上，LeapAlign 也保持最好或接近最好的结果：HPSv2.1 为 0.3915，HPSv3 为 15.5780，PickScore 为 23.6180，ImageReward 为 1.4736。这个补充实验说明它不是只适用于 FLUX.1-dev。

\subsection{10. 和 ReFL / DRaFT-LV / DRTune 的差别}

论文把 direct-gradient 方法的差别总结成四个维度：是否能更新早期步骤、是否保留 nested gradient、是否构造 leap trajectory、是否能每条轨迹更新多个 step。

\begin{center}
\resizebox{0.86\linewidth}{!}{
\begin{tabular}{lcccc}
\toprule
\textbf{方法} & \textbf{Early Steps} & \textbf{Nested Gradient} & \textbf{Leap Trajectory} & \textbf{Multi-Step}\\
\midrule
ReFL & \(\times\) & \(\times\) & \(\times\) & \(\times\)\\
DRaFT-LV & \(\times\) & \(\times\) & \(\times\) & \(\times\)\\
DRTune & \(\checkmark\) & \(\times\) & \(\times\) & \(\checkmark\)\\
LeapAlign & \(\checkmark\) & \(\checkmark\) & \(\checkmark\) & \(\checkmark\)\\
\bottomrule
\end{tabular}
}
\end{center}

\begin{itemize}[leftmargin=1.5em]
  \item \textbf{ReFL / DRaFT-LV：} 更偏向 late-step direct optimization，容易改善细节，但难以稳定改早期构图。
  \item \textbf{DRTune：} 能覆盖早期和多步，但会丢掉 nested gradient，即丢掉一部分 step interaction。
  \item \textbf{LeapAlign：} 通过 two-step leap 缩短路径，通过 discount 保留但控制 nested gradient，通过 similarity weighting 抑制不可信 leap。
\end{itemize}

\subsection{11. 和 Flow-GRPO / OP-GRPO / E-GRPO / MaskFocus 的关系}

\begin{center}
\resizebox{\linewidth}{!}{
\begin{tabular}{llll}
\toprule
\textbf{方法} & \textbf{路线} & \textbf{核心对象} & \textbf{解决的问题}\\
\midrule
Flow-GRPO & policy gradient & SDE transition logprob & Flow ODE 如何变成可做 GRPO 的 policy\\
OP-GRPO & off-policy GRPO & replay trajectory & 高质量历史轨迹能否安全复用\\
E-GRPO & entropy-aware GRPO & high-entropy / merged step & 哪些 denoising steps 值得随机探索和更新\\
MaskFocus & focused GRPO & MGM critical unmasking step & masked generation 哪些 step 值得更新\\
LeapAlign & direct gradient & two-step leap trajectory & reward gradient 如何稳定传到任意 flow step\\
\bottomrule
\end{tabular}
}
\end{center}

\begin{itemize}[leftmargin=1.5em]
  \item \textbf{和 Flow-GRPO：} Flow-GRPO 需要把 ODE 改成 SDE，利用 reward 形成 advantage，再通过 logprob ratio 更新；LeapAlign 不需要 logprob，而是直接反传 reward gradient。
  \item \textbf{和 OP-GRPO：} OP-GRPO 提高的是 policy-gradient 路线的样本复用效率；LeapAlign 提高的是 direct-gradient 路线的反传可达性和稳定性。
  \item \textbf{和 E-GRPO：} E-GRPO 通过 entropy-aware step selection 减少无效随机探索；LeapAlign 通过随机 two-step leap 让所有 timestep 都有机会接收 reward gradient。
  \item \textbf{和 MaskFocus：} 两者都承认 step importance 不均匀，但 MaskFocus 是 masked generation 里选择 critical update steps，LeapAlign 是 flow matching 里构造短反传路径。
\end{itemize}

\begin{warnbox}
\textbf{不要混淆：} LeapAlign 不是 ``少采样几步生成图像'' 的 fast sampler。它训练时仍然需要完整 rollout 得到真实 \(x_0\) 和中间 latents；它省的是 reward-gradient 反传路径，而不是 necessarily 省掉 inference / rollout 的全部中间步骤。
\end{warnbox}

\subsection{12. 我认为它的局限}

\begin{enumerate}[leftmargin=1.5em]
  \item \textbf{它依赖可微 reward。} 如果 reward 来自黑盒 VLM judge、规则系统、人类反馈或不可微 API，LeapAlign 就不如 GRPO / policy-gradient 路线自然。
  \item \textbf{仍然需要完整 rollout。} Latent connector 和真实 \(x_0\) reward 都要求先跑出完整采样轨迹，所以它主要减少的是反传成本，不是直接减少采样成本。
  \item \textbf{leap trajectory 是近似路径。} 如果 \(\hat{x}_{j\mid k}\) 或 \(\hat{x}_{0\mid j}\) 偏离真实轨迹太大，梯度会变得不可靠；trajectory-similarity weighting 只能缓解，不能完全消除这个问题。
  \item \textbf{超参数仍然敏感。} \(\alpha\)、\(\lambda\)、\(\tau\) 都需要根据 reward、模型和数据重新校准。
  \item \textbf{reward hacking 仍然存在。} Hinge loss 可以缓和过度优化，但 direct-gradient 对 reward model 的梯度依赖更强，一旦 reward model 有漏洞，模型可能沿着梯度方向更快钻空子。
\end{enumerate}

\subsection{13. 对我们当前主线的启发}

我觉得 LeapAlign 对前面关于 ``关键 step / 跳步 / credit assignment'' 的讨论有三个启发：

\begin{itemize}[leftmargin=1.5em]
  \item \textbf{跳步不一定用于推理加速，也可以用于训练反传加速。} 我们之前担心 ``先找 10 个关键点也要跑完整 50 步''，这个问题在 LeapAlign 里仍然存在；但它说明即使完整 rollout 不省，仍然可以通过短反传路径降低训练显存和稳定性成本。
  \item \textbf{早期 step 的价值不只是 entropy 高，也在于控制全局结构。} E-GRPO 用 entropy 判断探索价值；LeapAlign 从 direct-gradient 角度说明，能更新早期步骤会更影响 composition。
  \item \textbf{跨步骤依赖不能简单切断。} Dense / critical step credit 很容易只看单个 step 的贡献，但 LeapAlign 的 nested gradient 提醒我们：早期 step 的作用往往通过后续 step 间接体现。
\end{itemize}

\begin{summarybox}
\textbf{我最后会这样概括：} Flow-GRPO 让 flow model 可以做 policy-gradient RL；OP-GRPO 和 E-GRPO 分别提高轨迹复用效率和 timestep 探索效率；LeapAlign 则开辟另一条 direct-gradient 路线，用 two-step leap trajectory 把 reward gradient 稳定送到任意生成步骤，尤其是早期决定构图的步骤。
\end{summarybox}

\subsection{参考来源}

\begin{itemize}[leftmargin=1.5em]
  \item \textbf{论文主页：} \href{https://arxiv.org/abs/2604.15311}{LeapAlign: Post-Training Flow Matching Models at Any Generation Step by Building Two-Step Trajectories}。
  \item \textbf{arXiv HTML：} \href{https://arxiv.org/html/2604.15311v2}{LeapAlign HTML version}，用于核对方法公式、图表、表格和实验结果。
  \item \textbf{PDF：} \href{https://arxiv.org/pdf/2604.15311}{arXiv PDF}，用于核对论文正文、附录和算法描述。
  \item \textbf{项目主页：} \href{https://rockeycoss.github.io/leapalign/}{LeapAlign project page}，用于核对项目入口与可视化说明。
  \item \textbf{代码仓库：} \href{https://github.com/RockeyCoss/LeapAlign_Code}{RockeyCoss/LeapAlign\_Code}，来源于项目主页。
\end{itemize}

% ============================================================
% Paper deep reading: POPO
% Source: https://arxiv.org/abs/2605.06650
% ============================================================

\section{论文精读：POPO}

\begin{conceptbox}
\textbf{论文：} \href{https://arxiv.org/abs/2605.06650}{Beyond Negative Rollouts: Positive-Only Policy Optimization with Implicit Negative Gradients}\\
\textbf{作者/机构：} Mingwei Xu, Hao Fang（共同一作，University of Washington, Seattle）\\
\textbf{简称：} POPO。\\
\textbf{版本：} arXiv 2605.06650v1，COLM 2026 投稿版。\\
\textbf{代码：} \href{https://github.com/momo1443/colm2026-POPO}{momo1443/colm2026-POPO}。\\
\textbf{方向标签：} RLVR、LLM Reasoning、GRPO、Positive-Only Optimization、Implicit Negative Gradient、Siamese Policy、EMA。\\
\textbf{阅读日期：} 2026-05-09。\\
\textbf{阅读口径：} 这篇不在 Flow-GRPO 视觉生成主线上，但和 \key{egrpo / opgrpo / densegrpo} 共享一个上层问题——\textbf{RL 信号到底从哪里来、由谁来惩罚错误轨迹}。POPO 给的答案非常激进：在 RLVR 里干脆 \textbf{不显式使用负样本}，让负梯度从 softmax 概率守恒中自动涌现。
\end{conceptbox}

\subsection{0. 一句话结论}

POPO 的核心判断是：在二值可验证奖励（RLVR）下，GRPO 显式拿负 rollout 去算 advantage，其实是在用一个 \key{严重欠采样、且本身没有失败程度区分} 的负样本集合制造梯度噪声。它选择的替代路线是：

\[
\begin{aligned}
\text{POPO}
&= \text{仅在正 rollout 上做有界重要性采样}\\
&\quad + \text{Siamese 策略 + EMA 动量更新}\\
&\quad + \text{表征空间余弦相似度惩罚（替代 KL）}.
\end{aligned}
\]

它在 Qwen-Math、R1-Distill 1.5B/7B 和 Llama-3.1-8B 上都打过 GRPO/DAPO/Dr.~GRPO/SAPO，AIME 2025 在 Qwen-Math-7B 上从 30.00\% 提到 36.67\%（绝对 +6.67 pp，相对 +22\%）。

\begin{summarybox}
\textbf{我的判断：} 这篇值得精读。它的真正贡献不是 ``少采点负样本省钱'' 这种工程优化，而是 \key{从 softmax Jacobian 出发证明：只要在正样本集合内自归一化，负 logit 的梯度就严格为正}。换句话说，``罚错误'' 这件事并不需要显式去采错误样本，而是 softmax 概率守恒的直接副产品。这一点和 NGRPO 的 ``加强负样本'' 路线正好相反，但又在数学上自洽。
\end{summarybox}

\subsection{1. 为什么读这篇}

\begin{itemize}[leftmargin=1.5em]
  \item \textbf{和当前主线的关系：} 我在 post-training 主线上已经记了 DenseGRPO（reward 怎么分到每一步）、OP-GRPO（旧轨迹复用）、E-GRPO（哪些 step 值得探索）。POPO 是 \key{rollout 侧} 的另一个极端问题：到底要不要显式用错的 rollout？
  \item \textbf{它可能补上的知识缺口：} 之前我对 GRPO 的负样本一直有个隐约的疑问——失败答案的语义空间巨大，组内随机采到的几条 wrong rollouts 真的能代表 ``错误方向''吗？POPO 直接给了一个 ``根本不需要它们'' 的回答，并且配了证明。
  \item \textbf{我希望读完回答的问题：}
  \begin{itemize}[leftmargin=1.5em]
    \item 仅靠正样本，负方向梯度从哪里来？
    \item 没有 KL 怎么防 policy drift？
    \item 它和 NGRPO/SAPO/DAPO 这些 ``强化负样本'' 方法是不是真矛盾？
  \end{itemize}
\end{itemize}

\subsection{2. 问题定义与背景}

\textbf{问题：} RLVR 设定下，奖励是二值的（验证器判对/错），训练目标是用 RL 提升 LLM 在数学推理等任务上的 pass@k。

\textbf{输入/输出：} 输入是 prompt $x$，模型采样 $G$ 条 rollout $y_1,\ldots,y_G$，每条得到一个二值奖励 $R(x,y)\in\{0,1\}$；目标是更新 $\pi_\theta$。

\textbf{前作限制：} 论文重点批评了 GRPO 风格的负样本机制：

\begin{itemize}[leftmargin=1.5em]
  \item \textbf{失败程度无差异：} 二值奖励下 ``算错一步'' 和 ``完全跑题'' 都是 $R=0$，组内 advantage 把它们当成同一种负反馈。
  \item \textbf{负样本空间组合爆炸：} 错的 reasoning chain 数量远多于对的，组内 $G$ 条采样能覆盖的负方向极其稀疏。
  \item \textbf{KL 散度的副作用：} 在 token 概率空间约束 policy drift 会限制 exploration，且和 reasoning 这种 ``长 chain 多分支'' 任务不太契合。
\end{itemize}

\begin{intubox}
\textbf{直觉版：} GRPO 的负梯度是 ``我抽了几条错答案，告诉模型别像它们''。POPO 的逻辑是：``我只告诉模型对答案长什么样，softmax 自己会把概率从其它所有地方挪过来——这本身就是对错答案的惩罚，而且覆盖了所有错的方向，不只是我采到的那几条''。
\end{intubox}

\subsection{3. 方法拆解}

\subsubsection{3.1 总体框架}

POPO 维护两份策略：

\begin{itemize}[leftmargin=1.5em]
  \item \textbf{在线策略 $\pi_\theta$：} 接收梯度更新，负责采样和被优化。
  \item \textbf{Siamese 策略 $\pi_\xi$：} 不接收梯度，靠 EMA 动量更新，作为 ``自身锚点'' 用来做表征对齐。
\end{itemize}

每个 prompt $x$ 的 rollout 集合按奖励切成两半：

\[
\mathcal{S}^{+}(x)=\{y:R(x,y)=1\},
\qquad
\mathcal{S}^{-}(x)=\{y:R(x,y)=0\}.
\]

POPO 的全部训练梯度只来自 $\mathcal{S}^{+}(x)$；$\mathcal{S}^{-}(x)$ 中的样本不进入显式损失项，但论文会证明它们的 logit 仍然在被压低。

\subsubsection{3.2 关键机制}

\paragraph{机制一：有界重要性采样（bounded importance sampling）。}
不是和 PPO/GRPO 一样用 $\pi_\theta/\pi_{\theta_{\text{old}}}$ 这种跨 iteration 的 ratio，而是在当前正样本集合内自归一化：

\[
w_\theta(y\mid x)=\frac{\pi_\theta(y\mid x)}{Z^{+}(x)},
\qquad
Z^{+}(x)=\sum_{y'\in\mathcal{S}^{+}(x)}\pi_\theta(y'\mid x).
\]

满足 $\sum_{y\in\mathcal{S}^{+}}w_\theta(y\mid x)=1$。它有两个直接后果：

\begin{itemize}[leftmargin=1.5em]
  \item 高概率正样本得到更大权重，形成正样本之间的 ``自竞争''；
  \item 权重归一化让负 logit 的梯度行为变得干净（见 §3.3 的导数推导）。
\end{itemize}

\paragraph{机制二：Siamese 策略 + EMA 动量。}
Siamese 策略 $\pi_\xi$ 不参与梯度计算，只做指数滑动：

\[
\xi\leftarrow\tau\,\xi + (1-\tau)\,\theta,\qquad \tau=0.999.
\]

它的角色是给 ``当前策略不要漂得太远'' 提供一个稳定的参考点。

\paragraph{机制三：表征空间相似度惩罚替代 KL。}
论文不再用 token 概率空间的 KL，而是定义一个 BYOL/SimSiam 风格的不对称表征对齐损失：

\[
\mathcal{L}_{\text{sim}}(\theta,\phi,\xi)
=
-\mathbb{E}_x\!\!\sum_{y\in\mathcal{S}^{+}(x)}\!\!w_\theta(y\mid x)\,
\cos\!\Big(h_\phi\big(f_\theta(x,y)\big),\;\mathrm{sg}\big(f_\xi(x,y)+\epsilon\big)\Big),
\]

其中 $f_\theta,f_\xi$ 是两边网络的隐藏表征，$h_\phi$ 是只挂在 online 一侧的 2 层 MLP（hidden=4096）predictor，$\mathrm{sg}$ 是 stop-gradient，$\epsilon\sim\mathcal{N}(0,\sigma^2 I),\sigma=0.02$ 用来防表征塌缩。

\begin{warnbox}
\textbf{为什么不能 symmetric：} 如果两边都加 predictor、都要梯度，会塌缩到平凡解 $f_\theta\equiv f_\xi\equiv\text{const}$。BYOL/SimSiam 这条经验在这里被直接搬过来：\textbf{不对称 + stop-gradient + EMA target 才有非平凡解}。
\end{warnbox}

\subsubsection{3.3 关键公式与隐式负梯度证明}

\paragraph{总损失（Equation 5）。}
\[
\mathcal{L}_{\text{POPO}}(\theta)
=
-\mathbb{E}_{x}\!\!\sum_{y\in\mathcal{S}^{+}(x)}\!\!w_\theta(y\mid x)\log\pi_\theta(y\mid x)
\;+\;\alpha\,\mathcal{L}_{\text{sim}}
\;+\;\beta\,\mathcal{L}_{\text{ent}},
\]

其中 $\mathcal{L}_{\text{ent}}=-\mathbb{E}_x\,H(\pi_\theta(\cdot\mid x))$ 是熵正则化项。

注意：第一项是加权 NLL，但权重 $w_\theta$ 自己也是 $\theta$ 的函数。

\paragraph{核心定理（Theorem 3.1）：负 logit 的梯度严格为正。}
对任意错误响应 $y'\in\mathcal{S}^{-}(x)$，关于其 logit $z_{y'}$ 的梯度为：

\[
\frac{\partial \mathcal{L}_{\text{POPO}}}{\partial z_{y'}}
=
\pi_\theta(y'\mid x)\Big[1+\beta\big(\log\pi_\theta(y'\mid x)+H(\pi_\theta(\cdot\mid x))\big)\Big].
\]

当 $\pi_\theta(y'\mid x)>\exp\!\bigl(-\tfrac{1}{\beta}-H(\pi_\theta(\cdot\mid x))\bigr)$ 时，整体严格为正，即 \key{该负样本的 logit 一定被往下压}。

\paragraph{推导骨架（Appendix B.1）。}
关键步骤如下：

\begin{enumerate}[leftmargin=1.5em]
  \item Softmax Jacobian：因为 $y\in\mathcal{S}^{+},\,y'\in\mathcal{S}^{-}$ 不相交，$\delta_{y,y'}=0$，所以
  \[
  \frac{\partial \log\pi_\theta(y\mid x)}{\partial z_{y'}}=-\pi_\theta(y'\mid x).
  \]
  \item 权重 $w_\theta(y\mid x)$ 对负 logit 的偏导是 0：用商法则配合 $\partial Z^{+}/\partial z_{y'}=-\pi_\theta(y'\mid x)\,Z^{+}$，分子分母的因子精确抵消。\textbf{这是有界重要性采样的关键性质}：归一化不会削弱或扭曲负方向的隐式梯度。
  \item NLL 梯度（利用 $\sum_{y\in\mathcal{S}^{+}}w_\theta=1$）：
  \[
  \frac{\partial \mathcal{L}_{\text{NLL}}}{\partial z_{y'}}
  =\pi_\theta(y'\mid x)\cdot\sum_{y\in\mathcal{S}^{+}}w_\theta(y\mid x)
  =\pi_\theta(y'\mid x).
  \]
  论文把它叫做 \key{``概率税''}：无条件压低所有错答 logit，且 logit 越大压得越多。
  \item 熵项贡献 $\pi_\theta(y'\mid x)[\log\pi_\theta(y'\mid x)+H(\pi_\theta)]$：对高概率错答（$\log\pi$ 越大于 $-H$）施加最强抑制，自然偏向 ``打击当前最容易错下去的那条路''。
  \item $\mathcal{L}_{\text{sim}}$ 不依赖输出 logit，因此对 $z_{y'}$ 没有贡献。
\end{enumerate}

\begin{intubox}
\textbf{直觉：} softmax 是个零和游戏 $\sum_y \pi_\theta(y\mid x)=1$。任何对正样本概率的提升必然要从其它地方搬概率过来——而 $\mathcal{S}^{+}$ 之外的所有 token 序列就是被搬概率的 ``其它地方''。POPO 不需要点名打击某条具体的错答，softmax 自己会等比例打击所有错答，并且对当前最高概率的那条错答打得最狠。
\end{intubox}

\paragraph{EMA 参数有界性（Lemma 3.2）。}
在 $\|\nabla_\theta\mathcal{L}_{\text{POPO}}\|\le G_{\max}$ 的条件下，稳态时

\[
\|\theta_t-\xi_t\|\le \frac{\tau\eta G_{\max}}{1-\tau}.
\]

代入论文设定 $\tau=0.999,\,\eta=10^{-6},\,G_{\max}=1$，得 $\|\delta\|<10^{-3}$。这是论文 \key{放心去掉 KL 惩罚} 的依据：表征对齐 + EMA 动量已经在参数空间提供了一个有界 drift 保证，不需要再在 token 概率层面加 KL。

\begin{warnbox}
\textbf{这一步在数学上其实很关键：} 它把 ``policy 不能漂太远'' 这个约束从 \key{每条 token 概率都不能动太多}（KL）放松到了 \key{两份策略参数距离有上界}。前者把每一个分支都锁住，后者允许策略在任意单条序列上自由演化、只要整体行为不偏离 EMA 锚点。这对 long-chain reasoning 显然更友好。
\end{warnbox}

\subsection{4. 实验与证据}

\subsubsection{4.1 训练与评测设定}

\begin{itemize}[leftmargin=1.5em]
  \item \textbf{训练数据：} DeepScaleR-Preview-Dataset（约 40K 数学题，覆盖 AIME 1984--2023、AMC、Omni-MATH、Still）。
  \item \textbf{基础模型：} Qwen-Math-1.5B / 7B、R1-Distill-Qwen 1.5B / 7B、Llama-3.1-8B、DeepSeek-Math-7B。
  \item \textbf{评测：} MATH-500、AMC 2023、AIME 2024、AIME 2025、OlympiadBench；Pass@8（$n=128$ 样本）；temperature 0.7，top-p 0.95，max output 4096。
  \item \textbf{对比方法：} GRPO、Dr.~GRPO、DAPO、SAPO（同样是 RLVR 后训练，但都显式使用负 rollout 或非对称 ratio）。
  \item \textbf{POPO 关键超参：} $\tau=0.999$，$\alpha=0.1$，$\beta=0.01$，$\sigma=0.02$，predictor 2 层 $\times$ 4096 维。
  \item \textbf{硬件：} 4 $\times$ A100 80GB，1.5B 约 11h/epoch，7B 约 24h/epoch。
\end{itemize}

\subsubsection{4.2 主结果：数学专用与蒸馏模型}

\begin{center}
\resizebox{\linewidth}{!}{
\begin{tabular}{llcccccc}
\toprule
\textbf{Base} & \textbf{Algo} & \textbf{MATH-500} & \textbf{AMC 23} & \textbf{AIME 24} & \textbf{AIME 25} & \textbf{Olympiad} & \textbf{Avg.}\\
\midrule
Qwen-Math-1.5B & GRPO     & 86.20 & 75.00 & 23.33 & 16.25 & 50.30 & 50.22\\
               & SAPO     & 86.20 & 76.17 & 23.33 & 17.11 & 50.74 & 50.71\\
               & Dr.~GRPO & 86.00 & 76.62 & 25.18 & 17.39 & 50.89 & 51.22\\
               & DAPO     & 86.40 & 75.28 & 22.85 & 16.67 & 50.59 & 50.36\\
               & \textbf{POPO} & \textbf{86.60} & \textbf{77.50} & \textbf{26.67} & \textbf{23.33} & \textbf{51.19} & \textbf{53.06}\\
\midrule
R1-Distill-1.5B & GRPO    & 90.20 & 80.00 & 30.00 & 26.32 & 58.61 & 57.03\\
               & Dr.~GRPO & 90.40 & 83.46 & 33.33 & 25.84 & 58.90 & 58.39\\
               & \textbf{POPO} & \textbf{90.80} & \textbf{85.24} & \textbf{36.67} & \textbf{27.86} & 59.05 & \textbf{59.92}\\
\midrule
Qwen-Math-7B   & GRPO     & 90.80 & 85.50 & 43.33 & 30.00 & 56.38 & 61.20\\
               & DAPO     & \textbf{91.40} & 85.97 & 44.76 & 32.31 & 56.53 & 62.19\\
               & \textbf{POPO} & 90.80 & 85.75 & 45.13 & \textbf{36.67} & \textbf{57.27} & \textbf{63.12}\\
\midrule
R1-Distill-7B  & GRPO     & 93.60 & 87.50 & 46.67 & 33.33 & 66.91 & 65.60\\
               & SAPO     & \textbf{94.20} & \textbf{88.21} & 46.71 & 35.00 & 67.06 & \textbf{66.24}\\
               & \textbf{POPO} & 93.20 & 87.25 & 47.22 & \textbf{36.67} & 66.76 & 66.22\\
\bottomrule
\end{tabular}
}
\end{center}

\begin{summarybox}
\textbf{关键观察：} 任务越难，POPO 相对 GRPO 的相对提升越大——MATH-500 上只有 ~1.7\% 相对提升，到 AIME 2025 上就有 \textbf{17.6\%} 相对提升。这和论文的核心主张一致：\key{越是答案空间稀疏、负样本越无信息量的任务，POPO 的 ``不显式用负样本'' 策略收益越明显}。R1-Distill-7B 上 POPO 略低于 SAPO，作者归因于蒸馏后 base policy 的熵已经很低，on-policy 探索受限。
\end{summarybox}

\subsubsection{4.3 通用文本 LLM 上 POPO 也可迁移}

\begin{center}
\resizebox{\linewidth}{!}{
\begin{tabular}{llcccccc}
\toprule
\textbf{Base} & \textbf{Algo} & \textbf{MATH-500} & \textbf{AMC 23} & \textbf{AIME 24} & \textbf{AIME 25} & \textbf{Olympiad} & \textbf{Avg.}\\
\midrule
Llama-3.1-8B    & GRPO  & 56.40 & 35.22 & 10.00 & 6.67  & 25.32 & 26.72\\
                & \textbf{POPO}  & \textbf{59.80} & \textbf{37.78} & \textbf{13.33} & \textbf{10.00} & \textbf{28.29} & \textbf{29.84}\\
DeepSeek-Math-7B & GRPO & 80.20 & 73.38 & 20.00 & 16.67 & 48.81 & 47.81\\
                & \textbf{POPO}  & \textbf{82.60} & \textbf{74.53} & \textbf{23.33} & \textbf{20.00} & \textbf{50.29} & \textbf{50.15}\\
\bottomrule
\end{tabular}
}
\end{center}

POPO 在两个非数学专用 base 上对 GRPO 都是全面胜出，说明 ``正样本归一化 + 隐式负梯度'' 不依赖于特定 base 是否经过数学预训练。

\subsubsection{4.4 关键消融：哪一部分真的重要}

\paragraph{(a) Importance sampling 与权重再分配（Qwen-Math-1.5B）。}

\begin{center}
\begin{tabular}{lccc}
\toprule
\textbf{配置} & \textbf{AIME 25} & \textbf{Olympiad} & \textbf{Avg.}\\
\midrule
POPO（无负样本 + 再分配权重） & \textbf{23.33} & \textbf{51.19} & \textbf{37.26}\\
无负样本 + 均匀权重           & 16.67 & 49.11 & 32.89\\
有负样本 + 无再分配权重       & 13.33 & 48.07 & 30.70\\
有负样本 + 均匀权重           & 6.67  & 39.32 & 23.00\\
\bottomrule
\end{tabular}
\end{center}

\begin{warnbox}
\textbf{这组消融最值得记：} 单独 ``扔掉负样本''（无负样本 + 均匀权重）已经能从 23 涨到 33；再加 \key{有界重要性采样的权重再分配}，又从 33 涨到 37。这说明 POPO 的性能不全靠 ``少损害''，而是真的来自 \key{$w_\theta=\pi_\theta/Z^{+}$ 这个归一化引入的隐式负梯度结构}——这个证明在 §3.3 已经做过了。
\end{warnbox}

\paragraph{(b) Siamese 动量与表征对齐。}

\begin{center}
\begin{tabular}{lccc}
\toprule
\textbf{配置} & \textbf{AIME 25} & \textbf{Olympiad} & \textbf{Avg.}\\
\midrule
POPO（完整）            & \textbf{23.33} & \textbf{51.19} & \textbf{37.26}\\
无动量（静态锚 $\xi$）   & 13.33 & 48.37 & 30.85\\
无表征对齐（去掉 $\mathcal{L}_{\text{sim}}$） & 13.33 & 50.15 & 31.74\\
\bottomrule
\end{tabular}
\end{center}

去掉 EMA 或去掉表征对齐都会掉约 6 个点，二者都是 ``在没有 KL 的情况下保持训练稳定'' 必需的两个支柱。

\paragraph{(c) Predictor 与动量的超参扫描。}
2 层 predictor 优于 3/4 层，hidden=4096 优于 1024/2048，$\tau=0.999$ 是最佳动量值（0.99 太短导致漂移，0.9999 太慢失去自适应能力）。

\subsubsection{4.5 实验是否充分支撑论点}

\begin{itemize}[leftmargin=1.5em]
  \item ``正样本归一化引入隐式负梯度'' 既有 Theorem 3.1 的证明，也有消融 (a) 的数值支撑——\textbf{支撑充分}。
  \item ``EMA + 表征对齐能替代 KL'' 有 Lemma 3.2 的有界性 + 消融 (b) 的下降幅度——\textbf{支撑较强}。
  \item ``难任务收益更大'' 是论文最强的实证现象之一（AIME 25 比 MATH-500 提升大一个数量级）——\textbf{支撑显著}。
  \item 但 ``为什么 R1-Distill-7B 上反而打不过 SAPO''，作者只给了 ``蒸馏模型熵低、on-policy 探索受限'' 的解释，没有更深入的实验验证。
\end{itemize}

\subsection{5. 局限与风险}

\begin{itemize}[leftmargin=1.5em]
  \item \textbf{需要正样本存在：} 如果某个 prompt 的所有 $G$ 条 rollout 都失败（$\mathcal{S}^{+}=\emptyset$），整条梯度就消失。论文没有像 DAPO 那样专门处理 ``全错 group''——在更难的任务（IMO 级别）上这是真实风险。
  \item \textbf{熵下降模型不友好：} R1-Distill 系列已经被蒸馏得很 ``确定'' 了，正样本采样多样性不够，POPO 的 ``正样本之间自竞争'' 就退化成 ``一条主导路径吃掉全部权重''。这正是 R1-Distill-7B 上 POPO 略输 SAPO 的根因。
  \item \textbf{表征空间相似度的语义对齐性未必稳：} $\mathcal{L}_{\text{sim}}$ 用的是隐藏表征上的 cosine，而不是 token 概率上的 KL。这意味着 ``policy 没怎么变'' 是在表征几何上判断的，对于真正决定输出分布的最后几层 logit 头，约束相对宽松——可能在某些边角任务上表现为输出分布偏离较大但 sim 损失不报警。
  \item \textbf{只验证了数学推理：} 没有 code、agent、安全对齐这些 RLHF/RLVR 常见赛道上的实验，所以 ``正样本 + 隐式负梯度'' 是否在 reward 不那么稀疏二值的设置下还成立，仍然待证。
\end{itemize}

\subsection{6. 与已有工作的关系}

\begin{itemize}[leftmargin=1.5em]
  \item \textbf{直接前作：} GRPO（DeepSeek-Math 提出的组相对优势），以及它的几个后续 Dr.~GRPO、DAPO、SAPO。POPO 把 ``如何用好负样本'' 这条主线整体翻面，说 ``可以不用''。
  \item \textbf{相反路线：} NGRPO（Negative-enhanced GRPO，arXiv 2509.18851）走的是另一个极端——\key{加大对负样本的利用}。两篇都直觉上有合理性，差异在于：NGRPO 假设负样本含有 ``避免某类错误'' 的额外信息，POPO 假设这些信息在 softmax 守恒下已经被免费提供。
  \item \textbf{相近自监督思想：} Siamese + EMA + stop-gradient + predictor 几乎照搬 BYOL/SimSiam 的 anti-collapse 设计。POPO 的贡献不是发明这个结构，而是\key{发现它能用作 RLVR 里的 KL 替代品}。
  \item \textbf{真正差异：} 大部分 ``positive-only'' 工作（包括 PODPO、SAPO 中的 positive 加权变体）依赖正样本的绝对优势或某种锚点对比；POPO 的特别之处是把它直接还原成 \key{有界 importance sampling 的归一化结构}，并给出了负 logit 严格被压低的解析证明。
\end{itemize}

\subsection{7. 对我的可复用点}

\begin{itemize}[leftmargin=1.5em]
  \item \textbf{可复用 idea 1（视觉生成版 POPO）：} 在 Flow-GRPO / DanceGRPO 这类视觉 RL 里，``坏图'' 的语义空间同样巨大且无差异（一张全是噪点的图和一张轻微 misalign 的图都给 reward 0）。能不能在视觉 RL 里也只对 high-reward trajectory 做 $w_\theta=\pi_\theta/Z^{+}$ 自归一？这个想法值得放进 \key{post\_training/index} 的 ideas 列表。
  \item \textbf{可复用 idea 2（KL 替代）：} 表征空间余弦对齐 + EMA 动量，作为 KL 替代项，可能比 token-level KL 更适合 long horizon 视频生成 RL，因为视频策略更适合 ``整体行为别漂'' 而不是 ``每一帧概率别动''。
  \item \textbf{可复用证明套路：} Theorem 3.1 的 softmax Jacobian + 商法则证明负梯度涌现，是一种很干净的 ``利用 softmax 守恒导出隐式约束'' 范式。我以后写 implicit gradient / implicit constraint 类的论文可以借用这种结构。
  \item \textbf{可复用实验设计：} 4 个 base 模型 $\times$ 5 个 benchmark 的方阵 + 2 张消融表覆盖了 ``算什么 vs 怎么算''，结构非常工整，是 RLVR 类论文写法的好模板。
  \item \textbf{值得追的 follow-up：} POPO + DAPO 的全错 group 处理；POPO 在 code RL 上是否仍然成立；POPO 与 OP-GRPO 风格 replay 的兼容性（``只复用旧的正轨迹'' 是不是天然就是 POPO 的特例）。
\end{itemize}

\subsection{8. 待查问题}

\begin{itemize}[leftmargin=1.5em]
  \item Theorem 3.1 中 ``$\pi_\theta(y'\mid x)>\exp(-1/\beta-H(\pi_\theta))$ 时严格为正'' 这个条件在训练后期会不会失效？特别是当 $H$ 随训练下降时，它会反过来抬高门槛。
  \item 论文用 Pass@8（n=128）来评测，但训练时 group size 只有 8。Pass@8 评测是否对方法间公平？是否有 Pass@1 的对比？
  \item EMA $\tau=0.999$ + lr=$10^{-7}$ 的组合非常具体，换一个数据集量级（比如 100K+ 训练）后 Lemma 3.2 的 bound $\tau\eta G_{\max}/(1-\tau)$ 还能否保证训练稳定？
  \item POPO 在 R1-Distill-7B 上略输 SAPO，那是否意味着 ``越是 SFT/distill 后熵低的模型，越需要显式负样本''？如果是，POPO 的适用区间就比论文标题暗示的窄。
\end{itemize}

\subsection{参考来源}

\begin{itemize}[leftmargin=1.5em]
  \item \textbf{论文主页：} \href{https://arxiv.org/abs/2605.06650}{Beyond Negative Rollouts: Positive-Only Policy Optimization with Implicit Negative Gradients}。
  \item \textbf{arXiv HTML：} \href{https://arxiv.org/html/2605.06650v1}{arXiv HTML version}（Theorem 3.1、Lemma 3.2、Equations 5--8 和全部消融表均来自此版本）。
  \item \textbf{代码仓库：} \href{https://github.com/momo1443/colm2026-POPO}{momo1443/colm2026-POPO}。
  \item \textbf{对照阅读：} NGRPO（arXiv 2509.18851，反向路线）、SAPO / DAPO / Dr.~GRPO（GRPO 的负样本改良路线）。
\end{itemize}

% ============================================================
% Paper deep reading: DGPO
% Source: https://arxiv.org/abs/2510.08425
% ============================================================

\section{论文精读：DGPO}

\begin{conceptbox}
\textbf{论文：} \href{https://arxiv.org/abs/2510.08425}{Reinforcing Diffusion Models by Direct Group Preference Optimization}\\
\textbf{作者/机构：} Yihong Luo (HKUST), Tianyang Hu (CUHK-SZ), Jing Tang (HKUST-GZ / HKUST)\\
\textbf{简称：} DGPO。\\
\textbf{版本：} arXiv 2510.08425v1，\textbf{ICLR 2026} 已接收。\\
\textbf{代码：} \href{https://github.com/Luo-Yihong/DGPO}{Luo-Yihong/DGPO}（训练代码已开源，集成进 Flow-Factory）。\\
\textbf{方向标签：} Flow Matching、Diffusion RL、GRPO、DPO、Group Preference、ODE Sampler、Timestep Clip、SD3.5。\\
\textbf{阅读日期：} 2026-05-09。\\
\textbf{阅读口径：} 这篇是 Flow-GRPO 视觉后训练主线上第一次 \textbf{彻底抛弃 policy-gradient 框架} 的工作。它的问题不是 ``怎么更好地做 SDE rollout''，而是 \key{``我们真的需要 SDE 吗''}。
\end{conceptbox}

\subsection{0. 一句话结论}

DGPO 把 GRPO 成功的那部分拆成两块：``policy-gradient 形式''（可有可无）和 ``组内细粒度相对偏好''（这才是核心），然后只保留后者。它的结果是：

\[
\begin{aligned}
\text{DGPO}
&= \text{group-level preference 最大似然}\\
&\quad + \text{Diffusion-DPO 的 denoising-based reward 参数化}\\
&\quad + \text{advantage-based weighting（自动消 } Z(c)\text{）}\\
&\quad + \text{ODE rollout + timestep clip}.
\end{aligned}
\]

落到训练上就是：\key{同组里的正负样本，不再去做 SDE 采样、不再算整条轨迹的 logprob ratio}，而是沿 Diffusion-DPO 的做法，用随机 timestep 上的 denoising loss 差值近似 ``组内 reward 差''，再套一个 Bradley--Terry sigmoid 做最大化。

\begin{summarybox}
\textbf{我的判断：} 这篇是 Flow-GRPO 这条线里少见的 ``换路径'' 工作。Flow-GRPO / DanceGRPO / MixGRPO / E-GRPO / OP-GRPO 都在 \key{policy-gradient 框架内} 修 credit assignment / sample efficiency / exploration；DGPO 直接说 ``不用 policy gradient 也行，而且训练快约 20 倍''。值得精读，而且要和 OP-GRPO、E-GRPO 并列记：它们代表了 \key{同一症状的三种不同治法}。
\end{summarybox}

\begin{figure}[htbp]
\centering
\includegraphics[width=0.92\linewidth]{dgpo/fig1_speed_curve.png}\\[4pt]
\includegraphics[width=0.58\linewidth]{dgpo/fig1_ood_radar.png}
\caption{DGPO 总览（论文 Fig.~1）。上：在 GenEval 上 DGPO 的训练速度接近 Flow-GRPO 的 30 倍。下：同时在 DrawBench 的 OOD 指标（Aesthetic / DeQA / ImageReward / UnifiedReward / PickScore）上全部优于或持平 Flow-GRPO，说明速度收益不是以 OOD 质量为代价。图源：\href{https://arxiv.org/html/2510.08425v1}{arXiv HTML Figure 1}。}
\label{fig:dgpo-overview}
\end{figure}

\subsection{1. 为什么读这篇}

\begin{itemize}[leftmargin=1.5em]
  \item \textbf{和当前主线的关系：} Flow-GRPO 把 ODE 改写成 SDE 才能做 policy gradient，这是整条 Flow-GRPO 后训练主线的默认前提。DGPO 直接质疑这个前提。
  \item \textbf{它可能补上的知识缺口：}
    \begin{itemize}[leftmargin=1.5em]
      \item DPO 类方法（Diffusion-DPO）和 GRPO 类方法在数学上能否统一？
      \item 组级偏好怎样在数学上消掉 DPO 里那个难算的配分函数 $Z(c)$？
      \item ``用 ODE 采样去做 RL'' 实践上能否稳定训练？
    \end{itemize}
  \item \textbf{我希望读完回答的问题：}
    \begin{itemize}[leftmargin=1.5em]
      \item DGPO 是不是真的不依赖随机策略？
      \item 它能替代 Flow-GRPO 吗，还是只在某些任务上有优势？
      \item 它和 OP-GRPO / E-GRPO 是互斥还是可叠加？
    \end{itemize}
\end{itemize}

\subsection{2. 问题定义与背景}

\textbf{问题：} 给定预训练扩散 / flow-matching 模型 $f_\theta$、一个可计算 reward $r(c,x_0)$ 的奖励模型，用 RL 去后训练 $f_\theta$ 使生成结果对齐 reward，同时不破坏原模型的 OOD 质量（aesthetic / image-reward / picks-core）。

\textbf{前作限制（GRPO 线）：} 论文指出 GRPO 迁移到扩散模型有三个痛点：

\begin{enumerate}[leftmargin=1.5em]
  \item \textbf{SDE rollout 效率低。} GRPO 需要随机策略，逼迫使用 SDE 采样器取代更高效的 ODE，\key{相同 NFE 下样本质量更差}。
  \item \textbf{收敛慢。} 随机性来源是 ``model-agnostic 高斯噪声''，它在高维空间提供的学习信号非常弱。
  \item \textbf{训练昂贵。} 必须在整条采样轨迹上算 logprob ratio 和梯度，单次迭代成本高。
\end{enumerate}

\textbf{前作限制（DPO 线）：} Diffusion-DPO 已经做到 ``不需要随机策略''，但它只有 \key{成对偏好（pairwise）}，完全丢掉 GRPO 的 \key{组内细粒度相对信息}。

\begin{intubox}
\textbf{直觉版：} GRPO 有组信息但被迫用 SDE；DPO 用 ODE 友好但只认两两比较。DGPO 要的是 ``ODE 友好 + 组级偏好''，所以它把 DPO 从 pairwise 扩到 \key{pairwise of groups}。
\end{intubox}

\subsection{3. 方法拆解}

\subsubsection{3.1 关键一步：把 DPO 从 pair 扩到 group pair}

Bradley--Terry 偏好模型本来写在 \textbf{一对样本上}：

\[
p_{\text{BT}}(x_0^w\succ x_0^l\mid c) = \sigma\bigl(r(c,x_0^w)-r(c,x_0^l)\bigr).
\]

DGPO 把 ``一对样本'' 换成 ``一对样本组''。对一个 prompt $c$ 生成 $G$ 条样本 $\mathcal{G}=\{x_0^1,\ldots,x_0^G\}$，按 reward 切成 $\mathcal{G}^{+}$（偏好）和 $\mathcal{G}^{-}$（非偏好），目标（Eq.~8）：

\[
\max_\theta\ \mathbb{E}_{(\mathcal{G}^{+},\mathcal{G}^{-},c)\sim\mathcal{D}}
\log\sigma\bigl(R_\theta(\mathcal{G}^{+}\mid c)-R_\theta(\mathcal{G}^{-}\mid c)\bigr),
\]

组级 reward 定义为加权和（Eq.~9）：

\[
R_\theta(\mathcal{G}\mid c)=\sum_{x_0\in\mathcal{G}}w(x_0)\cdot r_\theta(c,x_0).
\]

\subsubsection{3.2 单样本 reward：沿用 Diffusion-DPO 的 denoising 近似}

DGPO 沿用 Diffusion-DPO 的 reward 参数化（Eq.~10）：

\[
r_\theta(c,x_0)\approx
\beta\,\mathbb{E}_{q(x_{1:T}\mid x_0)}\!
\log\frac{p_\theta(x_{0:T}\mid c)}{p_{\text{ref}}(x_{0:T}\mid c)}
\;+\;\beta\log Z(c).
\]

用 \textbf{前向扩散分布} $q(x_{1:T}\mid x_0)$ 近似反向链，把 ``要跑完整条采样轨迹'' 换成 ``随机采一个 $t$ 的 denoising loss''。

\begin{warnbox}
\textbf{这一步继承自 DPO：} 昂贵的反向链被近似成加噪前向链，DGPO 训练时 \key{只需要对随机 $t$ 做一次 denoising loss}，和 DPO 一样便宜。这是 DGPO 比 Flow-GRPO 快约 20 倍的主要来源——后者要对轨迹上每个 $t$ 都算 transition logprob。
\end{warnbox}

\subsubsection{\texorpdfstring{3.3 关键设计：用 advantage 做权重，让 $Z(c)$ 自动消掉}{3.3 关键设计：用 advantage 做权重，让 Z(c) 自动消掉}}

把 Eq.~9 + Eq.~10 代回 Eq.~8，$\beta\log Z(c)$ 在 $R_\theta(\mathcal{G}^{+})-R_\theta(\mathcal{G}^{-})$ 里留下 $\bigl(\sum_{\mathcal{G}^{+}}w-\sum_{\mathcal{G}^{-}}w\bigr)\log Z(c)$。要让 $Z(c)$ 彻底消去，权重需满足：

\begin{enumerate}[leftmargin=1.5em]
  \item 越好的正样本 / 越差的负样本应得更大权重（保留组内相对偏好信息）；
  \item $\sum_{x_0\in\mathcal{G}^{+}}w(x_0)=\sum_{x_0\in\mathcal{G}^{-}}w(x_0)$（两侧权重等质量）。
\end{enumerate}

它直接复用 GRPO 风格的组内归一化 advantage（Eqs.~12--14）：

\[
A(x_0^i)=\frac{r_i-\mathrm{mean}(\{r_j\})}{\mathrm{std}(\{r_j\})},
\quad
\mathcal{G}^{+}=\{A>0\},\;\mathcal{G}^{-}=\{A\le 0\},
\quad
w(x_0)=|A(x_0)|.
\]

因为 advantage 零均值（$\sum A=0$），所以 $|A|$ 在正负两侧自然满足 $\sum_{+}|A|=\sum_{-}|A|$，\textbf{$Z(c)$ 干净地抵消}。

\begin{intubox}
\textbf{直觉：} advantage 的 ``零均值'' 本来只是 GRPO 里做 variance reduction 的工程选择；在 DGPO 里它变成 \key{让 DPO 里讨厌的 $Z(c)$ 项自动消掉} 的数学条件。这是非常漂亮的复用。
\end{intubox}

\subsubsection{3.4 最终损失（Eq.~17）}

经 Jensen 不等式把期望挪到 sigmoid 外后，DGPO 的实际训练损失：

\[
\begin{aligned}
\mathcal{L}_{\text{DGPO}}(\theta)
=\;&-\mathbb{E}_{(\mathcal{G}^{+},\mathcal{G}^{-},c),t,\epsilon}
\log\sigma\!\Big(-\lambda_t\beta T\cdot\\
&\quad \Big[
\sum_{x_0\in\mathcal{G}^{+}}w(x_0)\bigl(\mathcal{L}_{\text{dsm}}^{\theta}-\mathcal{L}_{\text{dsm}}^{\theta_{\text{ref}}}\bigr)\\
&\qquad -\sum_{x_0\in\mathcal{G}^{-}}w(x_0)\bigl(\mathcal{L}_{\text{dsm}}^{\theta}-\mathcal{L}_{\text{dsm}}^{\theta_{\text{ref}}}\bigr)
\Big]\Big),
\end{aligned}
\]

其中 $\mathcal{L}_{\text{dsm}}^{\theta}(x,x_t,c)=\|f_\theta(x_t,t,c)-x\|_2^2$ 是 denoising score-matching loss，$\lambda_t$ 是 DPO 里常用的 timestep 加权，\textbf{同组样本共享一个 $\epsilon$} 以减小方差。

\begin{warnbox}
\textbf{形式上很 GRPO 但本质更像 DPO：} sigmoid 的自变量是 ``加权 denoising loss 差''——正样本要比 ref 模型 denoise 得更好（loss 差更负），负样本要比 ref 模型 denoise 得更差。这和 GRPO 的 clipped advantage 已经没有结构关系了，\key{它是 ``组内版的 Diffusion-DPO''}。
\end{warnbox}

\subsubsection{3.5 Timestep Clip：不让模型去学坏样本的伪影}

10 步 ODE rollout 采出来的训练样本会带 blurriness / 细节丢失。若从 $t\in[0,T]$ 均匀采，模型会在小 $t$ 处被迫拟合低质量细节，导致生成质量退化。

DGPO 的处理很直接：训练时 \key{只在 $[t_{\min},T]$ 内采样 $t$}，$t_{\min}>0$，把 ``对齐细节'' 的小 $t$ 区间从训练里排除。

\begin{intubox}
\textbf{直觉：} 用 10-step ODE 生成的样本，\key{大 $t$ 上是 ``语义/构图层面信息干净''，小 $t$ 上是 ``低层细节不干净''}。RL 只需要调语义 / 指令跟随，不需要调低层细节。这个设计和 E-GRPO 的 ``只训 high-entropy early steps'' 从不同角度指向同一结论：后期低 $t$ / 低熵步骤不是 RL 的好战场。
\end{intubox}

\begin{figure}[htbp]
\centering
\includegraphics[width=0.92\linewidth]{dgpo/fig4_timestep_clip_ablation.png}
\caption{Timestep Clip 的视觉消融（论文 Fig.~4）。左列为 SD3.5-M 基线，中列为没有 timestep clip 的 DGPO 变体——可以看到指令跟随仍然 OK，但视觉质量明显退化（画面变 blurry、颜色偏灰）；右列是完整 DGPO。说明 \key{低 $t$ 区间里 RL 会被 rollout 自带的伪影拖着走，必须排除}。图源：\href{https://arxiv.org/html/2510.08425v1}{arXiv HTML Figure 4}。}
\label{fig:dgpo-timestep-clip}
\end{figure}

\subsubsection{3.6 训练流程（Algorithm 1）}

一次 iteration 做的事：

\begin{enumerate}[leftmargin=1.5em]
  \item 从 prompt 集采 $c$；
  \item 用 $p_{\theta^{-}}$ 跑 \textbf{ODE} 采样得 $\mathcal{G}=\{x_0^1,\ldots,x_0^G\}$（$G=24$，10 步 Flow-DPM-Solver）；
  \item 查 reward $r_i=r_\phi(c,x_0^i)$；
  \item 组内归一化得 advantage $A_i$，划分 $\mathcal{G}^{+},\mathcal{G}^{-}$，权重 $w=|A|$；
  \item 采 $t\sim\mathcal{U}[t_{\min},T],\,\epsilon\sim\mathcal{N}(0,I)$，\textbf{同组共享 $\epsilon$}，构造 $x_t^i=\alpha_t x_0^i+\sigma_t\epsilon$；
  \item 在 online 与 ref 模型上算 denoising loss，代入 Eq.~17，一步梯度更新；
  \item 参考模型 $\theta^{-}$：前 200 步直接拷贝 $\theta$，之后走 EMA（decay=0.3）。
\end{enumerate}

\subsection{4. 实验与证据}

\subsubsection{4.1 设置}

\begin{itemize}[leftmargin=1.5em]
  \item \textbf{Backbone：} SD3.5-M（Stable Diffusion 3.5 Medium），LoRA rank=32，分辨率 512。
  \item \textbf{Rollout：} Flow-DPM-Solver，10 步 ODE，$G=24$；$\beta=100$，text dropout 0.05，A100 训练。
  \item \textbf{三任务 / 三 reward：}
    \begin{enumerate}[leftmargin=1.5em]
      \item 组合图像生成：GenEval（objects / counting / position / attributes）。
      \item 视觉文字渲染：``A sign that says `text''' 模板，reward 由 OCR 编辑距离给出 $r=\max(1-N_e/N_{\text{ref}},0)$。
      \item 人类偏好对齐：PickScore。
    \end{enumerate}
  \item \textbf{OOD 指标（DrawBench 上评估，防 reward hacking）：} Aesthetic、DeQA、ImageReward、UnifiedReward、PickScore。
\end{itemize}

\subsubsection{4.2 主结果：GenEval 上超越 GPT-4o}

\begin{center}
\resizebox{\linewidth}{!}{
\begin{tabular}{lccccccc}
\toprule
\textbf{模型} & \textbf{Overall} & \textbf{Single} & \textbf{Two Obj} & \textbf{Counting} & \textbf{Colors} & \textbf{Position} & \textbf{Attr.}\\
\midrule
GPT-4o          & 0.84 & 0.99 & 0.92 & 0.85 & 0.92 & 0.75 & 0.61\\
FLUX.1 Dev      & 0.66 & 0.98 & 0.81 & 0.74 & 0.79 & 0.22 & 0.45\\
SANA-1.5 4.8B   & 0.81 & 0.99 & 0.93 & 0.86 & 0.84 & 0.59 & 0.65\\
SD3.5-M (base)  & 0.63 & 0.98 & 0.78 & 0.50 & 0.81 & 0.24 & 0.52\\
+ Flow-GRPO     & 0.95 & 1.00 & 0.99 & 0.95 & 0.92 & 0.99 & 0.86\\
+ \textbf{DGPO} & \textbf{0.97} & \textbf{1.00} & \textbf{0.99} & \textbf{0.97} & \textbf{0.95} & \textbf{0.99} & \textbf{0.91}\\
\bottomrule
\end{tabular}
}
\end{center}

从 base 0.63 拉到 0.97，\key{比 GPT-4o 高 13 个点}；Position 从 0.24 拉到 0.99 是最戏剧的。

\begin{figure}[htbp]
\centering
\includegraphics[width=0.95\linewidth]{dgpo/fig2_qualitative_vs_flowgrpo.png}
\caption{与 Flow-GRPO 的定性对比（论文 Fig.~2，相同初始噪声）。DGPO 在指令跟随上和 Flow-GRPO 接近，但保留了更好的视觉质量，对应综合表里 \key{in-domain 提升但 OOD 不掉} 的现象。图源：\href{https://arxiv.org/html/2510.08425v1}{arXiv HTML Figure 2}。}
\label{fig:dgpo-qualitative}
\end{figure}

\subsubsection{4.3 三任务综合表：in-domain 赢 Flow-GRPO，OOD 也不掉}

\begin{center}
\resizebox{\linewidth}{!}{
\begin{tabular}{llcccccccc}
\toprule
\textbf{任务} & \textbf{方法} & \textbf{GenEval} & \textbf{OCR} & \textbf{Pick} & \textbf{Aes.} & \textbf{DeQA} & \textbf{ImgRwd} & \textbf{Pick-OOD} & \textbf{UniRwd}\\
\midrule
--        & SD3.5-M    & 0.63 & 0.59 & 21.72 & 5.39 & 4.07 & 0.87 & 22.34 & 3.33\\
组合生成  & Flow-GRPO  & 0.95 & --   & --    & 5.25 & 4.01 & 1.03 & 22.37 & 3.51\\
组合生成  & \textbf{DGPO} & \textbf{0.97} & -- & -- & \textbf{5.31} & \textbf{4.03} & \textbf{1.08} & \textbf{22.41} & \textbf{3.60}\\
文字渲染  & Flow-GRPO  & --   & 0.92 & --    & 5.32 & 4.06 & 0.95 & 22.44 & 3.42\\
文字渲染  & \textbf{DGPO} & -- & \textbf{0.96} & -- & \textbf{5.37} & \textbf{4.09} & \textbf{1.02} & \textbf{22.52} & \textbf{3.48}\\
偏好对齐  & Flow-GRPO  & --   & --   & 23.31 & 5.92 & 4.22 & 1.28 & 23.53 & 3.66\\
偏好对齐  & \textbf{DGPO} & -- & -- & \textbf{23.89} & \textbf{6.08} & \textbf{4.40} & \textbf{1.32} & \textbf{23.91} & \textbf{3.74}\\
\bottomrule
\end{tabular}
}
\end{center}

\begin{summarybox}
\textbf{这张表最值得记的点：} 不是 ``DGPO 在 in-domain 赢了 Flow-GRPO''，而是 \key{Flow-GRPO 在组合生成任务上为了刷 GenEval 牺牲了 Aesthetic（5.39\,→\,5.25）和 DeQA（4.07\,→\,4.01）}，DGPO 则在 GenEval 更高的同时 OOD 指标全部不降反升。这说明 DGPO 的 ``组级偏好 + denoising loss 差'' 损失不那么容易 reward-hack。
\end{summarybox}

\subsubsection{\texorpdfstring{4.4 速度：训练快约 20 倍（GenEval 上近 30$\times$）}{4.4 速度：训练快约 20 倍（GenEval 上近 30x）}}

速度收益来自三点：

\begin{enumerate}[leftmargin=1.5em]
  \item 不需要在整条轨迹上算 transition ratio，而是随机单个 $t$ 做 denoising loss；
  \item ODE rollout 比 SDE rollout 在相同 NFE 下更有信息量；
  \item 不需要 model-agnostic 高斯噪声作为探索源，真正的 ``信号'' 来自 reward 的组内相对差。
\end{enumerate}

\begin{figure}[htbp]
\centering
\includegraphics[width=0.95\linewidth]{dgpo/fig3_training_speed.png}
\caption{Flow-GRPO 与 DGPO 训练速度对比（论文 Fig.~3）。横轴是 wall-clock 时间 / 训练步数，纵轴是各任务的 in-domain reward。DGPO 在 GenEval 上达到 Flow-GRPO 同等水平只需约 1/30 时间，在 OCR 和 PickScore 上也有约 10--20$\times$ 加速。图源：\href{https://arxiv.org/html/2510.08425v1}{arXiv HTML Figure 3}。}
\label{fig:dgpo-training-speed}
\end{figure}

\subsubsection{4.5 消融：哪个设计真的重要}

\begin{center}
\begin{tabular}{ll}
\toprule
\textbf{消融项} & \textbf{结论}\\
\midrule
Timestep Clip       & 去掉 → OCR Acc 0.96\,→\,0.95（微降），但视觉质量显著退化\\
ODE vs.\ SDE rollout & ODE 在收敛速度与最终指标上都显著优于 SDE\\
Offline DGPO        & 用 $p_{\text{ref}}$ 预采数据集能赢 baseline，但显著弱于 online\\
vs.\ Diffusion DPO  & Online/offline 两种设置下 DGPO 都显著优于 DPO\\
\bottomrule
\end{tabular}
\end{center}

\begin{figure}[htbp]
\centering
\includegraphics[width=0.95\linewidth]{dgpo/fig5_text_rendering_variants.png}
\caption{视觉文字渲染任务上各变体的对比（论文 Fig.~5）。依次展示：SD3.5-M 基线、SDE rollout、offline DGPO、Diffusion-DPO 与完整 DGPO。可以看到 \key{ODE rollout 的字体清晰度一致更好}；offline 版本能改善 OCR 但丢视觉质量；Diffusion-DPO 的 pairwise 偏好不足以把文字 ``写准''。图源：\href{https://arxiv.org/html/2510.08425v1}{arXiv HTML Figure 5}。}
\label{fig:dgpo-text-variants}
\end{figure}

\begin{itemize}[leftmargin=1.5em]
  \item \textbf{ODE vs.\ SDE} 这个消融非常重要：它把 ``Flow-GRPO 必须用 SDE'' 的假设实验性地证伪了，\key{SDE 不是方法论必须，只是 policy-gradient 框架的被迫选择}。
  \item \textbf{vs.\ Diffusion DPO} 说明 DGPO 的提升不全靠 ``换成 DPO 式 objective'' 这一步，更来自 \key{``pair → group pair'' 的组级扩展}。
\end{itemize}

\subsection{5. 局限与风险}

\begin{itemize}[leftmargin=1.5em]
  \item \textbf{reward model 决定上限。} DGPO 的 advantage 权重只用了 $|A|$，没有做类似 PCPO / ChunkGRPO 的 multi-objective 组合；如果 reward 模型本身 bias 大，``组内自归一'' 会把这个 bias 等比例放大。
  \item \textbf{ODE rollout 的 ``坏样本伪影'' 仍是隐患。} Timestep Clip 只是缓解，没有根治。如果换一个更少步数的 ODE（比如 4~6 步），$t_{\min}$ 需要重新标定。
  \item \textbf{组划分是 hard 阈值（$A>0$ vs.\ $A\le 0$）。} 在 reward 方差极小的 group（几乎没有信号的 prompt）上，正负划分会变成纯噪声；论文没有像 DAPO 那样给 ``全同组'' 的降权策略。
  \item \textbf{只在 SD3.5-M 文生图上验证。} 作者自己也把 ``T2V'' 列为 future work；视频场景 ODE rollout 成本更高，能否保持 20$\times$ 优势仍待证。
  \item \textbf{KL / CFG 设定复杂。} README 里给了三种训练模式（带/不带 CFG、是否 EMA-ref），不同模式的 $\beta_{\text{dpo}}$ 差一个数量级。调参负担并不小。
\end{itemize}

\subsection{6. 与已有工作的关系}

\begin{center}
\resizebox{\linewidth}{!}{
\begin{tabular}{lcccc}
\toprule
\textbf{方法} & \textbf{需要随机策略} & \textbf{rollout} & \textbf{训练粒度} & \textbf{用到组信息}\\
\midrule
Flow-GRPO / DanceGRPO & 是（SDE）   & SDE（低效） & 整条轨迹 & 是\\
E-GRPO               & 是（SDE）   & ODE/SDE hybrid & consolidated step & 是\\
OP-GRPO              & 是（SDE）   & SDE + replay & 整条轨迹 & 是\\
Diffusion DPO        & 否          & 离线成对     & 随机 timestep & 否（只成对）\\
\textbf{DGPO}        & \textbf{否} & \textbf{ODE（高效）} & \textbf{随机 timestep} & \textbf{是（组级）}\\
\bottomrule
\end{tabular}
}
\end{center}

\begin{itemize}[leftmargin=1.5em]
  \item \textbf{直接前作：} Flow-GRPO（policy-gradient 路线的代表）、Diffusion-DPO（denoising-reward 近似的起点）。DGPO 的身份正好是 ``把两者优势合并''。
  \item \textbf{和 E-GRPO / OP-GRPO 的关系：} 三者都是在解 Flow-GRPO 的效率问题。E-GRPO 改 \key{探索发生在哪些 step}，OP-GRPO 改 \key{采样能否跨 iteration 复用}，DGPO 直接 \key{换整个优化框架}。逻辑上 DGPO 不需要 E-GRPO / OP-GRPO——它不做 SDE，所以 E-GRPO 的 entropy-aware step merging 没有意义；它不需要 logprob ratio，所以 OP-GRPO 的 sequence-level 纠偏也没有必要。
  \item \textbf{真正差异：} DGPO 是 Flow-GRPO 线上第一篇在数学层面 \key{完全脱离 policy gradient} 的方法，但仍然保留 GRPO 最有价值的那部分——\key{组内细粒度相对偏好}。
\end{itemize}

\subsection{7. 对我的可复用点}

\begin{itemize}[leftmargin=1.5em]
  \item \textbf{可复用 idea 1（advantage 零均值做 $Z(c)$ 消元）：} 任何带配分函数的 DPO-like 方法都可以复用这个技巧——只要把权重设计成 ``正负两侧等总权重''，$Z(c)$ 就自动消掉。这个套路可以迁到 video DPO、3D DPO 等场景。
  \item \textbf{可复用 idea 2（timestep clip）：} ``不让 RL 去学 rollout 低质量细节''，这个思路和 E-GRPO 的 high-entropy step training 是同一族观察。它可以成为 \key{所有 few-step ODE rollout 后训练方法的默认预处理}。
  \item \textbf{可复用 idea 3（组共享 $\epsilon$ 做方差缩减）：} 同组样本共享同一条 $\epsilon$，在 DGPO 里是为了让 $\mathcal{L}_{\text{dsm}}^{\theta}-\mathcal{L}_{\text{dsm}}^{\theta_{\text{ref}}}$ 在 prompt+noise 上配对。这个 trick 在 ``任何带 ref 模型的扩散后训练'' 里都可以白用。
  \item \textbf{值得追的 follow-up：}
    \begin{itemize}[leftmargin=1.5em]
      \item DGPO $\times$ 视频生成（T2V / I2V）：作者自己说的 future work。
      \item DGPO $\times$ multi-reward（PCPO / ChunkGRPO 风格的 credit 组合）：DGPO 的组级权重天然可以接多 reward aggregation。
      \item DGPO $\times$ 统一多模态生成后训练（BAGEL / Z-Image / Chameleon 这类）。
    \end{itemize}
\end{itemize}

\subsection{8. 待查问题}

\begin{itemize}[leftmargin=1.5em]
  \item Eq.~17 用 Jensen 不等式把期望移到 sigmoid 外，得到的是 \key{原目标的下界}。这个下界在训练末期是否足够紧？论文有无实证比较 ``用下界'' 和 ``用原始期望内形式''？
  \item $w(x_0)=|A(x_0)|$ 保证了 $\sum_{+}|A|=\sum_{-}|A|$，但如果组内正负样本数量严重失衡（比如 23 条负、1 条正），这一条正样本会拿走所有正权重质量。DGPO 有没有对失衡 group 做保护？
  \item DGPO 在 10 步 ODE 下成立；如果 base 模型本身只是 2\,--\,4 步蒸馏模型，$t_{\min}$ 的标定、ref 模型的 EMA 都可能要重新设计。
  \item 论文只用 LoRA 训了 SD3.5-M，若直接 full-finetune 更大模型（FLUX-12B / SD3.5-L），DGPO 的 20$\times$ 速度优势是否依然保持？
\end{itemize}

\subsection{9. 附录定性对比}

论文附录 B 给了两组完整的定性对比，分别对应 OCR Accuracy 和 PickScore 作为训练 reward 的情况：

\begin{figure}[htbp]
\centering
\includegraphics[width=0.95\linewidth]{dgpo/fig6_appendix_ocr_qual.png}
\caption{以 OCR Accuracy 为训练 reward 的定性对比（论文 Fig.~6，相同初始噪声）。可以看到 DGPO 在 ``写得对的同时画得也好''；Flow-GRPO 虽然写对但画质明显退化（笔画粗糙、背景纹理丢失）。图源：\href{https://arxiv.org/html/2510.08425v1}{arXiv HTML Figure 6}。}
\label{fig:dgpo-appendix-ocr}
\end{figure}

\begin{figure}[htbp]
\centering
\includegraphics[width=0.95\linewidth]{dgpo/fig7_appendix_pickscore_qual.png}
\caption{以 PickScore 为训练 reward 的定性对比（论文 Fig.~7，相同初始噪声）。在主观偏好对齐任务里，DGPO 的整体审美和构图更协调，且 OOD 风格不像 Flow-GRPO 那样收敛到 ``高 PickScore 模板化构图''。图源：\href{https://arxiv.org/html/2510.08425v1}{arXiv HTML Figure 7}。}
\label{fig:dgpo-appendix-pickscore}
\end{figure}

\subsection{参考来源}

\begin{itemize}[leftmargin=1.5em]
  \item \textbf{论文主页：} \href{https://arxiv.org/abs/2510.08425}{Reinforcing Diffusion Models by Direct Group Preference Optimization}（ICLR 2026）。
  \item \textbf{arXiv HTML：} \href{https://arxiv.org/html/2510.08425v1}{arXiv HTML version}（全部 equations、消融与图表均来自此版本）。
  \item \textbf{代码：} \href{https://github.com/Luo-Yihong/DGPO}{Luo-Yihong/DGPO}；已集成进 Flow-Factory。
  \item \textbf{对照阅读：} Flow-GRPO（policy-gradient 代表）、Diffusion-DPO（pairwise 代表）、E-GRPO / OP-GRPO（同样攻击 Flow-GRPO 效率但仍留在 policy gradient 框架内）。
\end{itemize}

% ============================================================
% Paper deep reading: UniSD
% Source: https://arxiv.org/abs/2605.06597
% ============================================================

\section{论文精读：UniSD}

\begin{conceptbox}
\textbf{论文：} \href{https://arxiv.org/abs/2605.06597}{UniSD: Towards a Unified Self-Distillation Framework for Large Language Models}\\
\textbf{作者/机构：} Yiqiao Jin, Yiyang Wang, Lucheng Fu, Haoxin Liu, B.\,A.\,Prakash, Josiah Hester, Srijan Kumar (Georgia Tech)；Yijia Xiao (UCLA)；Yinyi Luo (CMU)；Jindong Wang (William \& Mary)。\\
\textbf{简称：} UniSD / UniSD$^{*}$。\\
\textbf{版本：} arXiv 2605.06597v1（2026-05-07）。\\
\textbf{方向标签：} LLM Self-Distillation、Multi-teacher Agreement、EMA Teacher、Token Contrastive、Feature Matching、Divergence Clipping。\\
\textbf{阅读日期：} 2026-05-09。\\
\textbf{阅读口径：} 这篇不是 RL/RLVR 路线，而是 \textbf{无外部更强教师} 的 LLM adaptation 路线——和我已经记的 POPO 同属 ``rollout 侧靠正样本驱动'' 这条思路，但用的是蒸馏机制而不是策略梯度。
\end{conceptbox}

\subsection{0. 一句话结论}

UniSD 的目标是把已有 self-distillation（SD）方法里 \textbf{各自孤立的} 设计选项，整合进同一个目标函数，做一次系统性的因果研究：哪些组件真的有用、它们各自在哪类任务上贡献最大、能否复合叠加。它给出的统一目标（Eq.~1）覆盖三个关注维度：

\[
\begin{aligned}
\text{UniSD} &= \underbrace{\text{Multi-teacher Agreement}}_{\text{监督可靠性}}
              + \underbrace{\text{EMA Teacher}}_{\text{监督稳定性}}\\
&\quad + \underbrace{\text{Token Contrastive} + \text{Feature Matching}}_{\text{表征对齐}}\\
&\quad + \underbrace{\text{Divergence Clipping}}_{\text{训练稳定性}}.
\end{aligned}
\]

把它们组合成 \key{UniSD$^{*}$} 后，在 Qwen2.5-7B 上比 base 提升 \textbf{+5.4} 总分、比最强 baseline GKD 提升 \textbf{+2.8}，并在三个模型家族（Qwen / Llama / Gemma）的 18 个 model-dataset 配对里 15 个提升、2 个持平、仅 1 个 OOD 退步。

\begin{summarybox}
\textbf{我的判断：} 这篇的真正价值不在 ``又一个 SD 方法''，而在它把 SD 拆成 \key{监督可靠性 / 表征对齐 / 训练稳定性} 三个维度，并给出每个组件 ``主导哪个任务'' 的定性结论。这是一篇可以直接当 \textbf{设计手册} 用的论文。
\end{summarybox}

\begin{figure}[htbp]
\centering
\includegraphics[width=0.95\linewidth]{unisd/fig1_overview.png}
\caption{UniSD 总览（论文 Fig.~1）。把 self-distillation 的设计空间分成 supervision reliability / representation alignment / training stability 三块，对应五个互补组件，最终融合为 UniSD$^{*}$。图源：\href{https://arxiv.org/html/2605.06597v1}{arXiv HTML Figure 1}。}
\label{fig:unisd-overview}
\end{figure}

\subsection{1. 为什么读这篇}

\begin{itemize}[leftmargin=1.5em]
  \item \textbf{和当前主线的关系：} 我已记的 POPO 是 ``RLVR 下不显式用负样本''，UniSD 是 ``self-distillation 下不依赖外部更强教师''。两者都属于 ``LLM 后训练里如何在没有完美 supervision 时仍然学出更好策略'' 这个母问题，互为对照。
  \item \textbf{它可能补上的知识缺口：}
    \begin{itemize}[leftmargin=1.5em]
      \item Self-distillation 在自回归 LLM 上为什么 \textbf{容易不稳定}？
      \item Multi-teacher / EMA / contrastive / feature matching / divergence clipping 这五条 ``各家用过的招'' 在统一 setup 下到底谁有用？
      \item 各组件的收益是叠加的、互补的、还是会互相抵消？
    \end{itemize}
  \item \textbf{我希望读完回答的问题：}
    \begin{itemize}[leftmargin=1.5em]
      \item 在哪类任务（生成密集 / 多选 / 工具使用 / 代码）上，哪一个组件是必装项？
      \item 这个框架能否搬到 \textbf{视觉生成 self-distillation}（FLUX、SD3.5、Flow-matching distillation）？
    \end{itemize}
\end{itemize}

\subsection{2. 问题定义与背景}

\textbf{问题：} 给定一个已经 instruction-tuned 的 LLM $\pi_\theta$，没有更强的外部教师，如何 \textbf{用模型自己的轨迹} 进一步提升它在下游任务上的能力？

\textbf{Self-distillation 的核心难点：} 在自回归 LLM 上，自蒸馏面临三个原生问题：

\begin{enumerate}[leftmargin=1.5em]
  \item \textbf{自由形式轨迹。} 学生生成的 trajectory 没有 ground truth，正确性需要由任务度量决定；
  \item \textbf{表面合理 $\ne$ 可靠。} 看起来 plausible 的 rationale 可能仍然提供 ``似是而非'' 的不稳定监督；
  \item \textbf{孤立设计选择。} 已有 SD 方法（SDFT / GKD / SSD / OPSD 等）各自只调一个轴，没人系统比过。
\end{enumerate}

\textbf{统一目标（Eq.~1）：} UniSD 把所有 SD 变体都框进同一个 token 级散度损失：

\[
\mathcal{L} =
\mathbb{E}_x\,\mathbb{E}_{y\sim\pi_\theta(\cdot\mid x)}\Bigg[
\sum_{t=1}^{T}
\underbrace{m_t}_{\text{token mask}}\,
\underbrace{w_t}_{\text{reliability}}\,
D\!\bigl(\pi_\theta(\cdot\mid x,\hat y_{<t})\,\big\|\,\pi_*^{T}(\cdot\mid x,c,\hat y_{<t})\bigr)
+\lambda_{\text{aux}}\mathcal{L}_{\text{aux}}
\Bigg].
\]

不同 SD 方法的差异，本质上就是 ``如何选 $w_t$ / 怎么定 $\pi_*^{T}$ / 加什么 $\mathcal{L}_{\text{aux}}$''。这一步框架化是这篇论文最重要的贡献。

\begin{figure}[htbp]
\centering
\includegraphics[width=0.95\linewidth]{unisd/fig2_components.png}
\caption{UniSD 五大组件结构图（论文 Fig.~2）。左侧：multi-teacher agreement 通过多个辅助上下文 $c^k$ 给同一个学生 token 打分，统计 $\delta_t$ 当 reliability weight；中部：EMA teacher 给主蒸馏一个慢慢漂移的稳定目标分布；右侧：token-level contrastive、feature matching、divergence clipping 分别从对比学习、表征对齐和训练稳定性角度补充。图源：\href{https://arxiv.org/html/2605.06597v1}{arXiv HTML Figure 2}。}
\label{fig:unisd-components}
\end{figure}

\subsection{3. 方法拆解：五个组件}

\subsubsection{3.1 Multi-Teacher Agreement（监督可靠性）}

\textbf{思路：} 同一个学生轨迹，让 \key{同一个教师在 $K$ 个不同上下文 $c^k$ 下} 各打一次分；上下文之间分歧大的 token 不可靠，分歧小的 token 才是可靠监督。

每个上下文给 token-level log-prob（Eq.~2）：

\[
\ell_t^{k} = \log \pi_k^{T}(\hat y_t\mid x, c^k, \hat y_{<t}),\quad t\in[1,T].
\]

两种聚合方式：

\begin{itemize}[leftmargin=1.5em]
  \item \textbf{Token 级（局部）}：$\delta_t=A(\{\ell_t^{k}\}_k)$，其中 $A$ 是变异度统计量（方差 / 极差）；用 $\delta_t$ 反向作为 $w_t$。
  \item \textbf{Sequence 级（全局）}：$L^{k}=\sum_t m_t\ell_t^{k}/\sum_t m_t$，再做 $\delta_{\text{seq}}=A(\{L^k\}_k)$。
\end{itemize}

\textbf{三种辅助上下文 $c^k$：}

\begin{enumerate}[leftmargin=1.5em]
  \item retrieval-based few-shot（检索式）；
  \item random few-shot（随机式）；
  \item instruction induction（从 train set 归纳出的高层指令）。
\end{enumerate}

\begin{warnbox}
\textbf{为什么不直接训多个教师？} 一是成本，二是发散方向不可控。Multi-teacher agreement 的妙处在于 \key{只用同一个模型，但用上下文制造 diversity}，相当于 ``低成本 ensemble of views''。
\end{warnbox}

\begin{figure}[htbp]
\centering
\includegraphics[width=0.85\linewidth]{unisd/fig4_agree_views.png}
\caption{Token-level Agreement 在三种辅助上下文构造下的对比（论文 Fig.~4）：retrieval-based、random few-shot、instruction induction。可以看到不同任务偏好的 view 不一样：ScienceQA / GPQA / HumanEval 偏 retrieval，ToolAlpaca 偏 induction。图源：\href{https://arxiv.org/html/2605.06597v1}{arXiv HTML Figure 4}。}
\label{fig:unisd-agree-views}
\end{figure}

\subsubsection{3.2 EMA Teacher Stabilization（监督稳定性）}

学生模型自己做教师时，最大的风险是 ``教师的 distribution 跟着学生一起漂''——一旦学生进入 mode collapse，教师也跟着塌缩。EMA 给一个慢动量的稳定锚（Eq.~3）：

\[
\bar\theta_n = \beta\,\bar\theta_{n-1} + (1-\beta)\,\theta_n,\qquad \beta\in[0,1].
\]

教师分布换成 $\pi_{\bar\theta_n}(\cdot\mid x,c^*,\hat y_{<t})$。这个机制和 BYOL / POPO 的 EMA target 是同一族思想——\key{训练时所引用的目标，应该比当前 online 模型更慢变}。

\subsubsection{3.3 Token-Level Contrastive Learning（表征对齐 I）}

构造负样本 $y^-$ 三种方式：让 LLM 自己生成 ``似是而非'' 的错误版、破坏 $y^+$ 中的 reasoning、用 WordNet/PPDB/TextAttack 词级扰动。

打三个分（Eq.~4）：

\[
\ell_t^\theta = \log\pi_\theta(\hat y_t\mid x,\hat y_{<t}),\;
\ell_t^{+}=\log\pi^{T}(\hat y_t\mid x,y^{+},\hat y_{<t}),\;
\ell_t^{-}=\log\pi^{T}(\hat y_t\mid x,y^{-},\hat y_{<t}).
\]

Margin loss（Eq.~5）：

\[
\mathcal{L}_{\text{aux}}(\theta;x,y,c) =
\sum_{t=1}^{T}m_t\,\max\!\bigl(0,\;\gamma + d_t^{+} - d_t^{-}\bigr),
\quad d_t^{\pm}=|\ell_t^\theta-\ell_t^{\pm}|.
\]

要求 $\pi_\theta$ 对 token 的打分 \key{靠近正条件下教师的打分，远离负条件下教师的打分}，余量 $\gamma$。

\subsubsection{3.4 Feature Matching（表征对齐 II）}

直接对齐 hidden state（Eq.~6）：

\[
\mathcal{L}_{\text{feat}}
= \sum_{t=1}^{T}m_t \,\|\mathbf{f}_t^\theta - \mathbf{f}_t^{*}\|_2^{2}.
\]

实现里默认对齐 \key{completion tokens 上的最后一层 hidden states}。可选对象还包括 layer-wise 表征、self-attention 关系、attention-derived features。

\subsubsection{3.5 Divergence Clipping（训练稳定性）}

把 forward-KL 换成加权 JSD（Eqs.~7--8）：

\[
\mathcal{D}_t^{(\alpha)} = \alpha\,D(\pi_*^{T}\|M_t)+(1-\alpha)\,D(\pi_\theta\|M_t),
\quad
M_t=(1-\alpha)\pi_\theta+\alpha\pi_*^{T}.
\]

再硬裁剪：$\widetilde{\mathcal{D}}_t = \min(\mathcal{D}_t^{(\alpha)},\kappa)$，最终蒸馏目标 (Eq.~9):

\[
\mathcal{L}_{\text{clip}}
= \frac{\sum_{t=1}^{T}m_t w_t \widetilde{\mathcal{D}}_t}{\sum_{t=1}^{T}m_t w_t}.
\]

\begin{intubox}
\textbf{直觉：} 学生分布偏离教师分布过远的 token，是 ``超大 gradient + 极不可靠 supervision'' 的源头；clipping 不是在限制学习信号，而是在 \key{限制单 token 的最大破坏力}，这一招在 GRPO/PPO 的 ratio clipping 里也是同源思想。
\end{intubox}

\subsubsection{\texorpdfstring{3.6 完整算法 UniSD$^{*}$}{3.6 完整算法 UniSD*}}

把以上 5 个组件全开：

\begin{enumerate}[leftmargin=1.5em]
  \item 用学生 $\pi_\theta$ 采样轨迹 $\hat y$；
  \item 同一个 frozen 教师在 $K$ 个上下文下打分，得 $w_t$；
  \item 主蒸馏 loss：相对 \textbf{EMA teacher} 的 clipped JSD；
  \item 辅助 loss：token contrastive margin + feature matching；
  \item 总 loss = 主项 $+$ $\lambda_{\text{aux}}$($\mathcal{L}_{\text{contrast}}+\mathcal{L}_{\text{feat}}$)；
  \item 每步更新后做 EMA $\bar\theta\leftarrow\beta\bar\theta+(1-\beta)\theta$。
\end{enumerate}

\subsection{4. 实验与证据}

\subsubsection{4.1 设置}

\begin{itemize}[leftmargin=1.5em]
  \item \textbf{六个基准（4 类任务）：} ScienceQA / GPQA（科学推理；GPQA test-only 做 OOD）、CoS-E（常识）、MBPP / HumanEval（代码；HumanEval 做 OOD）、ToolAlpaca（工具使用）。
  \item \textbf{六个模型 / 三家族：} Qwen2.5-Instruct 0.5B / 1.5B / 3B / 7B；Llama-3.1-8B-Instruct；gemma-3-4b-it。主表用 Qwen2.5-7B。
  \item \textbf{训练：} LoRA rank=64, $\alpha=128$, dropout=0.05；AdamW, lr=2e-5, cosine, 10\% warmup, 1 epoch；bf16；grad accum=4；vLLM colocate sample, $\tau=0.7$；max prompt/completion 3072/1024；6$\times$A100 80GB。
  \item \textbf{评测：} 代码 pass@1（沙盒，10s 超时）；多选 accuracy；ToolAlpaca 完整 accuracy（动作名 + 全部参数 exact match）；解码贪心。
\end{itemize}

\subsubsection{4.2 主表（Qwen2.5-7B，Table~1）}

\begin{center}
\resizebox{\linewidth}{!}{
\begin{tabular}{lccccccc}
\toprule
\textbf{方法} & \textbf{ScienceQA} & \textbf{MBPP} & \textbf{CoS-E} & \textbf{ToolAlpaca} & \textbf{GPQA(OOD)} & \textbf{HumanEval(OOD)} & \textbf{Avg.}\\
\midrule
Raw            & 81.5 & 70.8 & 81.9 & 61.8 & 31.0 & 80.5 & 67.9\\
SFT            & 80.8 & 70.4 & 82.6 & 66.2 & 30.6 & 79.3 & 68.3\\
SDFT           & 81.6 & 71.6 & 81.2 & 73.5 & 34.2 & 78.7 & 70.1\\
GKD            & 81.2 & 72.8 & 81.8 & 72.1 & 33.0 & 82.3 & 70.5\\
SSD            & 80.8 & 72.4 & 79.9 & 55.9 & 33.7 & 81.1 & 67.3\\
OPSD           & 81.2 & 72.8 & 82.0 & 61.8 & 31.5 & 79.9 & 68.2\\
\midrule
Agree (Tok.)   & \textbf{85.2} & 71.2 & \textbf{82.2} & 75.0 & 36.2 & \textbf{83.5} & 72.2\\
Agree (Seq.)   & 84.4 & 73.2 & 81.9 & 76.5 & 35.7 & \textbf{83.5} & 72.5\\
EMA            & 84.3 & 73.2 & 81.4 & \textbf{77.9} & 35.3 & 82.9 & 72.5\\
Contrast       & 83.9 & 73.5 & 82.0 & 75.0 & 34.6 & 82.3 & 71.9\\
Match (Joint)  & 84.8 & 73.2 & 81.8 & 76.5 & 33.5 & 82.9 & 72.1\\
Match (Repr.)  & 83.7 & 73.2 & 81.7 & 72.1 & 35.9 & 82.3 & 71.5\\
Clip           & 82.8 & 73.2 & 81.7 & 70.6 & 31.5 & 81.7 & 70.3\\
\midrule
\textbf{UniSD$^{*}$} & 85.0 & \textbf{74.7} & \textbf{82.2} & \textbf{77.9} & \textbf{36.4} & \textbf{83.5} & \textbf{73.3}\\
\bottomrule
\end{tabular}
}
\end{center}

\textbf{核心数字：} UniSD$^{*}$ 比 Raw +5.4，比最强 baseline GKD +2.8；并且 \key{在所有 6 个 benchmark 上都不比当前最佳的单组件版本差}。

\subsubsection{4.3 跨家族泛化（Fig.~7 / Fig.~12）}

\begin{itemize}[leftmargin=1.5em]
  \item Qwen2.5-7B：+5.4；Llama-3.1-8B：+3.1；Gemma-3-4B：+2.2。
  \item 18 个 model-dataset 配对里，\textbf{15 提升 / 2 持平 / 1 OOD 退步}。
\end{itemize}

\begin{figure}[htbp]
\centering
\includegraphics[width=0.95\linewidth]{unisd/fig7_three_families.png}
\caption{三个家族（Qwen2.5 / Llama-3.1 / Gemma-3）下 UniSD$^{*}$ 相对各自 Raw 的提升（论文 Fig.~7）。Qwen 提升最大（+5.4），Gemma 最小（+2.2）；尺寸更大的模型反而提升更多，符合 ``更强 base 更能消化自蒸馏信号'' 的直觉。图源：\href{https://arxiv.org/html/2605.06597v1}{arXiv HTML Figure 7}。}
\label{fig:unisd-three-families}
\end{figure}

\begin{figure}[htbp]
\centering
\includegraphics[width=0.95\linewidth]{unisd/fig12_per_dataset_gain.png}
\caption{UniSD 各变体相对 Raw Qwen2.5-7B 的逐数据集提升（论文 Fig.~12）。可以看到 \key{ToolAlpaca 是受益最大的任务（EMA 单独就 +16.1）}，而代码 OOD（HumanEval）则是几乎所有变体都能稳定 +2$\sim$+3 的任务；CoS-E 几乎无提升空间。图源：\href{https://arxiv.org/html/2605.06597v1}{arXiv HTML Figure 12}。}
\label{fig:unisd-per-dataset-gain}
\end{figure}

\subsubsection{4.4 Scaling 与 retention（Fig.~3）}

\begin{figure}[htbp]
\centering
\begin{minipage}{0.62\linewidth}
\centering
\includegraphics[width=\linewidth]{unisd/fig3a_scaling.png}
\end{minipage}\hfill
\begin{minipage}{0.34\linewidth}
\centering
\includegraphics[width=\linewidth]{unisd/fig3b_retention.png}
\end{minipage}
\caption{论文 Fig.~3。左：Qwen2.5 在 0.5B/1.5B/3B/7B 四个尺寸下，UniSD 在 ScienceQA 与 GPQA(OOD) 上一致提升，且尺寸越大相对收益越大。右：base-distribution retention perplexity；UniSD$^{*}$ 比 SFT 更接近原始 base 分布（PPL 退化更小），证明它没有把 base 知识冲掉。图源：\href{https://arxiv.org/html/2605.06597v1}{arXiv HTML Figure 3}。}
\label{fig:unisd-scaling-retention}
\end{figure}

\subsubsection{4.5 训练时间 / 损失曲线（Fig.~5）}

\begin{figure}[htbp]
\centering
\begin{minipage}{0.32\linewidth}
\centering
\includegraphics[width=\linewidth]{unisd/fig5a_train_time.png}\\
{\footnotesize (a) 训练时间 vs 准确率}
\end{minipage}\hfill
\begin{minipage}{0.32\linewidth}
\centering
\includegraphics[width=\linewidth]{unisd/fig5b_components_eff.png}\\
{\footnotesize (b) 组件有效性}
\end{minipage}\hfill
\begin{minipage}{0.32\linewidth}
\centering
\includegraphics[width=\linewidth]{unisd/fig5c_loss_curve.png}\\
{\footnotesize (c) 训练 loss 曲线}
\end{minipage}
\caption{论文 Fig.~5（三联）。(a) 显示 UniSD$^{*}$ 用相近训练时间换更高准确率；(b) 各组件单独的有效性；(c) UniSD$^{*}$ 的训练 loss 比单组件方案更平滑稳定，说明 EMA + clipping 的稳定化机制确实生效。图源：\href{https://arxiv.org/html/2605.06597v1}{arXiv HTML Figure 5}。}
\label{fig:unisd-fig5}
\end{figure}

\subsubsection{4.6 分布保留（Fig.~6 / Table~5）}

\begin{itemize}[leftmargin=1.5em]
  \item EMA Teacher 比 SFT 让 retention PPL 降 \textbf{33.9\%}；
  \item UniSD$^{*}$ 的 token-level JSD 比 SFT 低（0.041 vs 0.054），\key{70.3\% 的样本上更接近 base 分布}。
\end{itemize}

\begin{figure}[htbp]
\centering
\includegraphics[width=0.95\linewidth]{unisd/fig6_jsd_distribution.png}
\caption{Base-scored perplexity 与 token-level JSD 分布（论文 Fig.~6）。UniSD$^{*}$ 的分布更靠近 base，说明它在 ``学新任务'' 与 ``不冲掉原有能力'' 之间取得了更好平衡。图源：\href{https://arxiv.org/html/2605.06597v1}{arXiv HTML Figure 6}。}
\label{fig:unisd-jsd}
\end{figure}

\subsubsection{4.7 各组件 ``主导哪个任务''}

\begin{center}
\resizebox{\linewidth}{!}{
\begin{tabular}{lll}
\toprule
\textbf{组件} & \textbf{主导/最佳任务} & \textbf{机理}\\
\midrule
Agree (Tok.)   & ScienceQA / HumanEval / CoS-E & 局部可靠性，峰值最强\\
Agree (Seq.)   & 总分稳健 (72.5)              & 全局保守、噪声更小\\
EMA            & ToolAlpaca (+16.1) / MBPP / HumanEval & 生成密集 / 格式严格任务最受益\\
Contrast       & 6 项均提升（最均衡）         & 区分 plausible-but-wrong\\
Match (Joint)  & 优于单 representation        & 输出 + 隐状态联合\\
Match (Repr.)  & 辅助正则                     & 单独使用收益有限\\
Clip           & 最轻量稳定器（+2.4）          & 防极端 token 破坏\\
\textbf{UniSD$^{*}$} & MBPP/ToolAlpaca/GPQA/HumanEval 最佳；ScienceQA/CoS-E 第二 & 最一致提升\\
\bottomrule
\end{tabular}
}
\end{center}

\begin{summarybox}
\textbf{这张表是这篇论文最值钱的输出。} 它意味着：
\begin{itemize}[leftmargin=1.5em]
  \item 工具使用 / 代码 / 严格格式任务 → \textbf{先开 EMA}；
  \item 多选 / 推理 → \textbf{先开 Agree (Tok.)} 配 retrieval 上下文；
  \item 怕 mode collapse → \textbf{先加 Clip}；
  \item 想要进一步 +1 个点 → \textbf{加 Contrast 或 Joint Match}。
\end{itemize}
这张组件--任务对应表可以直接当工程 checklist 用。
\end{summarybox}

\subsection{5. 局限与风险}

\begin{itemize}[leftmargin=1.5em]
  \item \textbf{超参未给具体值。} 论文没明确给出 $\lambda_{\text{aux}}$、$\beta$、$\alpha$、$\kappa$、$\gamma$ 的实际数值，复现需要照 Appendix A 的 Algorithm 1 反推或自调；
  \item \textbf{Agreement 对上下文构造敏感。} retrieval 在 ScienceQA/GPQA/HumanEval 最强，induction 在 ToolAlpaca 最强——\key{需要按任务挑 view，没有一招通用}；
  \item \textbf{Match (Repr.) 单独不够。} 仅对齐最后一层 hidden 收益有限，要和 contrastive 一起用才划算；
  \item \textbf{Gemma-3 上提升明显较小（+2.2）。} 论文没深入解释，可能跟 instruction-tuning 风格 / token izer / 模型容量都有关；
  \item \textbf{只测了相对小规模数据集。} 训练样本 100$\sim$12K 这一档次；当训练数据扩到 100K+ 时，retention 优势是否会被 ``更长训练'' 抵消还未知；
  \item \textbf{没和 RL/RLVR 路线直接比。} 这篇刻意避开了 RL，所以无法回答 ``UniSD vs.\ POPO / RLVR'' 谁更好。
\end{itemize}

\subsection{6. 与已有工作的关系}

\begin{itemize}[leftmargin=1.5em]
  \item \textbf{直接前作（baseline）：} SFT、SDFT、GKD、SSD、OPSD。UniSD 把它们都放进 Eq.~1 的统一框架里做对照实验；
  \item \textbf{相近思想：}
    \begin{itemize}[leftmargin=1.5em]
      \item \textbf{POPO：} 同样依赖 ``EMA / Siamese 锚 + 表征对齐 + 不显式负样本''，但走的是 RLVR 路线；UniSD 是蒸馏路线下的对偶版本；
      \item \textbf{BYOL / SimSiam：} EMA target + 不对称 predictor 是这套组件的祖师爷；
      \item \textbf{Contrastive Decoding / DoLa：} 这些走的是 inference-time，UniSD 把对比信号搬到 training-time。
    \end{itemize}
  \item \textbf{真正差异：} UniSD 是第一篇 \key{在统一 token-level 散度损失下} 系统比较所有 SD 组件、并给出 ``组件--任务'' 对应表的工作；它的贡献更像 ``benchmark + design guideline'' 而不是 ``又一个新方法''。
\end{itemize}

\subsection{7. 对我的可复用点}

\begin{itemize}[leftmargin=1.5em]
  \item \textbf{可复用 idea 1（Multi-view agreement 替代 multi-teacher ensemble）：} 用同一个模型 + 多种上下文制造 diversity，作为可靠性权重——这套做法可以直接搬到 \textbf{视觉生成 self-distillation}（用不同 prompt rewriting / 不同 CFG scale 当 view）；
  \item \textbf{可复用 idea 2（EMA teacher 防 mode collapse）：} 任何 ``学生即教师'' 的设置都可以加，\textbf{尤其在 ToolAlpaca 这种格式严格任务上 EMA 单独就 +16.1}，这是非常便宜的工程 trick；
  \item \textbf{可复用 idea 3（divergence clipping）：} 把 token 级 KL 换成 clipped JSD，几乎是 ``一个 if 语句'' 的代价但能稳定训练；这一招完全可以加进 \key{Flow-GRPO / DGPO / POPO} 的实现里；
  \item \textbf{可复用实验设计：} ``一张组件--任务热力图'' 的展示方式，直接当模板用；
  \item \textbf{值得追的 follow-up：}
    \begin{itemize}[leftmargin=1.5em]
      \item UniSD $\times$ Flow-Matching distillation（把它的五件套搬到 SD3.5 / FLUX 蒸馏）；
      \item UniSD $\times$ POPO（蒸馏 + 隐式负梯度叠加）；
      \item UniSD 在 100K+ 训练样本规模下的扩展行为。
    \end{itemize}
\end{itemize}

\subsection{8. 待查问题}

\begin{itemize}[leftmargin=1.5em]
  \item Eq.~1 里 $w_t$ 由 agreement 直接给出，但没有保证 $\sum w_t$ 在不同 batch 间稳定——是否做了 batch normalization？
  \item Token contrastive 的负样本 $y^-$ 通过 LLM 自己生成 / WordNet 扰动构造，\textbf{论文没给负样本质量校验}，这部分是否对结果敏感？
  \item EMA teacher 的 $\beta$ 在 1 epoch 训练里大小如何取？$\beta$ 太大 EMA 跟不上，太小又失去稳定作用；
  \item Gemma-3 上提升明显较小，这是 base 模型问题、tokenizer 问题，还是 retrieval-view 在 Gemma 上覆盖不够？
  \item UniSD$^{*}$ 是否会在 100K 训练样本规模下退化成 ``重新发现 SFT''？这是 self-distillation 类方法的共性疑问。
\end{itemize}

\subsection{参考来源}

\begin{itemize}[leftmargin=1.5em]
  \item \textbf{论文主页：} \href{https://arxiv.org/abs/2605.06597}{UniSD: Towards a Unified Self-Distillation Framework for Large Language Models}。
  \item \textbf{arXiv HTML：} \href{https://arxiv.org/html/2605.06597v1}{arXiv HTML version}（全部公式 Eq.~1--9、Table~1、Fig.~1--12 均来自此版本）。
  \item \textbf{对照阅读：} POPO（RLVR 下不显式用负样本）；BYOL / SimSiam（EMA target + 不对称 predictor 的祖师爷）；GKD / SDFT / SSD / OPSD（被它统一进 Eq.~1 的所有 SD 前作）。
\end{itemize}

% ============================================================
% Paper deep reading: OPSD
% Source: https://arxiv.org/abs/2601.18734
% ============================================================

\section{论文精读：OPSD}

\begin{conceptbox}
\textbf{论文：} \href{https://arxiv.org/abs/2601.18734}{Self-Distilled Reasoner: On-Policy Self-Distillation for Large Language Models}\\
\textbf{作者/机构：} Siyan Zhao (UCLA / Meta), Zhihui Xie (HKU), Mengchen Liu, Jing Huang, Guan Pang, Feiyu Chen (Meta Superintelligence Labs), Aditya Grover (UCLA)。\\
\textbf{简称：} OPSD / On-Policy Self-Distillation。\\
\textbf{版本：} arXiv 2601.18734v3；ICML 2026。\\
\textbf{代码：} \href{https://github.com/siyan-zhao/OPSD}{siyan-zhao/OPSD}。\\
\textbf{方向标签：} LLM Reasoning Post-training、On-policy Distillation、Self-distillation、Privileged Context、Dense Token-level Supervision、RLVR Alternative。\\
\textbf{阅读日期：} 2026-05-14。\\
\textbf{阅读口径：} 这篇不是在改 GRPO 的 advantage、sampling 或 reward，而是在 LLM reasoning 后训练里提出一条 \key{介于 SFT 和 RLVR 之间} 的路线：利用带标准解的数据集，把同一个模型拆成 ``能看标准解的 teacher'' 和 ``只能看题目的 student''，在 student 自己生成的 on-policy 轨迹上做 token-level soft distillation。
\end{conceptbox}

\subsection{0. 一句话结论}

OPSD 的核心想法可以概括成：\textbf{让模型带着标准答案做一次 ``会解题的自己''，去监督不看答案的 ``正常推理自己''}。形式上，它保留了 RLVR 的 on-policy 轨迹分布，又用 privileged teacher distribution 提供了比 binary reward 更密的 token-level 学习信号。

\[
\begin{aligned}
\text{OPSD}
&= \text{student on-policy rollout}\\
&\quad + \text{same-model privileged teacher}\\
&\quad + \text{full-vocabulary per-token divergence}\\
&\quad + \text{pointwise KL clipping for stability}.
\end{aligned}
\]

\begin{summarybox}
\textbf{我的判断：} OPSD 最重要的贡献不是 ``自蒸馏'' 这个名字，而是把 LLM reasoning 后训练里的三种信号重新组合：\key{SFT 的标准解信息}、\key{RLVR 的 on-policy 分布}、\key{distillation 的 dense soft targets}。它适合作为 RLVR 前的 reasoning warm-up，或在 RLVR reward variance / rollout cost 太高时作为更便宜的替代阶段。
\end{summarybox}

\subsection{1. 在 LLM 后训练中的定位与所处 stage}

\subsubsection{1.1 它处在什么 stage}

如果把一个 reasoning LLM 的训练流程粗略拆成：

\[
\text{Pretraining}
\rightarrow
\text{Instruction / SFT}
\rightarrow
\text{Reasoning specialization}
\rightarrow
\text{RLVR / preference alignment}
\rightarrow
\text{test-time scaling},
\]

OPSD 位于 \key{Reasoning specialization / post-SFT reasoning post-training} 阶段。它通常接在已经 instruction-tuned 的模型之后，例如论文使用 Qwen3 Instruct 系列作为 base，然后用 OpenThoughts 的数学 reasoning problem-solution pairs 做 LoRA 后训练。

更具体地说，它不是以下几类：

\begin{itemize}[leftmargin=1.5em]
  \item \textbf{不是预训练：} 不学通用语言建模分布，而是针对数学 reasoning 数据做能力增强。
  \item \textbf{不是传统 SFT：} 不直接 teacher-forcing 标准解轨迹，而是在 student 自己采样出的轨迹上对齐 teacher 分布，因此是 on-policy。
  \item \textbf{不是标准 RLVR：} 不用最终答案正确性给二值 reward，也不做 group advantage / clipped policy ratio；它的训练信号是 token-level distribution divergence。
  \item \textbf{不是无监督 self-improvement：} 需要数据集给出 ground-truth answer 或 reference chain-of-thought 作为 privileged information。
\end{itemize}

\begin{summarybox}
\textbf{Stage 定位：} OPSD 是一种 \key{有标准解数据驱动的 on-policy soft self-distillation stage}。在工程 pipeline 里，它可以放在 ``SFT 之后、RLVR 之前''，用于低成本提升 reasoning；也可以在一些任务上直接替代早期 RLVR，因为它避免了多 rollout、稀疏 reward 和 zero-variance group 的问题。
\end{summarybox}

\subsubsection{1.2 它和 SFT / RLVR / on-policy distillation 的关系}

\begin{center}
\resizebox{\linewidth}{!}{
\begin{tabular}{lcccc}
\toprule
\textbf{方法} & \textbf{on-policy 数据} & \textbf{dense token 信号} & \textbf{需要外部 teacher} & \textbf{需要 outcome reward}\\
\midrule
SFT / off-policy distillation & 否 & 是 & 视数据来源而定 & 否\\
GRPO / RLVR & 是 & 否 & 否 & 是\\
普通 on-policy distillation & 是 & 是 & 是 & 否\\
\textbf{OPSD} & \textbf{是} & \textbf{是} & \textbf{否} & \textbf{否}\\
\bottomrule
\end{tabular}
}
\end{center}

从这个表看，OPSD 的位置很清楚：它试图拿到 GRPO 的 on-policy 分布优势，同时保留 distillation 的 dense feedback，但不引入额外大 teacher。

\subsubsection{1.3 它在当前 post-training 版图里的角色}

在 LLM 后训练里，OPSD 可以理解为一条 ``非 RL 但 on-policy'' 的 reasoning alignment 路线：

\begin{itemize}[leftmargin=1.5em]
  \item \textbf{相对 SFT：} OPSD 解决 exposure bias，因为学生真的在自己的前缀上接受监督，而不是只在标准答案前缀上训练。
  \item \textbf{相对 GRPO：} OPSD 解决稀疏 reward 和 zero reward-std 问题，因为即使 student 最终答错，每个 token 仍有 teacher distribution 可学。
  \item \textbf{相对外部 teacher distillation：} OPSD 省掉更强 teacher 模型，但要求当前模型在看到标准解后能 rationalize，也就是 ``评估 / 理解答案'' 比 ``独立生成答案'' 更容易。
  \item \textbf{相对 STaR / ReST：} OPSD 不把成功轨迹筛出来再 SFT，而是在所有 student rollout 上做 soft token distribution matching，学习信号更密。
\end{itemize}

\begin{intubox}
\textbf{一句话定位：} OPSD 是 ``solution-conditioned self-teaching''。它不是让模型凭空自我提升，而是让模型借助 reference solution 形成一个 privileged policy，再把这个 policy 的 token distribution 蒸馏回正常 inference policy。
\end{intubox}

\subsection{2. 为什么读这篇}

\begin{itemize}[leftmargin=1.5em]
  \item \textbf{和当前主线的关系：} 当前 post-training 笔记里已有 GRPO/RLVR、positive-only rollout、自蒸馏统一框架等方向。OPSD 是 LLM reasoning 里 \textbf{用 on-policy distillation 替代稀疏 RLVR 信号} 的代表工作。
  \item \textbf{它可能补上的知识缺口：}
    \begin{itemize}[leftmargin=1.5em]
      \item 如何同时利用标准解数据和 on-policy rollout？
      \item 没有外部大 teacher 时，self-distillation 的 teacher 从哪里来？
      \item Dense token-level feedback 是否真的比 binary outcome reward 更 token-efficient？
    \end{itemize}
  \item \textbf{我希望读完回答的问题：}
    \begin{itemize}[leftmargin=1.5em]
      \item OPSD 和 GRPO 到底谁在学 ``更正确的 reasoning''，谁只是更高效？
      \item 它能不能作为 RLVR 的前置阶段？
      \item 它的 ``teacher = same model + privileged context'' 是否真的足够可靠？
    \end{itemize}
\end{itemize}

\subsection{3. 问题定义与背景}

\textbf{问题：} 给定 reasoning 数据集

\[
\mathcal{S}=\{(x_i,y_i^\star)\}_{i=1}^N,
\]

其中 $x_i$ 是题目，$y_i^\star$ 是标准答案或 reference chain-of-thought。目标是提升一个已经 instruction-tuned 的 LLM 在数学 reasoning benchmark 上的表现，同时降低 RLVR 需要多 rollout 和稀疏 reward 带来的 token cost。

\textbf{SFT 的限制：} SFT 只在 reference trajectory 上 teacher forcing，训练时前缀永远是标准解前缀；但推理时模型会遇到自己生成的错误或偏移前缀。这就是 exposure bias。论文还观察到，在 OpenThoughts 这类数据上，标准解风格较精简，直接 SFT 可能缩短测试时 reasoning length，导致性能下降。

\textbf{RLVR / GRPO 的限制：}

\begin{enumerate}[leftmargin=1.5em]
  \item \textbf{采样贵。} GRPO 对每个 prompt 采样一组 responses，例如论文设置里是 $8$ 条、每条最多 $16k$ tokens。
  \item \textbf{reward 稀疏。} 只有最终答案正确 / 错误，所有 token 共享同一个 sequence-level advantage。
  \item \textbf{reward std collapse。} 如果一组 responses 全对或全错，组内 reward 方差为零，advantage 消失，采样预算被浪费。
\end{enumerate}

\textbf{外部 teacher on-policy distillation 的限制：} 它解决了 on-policy 和 dense supervision，但需要另一个通常更强、更大的 teacher LLM；这在成本和闭源依赖上都不理想。

\begin{intubox}
\textbf{直觉版：} GRPO 像让学生反复考试，只告诉最后对错；SFT 像让学生照抄标准答案；OPSD 像让学生先自己写解法，再让 ``看过标准答案的自己'' 在每一步告诉他下一步分布应该更像什么。
\end{intubox}

\subsection{4. 方法拆解}

\subsubsection{4.1 同一个模型，两种上下文：student policy 与 teacher policy}

OPSD 从同一个 LLM $p_\theta$ 构造两个 policy：

\[
 p_S(\cdot\mid x) \triangleq p_\theta(\cdot\mid x),
\]

\[
 p_T(\cdot\mid x,y^\star) \triangleq p_\theta(\cdot\mid x,y^\star).
\]

其中 student 只看到题目 $x$，这与测试时一致；teacher 额外看到标准解 $y^\star$，因此是 privileged policy。两者参数相同，但 conditioning context 不同。

论文的 prompt 设计很关键：teacher prompt 不是简单把标准解塞进去要求复述，而是先给出 reference solution，然后要求模型 ``理解 reference solution 后用自己的方式再解一遍''。不过 teacher 实际训练时不需要生成完整答案；它只是在 student rollout 的每个 prefix 上做一次 forward pass，输出 next-token distribution。

\subsubsection{4.2 On-policy student rollout}

对每个题目 $x$，先从 student policy 采样一条回答：

\[
\hat{y}=(\hat{y}_1,\ldots,\hat{y}_{|\hat{y}|})\sim p_S(\cdot\mid x).
\]

这一步是 OPSD 区别于 SFT 的关键：训练轨迹来自 student 当前会访问的分布，而不是固定 reference solution。

在每个位置 $n$，student 和 teacher 都在同一个 student prefix $\hat{y}_{<n}$ 上给出 next-token distribution：

\[
 p_S(y_n\mid x,\hat{y}_{<n}),
\qquad
 p_T(y_n\mid x,y^\star,\hat{y}_{<n}).
\]

因此 teacher 并不是监督 ``标准答案的下一个 token''，而是监督 ``在 student 已经写到这里的情况下，如果我知道标准解，我会怎样分布式地修正下一步''。

\subsubsection{4.3 Full-vocabulary token-level divergence objective}

训练目标是在 student 自己的 trajectory 上最小化 teacher 与 student 的 token-level divergence：

\[
\mathcal{L}_{\mathrm{OPSD}}(\theta)
=
\mathbb{E}_{(x,y^\star)\sim\mathcal{S}}
\mathbb{E}_{\hat{y}\sim p_S(\cdot\mid x)}
\left[
\frac{1}{|\hat{y}|}
\sum_{n=1}^{|\hat{y}|}
D\!\left(
 p_T(\cdot\mid x,y^\star,\hat{y}_{<n})
 \Vert
 p_S(\cdot\mid x,\hat{y}_{<n})
\right)
\right].
\]

论文比较了 forward KL、reverse KL 和 generalized JSD。主实验最后采用 forward KL：

\[
D_{\mathrm{KL}}(p_T\Vert p_S)
=
\sum_{v\in\mathcal{V}}
 p_T(v)\log\frac{p_T(v)}{p_S(v)}.
\]

\begin{summarybox}
\textbf{为什么 forward KL 更合理：} teacher distribution 代表 ``看过标准解后，哪些 next tokens 都是合理的''。Forward KL 会惩罚 student 漏掉 teacher 认为可能的 token，因此更偏覆盖 teacher support；reverse KL 更 mode-seeking，容易只追一个高概率 token，在 reasoning 里可能丢掉替代推理路径。
\end{summarybox}

\subsubsection{4.4 梯度只走 student logits}

虽然定义上 teacher 和 student 来自同一个模型，训练时梯度只回传到 student logits，teacher distribution 被当作 target。论文还在实现细节里提到，为了稳定训练，他们把 teacher policy 固定为 initial policy，而不是持续更新的 current policy。

这个实现细节很重要，因为它把 ``同一模型自蒸馏'' 变成：

\[
\text{initial model + privileged solution context}
\quad\longrightarrow\quad
\text{updated student without privileged context}.
\]

\begin{warnbox}
\textbf{需要注意的细节：} OPSD 的 ``self'' 并不等于完全闭环自举。它仍然依赖外部标准解，并且主实验里 teacher 是固定 initial policy 的 privileged-context distribution。这更像是 ``同源模型 + 标准解上下文'' 的 privileged distillation，而不是 teacher/student 同步演化。
\end{warnbox}

\subsubsection{4.5 Per-token pointwise KL clipping}

论文发现 token-level divergence 非常 heavy-tailed：一些风格 token，例如 ``wait''、``think''、``therefore''、``maybe'' 等，可能产生远高于数学 token 的 KL 值。这样会导致训练信号被 style pattern 主导，而不是被真正的 math reasoning token 主导。

对一般 $f$-divergence，在位置 $n$ 和词表项 $v$ 上的 pointwise contribution 为：

\[
\ell^{(f)}_{n,v}
=
 p_T(v\mid\cdot)
 f\!\left(
 \frac{p_S(v\mid\cdot)}{p_T(v\mid\cdot)}
 \right).
\]

OPSD 对每个词表项的贡献做 clipping：

\[
D_{\mathrm{clip}}^{(f)}(p_T\Vert p_S)
=
\frac{1}{|\hat{y}|}
\sum_{n=1}^{|\hat{y}|}
\sum_{v\in\mathcal{V}}
\min(\ell^{(f)}_{n,v},\tau).
\]

\begin{intubox}
\textbf{直觉：} 如果不 clip，模型可能主要学会 ``看过答案的老师更喜欢哪些思考语气词''，而不是 ``哪一步数学推理该怎么转向''。Pointwise clipping 相当于限制每个单词能贡献的最大梯度，让信号更平均地落到数学 token 上。
\end{intubox}

\subsubsection{4.6 为什么 full-vocabulary logit distillation 优于 sampled-token policy gradient}

论文也讨论了一个替代目标：只在 student 实际采样出来的 token $\hat{y}_n$ 上计算 teacher/student log-prob 差，并把它当成 token-level advantage：

\[
A_n(x,\hat{y})
=
\log p_T(\hat{y}_n\mid x,y^\star,\hat{y}_{<n})
-
\log p_S(\hat{y}_n\mid x,\hat{y}_{<n}).
\]

再用 policy-gradient 风格更新：

\[
\mathcal{L}(\theta)
=
-
\mathbb{E}_{(x,y^\star),\hat{y}}
\left[
\frac{1}{|\hat{y}|}
\sum_{n=1}^{|\hat{y}|}
A_n(x,\hat{y})
\log p_S(\hat{y}_n\mid x,\hat{y}_{<n})
\right],
\]

其中 $A_n$ 停梯度。这个目标很像 dense-reward policy gradient，但它只监督被采样到的 token。Full-vocabulary objective 则监督整个 next-token distribution，因此信息更丰富，但显存更贵。

\begin{summarybox}
\textbf{关键取舍：} sampled-token objective 更像 RL，便宜但信息少；full-vocabulary distillation 更像 KL imitation，贵但更稳。论文实验显示 full-vocabulary 版本在 AIME25 / HMMT25 上更好。
\end{summarybox}

\subsection{5. 算法流程}

\subsubsection{5.1 用一道小算术题走完整个流程}

可以把 OPSD 想成一个很具体的教学场景：\key{学生先按自己的方式做题，老师手里有标准答案，但老师不直接让学生照抄标准解，而是在学生自己的每一步草稿上告诉他下一步 token 分布应该往哪里偏。}

假设训练集中有一道题：

\begin{quote}
小明买了 4 本笔记本，每本 3 元；又买了 5 支笔，每支 2 元。一共花了多少钱？
\end{quote}

标准解 $y^\star$ 是：

\begin{quote}
$4\times 3=12$，$5\times 2=10$，$12+10=22$，所以一共花了 22 元。
\end{quote}

在 OPSD 里，同一个模型被放进两种上下文：

\begin{itemize}[leftmargin=1.5em]
  \item \textbf{Student：} 只看到题目 $x$，这和真实测试时一致。
  \item \textbf{Teacher：} 看到题目 $x$ 和标准解 $y^\star$，因此它拥有 privileged information。
\end{itemize}

第一步不是让 student 抄标准解，而是先让 student 自己生成一条 on-policy 回答。比如 student 可能写出：

\begin{quote}
4 本笔记本是 $4\times 3=12$ 元。5 支笔是 $5\times 3=15$ 元。所以总共是 27 元。
\end{quote}

这条轨迹里 student 把笔的单价从 2 元错看成了 3 元。OPSD 的关键就在这里：它不会把这条错误轨迹直接丢掉，也不会只给一个 ``错了'' 的 sequence-level reward，而是在 student 自己真的走到的前缀上进行纠偏。

例如 student 已经生成到前缀：

\begin{quote}
5 支笔是 $5\times$
\end{quote}

此时 student 只看题目，next-token distribution 可能更偏向错误 token：

\begin{center}
\begin{tabular}{lcc}
\toprule
\textbf{候选 next token} & \textbf{Student 概率} & \textbf{Teacher 概率}\\
\midrule
$3$ & 0.55 & 0.08\\
$2$ & 0.25 & 0.78\\
$5$ & 0.08 & 0.04\\
其他 & 0.12 & 0.10\\
\bottomrule
\end{tabular}
\end{center}

Teacher 因为看过标准解，知道这里应该乘以 2，所以它在同一个 student prefix 上给出更偏向 $2$ 的分布。训练目标就是让 student 的分布更接近 teacher 的分布。注意 teacher 不是重新生成一条标准答案轨迹，而是在 student 已经写到的错误前缀上判断：\key{如果我知道标准解，那么下一步更合理的 token 分布是什么？}

同样地，当 student 后面写到：

\begin{quote}
所以总共是
\end{quote}

teacher 会让后续分布更靠近 22，而不是 27。于是 OPSD 的监督不是只发生在最终答案，也不是只发生在标准解 token 上，而是发生在 student rollout 的每一个 prefix 上。

\begin{intubox}
\textbf{这个例子的记忆点：} SFT 像老师把标准解发下来让学生照抄；GRPO 像考试后只告诉学生整题对错；OPSD 像老师拿着标准解，沿着学生自己的草稿逐步指出 ``下一步更应该往哪个 token 方向走''。所以它既保留了 on-policy 轨迹，又提供了 dense token-level soft correction。
\end{intubox}

用流程图表示就是：

\begin{center}
\resizebox{0.92\linewidth}{!}{%
\begin{tikzpicture}[
  node distance=0.9cm and 1.2cm,
  box/.style={draw, rounded corners, align=center, inner sep=6pt, font=\small},
  arrow/.style={-{Latex[length=2mm]}, thick}
]
  \node[box] (data) {训练数据\\题目 $x$ + 标准解 $y^\star$};
  \node[box, below left=of data] (student) {Student\\只看题目 $x$};
  \node[box, below right=of data] (teacher) {Teacher\\看题目 $x$ + 标准解 $y^\star$};
  \node[box, below=of student] (rollout) {Student 自己生成回答 $\hat y$\\例如走到错误前缀 $5\times$};
  \node[box, below=of rollout] (prefix) {取 $\hat y$ 的每个 prefix};
  \node[box, below left=of prefix] (sdist) {Student next-token distribution};
  \node[box, below right=of prefix] (tdist) {Teacher next-token distribution};
  \node[box, below=1.0cm of prefix] (loss) {比较两个分布\\forward KL + clipping};
  \node[box, below=of loss] (update) {更新 student\\下次更倾向正确推理方向};

  \draw[arrow] (data) -- (student);
  \draw[arrow] (data) -- (teacher);
  \draw[arrow] (student) -- (rollout);
  \draw[arrow] (rollout) -- (prefix);
  \draw[arrow] (prefix) -- (sdist);
  \draw[arrow] (prefix) -- (tdist);
  \draw[arrow] (teacher) |- (tdist);
  \draw[arrow] (sdist) -- (loss);
  \draw[arrow] (tdist) -- (loss);
  \draw[arrow] (loss) -- (update);
\end{tikzpicture}%
}
\end{center}

\subsubsection{5.2 抽象成训练 iteration}

一次 OPSD 训练 iteration 可以写成：

\begin{enumerate}[leftmargin=1.5em]
  \item 从 reasoning dataset 中采样 minibatch $B=\{(x,y^\star)\}$。
  \item 对每个题目，用 student prompt 采样 on-policy response $\hat{y}\sim p_S(\cdot\mid x)$。
  \item 对同一条 student trajectory，在每个 prefix $\hat{y}_{<n}$ 上分别计算：
    \begin{itemize}[leftmargin=1.5em]
      \item student distribution $p_S(\cdot\mid x,\hat{y}_{<n})$；
      \item privileged teacher distribution $p_T(\cdot\mid x,y^\star,\hat{y}_{<n})$。
    \end{itemize}
  \item 对每个 token position 计算 full-vocabulary forward KL，并做 pointwise clipping。
  \item 平均所有 token loss，梯度只更新 student / online model 参数。
\end{enumerate}

实际主实验配置：

\begin{itemize}[leftmargin=1.5em]
  \item \textbf{Backbone：} Qwen3-1.7B / 4B / 8B Instruct。
  \item \textbf{训练数据：} OpenThoughts 数学 reasoning 子集，最多 30K problem-solution pairs。
  \item \textbf{训练方式：} LoRA，rank 64，alpha 128，目标模块包括 attention projection 与 MLP projection。
  \item \textbf{OPSD rollout：} 每题 1 条 response，max completion length 1024，temperature 1.1。
  \item \textbf{GRPO baseline：} 每题 8 条 response，每条最多 16k tokens，temperature 1.2，最多 500 steps。
  \item \textbf{OPSD 训练步数：} 100 steps；README 中 Qwen3-1.7B 复现实验约 4×H100 上 15 分钟。
  \item \textbf{评测：} AIME 2024、AIME 2025、HMMT 2025；Avg@12；temperature 1.0，thinking mode enabled，max new tokens 38912。
\end{itemize}

\subsection{6. 实验与证据}

\subsubsection{6.1 主结果：OPSD 在 1.7B 上收益最大，整体匹配或超过 GRPO}

论文表 2 的核心结果如下：

\begin{center}
\resizebox{\linewidth}{!}{
\begin{tabular}{llcccc}
\toprule
\textbf{模型} & \textbf{方法} & \textbf{AIME24} & \textbf{AIME25} & \textbf{HMMT25} & \textbf{Average}\\
\midrule
Qwen3-8B & Base & 75.8 & 65.6 & 43.9 & 61.8\\
Qwen3-8B & SFT & 72.3 & 64.2 & 42.9 & 59.8\\
Qwen3-8B & GRPO & 76.4 & 68.9 & \textbf{46.7} & 64.0\\
Qwen3-8B & \textbf{OPSD} & \textbf{77.8} & \textbf{70.8} & 45.8 & \textbf{64.8}\\
\midrule
Qwen3-4B & Base & 74.9 & 66.4 & 42.2 & 61.2\\
Qwen3-4B & SFT & 70.2 & 62.3 & 43.4 & 58.6\\
Qwen3-4B & GRPO & 75.6 & 68.1 & 44.4 & 62.7\\
Qwen3-4B & \textbf{OPSD} & \textbf{76.4} & \textbf{68.3} & \textbf{46.1} & \textbf{63.6}\\
\midrule
Qwen3-1.7B & Base & 51.5 & 36.7 & 23.1 & 37.1\\
Qwen3-1.7B & SFT & 48.4 & 36.3 & 22.7 & 35.8\\
Qwen3-1.7B & GRPO & 51.1 & 38.3 & 23.7 & 37.7\\
Qwen3-1.7B & \textbf{OPSD} & \textbf{57.2} & \textbf{43.9} & \textbf{29.2} & \textbf{43.4}\\
\bottomrule
\end{tabular}
}
\end{center}

最有信息量的是 Qwen3-1.7B：OPSD 平均从 37.1 拉到 43.4，明显超过 GRPO 的 37.7。8B 上 OPSD 平均也超过 GRPO，但 HMMT25 单项略低于 GRPO。

\begin{summarybox}
\textbf{我会这样读结果：} OPSD 对较小模型特别有价值，因为小模型从 sparse binary reward 中很难稳定学到东西，但它在看到标准解后仍能产生有用的 privileged distribution。模型越大，GRPO 也能通过更多 rollout 找到有效 reward signal，两者差距缩小。
\end{summarybox}

\subsubsection{6.2 Token efficiency：OPSD 的核心优势}

论文强调 OPSD 只需要每题 1 条、最多 1024 token 的 rollout；GRPO 则需要每题 8 条、每条最多 16k token。即便 GRPO 生成更多 token，它得到的仍是最终 binary reward；OPSD 在每个 token 位置都有 teacher distribution。

更重要的是，GRPO 在训练后期会出现 reward diversity collapse：很多 batch 内所有 samples 都全对或全错，组内 reward 标准差为 0，advantage 归零，导致没有梯度。OPSD 不依赖组内 reward 方差，因此不会在这种情况下浪费采样。

\begin{warnbox}
\textbf{这里的比较要谨慎：} OPSD 使用了 reference solution 作为 privileged information；GRPO 使用的是 outcome verifier。两者不是完全同一种 supervision budget。OPSD 的 token efficiency 部分来自更强、更密的数据条件，而不只是算法更省。
\end{warnbox}

\subsubsection{6.3 SFT 为什么反而下降}

表 2 里 SFT 在三个模型尺度上平均都下降：8B 从 61.8 到 59.8，4B 从 61.2 到 58.6，1.7B 从 37.1 到 35.8。论文解释是 OpenThoughts 标准解较简洁，直接 SFT 会改变模型 test-time reasoning style，使它生成更短、更少探索的推理。

OPSD 不直接模仿标准解文本，而是让 teacher 在 student 自己轨迹上提供 distribution-level guidance，因此可以把 concise solution 转化成 dense learning signal，而不是直接把 concise style 强塞给 student。

\subsection{7. 消融：哪些设计真的重要}

\subsubsection{7.1 Forward KL 最好}

在 Qwen3-1.7B 的 AIME25 上，三种 divergence 的结果：

\begin{center}
\begin{tabular}{lccc}
\toprule
\textbf{Objective} & \textbf{Base} & \textbf{Step 50} & \textbf{Step 100}\\
\midrule
Forward KL $D_{\mathrm{KL}}(p_T\Vert p_S)$ & 36.7 & \textbf{43.9} & \textbf{41.1}\\
Reverse KL $D_{\mathrm{KL}}(p_S\Vert p_T)$ & 36.7 & 37.5 & 35.0\\
JSD $(\beta=0.5)$ & 36.7 & 36.9 & 39.0\\
\bottomrule
\end{tabular}
\end{center}

Forward KL 在 step 50 达到最佳，并在 step 100 仍高于 base；reverse KL 甚至会退化。

\subsubsection{7.2 Thinking mode 组合：student off / teacher on}

Qwen3 支持 thinking mode。论文比较四种 student/teacher 组合后发现：\key{student thinking mode off + teacher thinking mode on} 能让 math tokens 的 KL 信号最大，因此主实验采用这个组合。

这件事很有意思：student off 让 rollout 更短、更直接，teacher on 则让 privileged policy 有更强的 reasoning/rationalization 能力。于是 student 被迫在较简洁的前缀上对齐一个更会展开推理的 teacher distribution。

\subsubsection{7.3 Pointwise KL clipping 防止 collapse}

Style tokens 的 KL contribution 可能显著高于 math tokens，导致训练偏向风格模仿。Pointwise clipping 在 AIME24 曲线上防止了训练后期 performance collapse。

\begin{intubox}
\textbf{我认为这是 OPSD 最可复用的 trick：} 任何 full-vocabulary distillation 都可能被少数高 KL token 支配。把 divergence contribution 从 ``整个 token position 裁剪'' 改成 ``每个 vocabulary entry 裁剪''，更细粒度，也更适合 LLM 大词表。
\end{intubox}

\subsubsection{7.4 生成长度：更长不一定更好}

OPSD 的 loss 按 token 求和，看上去更长 rollout 会带来更多 supervision。但论文比较 1024 与 4096 token 后发现，更长并不稳定提升。

解释是：早期 tokens 往往是 reasoning 分叉点，对最终解题路径更关键；当 prefix 很长后，teacher 在看到大量 student prefix 后对下一 token 的纠偏空间变小，后续 token 的学习信号可能变弱甚至变噪。

\subsubsection{7.5 Full-vocabulary objective 优于 sampled-token distillation}

在 Qwen3-4B、2048 generation length 下：

\begin{center}
\begin{tabular}{lcc}
\toprule
\textbf{方法变体} & \textbf{AIME25 pass@8} & \textbf{HMMT25 pass@8}\\
\midrule
Full-vocabulary logit distillation & \textbf{84.1} & \textbf{60.0}\\
Sampled-token distillation & 82.1 & 57.3\\
\bottomrule
\end{tabular}
\end{center}

这说明 teacher 的完整 distribution 比单个 sampled token 上的 advantage 更有用。不过它也带来更高显存，因为需要存每个位置的 vocabulary-sized logits。

\subsection{8. 和已有工作的关系}

\begin{center}
\resizebox{\linewidth}{!}{
\begin{tabular}{lll}
\toprule
\textbf{方法} & \textbf{核心学习信号} & \textbf{和 OPSD 的差异}\\
\midrule
SFT & reference token hard label & off-policy；直接模仿标准解风格\\
STaR / ReST & self-generated correct rationale + filtering & sequence-level hard selection；错样本无学习信号\\
GRPO / RLVR & group-relative binary/verifiable reward & on-policy 但 sparse；多 rollout；token credit 粗\\
On-policy distillation / GKD & external teacher distribution & dense + on-policy，但需要外部 teacher\\
Context distillation & privileged context teacher outputs & 通常生成 teacher trajectory 后 SFT；不是 student rollout 上逐 token soft matching\\
\textbf{OPSD} & same-model privileged distribution & on-policy + dense + no external teacher，但需要 reference solution\\
\bottomrule
\end{tabular}
}
\end{center}

\subsubsection{8.1 和 GRPO/RLVR}

两者都使用 on-policy 数据，但优化对象不同：

\begin{itemize}[leftmargin=1.5em]
  \item \textbf{GRPO：} 最大化 outcome reward，学习信号是 group-normalized sequence advantage。
  \item \textbf{OPSD：} 最小化 privileged teacher 与 student 的 token distribution divergence，学习信号是每个 token 的 soft target。
\end{itemize}

因此 OPSD 可以看成 ``不直接优化正确率 reward 的 RLVR 替代物''。它更像 dense shaping，而不是 reward maximization。

\subsubsection{8.2 和 SFT}

OPSD 仍然依赖标准解，但不直接把标准解当 label。标准解只用于构造 teacher context；真正被监督的是 student rollout 上的 next-token distribution。这个区别解释了为什么同一数据集上 SFT 会退化，而 OPSD 能提升。

\subsubsection{8.3 和 UniSD}

在更大的 self-distillation 设计空间里，OPSD 是一个具体点：

\begin{itemize}[leftmargin=1.5em]
  \item teacher 由 same model + privileged answer context 构造；
  \item student 轨迹是 on-policy；
  \item loss 是 token-level divergence；
  \item 稳定性主要靠 pointwise KL clipping 和固定 teacher。
\end{itemize}

UniSD 后续把 self-distillation 拆成 teacher reliability、EMA teacher、contrastive / feature matching、divergence clipping 等组件；OPSD 可以被看作 privileged-context self-distillation 分支的代表。

\subsection{9. 我认为它的真正贡献}

\begin{enumerate}[leftmargin=1.5em]
  \item \textbf{把 reference solution 从 hard label 变成 privileged context。} 这使标准解不再直接决定学生输出文本，而是改变 teacher 的分布判断。
  \item \textbf{把 RLVR 的 on-policy 优势和 distillation 的 dense signal 合起来。} 这是它比 SFT 和 GRPO 都更 token-efficient 的核心。
  \item \textbf{证明小模型 reasoning 后训练不一定要先走昂贵 GRPO。} 在 1.7B 上，100 steps OPSD 的平均提升显著高于 GRPO baseline。
  \item \textbf{提出 pointwise KL clipping。} 这个技巧对所有 LLM logit distillation 都值得复用，尤其是 reasoning 中 style tokens 容易主导 KL 的场景。
\end{enumerate}

\subsection{10. 局限与风险}

\begin{itemize}[leftmargin=1.5em]
  \item \textbf{依赖高质量标准解。} OPSD 不是无监督 self-improvement；它需要 ground-truth answer 或 reference CoT。如果标准解错误、过短或风格偏置强，teacher distribution 也会被污染。
  \item \textbf{teacher 的能力有理解阈值。} 如果题目难到模型即使看到标准解也不能 rationalize，privileged teacher 就不能提供有效监督。论文也把 curriculum learning 作为 future direction。
  \item \textbf{没有显式使用 student answer correctness。} 训练目标只对齐 teacher distribution，没有把 verifier correctness 作为额外信号；这可能错过一些 outcome-level 校正。
  \item \textbf{规模只验证到 8B。} 论文受计算限制，尚不知道 32B/70B 级别是否仍然同样 token-efficient。
  \item \textbf{full-vocabulary loss 显存贵。} 它比 sampled-token objective 更好，但需要存储每个 token 的词表 logits，长序列和大词表下成本明显。
  \item \textbf{``same model teacher'' 的定义有实现 nuance。} 主实验固定 initial teacher，这增强稳定性，但也意味着 teacher 不随 student 进步而变强；如果同步更新，又可能出现 target drift。
\end{itemize}

\subsection{11. 对我的可复用点}

\begin{itemize}[leftmargin=1.5em]
  \item \textbf{可复用 idea 1：privileged-context teacher。} 任何有 reference solution / answer / tool trace / proof trace 的任务，都可以尝试把这些信息作为 teacher context，而不是直接 SFT label。
  \item \textbf{可复用 idea 2：on-policy prefix 上做 soft correction。} 对模型自己生成的错误前缀，teacher distribution 可能比标准解 token 更有价值，因为它提供 ``从错误状态如何回正'' 的局部信号。
  \item \textbf{可复用 idea 3：pointwise divergence clipping。} 对 LLM 蒸馏尤其重要，避免少数 style tokens 或模板 tokens 放大成训练主方向。
  \item \textbf{可复用 idea 4：teacher/student generation style 不对称。} Student 可以关 thinking 以降低 rollout 成本，teacher 可以开 thinking 以提高监督质量；这种不对称设计值得推广到其他模型家族。
  \item \textbf{可复用实验设计：} 一定要同时报 token budget、rollout 数、zero reward-std fraction，否则 ``OPSD 比 GRPO 高效'' 的结论不完整。
\end{itemize}

\subsection{12. 可能的 follow-up}

\begin{itemize}[leftmargin=1.5em]
  \item \textbf{OPSD + verifier：} 在 token-level KL 外，再加入 answer correctness 或 process verifier 的 sequence-level loss，形成 dense + outcome 双信号。
  \item \textbf{OPSD as RLVR warm-up：} 先用 OPSD 让小模型获得初始 reasoning policy，再用 GRPO/DAPO/VAPO 做 outcome optimization，可能降低早期 reward collapse。
  \item \textbf{Adaptive curriculum：} 选择模型 ``看答案后能理解、但不看答案还不会做'' 的题目，维持 teacher 有效性和 student 学习压力。
  \item \textbf{Moving / EMA teacher：} 在固定 initial teacher 和同步更新 teacher 之间，引入 EMA teacher，可能兼顾稳定性和持续改进。
  \item \textbf{Beyond math：} 代码、工具调用、proof assistant、agent trace 都有 privileged traces，可验证 OPSD 是否是通用 reasoning post-training stage。
\end{itemize}

\subsection{13. 待查问题}

\begin{itemize}[leftmargin=1.5em]
  \item Pointwise clipping 的阈值 $\tau$ 如何随模型大小、词表、temperature 和 thinking mode 变化？论文没有充分调参。
  \item Teacher 固定为 initial policy 是否限制了最终上限？EMA teacher 或 checkpoint ensemble 是否更好？
  \item OPSD 与 RLVR 组合时，最佳顺序是 OPSD $\rightarrow$ GRPO，还是交替训练？
  \item 如果 reference solution 是错误的但最终答案正确，teacher distribution 会不会学到伪推理？需要 verifier 或 PRM 来防。
  \item 4B/8B 上 SFT 下降明显，是否主要来自 reasoning length 缩短？能否通过 length control 或 SFT 数据重写解决？
\end{itemize}

\subsection{14. 读完后的记忆钩子}

\begin{summarybox}
\textbf{我最后会这样记 OPSD：} 它是 LLM reasoning 后训练里的 ``看答案自教'' 阶段。SFT 是照抄标准解，GRPO 是做题后只看对错，OPSD 是先自己做题，再让 ``看过标准解的自己'' 对每一步 next-token distribution 做软纠偏。它的 stage 应放在 post-SFT reasoning specialization，既可作为 RLVR 前置 warm-up，也可作为低成本 RLVR 替代。
\end{summarybox}

\subsection{参考来源}

\begin{itemize}[leftmargin=1.5em]
  \item \textbf{论文 PDF：} \href{https://siyan-zhao.github.io/assets/img/opsd/opsd_v3.pdf}{Self-Distilled Reasoner: On-Policy Self-Distillation for Large Language Models}。
  \item \textbf{arXiv：} \href{https://arxiv.org/abs/2601.18734}{arXiv:2601.18734v3}，包含摘要、版本信息与 DOI。
  \item \textbf{作者博客：} \href{https://siyan-zhao.github.io/blog/2026/opsd/}{Self-Distilled Reasoner: On-Policy Self-Distillation}，用于核对方法动机、训练范式对比和讨论。
  \item \textbf{代码仓库：} \href{https://github.com/siyan-zhao/OPSD}{siyan-zhao/OPSD}，用于核对训练脚本、OPSDTrainer、评测设置和 README 中的复现实验说明。
  \item \textbf{相关背景：} \href{https://arxiv.org/abs/1503.02531}{Knowledge Distillation}；\href{https://arxiv.org/abs/2402.03300}{DeepSeekMath / GRPO}；\href{https://arxiv.org/abs/2506.04178}{OpenThoughts}；\href{https://thinkingmachines.ai/blog/on-policy-distillation}{On-Policy Distillation}；\href{https://arxiv.org/abs/2605.06597}{UniSD}。
\end{itemize}

% ============================================================
% Idea pool: DenseGRPO improvement directions
% ============================================================

\section{Idea 池：DenseGRPO 之后怎么改}

\begin{conceptbox}
\textbf{核心判断：} 读完 Flow-GRPO、DenseGRPO、Chunk-GRPO、E-GRPO 和 LeapAlign 后，我现在更倾向于把后续工作理解成一个问题：\key{denoising timestep 是数值积分单位，不一定是最优 RL credit unit}。DenseGRPO 证明 step-wise improvement 有用，但太贵；Chunk-GRPO 证明 chunk 是更合理的中间粒度，但 advantage 仍偏粗；E-GRPO 证明不是所有 step 都值得随机探索；LeapAlign 证明跨步依赖和早期步骤很重要，但完整长轨迹反传也太贵。
\end{conceptbox}

\subsection{0. 总体判断：下一步不该只是 ``更 dense''}

DenseGRPO 的关键 reward 是：

\[
\Delta R_t = R_{t-1} - R_t.
\]

它回答了 Flow-GRPO 没有回答清楚的问题：

\[
\text{当前 denoising transition 到底让最终可读出的结果变好了多少？}
\]

但如果每个 timestep 都做完整 ODE readout，额外成本接近：

\[
T + (T-1) + \cdots + 1 = \frac{T(T+1)}{2}.
\]

所以后续真正值得做的不是简单把 reward 算得更密，而是：

\begin{quote}
如何用更少的 readout / 更少的 stochastic update / 更短的反传路径，保留 DenseGRPO 的 credit assignment 收益？
\end{quote}

我现在会把方向分成三层：

\begin{enumerate}[leftmargin=1.5em]
  \item \textbf{近端最稳：} 在原始采样步数不变的情况下，减少 readout 和 RL update 的无效位置，主打训练效率与 credit 稳定性。
  \item \textbf{中期更强：} 前验得到 sparse schedule 或 macro-step sampler，让 rollout 本身也变少，主打采样效率。
  \item \textbf{更激进：} 把 GRPO 的 group-relative credit 和 LeapAlign 的 short-gradient path 结合，形成 policy-gradient + direct-gradient hybrid。
\end{enumerate}

\begin{warnbox}
\textbf{一个必须写清楚的约束：} 如果 10 个关键点是先完整采样 50 步后再后验选出来，那它不省 rollout 成本，只能改善 credit assignment 或减少 update 成本。只有关键点来自前验 schedule、低成本 selector、解析 entropy/dynamics 指标，或者已经蒸馏好的 macro-step student，才真的能提升采样效率。
\end{warnbox}

\subsection{0.1 先把 readout 讲清楚：它不是 rollout}

这里最容易混的是 \textbf{readout} 和 \textbf{rollout}。

\begin{itemize}[leftmargin=1.5em]
  \item \textbf{Rollout：} 真正跑一条采样轨迹。比如从纯噪声一路降噪到最终图像：
  \[
  x_{50}\rightarrow x_{49}\rightarrow\cdots\rightarrow x_0.
  \]
  这条轨迹会产生训练样本，也就是 RL 里的 ``这次模型实际做了什么''。
  \item \textbf{Readout：} 不一定是在产生新样本，而是在问一个评估问题：\key{如果我现在停在某个中间 latent，比如 \(x_{30}\)，从这里继续按确定性 ODE 走到终点，它最后会长成什么图？这个图能得多少 reward？}
\end{itemize}

所以 DenseGRPO 里的 readout 通常是：

\[
x_t
\xrightarrow{\mathrm{deterministic\ ODE\ suffix}}
\hat{x}_{0|t}
\xrightarrow{\mathrm{Decode}}
\mathrm{image}_{t}
\xrightarrow{\mathrm{Reward\ Model}}
R_t.
\]

它的目的不是重新探索一条 trajectory，而是把 noisy / intermediate latent ``读成'' 一个 reward model 能评价的 clean image proxy。

\begin{intubox}
\textbf{更通俗地说：} rollout 像是真的把电影从开头拍到结尾；readout 像是在电影中间暂停，然后问：``如果从这一帧按默认剧情继续拍完，结局大概是什么样？这个结局打几分？''
\end{intubox}

这也是 DenseGRPO 贵的原因。普通 rollout 只需要跑一条：

\[
x_{50}\rightarrow\cdots\rightarrow x_0.
\]

但如果每个中间点都 readout，就还要额外跑：

\[
x_{49}\rightarrow x_0,
\quad
x_{48}\rightarrow x_0,
\quad
\ldots,
\quad
x_1\rightarrow x_0.
\]

所以成本会从普通 rollout 的 \(O(T)\) 变成接近 \(O(T^2)\)。

\begin{summarybox}
\textbf{本文后面所有 ``减少 readout'' 的意思：} 不是说一定少生成最终样本，而是少做这种 ``从中间 latent 额外读到 clean image 再打 reward'' 的昂贵评估。P0 主线首先省的是 readout / reward-evaluation / update 成本；只有到 Amortized Knot-GRPO 或 distilled macro sampler，才进一步声称 rollout 本身也变少。
\end{summarybox}

\subsection{1. 可排列组合的设计轴}

后续 idea 可以从下面几条轴组合：

\begin{enumerate}[leftmargin=1.5em]
  \item \textbf{credit unit：} step、chunk、entropy-merged step、macro-knot、two-step leap、region / component。
  \item \textbf{reward signal：} terminal reward、DenseGRPO improvement、chunk improvement、branch sensitivity、component-wise improvement、direct reward gradient。
  \item \textbf{readout 方式：} full ODE、boundary ODE、truncated ODE、leap readout、latent proxy、learned rank predictor。
  \item \textbf{step selection：} uniform、temporal dynamics、entropy mass、reward variance、trajectory similarity、prompt-conditioned selector。
  \item \textbf{rollout 策略：} full-step SDE、ODE/SDE hybrid、consolidated SDE step、sparse macro-step sampler、teacher-student distilled sampler。
  \item \textbf{optimization side：} advantage-side credit、ratio-side calibration、chunk-level ratio、macro-transition ratio、direct-gradient auxiliary loss。
\end{enumerate}

\subsection{2. 我当前最看好的主线：Entropy-Dynamic Chunk DenseGRPO}

\begin{conceptbox}
\textbf{组合：}
\[
\text{Chunk-GRPO 的 chunk unit}
+ \text{DenseGRPO 的 reward gain}
+ \text{E-GRPO 的 entropy-aware step selection}
\]
\textbf{一句话：} 把 50 个 denoising micro-steps 先合成少量 ``真正像决策阶段'' 的 chunks；只在 chunk 边界做 readout，再用相邻边界的 reward 差值给这个 chunk 分 credit。
\end{conceptbox}

\subsubsection{2.1 先用白话说这个 idea}

DenseGRPO 的做法像是：电影有 50 帧，我每一帧都暂停一次，然后从这一帧继续播放到结尾，看结局能打几分。这样当然能更细地知道每一帧有没有帮助，但成本太高。

我现在最看好的方向不是 ``每一帧都看''，而是：

\begin{quote}
先把 50 帧合成 4 到 8 个剧情段落；只在段落开头和结尾做 readout；如果这一段结束后的 readout 分数比这一段开始时高，就说明这个 chunk 对最终结果有正贡献。
\end{quote}

这就是 \textbf{Entropy-Dynamic Chunk DenseGRPO}。

\begin{summarybox}
\textbf{核心直觉：} denoising step 是数值积分为了稳定而切出来的 ``小格子''，但 RL credit assignment 不一定要跟着这些小格子走。更合理的 credit unit 应该是 ``语义变化明显、探索空间足够、reward 差异可区分'' 的 chunk。
\end{summarybox}

\subsubsection{2.2 三个词分别是什么意思}

这个名字里的三个部分分别对应三篇论文的启发：

\begin{itemize}[leftmargin=1.5em]
  \item \textbf{Entropy：} 来自 E-GRPO。不是所有 timestep 都值得 RL 探索；低熵 step 的 group 差异很小，强行更新容易学到 reward noise。因此 chunk boundary 应该避开大量低熵、低探索价值的位置，或者把这些位置合并掉。
  \item \textbf{Dynamic：} 来自 Chunk-GRPO。trajectory 不是匀速变化的：有些区间 latent / velocity 变化很剧烈，有些区间只是细节修补。chunk 应该尽量覆盖相对完整的动态阶段，而不是机械地每 10 步切一刀。
  \item \textbf{Chunk DenseGRPO：} 来自 DenseGRPO。真正的 advantage 不再只用 terminal reward，而是用 chunk boundary readout 的差值：这一段开始时 ``默认走完'' 得几分，这一段结束后 ``默认走完'' 又得几分。
\end{itemize}

可以把它理解成三句话：

\begin{enumerate}[leftmargin=1.5em]
  \item E-GRPO 告诉我们：哪些地方有探索价值。
  \item Chunk-GRPO 告诉我们：不要把每个微小 step 都当成独立决策。
  \item DenseGRPO 告诉我们：credit 最好来自 local reward gain，而不是整条轨迹共用一个 terminal reward。
\end{enumerate}

\subsubsection{2.3 具体怎么跑：rollout 还是照常，readout 只看边界}

设普通 rollout 是一条 50-step 采样轨迹：

\[
x_{50}\rightarrow x_{49}\rightarrow\cdots\rightarrow x_0.
\]

第一版不急着减少这条 rollout。先承认：我们仍然完整采样，得到一条真实 trajectory。真正先省的是 readout。

假设我们选出 \(K=6\) 个 chunks：

\[
50=\tau_6>\tau_5>\cdots>\tau_1>\tau_0=0.
\]

我们不再对 50 个中间点全部 readout，而只对这些 chunk boundaries readout：

\[
R_{\tau_j}
=
R\left(
\mathrm{Decode}\left(
\mathrm{ODE}(x_{\tau_j}\rightarrow x_0)
\right),c
\right).
\]

然后第 \(j\) 个 chunk 的贡献就是：

\[
\Delta R_{\mathrm{ch}_j}
=
R_{\tau_{j-1}}-R_{\tau_j}.
\]

\begin{intubox}
\textbf{白话解释：} \(R_{\tau_j}\) 是 ``进入这个 chunk 之前，如果默认走完，结局大概几分''；\(R_{\tau_{j-1}}\) 是 ``走完这个 chunk 之后，如果默认走完，结局大概几分''。两者相减，就是这个 chunk 把未来结局改善了多少。
\end{intubox}

用数字感受一下成本：如果 \(T=50\)，full DenseGRPO 近似要对很多中间点都做 suffix ODE readout；而 \(K=6\) 时，只需要对 7 个边界做 readout。

\begin{warnbox}
\textbf{这里不要夸大：} 这个最小版本减少的是 readout / reward scoring / advantage 计算成本，不是把 50-step rollout 直接变成 6-step rollout。要让 rollout 本身也变少，需要后面的前验 sparse schedule、cheap selector 或 macro-step distillation。
\end{warnbox}

\subsubsection{2.4 boundary 怎么选：不要只 uniform 切}

最朴素做法是 uniform chunks，但这只是 baseline。真正的方法应该用一个 boundary importance：

\[
I(t)
=
\alpha\,\widetilde{H}(t)
+\beta\,\widetilde{D}_{\mathrm{dyn}}(t)
+\gamma\,\widetilde{V}_{\Delta R}(t).
\]

其中：

\begin{itemize}[leftmargin=1.5em]
  \item \(H(t)\)：entropy proxy，表示这里有没有足够探索空间；
  \item \(D_{\mathrm{dyn}}(t)\)：trajectory dynamics，表示这里 latent / velocity 变化是否明显；
  \item \(V_{\Delta R}(t)\)：少量 calibration readout 得到的 reward-gain variance，表示这里的 credit 是否真的能区分好坏样本。
\end{itemize}

更通俗地说，chunk boundary 应该放在三类位置附近：

\begin{enumerate}[leftmargin=1.5em]
  \item \textbf{模型还真有选择空间的位置：} entropy 不能太低。
  \item \textbf{图像语义或结构正在明显变化的位置：} dynamics 不能太平。
  \item \textbf{不同样本在 reward 上会拉开差距的位置：} readout 差值不能全都差不多。
\end{enumerate}

\subsubsection{2.5 接到 GRPO 目标}

对同一个 prompt 采样 \(G\) 条 trajectories，第 \(i\) 条轨迹、第 \(j\) 个 chunk 的 dense chunk advantage 是：

\[
A^i_{\mathrm{ch}_j}
=
\frac{
\Delta R^i_{\mathrm{ch}_j}
-
\frac{1}{G}\sum_{g=1}^{G}\Delta R^g_{\mathrm{ch}_j}
}{
\mathrm{std}_{g}(\Delta R^g_{\mathrm{ch}_j})+\epsilon
}.
\]

它的含义非常直接：在同一个 prompt 的 group 里，第 \(i\) 条样本的第 \(j\) 个 chunk，是不是比其他样本的同一阶段更会提升最终 reward。

ratio 可以沿用 Chunk-GRPO 的几何平均：

\[
r^i_j(\theta)
=
\left(
\prod_{t\in \mathrm{ch}_j}
\frac{p_{\theta}(x^i_{t-1}|x^i_t,c)}{p_{\theta_{\mathrm{old}}}(x^i_{t-1}|x^i_t,c)}
\right)^{1/cs_j}.
\]

目标仍然是 PPO / GRPO 风格的 clipped objective：

\[
J(\theta)
=
\mathbb{E}
\left[
\frac{1}{GK}
\sum_{i=1}^{G}
\sum_{j=1}^{K}
\min\left(
 r^i_j(\theta)A^i_{\mathrm{ch}_j},
 \mathrm{clip}(r^i_j(\theta),1-\epsilon,1+\epsilon)A^i_{\mathrm{ch}_j}
\right)
\right].
\]

\subsubsection{2.6 为什么它比几个直接 baseline 更有吸引力}

\begin{itemize}[leftmargin=1.5em]
  \item \textbf{相比 DenseGRPO：} 不再每个 micro-step 都 readout，成本明显更低；而且 chunk-level gain 比 step-level gain 更平滑，噪声可能更小。
  \item \textbf{相比 Chunk-GRPO：} 不只是把 ratio 合成 chunk-level，advantage 也变成 chunk-specific improvement，因此 credit 更细。
  \item \textbf{相比 E-GRPO：} entropy 不只是用来决定哪些 step 更新，还参与决定 credit boundary；同时 DenseGRPO-style \(\Delta R\) 会检查这些 high-entropy 区域是否真的带来 reward improvement。
  \item \textbf{相比 LeapAlign：} 这条主线不依赖可微 reward，也不需要先处理 direct-gradient 的稳定性；但后续可以借 LeapAlign 的 similarity weighting 来判断某些 cheap readout 是否可信。
\end{itemize}

\begin{summarybox}
\textbf{为什么这是第一优先级：} 它最容易落地，因为不用一开始解决 macro-step logprob，也不要求 reward model 可微；但它已经把 DenseGRPO、Chunk-GRPO 和 E-GRPO 的最核心优点接到一起：少 readout、chunk credit、entropy-aware exploration。
\end{summarybox}

\subsubsection{2.7 最小可做版本和升级路线}

\textbf{Version 0：只省 readout，不声称省 rollout。}

\begin{enumerate}[leftmargin=1.5em]
  \item 完整跑 50-step SDE rollout。
  \item 固定 \(K=4\) 或 \(K=6\) 个 chunks。
  \item 比较 uniform boundary、dynamics boundary、entropy boundary、entropy+dynamics boundary。
  \item 只在 chunk boundaries 做 full ODE readout。
  \item 用 \(\Delta R_{\mathrm{ch}_j}\) 做 chunk-level advantage。
\end{enumerate}

\textbf{Version 1：进一步少读。}

只对 active chunks 做 readout：entropy 太低、dynamics 太平、或者历史 \(\Delta R\) variance 太低的 chunk，直接合并或用 terminal advantage 保底。

\textbf{Version 2：用 proxy / active learning 压 readout。}

训练 latent reward proxy 预测 boundary reward，但只要求 group ranking 对，不强求 reward MSE；当 proxy uncertainty 高时再调用 full ODE readout。

\textbf{Version 3：开始减少 rollout。}

把高价值 chunk boundaries 蒸馏成前验 sparse schedule 或 macro-step student。这个阶段才能严肃声称：rollout 本身也从 50 步降到更少步。

\begin{warnbox}
\textbf{论文写法要诚实：} P0 的第一阶段是 efficient credit assignment，不是 fast sampler。它先减少 DenseGRPO 最贵的 readout；后续如果结合 Amortized Knot-GRPO / distillation，才是 rollout acceleration。
\end{warnbox}

\subsection{3. Idea A：Entropy-Gain Early-Exit GRPO}

\begin{conceptbox}
\textbf{想法：}
完整 rollout 很贵，而且 E-GRPO / critical-step 类结果都在暗示：不是所有 denoising steps 都贡献同等有效的 RL 信号。更激进的做法是，不再默认每条 trajectory 都把 stochastic rollout 跑到最后；当一个 chunk 的有效探索价值进入 plateau，就停止随机探索，切到确定性的 ODE suffix 把当前 latent readout 成最终图像。

\textbf{一句话：} 前半段用 SDE rollout 做探索和 RL credit，后半段如果已经低 entropy、低 dynamics、低 reward-gain variance，就提前退出 stochastic rollout，用 cheap deterministic completion 得到最终 reward。
\end{conceptbox}

\subsubsection{3.1 白话版}

原始想法可以说成：

\begin{quote}
既然完整 rollout 很耗时，而很多 step 其实对最终 reward 没有明显贡献，那能不能一边 rollout 一边看 entropy / 收益变化；如果后面已经不怎么增长，就提前停掉？
\end{quote}

我觉得这个方向是可以做的，但要把 ``停掉'' 说准确：

\begin{warnbox}
\textbf{不是停在中间就结束。} 中间 latent 本身通常不能直接给 reward model 打分。真正该停的是 \textbf{stochastic rollout / RL exploration / RL update}，不是停掉最终图像生成。停下后还要用 deterministic ODE suffix 把 \(x_{\tau^*}\) 读到 \(x_0\)。
\end{warnbox}

也就是从：

\[
x_T\rightarrow x_{T-1}\rightarrow\cdots\rightarrow x_0
\]

变成：

\[
x_T\rightarrow\cdots\rightarrow x_{\tau^*}
\xrightarrow{\mathrm{ODE\ suffix}}
\hat{x}_{0|\tau^*}
\xrightarrow{\mathrm{Reward}}
R_{\tau^*}.
\]

\begin{intubox}
\textbf{通俗理解：} 前半段像真的继续拍电影，因为剧情还在分叉；后半段如果剧情基本定型，就不再反复随机拍不同版本，而是按默认剪辑把结尾补出来，然后看这个结尾打几分。
\end{intubox}

\subsubsection{3.2 停止条件：不要只看平均 reward 不涨}

最危险的版本是：

\begin{quote}
reward 均值不增长了，所以停。
\end{quote}

这不够稳。因为 GRPO 关心的是 group-relative signal：就算某个 chunk 的平均收益接近 0，只要 group 内不同样本的收益差异很大，它仍然有训练价值。

所以更合理的停止信号应该是组合分数：

\[
U_j
=
\alpha\,\widetilde{H}_j
+
\beta\,\widetilde{D}_j
+
\gamma\,\widetilde{V}_{\Delta R,j}.
\]

其中：

\begin{itemize}[leftmargin=1.5em]
  \item \(H_j\)：第 \(j\) 个 chunk 的 SDE transition entropy，表示这里还有没有探索空间；
  \item \(D_j\)：latent dynamics / velocity change，表示这里是不是还在发生结构性变化；
  \item \(V_{\Delta R,j}\)：group 内 reward-gain variance，表示这个 chunk 能不能区分好坏样本。
\end{itemize}

当连续 \(p\) 个 chunks 满足：

\[
U_j < \eta,
\]

就停止后续 stochastic rollout，切到 ODE suffix。

\begin{summarybox}
\textbf{核心判断：} 真正该停的不是 ``reward 不涨''，而是 ``entropy 低、dynamics 平、group 内 reward gain 也拉不开差距''。这三个条件同时出现，才说明后续随机探索的 RL 价值可能已经很低。
\end{summarybox}

\subsubsection{3.3 接到 GRPO：只更新 stop 前的 active chunks}

设自适应停止点为 \(\tau^*\)。训练时只对 \(\tau^*\) 之前已经真实执行过的 chunks 做 policy-gradient 更新：

\[
J(\theta)
=
\mathbb{E}
\left[
\frac{1}{G}\sum_{i=1}^{G}
\sum_{j\leq j^*(i)}
\min\left(
 r^i_j(\theta)A^i_j,
 \mathrm{clip}(r^i_j(\theta),1-\epsilon,1+\epsilon)A^i_j
\right)
\right].
\]

最终 reward 来自 early-exit 后的 deterministic completion：

\[
R^i_{\mathrm{exit}}
=
R\left(
\mathrm{Decode}\left(
\mathrm{ODE}(x^i_{\tau^*}\rightarrow x_0)
\right),c
\right).
\]

如果想显式鼓励更早停止，可以加 cost-aware objective：

\[
\tilde{R}^i
=
R^i_{\mathrm{exit}}
-
\lambda\,\mathrm{Cost}(\tau^{*,i}).
\]

但第一版最好先不要让停止策略太自由，否则模型可能学会 ``骗 stop rule''，让 entropy / dynamics 看起来低，从而逃避后续惩罚。

\subsubsection{3.4 最小可做实验}

\textbf{Version 0：固定 cutoff baseline。}

先不要动态停，先比较几个固定停止点：

\[
\tau^*\in\{40,30,20,10\}.
\]

每条样本只 rollout 到 \(x_{\tau^*}\)，然后用 ODE suffix readout 到最终图像。报告：

\begin{itemize}[leftmargin=1.5em]
  \item early-exit reward 和 full rollout reward 的差距；
  \item early readout ranking 和 terminal reward ranking 的 correlation；
  \item 节省的 rollout forward calls；
  \item 哪些 prompt / reward 类型最容易因为早停受损。
\end{itemize}

\textbf{Version 1：entropy cutoff。}

用 E-GRPO 风格的 entropy proxy：

\[
H_j<\eta_H
\]

连续 \(p\) 个 chunks 后停止。

\textbf{Version 2：entropy + reward variance cutoff。}

只有同时满足：

\[
H_j<\eta_H,
\qquad
\mathrm{Var}_{G}(\Delta R_j)<\eta_V,
\]

才允许 early exit。这样可以避免 ``平均 reward 不涨但 group 内仍有强区分度'' 的情况被误停。

\textbf{Version 3：bucketed stop-time。}

不要让每条样本任意停，而是先把停止点限制在少量 buckets：

\[
\tau^*\in\{40,30,20,10\}.
\]

这样 advantage normalization、ratio statistics 和 KL 统计都会更稳。

\subsubsection{3.5 和当前 P0 主线的关系}

这条 idea 比前面的 Entropy-Dynamic Chunk DenseGRPO 更激进。

\begin{itemize}[leftmargin=1.5em]
  \item \textbf{P0 主线：} 完整 rollout 仍然跑完，只减少 readout / reward scoring / update 的无效位置。它更稳，目标是 efficient credit assignment。
  \item \textbf{Early-Exit GRPO：} 当后续 chunks 没有有效探索价值时，连 stochastic rollout 本身也提前停止，再切 deterministic completion。它更像 rollout-cost reduction。
\end{itemize}

\begin{summarybox}
\textbf{定位：} 这是 P0 之后最自然的 P0.5 / P1 升级版。P0 先证明 ``少 readout 也能学到好 credit''；这个 idea 再进一步证明 ``低价值后缀甚至可以不做 stochastic rollout''。
\end{summarybox}

\subsection{4. Idea B：Sparse Readout Dense Reward}

\textbf{想法：} DenseGRPO 不必每个 timestep 都 readout。先选择少量关键 timesteps：

\[
T = \tau_K > \tau_{K-1} > \cdots > \tau_1 > \tau_0 = 0,
\]

只计算这些点的 reward：

\[
R_{\tau_K}, R_{\tau_{K-1}}, \ldots, R_{\tau_0},
\]

再用区间 reward gain：

\[
\Delta R_{[\tau_j,\tau_{j-1}]}
= R_{\tau_{j-1}} - R_{\tau_j}.
\]

这个 idea 是前面主线的基础版本。它只回答 ``readout 成本怎么降''，还没有回答 ``哪些点值得选''。

\textbf{选点策略可以从弱到强排列：}

\begin{itemize}[leftmargin=1.5em]
  \item \textbf{Uniform sparse：} 简单 baseline，不该作为最终方法。
  \item \textbf{Entropy-mass sparse：} 按 E-GRPO entropy proxy 的累计质量切分。
  \item \textbf{Dynamics sparse：} 按 latent velocity / relative \(L1\) 曲线变化切分。
  \item \textbf{Reward-variance sparse：} 用少量 calibration readout 找 group 内 \(\Delta R\) variance 高的位置。
\end{itemize}

\begin{summarybox}
\textbf{定位：} Sparse Readout 是一个工程基线；真正有论文味的是把 sparse boundary 和 ``credit informativeness'' 绑定起来，而不是单纯少读几个点。
\end{summarybox}

\subsection{5. Idea C：Chunk-level Improvement Reward}

\textbf{想法：} 不估计每一个 step 的贡献，而是估计一个 chunk 的边际贡献。

第 \(j\) 个 chunk 为：

\[
\mathrm{ch}_j=(\tau_j,\tau_{j-1}],
\]

只在 chunk boundary 上做 readout：

\[
\hat{x}_{0|\tau_j}
=
\mathrm{ODE}(x_{\tau_j},c;\tau_j\rightarrow 0),
\qquad
R_{\tau_j}=R(\mathrm{Decode}(\hat{x}_{0|\tau_j}),c).
\]

chunk-level improvement：

\[
\Delta R_{\mathrm{ch}_j}
=
R_{\tau_{j-1}}-R_{\tau_j}.
\]

\begin{warnbox}
\textbf{和 Chunk-GRPO 的区别：} Chunk-GRPO 主要把 step-level ratio 合成 chunk-level ratio，但多数设置下仍用 terminal advantage。这个 idea 的变化是 advantage 也变成 chunk-specific improvement，所以它才真正把 DenseGRPO 的 reward gain 搬到了 chunk 粒度。
\end{warnbox}

\textbf{可以继续增强：}

\begin{itemize}[leftmargin=1.5em]
  \item 用 E-GRPO entropy 合并低熵 steps，避免 chunk boundary 落在无效 refinement 区。
  \item 用 DenseGRPO 的 \(\Delta R\) variance 反过来动态细分 chunk。
  \item 用 terminal advantage 做保底，避免早期 chunk readout 噪声过大。
\end{itemize}

\subsection{6. Idea D：E-DenseGRPO：只在高信息量 step 上算 dense credit}

\begin{conceptbox}
\textbf{组合：}
\[
\text{DenseGRPO step-wise improvement}
+ \text{E-GRPO entropy-aware active steps}
\]
\textbf{一句话：} DenseGRPO 不应该默认每个 step 都有同等 credit 价值；只在 high-entropy 或 consolidated SDE step 上做 \(\Delta R\)，低熵段合并或跳过。
\end{conceptbox}

E-GRPO 的结论是：low-entropy steps 的随机扰动太小，group 内样本差异不明显，继续做 RL 更新容易放大 reward noise。因此可以把 DenseGRPO 的 readout 点限制在：

\[
\mathcal{S}_{\mathrm{active}}
=
\{t: e^{h(t)}\ge \tau\}
\cup
\{\text{merged low-entropy segments}\}.
\]

然后只计算：

\[
\Delta R_t=R_{t^-}-R_t,
\qquad t\in \mathcal{S}_{\mathrm{active}}.
\]

低熵连续段不逐步 readout，而是合并成一个 consolidated improvement：

\[
\Delta R_{[t_m,t_{m-l}]}
=R_{t_{m-l}}-R_{t_m}.
\]

\begin{summarybox}
\textbf{优势：} 这比单纯 Sparse Readout 更有理论动机，因为选择依据来自 SDE transition entropy；也比 E-GRPO 原版更强，因为 advantage 不再只靠 terminal group reward，而有了 DenseGRPO 风格的 local improvement。
\end{summarybox}

\subsection{7. Idea E：Amortized Knot-GRPO：前验关键跳点，而不是后验挑点}

\begin{conceptbox}
\textbf{核心约束：} 如果每条 trajectory 都要先完整采 50 步再选 10 个 knot，那么采样效率没有提升。真正有意义的 Knot-GRPO 必须让 knots 在 rollout 前就确定，或者由低成本 selector 在线确定。
\end{conceptbox}

设原始有 50 个 denoising steps。目标不是后验选出 10 个点，而是学习一个 sparse schedule：

\[
\mathcal{K}=\{k_{10},k_9,\ldots,k_1,0\},
\]

之后 rollout 直接走：

\[
t_{k_{10}}\rightarrow t_{k_9}\rightarrow\cdots\rightarrow t_{k_1}\rightarrow t_0.
\]

\textbf{这个 idea 可以有三个版本：}

\begin{enumerate}[leftmargin=1.5em]
  \item \textbf{固定 schedule：} 用少量 calibration prompts 离线搜索一个通用 sparse schedule，然后所有 RL rollout 直接使用。
  \item \textbf{prompt-conditioned schedule：} 用 prompt embedding、initial noise statistics 或 early cheap features 预测 schedule。
  \item \textbf{teacher-student macro sampler：} 用 50-step teacher 生成关键状态，蒸馏出 10-step student，RL 阶段直接训练 student 的 macro transitions。
\end{enumerate}

\textbf{可以优化的目标：}

\[
\mathcal{K}^{*}
=
\arg\max_{|\mathcal{K}|=K}
\mathbb{E}_{c}
\left[
R(x^{\mathcal{K}}_0,c)
-\lambda\,\mathrm{Cost}(\mathcal{K})
+\eta\,\mathrm{CreditInfo}(\mathcal{K})
\right].
\]

其中 \(\mathrm{CreditInfo}\) 可以来自 entropy mass、reward variance、semantic divergence 或 chunk improvement correlation。

\begin{warnbox}
\textbf{最大难点：} macro-transition 的 logprob / ratio 怎么算。如果直接跨很大一步，\(p_{\theta}(x_{t_j}|x_{t_i},c)\) 不一定好精确建模。保守版本可以先内部用 deterministic solver，RL 只在 macro boundary 上做 credit；激进版本再训练 macro-step student。
\end{warnbox}

\subsection{8. Idea F：Leap-Readout Dense Reward}

\begin{conceptbox}
\textbf{组合：}
\[
\text{DenseGRPO 的 ODE readout}
+ \text{LeapAlign 的 two-step leap / similarity weighting}
\]
\textbf{一句话：} DenseGRPO 的 full ODE readout 太贵，可以尝试用 two-step leap readout 近似 \(R_t\)，再用 trajectory-similarity weight 控制不可信 readout。
\end{conceptbox}

DenseGRPO 为了给 \(x_t\) 打分，需要从 \(x_t\) 走到 \(x_0\)。LeapAlign 提醒我们，可以构造短路径：

\[
x_k
\Rightarrow
\hat{x}_{j|k}
\Rightarrow
x_j
\Rightarrow
\hat{x}_{0|j}
\Rightarrow
x_0.
\]

对应到 dense reward readout，可以用近似：

\[
\hat{x}_{0|t}^{\mathrm{leap}}
=
\mathrm{LeapReadout}(x_t,t,c),
\qquad
\hat{R}_t^{\mathrm{leap}}
=R(\mathrm{Decode}(\hat{x}_{0|t}^{\mathrm{leap}}),c).
\]

但 leap approximation 会偏离真实 ODE readout，所以需要类似 LeapAlign 的可信度权重：

\[
w_t
=
\frac{1}{\max(d_t,\epsilon)+\max(d_0,\epsilon)}.
\]

最终 dense gain 可以写成：

\[
\Delta \tilde{R}_t
=
w_t\left(\hat{R}_{t-1}^{\mathrm{leap}}-\hat{R}_{t}^{\mathrm{leap}}\right).
\]

\begin{summarybox}
\textbf{价值：} 这个方向不是把 LeapAlign 直接变成 GRPO，而是借它的 ``short path + similarity weighting'' 来降低 DenseGRPO readout 成本。它适合做成 DenseGRPO 的 cheap readout variant。
\end{summarybox}

\textbf{风险：} reward model 通常对真实 clean image 分布更可靠；如果 leap readout 图像有伪影，reward 会被误导。因此需要少量 full ODE readout 做校准，并报告 \(\Delta \hat{R}\) 与 \(\Delta R\) 的 ranking correlation。

\subsection{9. Idea G：Rank-Corrected Latent Reward Proxy}

\textbf{想法：} 用少量 full ODE readout 作为 teacher，训练一个 latent reward proxy：

\[
\hat{R}_{\phi}(x_t,t,c)\approx R_t.
\]

在线 RL 时直接预测：

\[
\Delta \hat{R}_t
=
\hat{R}_{\phi}(x_{t-1},t-1,c)-\hat{R}_{\phi}(x_t,t,c).
\]

但 GRPO 主要依赖 group-relative ranking，所以 proxy 不一定要拟合绝对 reward，重点是拟合排序：

\[
\mathrm{rank}(\Delta \hat{R}_t)
\approx
\mathrm{rank}(\Delta R_t).
\]

可以加入 rank loss：

\[
\mathcal{L}_{\mathrm{rank}}
=
\sum_{i,j}
\max\left(0,
 m - s_{ij}(\Delta \hat{R}^i_t-\Delta \hat{R}^j_t)
\right),
\]

其中 \(s_{ij}=\mathrm{sign}(\Delta R^i_t-\Delta R^j_t)\)。

\begin{summarybox}
\textbf{为什么比普通 proxy 更适合 GRPO：} GRPO 不太在意 reward 绝对值，只在意同 prompt group 内谁更好。因此 proxy 的第一目标应该是 group ranking，而不是 MSE。
\end{summarybox}

\subsection{10. Idea H：Adaptive Readout as Active Learning}

\textbf{想法：} 把 full ODE readout 看成昂贵 oracle。训练过程中不是固定 readout，而是根据不确定性主动查询。

触发 full readout 的条件可以包括：

\begin{itemize}[leftmargin=1.5em]
  \item proxy uncertainty 高；
  \item leap readout 和 truncated readout 分歧大；
  \item group 内 predicted \(\Delta R\) ranking 接近，容易翻转；
  \item policy ratio / KL 在某个 chunk 异常；
  \item entropy 高但 reward variance 低，说明探索可能无效。
\end{itemize}

可以定义查询分数：

\[
Q(t)
=
\lambda_1\,\mathrm{Unc}_{\phi}(t)
+\lambda_2\,\mathrm{Disagree}(t)
+\lambda_3\,\mathrm{Var}_{G}(\Delta \hat{R}_t)
+\lambda_4\,\mathrm{KL}(t).
\]

只对 top budget 的 \(t\) 做 full readout，其余用 proxy / truncated / chunk interpolation。

\begin{summarybox}
\textbf{定位：} 这是把 DenseGRPO 的 readout 成本变成主动学习问题。它很适合和 Idea F 的 proxy 组合：proxy 负责便宜估计，active readout 负责不断校准 proxy。
\end{summarybox}

\subsection{11. Idea I：Ratio-Calibrated Dense Credit}

\textbf{想法：} Flow-GRPO 的实际更新强度不只是 advantage：

\[
\Delta \theta_t
\propto
\nabla_{\theta}\log p_{\theta}(x_{t-1}|x_t,c)
\cdot A_t
\cdot w_{\rho}(t).
\]

DenseGRPO 主要修了 \(A_t\)，但如果不同 timestep / chunk 的 ratio scale 天然不同，那么同样的 \(A_t\) 乘到 objective 里仍然可能被放大或压小。

因此可以对 chunk-level dense advantage 再做 ratio-side normalization：

\[
\tilde{A}_{\mathrm{ch}_j}
=
\frac{A_{\mathrm{ch}_j}}{\sqrt{\mathrm{Var}_{G}(\log r_j)}+\epsilon},
\]

或者对 clipping range 做 timestep / chunk dependent calibration：

\[
\epsilon_j
=\epsilon_0\cdot g\left(\mathrm{Entropy}(\mathrm{ch}_j),\mathrm{Var}(\log r_j)\right).
\]

\begin{summarybox}
\textbf{价值：} 这个方向强调 ``credit assignment = advantage side + ratio side''。如果只做 DenseGRPO-style advantage，但忽略 ratio scale，训练信号仍可能在 objective 里偏掉。
\end{summarybox}

\subsection{12. Idea J：Credit-aware Exploration Noise}

DenseGRPO 已经用 dense reward 的正负分布校准 step-wise noise；E-GRPO 则用 entropy 判断哪些 steps 值得探索。可以把两者合成 chunk-level exploration controller。

对每个 active chunk，维护：

\[
\sigma_j
=f\left(
H_j,
\mathrm{Var}_{G}(\Delta R_{\mathrm{ch}_j}),
\mathrm{Balance}_{G}(\Delta R_{\mathrm{ch}_j}>0)
\right).
\]

直觉是：

\begin{itemize}[leftmargin=1.5em]
  \item entropy 低：不要强行加大随机性，优先 merge；
  \item reward variance 高：这个 chunk 的探索有区分度，可以保留或略增探索；
  \item reward 全正 / 全负：说明探索尺度可能不合适，要缩小或重新校准；
  \item ratio 波动过大：探索过强，容易触发 clipping。
\end{itemize}

\begin{summarybox}
\textbf{定位：} E-GRPO 告诉我们哪里值得探索；DenseGRPO 告诉我们探索结果是否真的带来 improvement。两者合起来可以形成 reward-aware entropy calibration。
\end{summarybox}

\subsection{13. Idea K：Hybrid Terminal + Chunk + Dense Reward}

\textbf{想法：} 不要完全相信某一种 reward credit。早期 dense / chunk readout 噪声大，后期 terminal reward 更稳定；因此可以做 hybrid advantage：

\[
A^i_{j}
=
\lambda_j\widehat{\Delta R^i_{\mathrm{ch}_j}}
+
\mu_j\widehat{R^i_{\mathrm{terminal}}}
+
\nu_j\widehat{S^i_{\mathrm{branch},j}}.
\]

其中：

\begin{itemize}[leftmargin=1.5em]
  \item \(\Delta R_{\mathrm{ch}_j}\)：DenseGRPO / Chunk improvement；
  \item \(R_{\mathrm{terminal}}\)：最终 reward 保底；
  \item \(S_{\mathrm{branch},j}\)：从 chunk boundary 分叉后的 reward sensitivity，可选；
  \item \(\lambda_j,\mu_j,\nu_j\)：由 entropy、readout uncertainty、trajectory stage 控制。
\end{itemize}

\begin{warnbox}
\textbf{关键点：} hybrid 不是为了堆公式，而是为了处理 DenseGRPO 的两个失败模式：早期 readout 不稳定，以及局部 \(\Delta R\) 可能暂时为负但对后续结构有帮助。
\end{warnbox}

\subsection{14. Idea L：GRPO + LeapAlign 的双通道后训练}

\begin{conceptbox}
\textbf{组合：}
\[
\text{GRPO for black-box / non-differentiable reward}
+ \text{LeapAlign for differentiable reward}
\]
\textbf{一句话：} 对不可微偏好模型走 GRPO；对可微 reward / CLIP-like reward 走 LeapAlign-style short-gradient path，让同一个 flow model 同时吃到 group-relative ranking signal 和 direct gradient signal。
\end{conceptbox}

可以写成双损失：

\[
\mathcal{L}
=
\mathcal{L}_{\mathrm{GRPO}}
+
\lambda_{\mathrm{DG}}\mathcal{L}_{\mathrm{LeapDG}}.
\]

其中 \(\mathcal{L}_{\mathrm{GRPO}}\) 使用 chunk / entropy-aware advantage，\(\mathcal{L}_{\mathrm{LeapDG}}\) 使用 LeapAlign 的 two-step leap、gradient discounting 和 similarity weighting。

\textbf{为什么可能有用：}

\begin{itemize}[leftmargin=1.5em]
  \item GRPO 不需要可微 reward，适合黑盒 VLM judge / human preference proxy。
  \item LeapAlign 能直接把可微 reward 梯度送到早期 steps。
  \item 两者的信号不同：GRPO 是相对排序，LeapAlign 是局部梯度方向。
\end{itemize}

\textbf{主要风险：} 两个 loss 可能互相冲突，尤其是 direct gradient 可能更容易 reward hacking。因此需要 gradient norm balancing、loss threshold 或 objective conflict filtering。

\subsection{15. Idea M：Component-wise / Region-wise Improvement for Editing and Video}

如果任务转向图像编辑或视频生成，单一 scalar reward 不够。更合理的是拆成多个 component：

\[
R
=\lambda_{\mathrm{edit}}R_{\mathrm{edit}}
+\lambda_{\mathrm{preserve}}R_{\mathrm{preserve}}
+\lambda_{\mathrm{id}}R_{\mathrm{id}}
+\lambda_{\mathrm{motion}}R_{\mathrm{motion}}
+\lambda_{\mathrm{aesthetic}}R_{\mathrm{aesthetic}}.
\]

然后分别计算 component-wise improvement：

\[
\Delta R^{(m)}_{\mathrm{ch}_j}
=
R^{(m)}_{\tau_{j-1}}-R^{(m)}_{\tau_j}.
\]

对图像编辑，还可以区分 target / non-target region：

\[
\Delta R_{\mathrm{target}}
\quad \text{vs.} \quad
\Delta R_{\mathrm{non\text{-}target}}.
\]

对视频任务，可以把 chunks 对齐到不同阶段：

\begin{itemize}[leftmargin=1.5em]
  \item early chunks：layout、camera、motion plan；
  \item middle chunks：object identity、interaction、temporal consistency；
  \item late chunks：texture、style、aesthetic。
\end{itemize}

\begin{summarybox}
\textbf{价值：} 这是最适合从 text-to-image 往 editing / video 扩展的方向。Dense reward 的优势不只是时间维度更细，也可以在 objective / region / component 维度更细。
\end{summarybox}

\subsection{16. 我当前的优先级排序}

\begin{center}
\resizebox{\linewidth}{!}{
\begin{tabular}{p{0.23\linewidth}p{0.34\linewidth}p{0.34\linewidth}}
\toprule
\textbf{优先级} & \textbf{方向} & \textbf{原因}\\
\midrule
P0 & Entropy-Dynamic Chunk DenseGRPO & 最容易接到现有 GRPO 框架，能同时回应 DenseGRPO 成本、Chunk-GRPO advantage 粗、E-GRPO step selection 三个问题。\\
P1 & Rank-Corrected Latent Reward Proxy + Active Readout & 如果 P0 有效，下一步就要继续压 readout 成本；rank correlation 比 reward MSE 更贴合 GRPO。\\
P1 & Amortized Knot-GRPO / sparse schedule & 论文潜力更大，因为它可能真的减少 rollout 步数；但 macro-transition ratio 和 10-step 质量保持难度更高。\\
P2 & Leap-Readout Dense Reward & 有新意，但需要证明 leap readout 的 reward ranking 足够可靠，否则容易被近似伪影误导。\\
P2 & GRPO + LeapAlign hybrid & 想象空间大，但实现复杂，且 reward hacking / gradient conflict 风险更高。\\
P2 & Component-wise editing / video extension & 任务辨识度强，适合后续换到 editing 或 video 时做。\\
\bottomrule
\end{tabular}
}
\end{center}

\subsection{17. 实验应该怎么证明}

不要只报告最终 reward，要显式报告 cost-quality tradeoff：

\begin{itemize}[leftmargin=1.5em]
  \item \textbf{效果：} HPS / PickScore / ImageReward / CLIP / GenEval / human preference。
  \item \textbf{成本：} rollout forward calls、ODE readout calls、VAE decode calls、reward model calls、GPU hours。
  \item \textbf{效率：} reward improvement per GPU hour、per sampled image、per reward call。
  \item \textbf{credit 质量：} chunk / proxy / leap \(\Delta R\) 与 full DenseGRPO \(\Delta R\) 的 group ranking correlation。
  \item \textbf{稳定性：} advantage variance、ratio clipping rate、KL per chunk、reward hacking 指标。
  \item \textbf{消融：} uniform boundary vs. dynamics boundary vs. entropy boundary vs. reward-variance boundary。
\end{itemize}

\begin{warnbox}
\textbf{最关键 baseline：} 必须比较 fixed sparse schedule / uniform chunks。否则 reviewer 会问：你的复杂 entropy-dynamics-credit selector，是不是只是因为少训了一些低价值 steps？
\end{warnbox}

\subsection{18. 最后形成的论文切入点}

我会把最强主张写成下面这句：

\begin{quote}
Denoising timesteps are numerical discretization units, but GRPO needs reward-informative decision units. We therefore compress dense credit assignment from all micro-steps into entropy- and dynamics-aware chunks / knots, preserving DenseGRPO-style improvement signals while reducing readout and update cost.
\end{quote}

对应的中文记忆钩子是：

\begin{summarybox}
\textbf{DenseGRPO 之后最自然的方向不是更细，而是更会挑。} 该细的位置做 dense credit，该合的位置做 entropy merge，该跳的位置用前验 sparse schedule 或 distillation 直接跳过去；同时别忘了 ratio side 和跨步依赖。
\end{summarybox}

\subsection{参考来源}

\begin{itemize}[leftmargin=1.5em]
  \item \textbf{DenseGRPO：} \href{https://arxiv.org/abs/2601.20218}{DenseGRPO: From Sparse to Dense Reward for Flow Matching Model Alignment}，用于本文关于 step-wise reward gain、ODE readout 成本和 exploration calibration 的讨论。
  \item \textbf{Chunk-GRPO：} \href{https://arxiv.org/abs/2510.21583}{Sample By Step, Optimize By Chunk: Chunk-Level GRPO For Text-to-Image Generation}，用于本文关于 chunk-level ratio、temporal dynamics chunking 和 chunk 粒度 credit 的讨论。
  \item \textbf{E-GRPO：} \href{https://arxiv.org/abs/2601.00423}{E-GRPO: High Entropy Steps Drive Effective Reinforcement Learning for Flow Models}，用于本文关于 entropy-aware active steps、low-entropy merging 和 ODE/SDE hybrid rollout 的讨论。
  \item \textbf{LeapAlign：} \href{https://arxiv.org/abs/2604.15311}{LeapAlign: Post-Training Flow Matching Models at Any Generation Step by Building Two-Step Trajectories}，以及 \href{https://rockeycoss.github.io/leapalign/}{LeapAlign project page}，用于本文关于 two-step leap、latent connector、gradient discounting 和 trajectory-similarity weighting 的讨论。
  \item \textbf{本目录主线总结：} \texttt{2026\_04\_28\_flow\_grpo\_credit\_summary.tex}，用于本文关于 Flow-GRPO credit assignment = advantage side + ratio side 的统一视角。
\end{itemize}


% Image generation and editing deep-reading notes.

\section*{图像生成与编辑模型}
\phantomsection
\addcontentsline{toc}{section}{图像生成与编辑模型}

\begin{conceptbox}
\textbf{收录范围：} text-to-image、image editing、local editing、visual text generation、identity preservation、controllable generation。\\
\textbf{归档规则：} 如果论文主要贡献是 RL 后训练，即使实验对象是图像生成，也优先放入“后训练与 RL 对齐”；如果主要贡献是图像架构、编辑控制或生产工作流，则放在这里。
\end{conceptbox}

% ============================================================
% Paper deep reading: Wan-Image
% ============================================================

\section{论文精读：Wan-Image}

\begin{conceptbox}
\textbf{论文：} \href{https://arxiv.org/abs/2604.19858}{Wan-Image: Pushing the Boundaries of Generative Visual Intelligence}\\
\textbf{作者/机构：} Wan Team, Alibaba Group\\
\textbf{时间：} arXiv v2, 2026-04-23；报告日期：2026-04-27；阅读日期：2026-04-29\\
\textbf{方向标签：} Image Generation、Image Editing、MLLM Planner、DiT Visualizer、4-channel VAE、Cascade RL、4K Refiner\\
\textbf{模型卡：} Wan2.7-Image；\href{https://create.wan.video/generate/image/generate?model=wan2.7-pro}{create.wan.video}
\end{conceptbox}

\subsection{0. 一句话结论}

Wan-Image 不是一篇“提出单个新模块”的论文，而是一份完整的工业级图像生成系统报告。它的目标不是证明某个纯架构假设，而是把图像模型从 casual image synthesizer 推向 professional productivity tool：复杂文字、身份保持、多图编辑、连续图组、透明 alpha、极端长宽比和 4K 输出都被当成一等能力来设计。

\begin{summarybox}
\textbf{我的判断：} Wan-Image 的核心价值在于系统工程路线：强 MLLM Planner 负责理解、路由和 prompt planning，DiT Visualizer 负责高保真生成，4-channel VAE 解决文字/透明图层/高频细节，Prompt Enhancer 和 Image Refiner 把用户意图与最终生产质量再拉一遍，Cascade RL 负责偏好和设计质量对齐。这篇不适合只读成“又一个 DiT”，而应该读成“专业视觉生产模型需要哪些基础设施”。
\end{summarybox}

\subsection{1. 它想解决什么问题}

论文开场的判断很直接：当前图像生成模型已经能做高审美图片，但还经常卡在生产工作流中的硬约束：

\begin{itemize}[leftmargin=1.5em]
  \item \textbf{复杂 typography：} 长文本、多行排版、海报和信息图中的文字仍容易错字、粘连或布局不稳。
  \item \textbf{严格 controllability：} 真实设计任务要求空间、逻辑、颜色、风格、身份和局部编辑都可控。
  \item \textbf{identity preservation：} 多参考主体、多图编辑、图组生成中，人物和物体身份不能漂移。
  \item \textbf{生产格式：} 电商、游戏、设计素材需要透明 alpha、极端长宽比和 2K--4K 质量，而不是单张 1K 美图。
\end{itemize}

因此 Wan-Image 的问题定义不是“如何提升文生图平均美感”，而是：

\begin{intubox}
如何把 MLLM 的语义理解、任务规划能力，与 DiT 的高质量像素合成能力，组合成一个面向专业生产的统一视觉生成系统？
\end{intubox}

\begin{figure}[htbp]
  \centering
  \includegraphics[width=0.82\linewidth]{wan_image/wan_image_fig1_eval_radar.pdf}
  \caption{Wan-Image 的综合人评雷达图。论文用 T2I/general editing 的 Win Score，以及 interactive editing/image series 的 Pass Rate 来展示能力覆盖。读这张图时要注意：Wan-Image 的卖点不是单项 benchmark，而是专业任务面上的覆盖，包括文字、身份、编辑、图组和设计场景。图源：arXiv source。}
\end{figure}

\subsection{2. 系统总览：Planner + Visualizer}

Wan-Image 的整体结构可以简化成：

\[
\begin{aligned}
\text{User Intent}
&\rightarrow \text{MLLM Planner}
\rightarrow \text{Prompt / Plan}\\
&\rightarrow \text{DiT Visualizer}
\rightarrow \text{Image Refiner}
\end{aligned}
\]

其中 Planner 与 Visualizer 不是简单串两个黑盒。论文称其采用 unified Transformer architecture：理解分支是 decoder-only Transformer，生成分支是 Rectified Flow 下训练的 DiT；二者在统一框架里对齐表示，并在训练上采用 decoupled strategy：先训练理解分支，再用理解分支初始化生成分支的 attention weights，新引入模块随机初始化。

\begin{figure}[htbp]
  \centering
  \includegraphics[width=\linewidth]{wan_image/wan_image_fig3_framework.pdf}
  \caption{Wan-Image 架构。重点看三层耦合：MLLM Planner 负责理解、推理和任务路由；DiT Visualizer 负责生成；VAE token、ViT token 和 text token 通过结构化 attention 组织。它不是 Tuna-2 那类 encoder-free 设计，而是更偏工业系统的 MLLM + ViT + VAE + DiT 组合。图源：arXiv source。}
\end{figure}

\subsubsection{2.1 Attention 与 token 组织}

论文里有几个关键设计：

\begin{itemize}[leftmargin=1.5em]
  \item \textbf{文本、ViT token、VAE token 之间用结构化 attention mask：} 文本侧保持 auto-regressive 生成，视觉 token 内部使用 bidirectional attention。
  \item \textbf{图组生成时，多个图片的视觉 token 之间也启用 bidirectional attention：} 这服务于 image series 的身份一致性和风格连贯性。
  \item \textbf{编辑任务里输入图像的 clean VAE tokens 作为条件输入：} clean tokens 的 timestep 显式设为 0，类似 teacher-forcing 的确定性条件；noised VAE tokens 则走生成分支。
  \item \textbf{位置编码分流：} 理解分支使用 MRoPE，生成分支使用 3D-RoPE，把 \(T,H,W\) 三个维度分开编码，并在 text/image segments 之间加固定 temporal offset 作为语义边界。
\end{itemize}

这里的重点是：Wan-Image 并不追求“把所有东西压成一个最纯的 token stream”。它更像把每种生产能力所需的 inductive bias 都保留下来，再通过统一 Transformer 和训练流程让它们协同。

\subsection{3. 4-channel VAE：为什么它还需要 VAE}

这篇和 Tuna-2 的立场几乎相反。Tuna-2 问的是“能否去掉 VAE 和 vision encoder”；Wan-Image 则认为，在生产级图像系统里，VAE 的质量直接决定生成上限，尤其是：

\begin{itemize}[leftmargin=1.5em]
  \item 小字号文字、密集排版、文档区域；
  \item 复杂纹理、细边界和高频局部结构；
  \item 透明 PNG、半透明边缘、图层素材和电商资产。
\end{itemize}

Wan-Image 提出 high-fidelity four-channel VAE，同时建模 RGB 和 alpha。结构上先做 \(2\times2\) patchify，再经过 \(8\times8\) spatial downsampling，总压缩率是 \(16\times16\)，latent dim 为 48。训练上采用三阶段：

\begin{enumerate}[leftmargin=1.5em]
  \item \(256\times256\) 低分辨率训练，先稳定透明内容和基本结构。
  \item mixed-resolution 训练，增强多尺度泛化。
  \item 加入 4-channel GAN discriminator，提升锐度、透明边界和整体一致性。
\end{enumerate}

此外它还加入 semantic distillation，让 latent 不只保重建像素，也携带更强语义信息。论文报告的重建结果：

\begin{center}
\small
\begin{tabular}{p{0.36\linewidth} p{0.20\linewidth} p{0.16\linewidth} p{0.16\linewidth}}
\toprule
\textbf{设置} & \textbf{PSNR} & \textbf{SSIM} & \textbf{LPIPS}\\
\midrule
RGB, 512p, 5K images & 35.429 & 0.958 & 0.010\\
RGBA, 1080p, 500 images & 38.274 & 0.970 & 0.018\\
\bottomrule
\end{tabular}
\end{center}

\begin{warnbox}
\textbf{和 Tuna-2 的区别：} Tuna-2 去掉 VAE 是为了验证 pixel-space unified modeling 的可行性；Wan-Image 保留并强化 VAE，是因为专业生产要可控、高清、可编辑、可透明合成。二者回答的问题不同，不能简单说谁路线更先进。
\end{warnbox}

\subsubsection{补充：VAE 是否应该保留}

如果把 Wan-Image 和近期一些去 VAE 的路线放在一起看，例如 Tuna-2、SenseNova 这类更偏 encoder-free / pixel-space 的方向，我的判断不是“VAE 必须保留”或“VAE 已经过时”，而是：\textbf{这取决于你到底想优化什么。}

\begin{itemize}[leftmargin=1.5em]
  \item \textbf{更倾向保留 VAE 的场景：} 目标里明确包含小字号文字、密集排版、透明 alpha、局部编辑、身份保持、素材合成、2K--4K 交付，或者你更在意商业工作流里的稳定边界。这类任务本质上都在惩罚“高频细节丢失”和“重建不可控”，而一个足够强的专用 VAE 可以把压缩、重建和语义先验先固定住，降低后端生成器需要同时解决的难题数量。
  \item \textbf{更倾向去掉 VAE 的场景：} 你更在意统一建模纯度、端到端 token 表示、理解与生成共享同一空间、避免 latent bottleneck，以及继续把 scaling law 往 pixel-space 推进。在这种目标下，VAE 本身反而可能变成中间瓶颈：它的压缩误差、频率偏置和重建上限，会提前决定系统能达到的视觉边界。
\end{itemize}

Wan-Image 代表的立场可以概括成一句话：\textbf{它愿意接受 latent bottleneck，换取更强的重建确定性和工程可控性。} 这对专业生产系统是合理的，因为它服务的不是“证明 unified modeling 能否完全去中介”，而是“让文字、透明图层、复杂编辑和高分辨率输出尽量稳定落地”。相对地，Tuna-2 / SenseNova 一类路线更像是在赌另一件事：随着模型规模、算力和 pixel-space 训练方法继续推进，VAE 这个中间压缩器本身，也许才是限制统一能力和端到端扩展性的地方。

\begin{intubox}
所以我的结论不是“VAE 一定要保留”，而是：如果目标是 \textbf{生产交付稳定性}，我会先保留强 VAE；如果目标是 \textbf{统一建模与长期 scaling 潜力}，去 VAE 是值得下注的路线，但要准备为文字、alpha、超高分辨率和编辑稳定性补更多额外工程。
\end{intubox}

\subsection{4. 数据引擎：真正的护城河}

Wan-Image 把数据系统写得很重。数据覆盖理解和生成两大部分。

\subsubsection{4.1 理解数据}

\begin{itemize}[leftmargin=1.5em]
  \item \textbf{General understanding data：} 文本和图文 QA，过滤、平衡、用 LLM/MLLM 重写低质量答案，并构造带 \texttt{<think>} 的 reasoning supervision。
  \item \textbf{Text-proxy data：} 为 interleaved generation 构造密集视觉 prompt。每个要生成的图片位置用 \texttt{<BOI>} 标记，视觉描述放在 \texttt{<imagine>} 与 \texttt{</imagine>} 内。
  \item \textbf{User-aligned query set：} 用 MLLM 按任务类型、难度、是否需要数值/逻辑推理、是否纯 VQA 等属性标注真实用户 query，再按规则采样，提升实用性。
\end{itemize}

\subsubsection{4.2 生成数据}

生成任务覆盖：

\begin{itemize}[leftmargin=1.5em]
  \item T2I：支持长文本渲染和最高 1:8 的极端长宽比。
  \item I2I：编辑、参考生成、多图条件，最多支持 9 张输入图。
  \item T2S / TI2S：一次生成最多 12 张图，强调视觉一致性和逻辑进展。
\end{itemize}

数据获取不是普通 web scrape，而是 multi-modal retrieval + annotation + cleaning 的闭环：支持 image-to-image、多图检索、text-to-image、图文混合和 batch retrieval，并用 cluster-based diversity re-ranking 控制多样性。

\begin{figure}[htbp]
  \centering
  \includegraphics[width=0.88\linewidth]{wan_image/wan_image_fig2_data_taxonomy.pdf}
  \caption{Wan-Image 的 structured data taxonomy。它先把图片路由到 photorealistic、non-photorealistic、text-centric、charts、multi-image compositions 五类，再按类别挂载不同属性维度。这个设计服务于两个目标：减少数据分布偏斜，并让模型学会专业任务需要的细粒度约束。图源：arXiv source。}
\end{figure}

\subsubsection{4.3 Annotation 与过滤}

论文的数据过滤 operator 包括五类：image feature extraction、aesthetic quality、AI-generated content detection、low-level information、overall image quality。具体例子包括：

\begin{itemize}[leftmargin=1.5em]
  \item compression artifact ratio：用理论未压缩大小和实际文件大小的比例判断过压缩。
  \item edge pixel variance：过滤大面积纯色背景和有效信息低的图片。
  \item JPEG bits-per-pixel：作为结构复杂度 proxy，过滤低信息密度样本。
  \item artificiality score：过滤明显 AI 生成或不真实图。
  \item watermark / greasy texture detection：提升数据可用性和真实感。
\end{itemize}

caption 也分三层：

\begin{itemize}[leftmargin=1.5em]
  \item raw captions：保留网页 alt-text、tag、标题等噪声弱监督。
  \item natural language captions：用 MLLM 生成多粒度自然语言描述。
  \item structured JSON captions：按类别填 25 个左右属性维度，再由 MLLM 重写成细节丰富的自然语言。
\end{itemize}

\subsection{5. Prompt Enhancer 与 Image Refiner}

\subsubsection{5.1 Prompt Enhancer}

Prompt Enhancer 是 Wan-Image 很工程化但很关键的一层。它不是单纯扩写 prompt，而是把用户意图翻译成更适合生成模型执行的设计 brief。

PE 有两个版本：

\begin{itemize}[leftmargin=1.5em]
  \item \textbf{Non-CoT：} 基于 Qwen3-VL-2B，输出约 400--600 tokens，面向低延迟，推理约 1 秒。
  \item \textbf{CoT：} 基于 Qwen3-VL-30B-A3B，输出约 1500--2000 tokens，显式做设计推理、元素扩展和空间布局规划。
\end{itemize}

它覆盖 T2I、I2I、T2S、TI2S。对 reference-guided I2I/TI2S，PE 输出 common 和 difference 两个字段：common 描述应保持不变的主体、身份或风格，difference 描述要修改的部分。这个拆分很重要，因为真实编辑任务的难点往往不是“改什么”，而是“什么绝对不能改”。

PE 的训练是 SFT + RL，RL 阶段使用 GRPO，reward 评估 intent fidelity、descriptive richness、linguistic accuracy。论文还提到 prompt rewriting 不是所有任务都强制开启，而是对文字渲染、海报、信息图、UI、storyboard、技术图解、逻辑推理等复杂类别选择性启用。

\subsubsection{5.2 Image Refiner}

Image Refiner 用在 DiT 输出之后，目标是把输出推到 2K--4K，并修复局部缺陷。它有两种模式：

\begin{itemize}[leftmargin=1.5em]
  \item \textbf{repair mode：} 输入生成图和详细 prompt，修复错字、噪声、局部 artifact。
  \item \textbf{high-fidelity mode：} 用在 identity preservation，尽量减少结构和身份漂移。
\end{itemize}

训练时用 blur、noise、lossy compression，以及 diffusion forward-noising/denoising 构造退化输入；之后再用 distillation 和 RL 优化效率与质量。

\subsection{6. 训练流程}

Wan-Image 的训练分为理解、生成、RL、distillation 四段。

\subsubsection{6.1 理解训练}

理解部分从 Qwen-VL base model 初始化，三阶段：

\begin{itemize}[leftmargin=1.5em]
  \item \textbf{Understanding CT：} 46K steps，0.1T tokens，混合图文理解和纯文本数据，保留语言推理能力。
  \item \textbf{Understanding \& Multi-task SFT：} 10K steps，0.033T tokens，用 understanding data 和 text-proxy data 联合训练，让模型能把用户输入转成生成指导。
  \item \textbf{On-policy multi-teacher distillation：} 0.4K steps，面向真实用户 query，用不同 teacher 迁移不同输入类别的处理能力。
\end{itemize}

\subsubsection{6.2 生成训练}

\begin{center}
\footnotesize
\setlength{\tabcolsep}{3pt}
\begin{tabular}{p{0.13\linewidth} p{0.13\linewidth} p{0.15\linewidth} p{0.20\linewidth} p{0.24\linewidth}}
\toprule
\textbf{阶段} & \textbf{Steps} & \textbf{Tokens} & \textbf{Resolution} & \textbf{Ratio}\\
\midrule
PT & 713K & 13.27T & 192 / 320 / 640 & 7:3, T2I/I2I\\
CT & 223K & 8.85T & 512--2048 & 7:2:0.5:0.5\\
SFT & 13K & 0.62T & 512--2048 & curated mix\\
\bottomrule
\end{tabular}
\end{center}

读法很清楚：PT 先靠大规模低中分辨率数据学通用生成；CT 把分辨率拉到 512--2048，并加入 T2S/TI2S，比例为 T2I:I2I:T2S:TI2S = 7:2:0.5:0.5；SFT 用更少但更干净的数据提升审美、细节和指令遵循。

\begin{figure}[htbp]
  \centering
  \includegraphics[width=0.92\linewidth]{wan_image/wan_image_fig5_training_stages.pdf}
  \caption{PT、CT、SFT 阶段输出对比。论文想表达的是训练阶段的职责分工：PT 学基本生成，CT 补高分辨率和复杂条件，SFT 强化美学和局部细节。图源：arXiv source。}
\end{figure}

\subsubsection{6.3 Cascade RL}

RL 阶段是这篇和你后训练专区连接最强的部分，而且它的 framing 很值得注意：Wan-Image 并不把 RL 写成“最后上一遍同一种偏好优化”，而是把它做成一个 \textbf{按 domain 组织、按 reward 形态选工具、按风险逐层推进} 的级联系统。这里的关键词不是某一个算法名，而是 \textbf{Cascade}。

先说论文明写的两条基础原则。第一，Wan-Image 构造 \textbf{domain-wise preference data}，说明偏好数据不是混成一个大池子，而是围绕不同生成场景分别采集和筛选。第二，\textbf{RL prompts 与 SFT prompts 互斥}，这是很重要的反过拟合设计：如果 RL 继续在 SFT 看过的 prompt 上做，很容易把 reward 学成“记住训练分布里什么图更像标注者喜欢”，而不是学到可迁移的偏好对齐能力。

再看图 4，论文把 RL 前面的 data/reward pipeline 单独画了出来：\textbf{Data Collection}、\textbf{Data Pair}、\textbf{Aesthetic Expert}、\textbf{Data Retrieval}、\textbf{Quality Monitoring}、\textbf{Reward Benchmark}、\textbf{Reward Model}。这其实透露了 Wan-Image 的系统思路：\textbf{先做偏好数据工程，再做 reward 工程，最后才是 RL 工程}。也就是说，它并不把 reward model 当成一个“训完就用”的静态打分器，而是持续用更广的 benchmark 和人工质检去检查 reward 是否偏离真实偏好。

从图里的维度名还能看出 reward 不是单一 aesthetic score，而是至少覆盖：\textbf{Realism、Text、Harmony、Structure、Aesthetics、Consistency}。这和前面模型的产品目标是对齐的：

\begin{itemize}[leftmargin=1.5em]
  \item \textbf{Realism：} 针对写实感、材质和整体自然度；
  \item \textbf{Text：} 针对文字可读性、拼写与排版；
  \item \textbf{Structure：} 针对空间关系、布局和物体结构正确性；
  \item \textbf{Aesthetics：} 针对构图、色彩和设计感；
  \item \textbf{Consistency / Harmony：} 更偏编辑、多图和组合场景里的身份保持、风格一致与局部协调。
\end{itemize}

这个设计的直接含义是：\textbf{Wan-Image 的 RL 目标不是把所有样本都往“更漂亮”一个方向推，而是按任务暴露出的短板做多维对齐。} 这也解释了为什么它最后不是单一 RL，而是 DPO、ReFL、DenseGRPO、ReFMA 的组合。

Cascade RL 组合了：

\begin{itemize}[leftmargin=1.5em]
  \item DPO；
  \item ReFL；
  \item DenseGRPO；
  \item in-house ReFMA，即 reinforcement learning with flow matching anchors。
\end{itemize}

我的理解是，这四类方法并不是互相替代，而是分别覆盖不同风险/成本/信用分配条件下的优化位置。

\textbf{1. DPO：先用离线 pairwise preference 把方向拉正。}

DPO 的优点是稳定、便宜、工程风险低。给定同一 prompt 下的 winner/loser 对，模型直接学“让 winner 相对 loser 更可能”，同时参考模型提供一个 KL-like anchor，避免分布漂太远。放到 Wan-Image 这里，DPO 很适合作为级联里的低风险起点：\textbf{先把明显的人类偏好方向学进去，例如更好的排版、更合理的构图、更顺眼的整体风格，而不急着上高成本在线探索。}

但 DPO 的边界也很明确：它主要吃现有偏好对，\textbf{不会自己探索新轨迹}，也不能细致回答“到底是哪几个 denoising steps 把图像变好了”。所以它更像是偏好对齐的第一层，不太像最终解法。

\textbf{2. ReFL：当 reward 可微时，直接把 reward 梯度压回生成过程。}

从 Wan-Image 把 ReFL 放进同一工具箱来看，它显然想利用一类和 policy-ratio 不同的优化信号：\textbf{如果 reward model 可微，就可以不只靠采样后比较 winner/loser，也可以直接沿 reward gradient 优化生成轨迹。} 这类方法的直觉通常是：对那些 reward 能稳定评估、而且修改更像局部连续修正的目标，直接梯度会比纯采样式 RL 更高效。

结合你前面后训练主线里的经验，可以把 ReFL 粗略记成“更偏 direct reward optimization 的路线”。它通常更容易先改善后期可见的质量问题，比如局部细节、审美微调、文字边缘和表面纹理。但它的前提也更苛刻：\textbf{reward 要可微，训练稳定性也更依赖 reward 本身的质量。}

\textbf{3. DenseGRPO：解决 flow/diffusion 里的时间步信用分配。}

这部分最值得和你已经读过的 DenseGRPO 对照。标准 Flow-GRPO/GRPO 的一个硬伤是：最终图像只有一个 terminal reward，但整条生成轨迹有很多 denoising steps。如果把同一个终局 reward 平均分给所有 timestep，就等于默认“最终图像好，说明每一步都好”，这通常不成立。

DenseGRPO 的价值正是在这里。它不满足于 trajectory-level reward，而是试图估计每一步到底带来了多少 reward gain。具体做法是：从中间 latent 通过 ODE readout 还原出可打分的 clean prediction，再看相邻 step 的 reward 提升量。这样得到的是 step-wise dense reward，而不是把同一个 terminal reward 广播到全轨迹。

放到 Wan-Image 的语境里，DenseGRPO 很适合那些\textbf{明显依赖生成过程信用分配}的目标，例如结构稳定性、文字渲染和多对象空间关系。因为这些问题往往不是“最后一张图总体不好看”这么简单，而是某些关键步把 layout、shape 或字符骨架画错了。如果 credit assignment 太粗，模型只知道“这张图分低”，却不知道应该改前期布局还是后期 refinement。

当然，DenseGRPO 的代价也很实在：中间步读出和打分很贵，工程上明显重于 DPO。Wan-Image 愿意把它放进 Cascade RL，本身就说明他们认为单靠离线偏好或整图 reward 已经不够解决生产场景里的细粒度问题。

\textbf{4. ReFMA：用 flow-matching anchor 稳住 reward-driven optimization。}

论文对 ReFMA 的公开细节没有 DenseGRPO 那么充分，但从名字和图示可以读出两个关键信号。第一，它是 \textbf{flow matching anchors}，说明它不是完全放飞的 reward maximization，而是保留一个和原始 flow trajectory/teacher policy 相关的锚点。第二，现有描述强调 \textbf{在不同 noise levels 之间平衡优化目标}，并使用 \textbf{reward model 的梯度信号} 指导优化。

这意味着 ReFMA 很可能是在解决一个非常实际的问题：\textbf{直接用 reward 去推 flow model，容易在某些噪声区间过度优化，导致分布漂移、结构塌陷或者局部 reward hacking；anchor 的作用就是把模型拴回原始可用分布附近。} 如果这么理解，ReFMA 扮演的是“高强度 reward optimization 的安全带”。

我会把它看成 Wan-Image 在工业落地里的关键补件：\textbf{既想吃到 reward-gradient 带来的强对齐收益，又不愿意让模型因为几类 reward 被拉坏基础生成分布。} 这和纯研究论文常见的“只报 reward 提升”不同，Wan-Image 更在乎长期可用性和跨任务稳定性。

\textbf{把四者连起来看，Cascade RL 背后的逻辑就清楚了。}

如果只用 DPO，优点是稳，但缺少在线探索和 step-level credit；如果只用 DenseGRPO，credit 更准，但成本高；如果只用 ReFL 一类 direct-gradient 方法，优化效率高，但受 reward 可微性和稳定性约束；如果只有强 reward gradient 而没有 anchor，又容易过冲。Wan-Image 把它们级联，本质上是在做下面这件事：

\begin{itemize}[leftmargin=1.5em]
  \item 先用更稳的离线偏好方法把大方向拉正；
  \item 再用更强的 reward-based 方法补足在线探索和细粒度 credit；
  \item 最后用 anchor/control 机制把 reward optimization 限制在可用分布附近。
\end{itemize}

所以这部分最值得记住的，不是“Wan-Image 也用了 DPO/DenseGRPO”，而是：

\begin{intubox}
对于专业图像系统，RL 不是一个统一算法，而是一套要和数据分区、reward 维度、生成轨迹结构、稳定性约束一起设计的后训练工具链。
\end{intubox}

论文报告 Cascade RL 相比 SFT baseline 在不同场景下 Win Score 提升约 15\%--50\%。结合图 4 的 reward 维度，这个提升不应理解成单纯 aesthetic bump，而更像是对 \textbf{photorealism、structural correctness、text rendering、aesthetic quality，以及 I2I consistency/harmony} 这些生产敏感指标的系统性补强。

\begin{figure}[htbp]
  \centering
  \includegraphics[width=0.92\linewidth]{wan_image/wan_image_fig4_rl_framework.pdf}
  \caption{Cascade RL 框架和相对 SFT baseline 的提升。图里除了 DPO、ReFL、DenseGRPO、ReFMA 的级联关系，还明确画出了 reward system 的前置基础设施：data pair、aesthetic expert、data retrieval、quality monitoring、reward benchmark、reward model，以及 realism/text/harmony/structure/aesthetics/consistency 六个对齐维度。图源：arXiv source。}
\end{figure}

\subsubsection{6.4 Distillation}

蒸馏在 Wan-Image 里不是泛指的“模型压缩”，而是非常具体的 \textbf{inference step distillation}：目标不是再造一个更小的基座模型，而是在尽量不破坏生成能力的前提下，把原始采样流程压缩到更少 step，从而满足真实交互应用对低延迟的要求。论文特别指出，尽管 base model 的参数规模已经相对克制，但其原始的 iterative generation scheme 再叠加 classifier-free guidance (CFG)，仍然会带来显著计算负担；部署时真正昂贵的，不只是参数量本身，而是一次生成要做多少次 denoising/function evaluation，也就是 NFE。

这里最重要的理解是：\textbf{Wan-Image 要解决的不是一般意义上的“少步采样会降质”，而是少步采样会优先破坏那些最有产品价值的能力。} 论文明确点出，已有 step distillation 方法虽然可以提升速度，但往往会在 \textbf{high-frequency facial details} 和 \textbf{text rendering capability} 上出现明显退化。这两个点之所以值得作者单独拿出来说，是因为它们恰好对应 Wan-Image 想重点建立优势的使用场景：真实人像、广告图文、海报排版、品牌素材，以及需要多轮编辑的商业工作流。也就是说，Wan-Image 的蒸馏目标不是“把图粗略生成出来”就够了，而是要求在提速之后，仍尽量保住脸部微细节、文字可读性、局部纹理锐度，以及整体结构稳定性。

为了解这个矛盾，论文没有采用单一蒸馏技巧，而是设计了一个 \textbf{multi-stage progressive distillation} 框架，把两类互补思路组合起来：trajectory distillation 和 distribution matching distillation。虽然这一节没有公开每个 stage 的精确步数、loss 权重和 scheduler 细节，但从作者对目标的描述可以比较清楚地看出，两者扮演的角色并不一样。

\begin{itemize}[leftmargin=1.5em]
  \item \textbf{Trajectory distillation：} 让 student 用更少的采样步去逼近 teacher 的 denoising trajectory。按我的理解，它更偏向“过程对齐”，核心作用是保住全局布局、主体关系、结构演化，以及多步去噪过程中逐步收敛的动态。对于高容量 teacher 来说，这一步尤其重要，因为 teacher 的去噪轨迹通常比小模型更复杂，如果直接从多步 teacher 硬压到极少步 student，学生模型很容易只学到一个粗糙终态，结果就是构图漂移、局部关系不稳，或者 prompt 条件在中间过程里被过早丢失。
  \item \textbf{Distribution matching distillation：} 进一步把 student 的输出分布往 teacher 的高质量分布上拉齐。它更像“结果对齐”，关注的不是每一步都长得像 teacher，而是最后生成出来的样本在统计意义上仍然落在 teacher 擅长的那类结果附近。结合论文强调的退化点，这一步可以理解为对 trajectory imitation 的补充：前者负责守住结构与轨迹，后者负责补 realism、纹理锐度和文本 fidelity，尽量避免少步近似之后出现的过平滑、字形塌陷、局部细节发糊等问题。
\end{itemize}

为什么还要做成 \textbf{progressive} 的多阶段蒸馏？一个合理解释是，作者并不想让 student 一次性从“高 NFE teacher”直接跳到“极少步部署模型”，而是把压缩过程做成一种 curriculum。这样每个阶段只承担一部分步数压缩和误差吸收，训练更容易稳定，也更容易在速度和保真之间找到可接受的折中。论文虽然没有展开完整训练配方，但从问题设定上看，这种 staged compression 本质上是在规避一个老问题：单次激进压步通常能换来更快采样，却很容易把 fidelity 尤其是细节 fidelity 一起压坏。

所以这部分最值得记住的，不是“Wan-Image 也做了蒸馏”，而是：

\begin{intubox}
Wan-Image 把蒸馏当成“保持能力边界的采样压缩”，而不是单纯追求更少 step。真正要守住的是 \textbf{全局结构、高频细节和文本对齐} 这三件事。
\end{intubox}

把它放回整篇论文的系统视角里看，这一步其实是前面所有能力模块的产品化收口：VAE 负责高压缩下的视觉重建，thinking/instruct 负责复杂理解与编辑控制，Cascade RL 负责把生成偏好往高质量分布推，而 distillation 则负责把这些能力压进真实可部署的延迟预算里。论文给出的结论是，经过这一套 multi-stage progressive distillation 后，推理流程获得了超过 \(8\times\) 的 speedup，同时仍能较好保持 generation fidelity。这里的“保真”更应理解为系统级保真，而不是所有子任务严格无损；但从作者特意强调脸部高频细节、文本渲染，以及 Figure 19 中 teacher vs. distilled model 的对比可以看出，他们真正担心的，就是那些最容易在少步采样里首先塌掉、却又最影响商业可用性的能力。
\subsection{7. 结果怎么读}

论文报告的核心结果包括：

\begin{itemize}[leftmargin=1.5em]
  \item \textbf{人评：} 整体超过 Seedream 5.0 Lite 和 GPT Image 1.5，在困难任务上接近 Nano Banana Pro。
  \item \textbf{交互编辑和图组生成：} pass rate 大约 80\%，强调真实工作流可用性。
  \item \textbf{理解能力：} Ours Thinking 在 7 个 image understanding benchmark 的平均分为 76.3，高于很多 unified models；text understanding 平均也高于其 instruct mode。
  \item \textbf{VAE：} RGB/RGBA 重建指标优于或接近强 baseline，尤其强调 \(16\times16\) 高压缩下的细节与 alpha。
  \item \textbf{RL：} Cascade RL 对 SFT baseline 有 15\%--50\% Win Score 提升。
\end{itemize}

但要谨慎：这篇没有给出完全可复现的开放训练配方，也没有把所有模块做严格公开 ablation。很多结论来自系统级人评和质性展示，适合判断产品能力和工程方向，不适合直接当作单一算法的严格因果证明。

\subsection{8. 与 Tuna-2 / BAGEL 的关系}

\begin{center}
\footnotesize
\setlength{\tabcolsep}{3pt}
\begin{tabular}{p{0.15\linewidth} p{0.25\linewidth} p{0.23\linewidth} p{0.25\linewidth}}
\toprule
\textbf{模型} & \textbf{视觉表示} & \textbf{生成路径} & \textbf{核心取舍}\\
\midrule
Tuna-2 & pixel embeddings；无 VAE/ViT & pixel-space flow & 架构假设更激进，证明 encoder-free scaling；生成侧未必最强。\\
BAGEL & SigLIP + VAE latent & Flow Matching + MoT & 理解/生成分路，强调 native unified。\\
Wan-Image & ViT + VAE + MLLM Planner & DiT + Refiner & 工业系统路线，保留并强化视觉先验，服务专业生产约束。\\
\bottomrule
\end{tabular}
\end{center}

Wan-Image 说明一个很现实的点：即使 encoder-free/pixel-space 是很有研究价值的方向，生产级图像系统仍可能选择强 VAE、强 ViT、强 MLLM、强 refiner 的组合。因为业务目标不是架构纯度，而是文字、身份、透明图层、编辑和 4K 交付的稳定性。

\subsection{9. 局限与风险}

\begin{itemize}[leftmargin=1.5em]
  \item \textbf{不可复现：} 数据规模、过滤器、reward model、teacher model 和私有标注流程都高度闭源。
  \item \textbf{系统归因困难：} 最终质量来自 PE、Planner、VAE、DiT、Refiner、RL、distillation 多模块叠加，很难知道单个模块贡献。
  \item \textbf{人评标准不透明：} Win Score 和 Pass Rate 有参考意义，但任务分布、标注一致性和失败样例覆盖仍需更细披露。
  \item \textbf{PE 可能改变用户意图：} prompt rewriting 提升美学和细节，但也可能把原始需求过度设计化。
  \item \textbf{RL 可能强化审美模式：} Cascade RL 会提升人类偏好分数，但也可能收缩风格多样性或偏向 reward model 的偏好。
  \item \textbf{4K 质量部分依赖 refiner：} 需要区分 base generation capability 和 post-refinement capability。
\end{itemize}

\subsection{10. 对我的可复用点}

\begin{itemize}[leftmargin=1.5em]
  \item \textbf{生产级指标先行：} 文字、身份、alpha、图组、局部编辑、极端比例，比普通 aesthetic prompt 更能暴露模型短板。
  \item \textbf{common/difference 编辑表示：} 对 reference editing 很实用，明确“保留什么”和“修改什么”。
  \item \textbf{数据 taxonomy + structured JSON caption：} 对图像生成训练比单纯自然语言 caption 更可控，尤其适合长尾任务。
  \item \textbf{RL 不是单算法：} 可以把 DPO、ReFL、DenseGRPO、flow-matching anchor 类方法按领域级联，而不是期待一个 reward objective 解决所有偏好。
  \item \textbf{VAE 是否保留取决于目标：} 如果目标是透明图层、小字、文档和 4K，专用 VAE 的工程价值很高；如果目标是追求 encoder-free 的统一建模与长期 scaling，可以继续关注 pixel-space 路线。
\end{itemize}

\subsection{11. 待查问题}

\begin{itemize}[leftmargin=1.5em]
  \item ReFMA 的具体 objective、anchor 设计和与 DenseGRPO 的衔接，论文没有完全展开，需要等技术细节或代码。
  \item unified Transformer 中 dedicated experts 的参数共享边界、attention 共享方式和推理时调度细节，需要进一步核对。
  \item Prompt Enhancer 的选择性触发规则是否会导致用户意图漂移，最好看更多实际案例。
  \item Refiner 对身份保持和文字修复的失败模式还不清楚，尤其是是否会在高分辨率增强时改动关键内容。
  \item 与 Nano Banana Pro、GPT Image 1.5、Seedream 的公平对比需要更透明的 prompt 集和人评协议。
\end{itemize}

\subsection{12. 后续阅读顺序}

\begin{itemize}[leftmargin=1.5em]
  \item \textbf{先对照 Wan2.2 / Wan video：} 看 Wan 系列的 VAE、DiT 和视频生成技术如何迁移到 image。
  \item \textbf{再读 DenseGRPO / Flow-GRPO / ReFL：} 补齐 Cascade RL 背后的偏好对齐工具。
  \item \textbf{对照 Tuna-2：} 前者看 encoder-free pixel scaling；后者看 VAE、ViT、MLLM 的生产系统栈。
  \item \textbf{对照 BAGEL：} 比较 native unified multimodal architecture 与 Wan-Image 的 planner-visualizer 工程路线。
\end{itemize}


% Video generation and editing deep-reading notes.

\section*{视频生成与编辑模型}
\phantomsection
\addcontentsline{toc}{section}{视频生成与编辑模型}

\begin{conceptbox}
\textbf{收录范围：} text-to-video、image-to-video、video editing、long-video DiT、controllable video、physical plausibility、audio-video generation。\\
\textbf{归档规则：} 如果论文主要贡献是 video model architecture、video editing control 或视频生成评测，放在这里；如果核心是用 RL 优化视频模型，则优先放入“后训练与 RL 对齐”。
\end{conceptbox}

\note{当前暂无单篇精读。候选：World-R1、FlowAnchor、FreqFormer、PhysLayer、Talker-T2AV。}

% Unified understanding and generation deep-reading notes.

\section*{统一理解与生成模型}
\phantomsection
\addcontentsline{toc}{section}{统一理解与生成模型}

\begin{conceptbox}
\textbf{收录范围：} unified understanding \& generation、any-to-any、interleaved multimodal generation、shared/decoupled visual representation、LLM + diffusion/flow hybrid。\\
\textbf{比较维度：} 视觉表示是 pixel embedding、VAE latent、VQ token 还是视觉编码器特征；理解和生成是共享主干、双编码器还是 MoT/adapter 分路。
\end{conceptbox}

% ============================================================
% Paper deep reading: Tuna-2
% ============================================================

\section{论文精读：Tuna-2}

\begin{conceptbox}
\textbf{论文：} \href{https://arxiv.org/abs/2604.24763}{Tuna-2: Pixel Embeddings Beat Vision Encoders for Multimodal Understanding and Generation}\\
\textbf{作者/机构：} Zhiheng Liu, Weiming Ren, Xiaoke Huang, Shoufa Chen, Tianhong Li, Mengzhao Chen, Yatai Ji, Sen He, Jonas Schult, Belinda Zeng, Tao Xiang, Wenhu Chen, Ping Luo, Luke Zettlemoyer, Yuren Cong；Meta AI / The University of Hong Kong / University of Waterloo\\
\textbf{时间：} arXiv v1, 2026-04-27；阅读日期：2026-04-28\\
\textbf{方向标签：} Unified Multimodal Model、Pixel Embedding、Encoder-Free、Pixel-space Flow Matching、Image Understanding / Generation / Editing\\
\textbf{项目页 / 代码：} \href{https://tuna-ai.org/tuna-2/}{tuna-ai.org/tuna-2} · \href{https://github.com/facebookresearch/tuna-2}{facebookresearch/tuna-2}
\end{conceptbox}

\subsection{0. 一句话结论}

Tuna-2 的核心不是又做一个 unified multimodal model，而是直接挑战当前统一理解与生成模型里的默认假设：\textbf{图像一定要先经过 VAE、SigLIP/CLIP 或其它 pretrained vision encoder 才能进入统一模型}。作者把视觉侧压到最简：图片只经过 patch embedding，后续理解和生成都交给同一个 LLM decoder + pixel-space flow head。

\begin{summarybox}
\textbf{我的判断：} 这篇值得精读，因为它把 unified model 的争论从“用哪种视觉编码器”推进到“是否还需要视觉编码器”。但结论不能简单读成 encoder 已经过时。更准确地说：在足够大规模 end-to-end pretraining 下，encoder-free pixel-space UMM 可以在理解侧超过 encoder-based variant；生成和编辑侧目前仍能从 representation encoder 的语义先验中受益。
\end{summarybox}

\subsection{1. 它在统一理解生成版图里的位置}

你已有的 BAGEL / BLIP3o / Transfusion / Janus 这条主线，核心矛盾都是：

\[
\text{理解需要语义表示} \quad \text{vs.} \quad \text{生成需要可重建表示}
\]

常见解法有三类：

\begin{itemize}[leftmargin=1.5em]
  \item \textbf{解耦表示：} 理解走 CLIP/SigLIP，生成走 VAE/VQ，代表 Janus 一类。
  \item \textbf{统一 latent / tokenizer：} 尝试训练一个同时适合理解与生成的连续或离散视觉表示。
  \item \textbf{LLM + diffusion/flow hybrid：} 文本和理解由 LLM 处理，图像生成由 diffusion/flow head 处理，代表 Transfusion、BAGEL、BLIP3o 的不同变体。
\end{itemize}

Tuna-2 的切入点更激进：\textbf{不要专门的视觉编码器，也不要 VAE latent，直接在 pixel patch 上做统一建模。}

\begin{center}
\begin{tabular}{p{0.22\linewidth} p{0.33\linewidth} p{0.34\linewidth}}
\toprule
\textbf{模型} & \textbf{视觉表示} & \textbf{关键取舍}\\
\midrule
Tuna & VAE encoder + representation encoder & 统一表示，但仍依赖两个视觉编码模块。\\
Tuna-R & representation encoder + pixel-space generation & 去掉 VAE，但保留 SigLIP 2 语义先验。\\
Tuna-2 & patch embedding / pixel embeddings & 去掉 VAE 和 representation encoder，把视觉建模压力全部交给统一 transformer。\\
\bottomrule
\end{tabular}
\end{center}

\begin{figure}[htbp]
  \centering
  \includegraphics[width=\linewidth]{tuna2/tuna2_fig1_architecture.pdf}
  \caption{Tuna-2 的架构演进图。阅读这张图时重点看两次“删模块”：Tuna-R 去掉 VAE，把生成从 latent space 推到 pixel space；Tuna-2 再去掉 representation encoder，只保留 patch embedding。这个设计让论文的问题变得很干净：视觉侧先验越少，统一 transformer 是否还能自己学出理解和生成所需的表示？图源：arXiv source。}
\end{figure}

\subsection{2. 论文真正要回答的问题}

\textbf{问题不是：} pixel-space generation 能不能做。这个方向已有 PixelFlow / DiP / PixelDiT / JiT 等工作。

\textbf{问题也不是：} encoder-free VLM 能不能做。Fuyu、Chameleon、Mono-InternVL、NEO 等已经探索过。

\textbf{Tuna-2 真正问的是：}

\begin{intubox}
能不能把 encoder-free visual understanding 和 pixel-space visual generation 合到一个 native UMM 里，并且在同一个训练框架下证明：不用 pretrained vision encoder 也能同时做好理解、生成和编辑？
\end{intubox}

这个问题重要，因为当前统一理解与生成模型经常把视觉编码器当成先验基础设施。Tuna-2 如果成立，说明统一模型的 scaling path 可以更接近语言模型：输入侧尽量简单，表示能力从大规模端到端训练中长出来。

\subsection{3. 方法拆解}

\subsubsection{3.1 架构：从 Tuna 到 Tuna-R 再到 Tuna-2}

作者不是直接宣称 encoder-free 最好，而是做了一个受控阶梯：

\begin{enumerate}[leftmargin=1.5em]
  \item \textbf{Tuna：} 图片先经过 VAE encoder，再经过 representation encoder，形成统一视觉表示；LLM decoder 负责图文建模；语言输出走 LM head，图像生成走 flow matching head。
  \item \textbf{Tuna-R：} 去掉 VAE，只保留 representation encoder。理解类似常规 LMM；生成改成 pixel-space flow matching。
  \item \textbf{Tuna-2：} 连 representation encoder 也去掉，只用简单 patch embedding 把原始图片变成视觉 token，再交给 LLM decoder。
\end{enumerate}

Tuna-2 的结构可以简化成：

\[
\begin{aligned}
\text{Image Pixels}
&\rightarrow \text{Patchify Layer}
\rightarrow \text{LLM Decoder} \\
&\rightarrow
\begin{cases}
\text{Language Head} & \text{理解/文本输出}\\
\text{Unpatchify + Flow Head} & \text{图像生成/编辑}
\end{cases}
\end{aligned}
\]

这和 BAGEL 的差异很明显：

\begin{itemize}[leftmargin=1.5em]
  \item BAGEL 仍有 SigLIP 理解侧和 VAE latent 生成侧；Tuna-2 把两者都拿掉。
  \item BAGEL 用 MoT 把理解/生成参数分路；Tuna-2 更像单一 decoder 承担视觉语言表示学习。
  \item BAGEL 的生成发生在 VAE latent / Flow Matching 空间；Tuna-2 的 flow matching 发生在 pixel space。
\end{itemize}

\subsubsection{3.2 Pixel-space Flow Matching}

去掉 VAE 后，不能再沿用 latent diffusion / latent flow 的老路径。Tuna-2 借鉴 JiT 的 x-prediction + v-loss 范式，直接在像素空间构造 rectified flow。

给定 clean image \(x_1\)、高斯噪声 \(x_0 \sim \mathcal{N}(0,I)\)、时间 \(t \in [0,1]\)，构造：

\[
x_t = t x_1 + (1-t)x_0
\]

模型直接预测 clean image：

\[
x_\theta = \pi_\theta(x_t, c, t)
\]

其中 \(c\) 是条件信号：T2I 时是文本，编辑时是文本 + 图像。虽然模型输出是 \(x_\theta\)，训练目标仍转成 velocity：

\[
v_\theta = \frac{x_\theta - x_t}{1-t}, \qquad v = x_1 - x_0
\]

\[
\mathcal{L}_{\mathrm{flow}} =
\mathbb{E}_{t,c,x_1,x_0}\left[\|v_\theta - v\|_2^2\right]
\]

这里容易误读的一点是：\textbf{加噪发生在完整图像 \(x_1\) 上，而不是只发生在某些 patch 上}。更具体地说，生成侧可以按下面的顺序理解：

\begin{enumerate}[leftmargin=1.5em]
  \item 先在 pixel space 对全图采样噪声，得到完整 noisy image \(x_t\)。
  \item 再把 \(x_t\) patchify 成 noisy visual tokens。
  \item 如果启用 masking-based feature learning，就在 token 输入层随机挑一部分 noisy visual tokens，用 learnable mask token 替换，得到送入 LLM decoder 的 masked visual input。
  \item 模型预测 clean image patches，经 unpatchify 回到图像空间，或在 patch tensor 上做等价计算。
  \item flow matching loss 仍然使用原始全图的 \(x_t,x_1,x_0\) 来构造 \(v_\theta\) 和 \(v\)，并对 masked 与 unmasked regions 一起回归。
\end{enumerate}

\begin{intubox}
\textbf{关键区分：} noisy visual token 来自 rectified flow 的连续噪声路径；learnable mask token 只是 patch/token 级输入遮挡正则。Tuna-2 不是同时训练一条“加噪 flow”和一条“mask flow”，而是在同一个 pixel-space flow matching 目标上加入 masked prediction 难度。
\end{intubox}

推理时用 Euler solver，根据 velocity 从 noisier image 逐步走向 denoised image。

\textbf{关键理解：} Tuna-2 不是把图片压到低维 latent 再生成，而是把高维 pixel denoising 的难题直接交给大模型。它的代价是训练更难、token/计算压力更大；收益是低层细节没有被 VAE 或 representation encoder 提前丢掉。

\subsubsection{3.3 Masking-based Feature Learning}

作者认为 pixel space 冗余更高，模型容易学 shortcut，而不是学真正对理解和生成都有用的视觉线索。因此引入 masking-based visual feature learning：

\begin{itemize}[leftmargin=1.5em]
  \item 随机选择一部分 image patches。
  \item 用 learnable mask token 替换这些视觉 token。
  \item generation examples 中，模型基于部分可见的 noisy visual tokens 和 mask tokens，仍要预测完整 clean image patches；loss 覆盖 masked 和 unmasked regions。
  \item understanding examples 中，模型要基于被遮挡的视觉输入预测文本答案。
\end{itemize}

\begin{intubox}
这一步可以理解成：对生成任务，它让 denoising 问题更难；对理解任务，它迫使模型不要只依赖局部显著纹理，而要在部分观测下建立更稳的视觉语义。
\end{intubox}

\begin{figure}[htbp]
  \centering
  \includegraphics[width=0.96\linewidth]{tuna2/tuna2_fig3_masking.pdf}
  \caption{Masking-based feature learning。左侧是生成样本：在带噪视觉 token 和 learnable mask token 条件下预测完整 clean image；右侧是理解样本：遮挡一部分视觉 token 后继续预测新的文本 token。我的读法是，这不是单纯的 MAE trick，而是给 encoder-free pixel embedding 补一个“被迫建模全局语义”的正则。图源：arXiv source。}
\end{figure}

论文的 ablation 显示 masking 对 Tuna-2 的提升比 Tuna-R 更明显。作者的解释是 Tuna-R 的 SigLIP 2 encoder 本身已经带有类似 masked prediction 的预训练先验，而 Tuna-2 需要在端到端训练中自己学出来。

\subsubsection{3.4 两阶段训练}

\begin{itemize}[leftmargin=1.5em]
  \item \textbf{Stage 1 full model pretraining：} 联合训练 image captioning 和 text-to-image generation，目标是建立 flow matching head 的初始化，并让 pixel-space visual input 适配统一理解/生成。
  \item \textbf{Stage 2 SFT：} 用 image instruction-following、image editing、高质量图像生成数据做全模型低学习率微调，提升编辑、指令跟随和生成质量。
\end{itemize}

具体训练配方：

\begin{itemize}[leftmargin=1.5em]
  \item LLM decoder：Qwen2.5-7B-Instruct。
  \item Stage 1：550M in-house image-text pairs，其中 70\% image captioning、30\% text-to-image；另有 Nemotron text-only data，占总预训练数据 20\%。
  \item Stage 1 训练：300k steps，64 nodes，AdamW，learning rate \(1 \times 10^{-4}\)，每 GPU sequence length pad 到 16k tokens。
  \item Stage 2：13M FineVision 对话样本、约 2M OmniEdit 图像编辑样本，以及高质量图像生成数据。
  \item Stage 2 训练：50k steps，AdamW，learning rate \(2 \times 10^{-5}\)。
  \item Tuna-R 额外有 connector alignment stage：3k steps，learning rate \(5 \times 10^{-4}\)。
\end{itemize}

\subsection{4. 主要实验结论}

\subsubsection{4.1 图像理解：Tuna-2 强于 Tuna-R 的地方最多}

Table 1 的关键信息是：Tuna-2 在 7B native UMM 中表现很强，尤其在 pixel-centric benchmarks 上优势明显。

\begin{center}
\begin{tabular}{p{0.25\linewidth} p{0.18\linewidth} p{0.18\linewidth} p{0.28\linewidth}}
\toprule
\textbf{Benchmark} & \textbf{Tuna-R} & \textbf{Tuna-2} & \textbf{变化}\\
\midrule
GQA & 63.5 & 65.0 & Tuna-2 更好\\
MMVet & 46.7 & 51.7 & 明显更好\\
MMVP & 74.7 & 77.3 & 更好\\
SEED-Bench2+ & 58.4 & 61.1 & 更好\\
OCRBench & 78.3 & 79.7 & 更好\\
V* & 57.6 & 59.2 & 更好\\
CountBench & 77.8 & 81.7 & 明显更好\\
VisuLogic & 26.2 & 28.8 & 更好\\
\bottomrule
\end{tabular}
\end{center}

不是所有理解指标都赢：RealWorldQA、MMMU 上 Tuna-R 略高，ChartQA 持平。但总体上，Tuna-2 说明 encoder-free pixel embedding 在大规模训练后并不弱，尤其适合低层细节、计数、小目标、视觉逻辑等 fine-grained perception。

\subsubsection{4.2 图像生成：Tuna-2 竞争力强，但不是最强}

在 GenEval / DPG-Bench：

\begin{itemize}[leftmargin=1.5em]
  \item Tuna：GenEval overall 0.90，DPG overall 86.76。
  \item Tuna-R：GenEval overall 0.88，DPG overall 86.35。
  \item Tuna-2：GenEval overall 0.87，DPG overall 86.54。
\end{itemize}

所以生成侧结论要谨慎：Tuna-2 可以接近强 unified models，但并没有全面超过 Tuna / Tuna-R。论文自己的分析也承认 representation encoder 的语义先验对 generation 仍有帮助。

\begin{figure}[htbp]
  \centering
  \includegraphics[width=0.92\linewidth,height=0.68\textheight,keepaspectratio]{tuna2/tuna2_fig2_generation_samples.pdf}
  \caption{Tuna-2 的文生图与图像编辑样例。它的意义主要是定性证明 encoder-free 路线不等于“只能理解、不能生成”：没有 VAE 和视觉表征 encoder，模型仍能生成高保真图片并做指令编辑。但这张图只能支持能力展示，不能单独推出 Tuna-2 在生成侧全面优于 Tuna/Tuna-R；定量表里 Tuna-2 的生成优势并不如理解侧稳定。图源：arXiv source。}
\end{figure}

不过 LLM-judge evaluation 给了另一个角度：在 1.5K prompts、每 prompt 4 张图的评测中，Tuna-2 在 diversity 上被 GPT-5.4 和 Claude Opus 4.7 明显偏好；quality 上和 Tuna-R 接近，高于 Tuna。这说明 encoder-free 设计可能减少某些视觉先验带来的模式收缩，让输出多样性更强。

\subsubsection{4.3 图像编辑：强，但仍落后于 Tuna / Tuna-R}

ImgEdit 总分：

\begin{itemize}[leftmargin=1.5em]
  \item Tuna：4.31。
  \item Tuna-R：4.18。
  \item Tuna-2：4.09。
  \item GPT-Image：4.20。
  \item BAGEL：3.20，OmniGen2：3.44。
\end{itemize}

Tuna-2 超过早期 unified editing baselines，但和 Tuna / Tuna-R 有差距。这很合理：编辑既需要低层保真，也需要对原图语义和局部区域的稳健理解；pretrained representation encoder 仍可能提供更强的编辑先验。

\subsubsection{4.4 图像重建：pixel-space unified representation 不是空话}

作者做了轻量 reconstruction finetuning，在 ImageNet validation 上比较 visual tokenizers。Tuna-2 的结果：

\begin{itemize}[leftmargin=1.5em]
  \item Tuna-R：rFID 0.12，PSNR 32.22，SSIM 0.93。
  \item Tuna-2：rFID 0.15，PSNR 32.80，SSIM 0.93。
  \item FLUX.1 VAE：rFID 0.06，PSNR 33.65，SSIM 0.93。
\end{itemize}

这说明 Tuna-2 的 pixel representation 在 reconstruction 上接近强 VAE，但还没有超过专用 VAE。它证明的是“统一模型中的 pixel embedding 可用”，不是“专用 tokenizer 已经不需要”。

\begin{figure}[htbp]
  \centering
  \includegraphics[width=\linewidth]{tuna2/tuna2_fig4_reconstruction.pdf}
  \caption{不同 visual tokenizer 的重建质性对比。Tuna-2 的重建已经很接近 Tuna-R，说明直接 patch/pixel embedding 可以保留足够多的低层细节；但 FLUX.1 VAE 仍然更像专门为重建优化的 tokenizer。因此这张图更适合读成“pixel-space UMM 可行”，而不是“专用 VAE 已经失去价值”。图源：arXiv source。}
\end{figure}

\subsection{5. Ablation 与训练动态}

\subsubsection{5.1 数据配比：generation objective 对数据比例更敏感}

作者比较了 generation-to-understanding 采样比例 \(xg yu\)。结论：

\begin{itemize}[leftmargin=1.5em]
  \item 提高某一类数据比例，会降低对应任务的训练 loss。
  \item flow matching 的 MSE 对比例变化更敏感；language CE 变化相对小。
  \item 论文最终采用 \(7:3\) 的 generation-to-understanding ratio。
\end{itemize}

\textbf{读法：} 统一训练不是简单把理解数据和生成数据混在一起。生成目标对数据量和采样比例更脆弱，训练 recipe 会显著影响最终结论。

\begin{figure}[htbp]
  \centering
  \includegraphics[width=0.98\linewidth]{tuna2/tuna2_fig5_training_loss.png}
  \caption{不同 understanding-generation 数据配比下的训练 loss。左侧 generation MSE 对采样比例更敏感，右侧 language CE 相对平滑；这解释了为什么 Tuna-2 的训练 recipe 不能只按理解任务配数据。对统一模型来说，generation objective 往往更吃数据比例和训练稳定性。图源：arXiv source。}
\end{figure}

\subsubsection{5.2 Masking 对 Tuna-2 很关键}

小模型 ablation 中，Tuna-2 加 masking 后：

\begin{center}
\begin{tabular}{p{0.22\linewidth} p{0.19\linewidth} p{0.19\linewidth} p{0.20\linewidth}}
\toprule
\textbf{指标} & \textbf{w/o Masking} & \textbf{w/ Masking} & \textbf{变化}\\
\midrule
OCRBench & 55.4 & 56.8 & +1.4\\
MMVP & 52.3 & 55.7 & +3.4\\
CountBench & 53.4 & 57.6 & +4.2\\
GenEval & 47.6 & 48.2 & +0.6\\
\bottomrule
\end{tabular}
\end{center}

论文在最终训练的后 40\% 阶段使用 masking-based feature learning。这个设计很值得记：不是从一开始就遮挡，而是在模型已有基础多模态知识后，用 masking 强化视觉表示。

\subsubsection{5.3 Tuna-R vs Tuna-2：scale 之后角色反转}

训练曲线给出的判断：

\begin{itemize}[leftmargin=1.5em]
  \item 早期 Tuna-R 理解更强，因为 SigLIP 2 representation encoder 提供语义先验。
  \item 随着预训练 token 规模增大，Tuna-2 追上并超过 Tuna-R。
  \item 生成侧 Tuna-R 基本全程略强，说明 representation encoder 的语义先验仍然帮助 generation。
  \item SFT 后两者生成差距接近消失。
\end{itemize}

\begin{figure}[htbp]
  \centering
  \includegraphics[width=0.98\linewidth]{tuna2/tuna2_fig6_scaling_curves.png}
  \caption{Tuna-R 与 Tuna-2 随训练 token 增长的 accuracy 曲线。最关键的现象是：Tuna-R 早期受益于 SigLIP 2 语义先验，起步更快；Tuna-2 随 scale 增大追上并在理解侧反超。生成侧 Tuna-R 仍略稳，说明 representation encoder 的语义先验对 generation/editing 仍有价值。图源：arXiv source。}
\end{figure}

\begin{warnbox}
\textbf{不要过度解读：} Tuna-2 不是在低数据、低算力下胜过 encoder-based model。它的优势依赖大规模 end-to-end vision pretraining。小数据场景下，pretrained encoder 仍然可能是更稳的工程选择。
\end{warnbox}

\subsection{6. Attention 可视化的意义}

论文用 attention map 比较 Tuna、Tuna-R、Tuna-2 和 LLaVA-OneVision、Qwen2.5-VL、Penguin-VL 等模型。重点案例包括：

\begin{itemize}[leftmargin=1.5em]
  \item 问“shining”时，Tuna-2 更准确关注发光区域。
  \item 问“purple object”时，Tuna-2 更聚焦目标物体。
  \item “cat cafe”语言先验和图像中真实动物冲突时，Tuna-2 更少被文本先验误导。
  \item football 场景中，问题真实目标是 glass cup，而图中 football 是显著干扰物，Tuna-2 更能定位问题语义对应的物体。
\end{itemize}

\begin{figure}[htbp]
  \centering
  \includegraphics[width=0.98\linewidth]{tuna2/tuna2_fig7_attention.pdf}
  \caption{Attention map 可视化。Tuna-2 在这些案例里更容易把注意力放到问题真正指向的图像区域，而不是被显著物体或文本先验带偏。它给论文主张提供了直观证据：encoder-free pixel embedding 可能保留更多细粒度视觉线索，帮助 grounding；但 attention map 仍是定性分析，不能替代 localization benchmark。图源：arXiv source。}
\end{figure}

这些可视化支撑了作者的解释：encoder-free pixel embedding 可能保留了更多低层视觉细节，使模型在细粒度 grounding 上更强。不过 attention map 只能作为定性证据，不能替代系统性 localization benchmark。

\subsection{7. 与 BAGEL / Janus / BLIP3o 的关键区别}

\begin{center}
\small
\begin{tabular}{p{0.16\linewidth} p{0.24\linewidth} p{0.21\linewidth} p{0.20\linewidth}}
\toprule
\textbf{模型} & \textbf{理解表示} & \textbf{生成表示} & \textbf{统一方式}\\
\midrule
Janus / Janus-Pro & semantic encoder & VQ/VAE-like generation representation & 解耦编码，避免理解/生成冲突。\\
BAGEL & SigLIP visual features & VAE latent + Flow Matching & MoT 分路，attention 交互。\\
BLIP3o & vision encoder / LLM condition & diffusion / flow decoder & LLM 条件 + 独立生成头。\\
Tuna & VAE + representation encoder & 统一 continuous visual representation & 仍依赖视觉编码模块。\\
Tuna-2 & patch embedding & pixel-space flow matching & encoder-free 单 transformer + flow head。\\
\bottomrule
\end{tabular}
\end{center}

Tuna-2 给这个领域增加了一条清晰路线：

\[
\text{representation engineering} \rightarrow \text{end-to-end pixel-space scaling}
\]

但这条路线的门槛是数据和算力。它不像 Janus 那样模块化，也不像 BAGEL 那样通过结构分路减轻任务冲突，而是把冲突交给大规模训练解决。

\subsection{8. 局限与风险}

\begin{itemize}[leftmargin=1.5em]
  \item \textbf{数据不可复现：} Stage 1 使用 550M in-house image-text pairs，核心结论高度依赖私有数据和训练规模。
  \item \textbf{算力门槛高：} 300k steps、64 nodes、16k tokens/GPU，不是普通实验室容易复现的路线。
  \item \textbf{生成侧没有全面胜出：} Tuna-2 在 GenEval、ImgEdit 上仍略低于 Tuna / Tuna-R。
  \item \textbf{只覆盖图像：} 虽然统一模型长期会走向视频、音频、any-to-any，但 Tuna-2 当前主实验是 image understanding / T2I / image editing。
  \item \textbf{LLM judge 评价有主观性：} quality/diversity 的模型裁判结果有参考价值，但需要和人评、自动指标、更多 prompt 分布交叉验证。
  \item \textbf{pixel token 成本：} 直接 pixel-space 建模通常会带来更高序列长度和训练成本，论文没有把效率问题作为主要贡献来解决。
\end{itemize}

\subsection{9. 对我的可复用点}

\begin{itemize}[leftmargin=1.5em]
  \item \textbf{受控模型阶梯：} Tuna \(\rightarrow\) Tuna-R \(\rightarrow\) Tuna-2 是很好的实验设计。后续比较架构时，不要只做最终模型对比，要构造“逐步去掉模块”的路径。
  \item \textbf{pixel-centric benchmark：} V*、CountBench、VisuLogic 很适合检验统一模型是否真的保留细粒度视觉信息。
  \item \textbf{训练动态分析：} 早期 encoder-based 更快、后期 encoder-free 更强这个现象很重要。未来看模型时要问“比较发生在训练哪个阶段”。
  \item \textbf{masking 作为后期表示增强：} 不是从一开始强行 MAE，而是在最终 40\% pretraining 引入 masking，这个 schedule 值得借鉴。
  \item \textbf{生成/理解分开读：} Tuna-2 的理解侧胜利比生成侧更强。统一模型不能只看一个综合分。
\end{itemize}

\subsection{10. 待查问题}

\begin{itemize}[leftmargin=1.5em]
  \item Tuna-2 的 patch size、视觉 token 数、不同分辨率下的 token budget 和 attention 成本，需要进一步从代码/config 里核对。
  \item pixel-space flow head 的具体参数规模、是否共享 transformer 层、unpatchify 细节，需要看官方实现。
  \item 如果把 Tuna-2 接上 RL 后训练，例如 visual reward / GenEval reward / editing reward，pixel-space trajectory 的 credit assignment 会不会比 latent-space 更难？
  \item 视频版 Tuna-2 是否可行：pixel-space video token 会非常贵，可能必须结合 patch compression、streaming attention 或 learned downsampling。
  \item 与 BAGEL 的真正公平对比需要控制参数量、数据量、训练 compute；当前表里 BAGEL 是 14B，Tuna-2 是 7B，不能直接下“谁更强”的结论。
\end{itemize}

\subsection{11. 后续阅读顺序}

\begin{itemize}[leftmargin=1.5em]
  \item \textbf{先读 Tuna 原文：} 理解 VAE + representation encoder 的统一表示设计，再看 Tuna-2 去模块化的动机。
  \item \textbf{再读 JiT / PixelFlow：} 补齐 pixel-space flow matching 为什么可行。
  \item \textbf{对照 BAGEL：} 重点比较 MoT 分路 vs encoder-free 单 transformer，二者对任务冲突的处理完全不同。
  \item \textbf{看官方代码：} 核对 patch embedding、masking schedule、flow head、resolution bucket 和 inference 参数。
\end{itemize}


\note{候选：BAGEL、Janus/Janus-Pro、BLIP3o、Transfusion。}

% World model and 3D deep-reading notes.

\section*{世界模型与 3D}
\phantomsection
\addcontentsline{toc}{section}{世界模型与 3D}

\begin{conceptbox}
\textbf{收录范围：} world model、agentic world modeling、3D/4D generation、spatial intelligence、scene simulation、CAD/asset generation、robotic policy evaluation by world model。\\
\textbf{核心问题：} 模型到底是在预测下一帧、模拟可交互环境、生成可编辑 3D 世界，还是作为策略评估器。
\end{conceptbox}

\note{当前暂无单篇精读。候选：Agentic World Modeling、dWorldEval、ReVSI、MetaEarth3D、Zero-to-CAD、GenAssets。}

% Embodied AI and VLA deep-reading notes.

\section*{具身智能与 VLA}
\phantomsection
\addcontentsline{toc}{section}{具身智能与 VLA}

\begin{conceptbox}
\textbf{收录范围：} vision-language-action、机器人后训练、低数据 SFT、online RL、具身安全、仿真环境、触觉/力觉多模态、长程操作。\\
\textbf{阅读重点：} 数据来源、动作表示、策略学习方式、真实机器人验证、泛化和安全失败模式。
\end{conceptbox}

\note{当前暂无单篇精读。候选：DeLock、RL Token、MoSS、VLA Safety、RedVLA、GazeVLA、CodeGraphVLP。}

% Style, identity, and design workflow deep-reading notes.

\section*{Style / Identity / Design}
\phantomsection
\addcontentsline{toc}{section}{Style / Identity / Design}

\begin{conceptbox}
\textbf{收录范围：} 图像风格、身份保持、voice/style preservation、设计工作流、审美奖励、版式生成、品牌一致性。\\
\textbf{阅读重点：} 风格如何定义、奖励/评测如何构造、是否能保持身份或作者 voice、是否贴近真实生产工作流。
\end{conceptbox}

\note{当前暂无单篇精读。候选：Can You Make It Sound Like You?、Voice Under Revision、AeSlides、POCA、Wan-Image 的 palette/identity/design workflow 部分。}

% Positional encoding topic notes.

\section*{位置编码专题}
\phantomsection
\addcontentsline{toc}{section}{位置编码专题}

\begin{conceptbox}
\textbf{收录范围：} absolute/sinusoidal positional encoding、RoPE、ALiBi、xPos、2D/3D RoPE、MRoPE、多模态 segment/offset 设计，以及图像/视频/统一理解生成模型中的位置组织方式。\\
\textbf{学习目标：} 不只记公式，而是搞清楚位置编码在不同模型里解决什么问题：语言顺序、图像空间位置、视频时间维、图组一致性、多模态 text/image segment 边界，以及长上下文外推。\\
\textbf{归档规则：} 单篇论文如果核心贡献是位置编码或 attention 位置机制，就放在这里；如果只是某篇模型论文里的一个模块，则在原论文笔记里讲一遍，并在这里做交叉索引。
\end{conceptbox}

% ============================================================
% Topic overview: Positional Encoding
% ============================================================

\subsection{位置编码学习总览}

这个专题先作为“位置编码学习入口”。后续每次在论文里遇到 RoPE、MRoPE、3D-RoPE、position offset、segment boundary、long-context extrapolation，都可以回到这里补一块。

\begin{intubox}
\textbf{先记住一个总问题：}Transformer 的 attention 本身只看 token 间内容相似度，不天然知道“谁在前、谁在后、谁在左上角、谁是第几帧”。位置编码就是给 token 补上可计算的位置关系，让模型知道顺序、空间、时间和模态边界。
\end{intubox}

\textbf{第一层：语言模型的位置。}

语言里最基础的问题是顺序。比如“狗咬人”和“人咬狗”token 集合相近，但顺序完全改变语义。早期 absolute/sinusoidal positional encoding 直接给每个位置加一个位置向量；RoPE 则把位置信息放进 query/key 的旋转关系里，让 attention score 能感知相对位移。

\textbf{第二层：图像模型的位置。}

图像 token 不只有一维顺序，还要知道二维空间。一个 patch 在左上角、右下角、人物脸部还是背景天空，对生成和编辑都很关键。因此图像模型常常需要 2D 位置编码，或者把 RoPE 扩展到 $H,W$ 两个空间维度。

\textbf{第三层：视频和图组的位置。}

视频、图像序列、multi-image generation 又多了时间/序列维。此时位置编码要同时表示第几张图/第几帧 $T$，以及图像内部第几行第几列 $H,W$。Wan-Image 里提到的 3D-RoPE 就是这一类：它让生成侧知道 latent token 在 $T,H,W$ 中的位置。

\textbf{第四层：多模态边界。}

统一理解与生成模型里，文本 token、ViT visual tokens、VAE latent tokens 可能放在同一个上下文里。如果只给它们普通位置，模型可能分不清“这是文本段的位置”还是“这是图像 latent 的位置”。所以很多系统会引入 segment boundary、positional offset 或 modality-specific encoding。Wan-Image 在 text segment 和 image segment 之间沿 $T$ 轴加入 fixed positional offset，就是这个方向的设计。

\begin{lstlisting}[language={},basicstyle=\ttfamily\small]
后续学习路线：
1. Absolute / sinusoidal PE：最基础的位置向量。
2. RoPE：为什么旋转能表达相对位置。
3. ALiBi / xPos：长上下文和外推问题。
4. 2D RoPE：图像 patch 的 H,W 位置。
5. MRoPE：多模态 text + image token 的位置组织。
6. 3D-RoPE：视频/图组/latent 的 T,H,W 位置。
7. Segment offset：文本段、图像段、latent 段如何隔开。
\end{lstlisting}

\textbf{和 Wan-Image 的连接。}

Wan-Image 3.1 里最值得回头细学的是这句话：理解侧使用 MRoPE，生成侧使用 3D-RoPE，并在 text/image segments 之间插入 fixed positional offset。它背后的直觉是：

\begin{itemize}[leftmargin=1.5em]
  \item \textbf{MRoPE：}帮助理解侧把文本描述和图像区域对齐，比如“左下角 logo”“右侧人物脸部”。
  \item \textbf{3D-RoPE：}帮助生成侧知道 VAE latent token 属于第几张图/第几帧、图像第几行、第几列。
  \item \textbf{fixed offset：}把 text segment 和 image segment 在位置空间里隔开，避免不同模态的位置混在一起。
\end{itemize}

后面精读位置编码时，不要只停在公式。每种方法都按三个问题学：

\begin{enumerate}[leftmargin=1.8em]
  \item 它给 token 注入了哪种位置信息？
  \item 它改变的是 embedding、attention score，还是 query/key 的几何关系？
  \item 它适合语言、图像、视频、多模态里的哪一种场景？
\end{enumerate}


% Research idea notes.

\section*{Idea 专区}
\phantomsection
\addcontentsline{toc}{section}{Idea 专区}

\begin{conceptbox}
\textbf{用途：} 这里专门记录读论文、做实验、讨论模型系统时冒出来的 idea。它不要求一开始就完整，只要能把动机、核心假设、可能实验和相关论文留下来即可。\\
\textbf{记录方式：} 先放进 Idea Inbox；如果某个 idea 变得清晰，再单独拆成一个文件，补充背景、技术路线、最小验证实验、风险和下一步。\\
\textbf{建议原则：} 不急着写成 proposal，先保留原始问题意识。真正重要的是：这个 idea 想解决什么痛点，为什么现在可能可做，最小实验怎么验证。
\end{conceptbox}

\subsection*{Idea Inbox}
\phantomsection
\addcontentsline{toc}{subsection}{Idea Inbox}

\begin{summarybox}
这个文件先作为所有研究想法的临时入口。后面如果某个 idea 需要深入推演，再从这里移动到单独的专题笔记里。
\end{summarybox}

\subsubsection*{记录模板}

\begin{lstlisting}[language={},basicstyle=\ttfamily\small]
日期：
标题：
来源：哪篇论文 / 哪次讨论 / 哪个实验现象

我真正想解决的问题：

核心想法：

为什么可能有效：

相关论文或系统：

最小验证实验：

可能风险：

下一步：
\end{lstlisting}

\subsubsection*{待记录}

\begin{itemize}[leftmargin=1.5em]
  \item \textbf{2026-05-02：} 从 Latent Forcing 出发，研究 representation latents 如何加速 $1K+$ pixel diffusion 训练。已拆成独立笔记：\textbf{Idea：Latent Forcing 如何加速高分辨率 Pixel Diffusion}。
\end{itemize}

\subsection*{Idea：Latent Forcing 如何加速高分辨率 Pixel Diffusion}
\phantomsection
\addcontentsline{toc}{subsection}{Idea：Latent Forcing 如何加速高分辨率 Pixel Diffusion}

\begin{conceptbox}
\textbf{核心问题：} 能不能利用 deterministic representation latents 加速 $1K+$ pixel diffusion 的训练，但最终仍然直接生成 raw pixels，不依赖 VAE decoder，也不走低分辨率到高分辨率的 cascade？\\
\textbf{当前判断：} 这个方向真正要抓的是 \textbf{generation order}，不是 resolution pyramid。Latent Forcing 的价值在于：把 representation latent 放进 diffusion trajectory，让模型先生成更容易的结构化表示，再生成高维像素细节。latent 是 inference-time scratchpad，最后丢掉。
\end{conceptbox}

\begin{figure}[htbp]
  \centering
  \includegraphics[width=0.82\linewidth]{latent_forcing/fig1_concept.jpg}
  \caption{Latent Forcing 的核心直觉。横轴是 pixel timestep，纵轴是 latent timestep，图里比较的不是“低分辨率图到高分辨率图”，而是不同变量在 diffusion trajectory 里变 clean 的顺序。绿色轨迹先让 latent 变得有信息，再让 pixels 生成细节；红色轨迹是纯 pixel diffusion，模型要直接从像素噪声里学结构。图源：arXiv HTML。}
\end{figure}

\subsubsection*{1. 先把问题摆正}

这个 idea 不是“先生成低分辨率图，再生成高分辨率图”。这条路已经有很多工作做过，比如 cascaded diffusion、Matryoshka Diffusion Models、PixelFlow 里的 cascade flow，以及各种 high-resolution pixel-space architecture。它们解决的是高分辨率生成的分解问题：先低频、后高频；先小图、后大图；先全局、后局部。

我真正想研究的是另一个问题：\textbf{pixel diffusion 直接在 $1K$ 甚至更高分辨率 raw pixels 上训练时，能不能借助 representation latents 改善优化轨迹，让训练更快、更稳？}

也就是说，最终输出仍然是 pixels：
\begin{lstlisting}[language={},basicstyle=\ttfamily\small]
final output: raw pixels
not output: VAE latent decoded by a decoder
not output: low-res image upsampled / refined to high-res
\end{lstlisting}

latent 的角色不是图像压缩格式，而是中间计算变量：
\begin{lstlisting}[language={},basicstyle=\ttfamily\small]
x = high-resolution raw pixels
y = f(x), deterministic representation latent

train joint diffusion over (x, y)
generate y earlier than x
discard y at the end
\end{lstlisting}

所以这张图对我的 idea 很重要：它把“加速”画成了\textbf{轨迹难度下降}，不是把图像空间换成低维 latent space。最终图像仍然由 pixel branch 产生；latent 只是让模型在早期先拿到更稳定的结构坐标。

\subsubsection*{2. 过往工作怎么走到 Latent Forcing}

\textbf{Latent Diffusion / VAE tokenizer 路线。}

Latent Diffusion 的基本逻辑是：先用 VAE/VQ tokenizer 把图像压到 latent space，在 latent space 训练 diffusion，最后通过 decoder 回到 pixel。它的优势很明显：token 数少、训练便宜、收敛快。但代价也很明显：encoder 会丢信息，decoder 是单独训练的，模型最终优化的是辅助 latent distribution，不是 raw image distribution。

这个路线长期存在一个 reconstruction-generation tradeoff：tokenizer 越保真，latent 越难生成；tokenizer 越 lossy，diffusion 越容易，但图像上限被 decoder 和重建误差限制。小字、人脸、高频纹理、透明边界这类细节特别容易被 tokenizer 阶段提前牺牲。

\textbf{Pixel Diffusion 回归路线。}

Pixel diffusion 想绕开这个问题：不压缩，不解码，直接在 raw pixels 上建模。过去它难，是因为像素空间维度太高，训练容易失败，尤其是 Transformer 直接预测 noise 或 velocity 时，目标不在自然图像流形上。

JiT 的关键贡献是把目标改成 clean data prediction，也就是让模型直接预测 denoised image，再用 v-loss weighting 平衡训练。这样 plain Transformer 可以直接在 pixel space 工作。PixelFlow、Matryoshka Diffusion、HDiT、SiD2 等工作则从不同角度证明：pixel-space generation 不是不可行，关键在 architecture、schedule 和训练目标。

\textbf{Representation Alignment 路线。}

REPA 提供了另一个重要线索：diffusion model 训练慢，部分原因是它要自己学出好的 visual representation。REPA 用 DINOv2 等 pretrained visual encoder 的 clean representation 去对齐 diffusion hidden states，大幅提升训练效率。

但 REPA 的 representation 主要是 training-time auxiliary signal。推理时，模型并不会显式生成这些 representation。换句话说，REPA 告诉模型“训练时你应该像 DINO 那样组织特征”，但没有把 representation 放进生成轨迹本身。

\textbf{RAE / REG / Latent Forcing 的位置。}

RAE 把 pretrained representation encoder 和 decoder 组成新的 autoencoder，让 diffusion 在更语义化的 latent space 里训练。它比传统 VAE 更强，但仍然需要 decoder，所以仍是 latent diffusion 家族。

REG 更进一步，把高层语义 token 放进 denoising，让 semantic token 在推理时也参与生成。Latent Forcing 则更系统地把这个想法写成 multi-time diffusion：pixels 和 representation latents 同时作为 diffusion variables，每个变量有自己的 time schedule。

\Needspace{18\baselineskip}
\begin{intubox}
这条线可以这样理解：REPA 是 training-time representation alignment；RAE 是 representation autoencoding；REG 是 semantic token entanglement；Latent Forcing 是 trajectory-level representation generation。
\end{intubox}

\subsubsection*{3. Latent Forcing 的核心机制}

Latent Forcing 设定里有两个变量：
\begin{lstlisting}[language={},basicstyle=\ttfamily\small]
x = pixels
y = f(x), e.g. DINOv2 / Data2Vec2 representation
\end{lstlisting}

注意，$y$ 是从 $x$ 确定性算出来的，不是人工标签，也不是另一个生成模型采样得到的东西。因此联合建模 $(x,y)$ 并没有改变最终目标分布。最后只要保留 $x$，丢掉 $y$，模型仍然是在生成 raw pixels。

训练时分别对两种变量加噪：
\[
x_{t_x} = t_x x + (1-t_x)\epsilon_x,
\qquad
y_{t_y} = t_y y + (1-t_y)\epsilon_y .
\]

模型输入是：
\begin{lstlisting}[language={},basicstyle=\ttfamily\small]
noised pixels, noised latents, pixel timestep, latent timestep
\end{lstlisting}

模型输出则同时预测 clean pixels 和 clean latents。关键不是“多一个 loss”，而是 $t_x$ 和 $t_y$ 可以不同。Latent Forcing 让 representation latent 比 pixels 更早 denoise：
\begin{lstlisting}[language={},basicstyle=\ttfamily\small]
early trajectory:   latent becomes informative first
later trajectory:   pixels use latent structure to generate detail
final:              discard latent, keep pixels
\end{lstlisting}

这就是它和 REPA 的根本区别。REPA 是把 representation 当监督信号；Latent Forcing 是把 representation 当生成过程的一部分。

\begin{figure}[htbp]
  \centering
  \includegraphics[width=0.94\linewidth]{latent_forcing/fig2_architecture.jpg}
  \caption{Latent Forcing 对 DiT 的改动其实很克制。左侧把 noised pixels 和 noised latents 都 patchify、embed，然后把同位置的 embedding 相加，所以 transformer token 数没有翻倍；中间把 $t_{\text{pixel}}$ 和 $t_{\text{latent}}$ 一起送入 adaLN；右侧可以把最后几层拆成 pixel output expert 和 latent output expert。图源：arXiv HTML。}
\end{figure}

这张架构图也解释了为什么它适合作为 high-resolution pixel diffusion 的加速思路：它不是额外接一个大 decoder，也不是在 latent space 里替代 pixel generation，而是在同一个 DiT 里把 representation latent 作为一个额外的生成变量。对高分辨率版本来说，真正需要重新设计的地方会是 latent resolution、patch 对齐方式、两个 timestep 的 schedule，以及 output head / expert 是否足够承载 pixel 和 latent 两个预测目标。

\subsubsection*{4. 为什么这可能加速高分辨率 Pixel Diffusion}

$1K+$ pixel diffusion 的难点不是模型不能表达像素，而是优化路径太难。早期 timestep 里，模型面对的是巨大维度的 pixel noise。它既要猜全局布局，又要猜对象结构，还要学局部纹理和高频细节。直接从 raw pixel noise 中学出语义结构，样本效率很低。

representation latent 的优势是：它通常更接近结构化语义空间。比如 DINOv2 patch tokens 会保留对象、部件和空间结构；SigLIP 更贴近 text-image alignment；MAE 或 reconstruction-oriented features 可能更保留局部结构。Latent Forcing 的想法是先让这些 representation 变得 informative，再让 pixel branch 生成细节。

这里的关键不是降低 pixel 分辨率，而是改变信息出现顺序：
\[
\text{semantic/spatial representation first}
\quad \rightarrow \quad
\text{raw pixels later}.
\]

\begin{figure}[htbp]
  \centering
  \includegraphics[width=0.94\linewidth]{latent_forcing/fig4_same_psnr.jpg}
  \caption{同样的 pixel PSNR 下，DINOv2 latent 预测比 pixel-only 预测更早保住了对象轮廓和空间结构。上排 pixel reconstruction 在低 PSNR 时更像局部颜色块；下排 DINO reconstruction 即使模糊，也能更早呈现狗的姿态和脸部位置。图源：arXiv HTML。}
\end{figure}

\Needspace{18\baselineskip}
\begin{intubox}
\textbf{PSNR 是什么。} PSNR 全称是 Peak Signal-to-Noise Ratio，通常用来衡量一张重建图 $\hat{x}$ 和真实图 $x$ 的 pixel-level reconstruction fidelity。先算均方误差：
\[
\operatorname{MSE}(x,\hat{x})
=
\frac{1}{N}\sum_{i=1}^{N}(x_i-\hat{x}_i)^2 .
\]
再把误差转成对数尺度：
\[
\operatorname{PSNR}
=
10\log_{10}\frac{\operatorname{MAX}^2}{\operatorname{MSE}}
=
20\log_{10}\operatorname{MAX}-10\log_{10}\operatorname{MSE}.
\]
如果图像归一化到 $[0,1]$，通常 $\operatorname{MAX}=1$。PSNR 越高，说明 pixel MSE 越小。但它不直接衡量语义结构：两张图可以有相同 PSNR，一张保住对象轮廓，另一张只是局部颜色更接近。所以上面这张图的重点是：\textbf{同样 pixel 误差水平下，DINO latent 更早保住了结构信息。}
\end{intubox}

从 optimization 角度看，这会把 pixel branch 的任务从“在纯噪声中自己发现结构”，变成“在已经形成的 representation structure 上生成像素”。这就是我认为它能加速 high-resolution pixel diffusion training 的原因。

\begin{warnbox}
不能把这个 idea 讲成 coarse-to-fine resolution training。正确说法是：\textbf{用 representation latents 重排 diffusion trajectory 中的信息揭示顺序，从而降低 high-dimensional pixel denoising 的优化难度。}
\end{warnbox}

\subsubsection*{5. 可以做创新的几个切口}

\textbf{第一，什么 latent 最适合加速 pixel diffusion？}

Latent Forcing 默认用 DINOv2，也试了 Data2Vec2 和 downsampled pixels。但如果目标是 $1K+$ pixel diffusion，latent 不一定越语义越好。iREPA 这类工作提醒我们，generation 里很重要的可能不是 global semantic accuracy，而是 patch-level spatial structure。

所以可以系统比较：
\begin{itemize}[leftmargin=1.5em]
  \item DINOv2 deep features：强对象语义和空间结构；
  \item SigLIP features：更适合 text-conditioned generation；
  \item MAE features：更偏 reconstruction，可能更适合局部结构；
  \item shallow / mid / deep 多层 features：把边缘、部件、语义拆开；
  \item spatial-structure-preserving features：不追求分类强，而追求 patch relation 稳定。
\end{itemize}

\textbf{第二，schedule 不应该照搬 256 ImageNet。}

Latent Forcing 主要验证在 ImageNet $256$。到 $1K$ 以后，pixel dimension、patch token 数、低频/高频比例、文本条件复杂度都变了。可能需要重新设计：
\begin{itemize}[leftmargin=1.5em]
  \item latent 要提前多少；
  \item 是否完全 cascaded；
  \item latent 和 pixel 是否需要 overlap；
  \item 多层 latent 是否应该有不同 schedule；
  \item 高分辨率细节阶段是否还需要 latent dropout / latent noise。
\end{itemize}

\begin{figure}[htbp]
  \centering
  \includegraphics[width=0.94\linewidth]{latent_forcing/fig3_trajectories.jpg}
  \caption{不同 $t_{\text{DINO}}$ 与 $t_{\text{pixel}}$ 轨迹的 FID-10K。颜色越黄 FID 越低；图里表现好的路径都有一个共同点：DINO latent 走在 pixel 前面，先变 clean。左图像 variance-shifted schedule，右图像 linear-offset schedule。图源：arXiv HTML。}
\end{figure}

\begin{intubox}
\textbf{FID 是什么。} FID 全称是 Fréchet Inception Distance，用来比较生成图像集合和真实图像集合的 distribution-level distance。做法不是逐张图对齐，而是先用一个固定的 Inception 网络把真实图和生成图都映射到 feature space，然后分别估计两个高斯分布：
\[
(\mu_r,\Sigma_r)
\quad\text{and}\quad
(\mu_g,\Sigma_g).
\]
FID 定义为：
\[
\operatorname{FID}
=
\|\mu_r-\mu_g\|_2^2
+
\operatorname{Tr}\left(
\Sigma_r+\Sigma_g-2(\Sigma_r\Sigma_g)^{1/2}
\right).
\]
FID 越低，说明生成分布在 Inception feature space 里越接近真实分布。FID-10K 通常表示用 $10K$ 张生成样本来估计这个距离。它比 PSNR 更适合看 generative quality，但也不是完美指标：它依赖 feature extractor、样本数和真实集统计，而且不能解释某一张图为什么好或坏。
\end{intubox}

这张图给我的启发是：schedule 不是形式问题，而是方法本体。Latent Forcing 的收益不只是“加了 DINO loss”，而是让模型在采样轨迹中先形成 DINO 结构，再让 pixel 追上来。因此做 $1K+$ 时，最值得系统扫的不是单个 latent 类型，而是 latent 提前量、overlap 区间、以及 pixel 什么时候开始强依赖 latent。

\begin{figure}[htbp]
  \centering
  \includegraphics[width=0.94\linewidth]{latent_forcing/fig5_psnr_heatmap.jpg}
  \caption{PSNR heatmap 用来观察两个变量互相提供信息的区域。一个值得注意的细节是：pixel 信息几乎总能帮助 DINO latent 预测；但当 pixel 已经比较 clean 时，DINO 对 pixel 的额外帮助会变小。这提示 high-res 版本不一定要让 latent 一直强控到最后，后段可能应该逐步把自由度还给 pixels。图源：arXiv HTML。}
\end{figure}

\textbf{第三，把 REPA 从 auxiliary loss 推到 trajectory variable。}

可以设计一个很清楚的对照：
\begin{lstlisting}[language={},basicstyle=\ttfamily\small]
Baseline:        pure pixel diffusion
REPA:            pixel diffusion + training-time representation alignment
Latent Forcing:  pixel diffusion + generated representation latent
Hybrid:          generated representation latent + light REPA alignment
\end{lstlisting}

如果 Latent Forcing 在同等 compute 下比 REPA 更持久地加速，说明 representation 不只是训练监督，而是应该进入生成轨迹。

\textbf{第四，多层 representation ordering。}

单一 DINOv2 latent 可能太粗。更自然的 high-res 版本是：
\begin{lstlisting}[language={},basicstyle=\ttfamily\small]
deep semantic latent      denoise earliest
mid-level structure       denoise next
local/spatial latent      denoise next
raw pixels                denoise last
\end{lstlisting}

这不是 image pyramid，因为每一层不是不同分辨率图像，而是不同抽象层级的 representation。它更像 continuous latent chain-of-thought for image generation。

\subsubsection*{6. 最小实验设计}

我觉得第一步不该直接冲 $1K$ 大模型，而是先证明机制：

\begin{enumerate}[leftmargin=1.8em]
  \item 在 ImageNet $256$ 或一个小规模高质量数据集上复现 JiT-style pixel diffusion baseline；
  \item 加 REPA，得到 training-time alignment baseline；
  \item 加 Latent Forcing：DINOv2 latent 作为 generated variable，latent 先于 pixel denoise；
  \item 比较 convergence speed：同等 step 下 FID / IS / precision-recall；
  \item 比较结构指标：对象布局、patch relation、低频结构稳定性；
  \item 再把分辨率推到 $512$，看收益是否更明显；
  \item 最后再到 $1K$，重点观察 sample efficiency 和高频细节质量。
\end{enumerate}

如果资源有限，先做一个很小的消融也有价值：
\begin{lstlisting}[language={},basicstyle=\ttfamily\small]
DINOv2 latent earlier
DINOv2 latent same schedule as pixels
DINOv2 latent later than pixels
REPA only
no latent / pure pixel
\end{lstlisting}

这组实验能直接回答：收益到底来自 representation 本身，还是来自 ordering。

\subsubsection*{7. 当前最清晰的研究假设}

\begin{summarybox}
在 $1K+$ pixel-space diffusion 中，显式生成 deterministic representation latents，并让这些 latents 在 diffusion trajectory 中先于 raw pixels 变得 informative，可以在不引入 VAE decoder 和 reconstruction ceiling 的前提下，加速 pixel diffusion 训练，并提升结构一致性。
\end{summarybox}

这个方向可以暂时命名为：
\begin{itemize}[leftmargin=1.5em]
  \item \textbf{Representation-Ordered Pixel Diffusion}
  \item \textbf{High-Resolution Latent Forcing}
  \item \textbf{Trajectory-Level Representation Guidance}
\end{itemize}

\subsubsection*{8. 相关论文}

\begin{itemize}[leftmargin=1.5em]
  \item \href{https://arxiv.org/abs/2602.11401}{Latent Forcing: Reordering the Diffusion Trajectory for Pixel-Space Image Generation}
  \item \href{https://arxiv.org/abs/2511.13720}{JiT / Back to Basics: Let Denoising Generative Models Denoise}
  \item \href{https://arxiv.org/abs/2410.06940}{REPA: Representation Alignment for Generation}
  \item \href{https://arxiv.org/abs/2510.11690}{Diffusion Transformers with Representation Autoencoders}
  \item \href{https://arxiv.org/abs/2512.10794}{What matters for Representation Alignment: Global Information or Spatial Structure?}
  \item \href{https://arxiv.org/abs/2507.01467}{Representation Entanglement for Generation}
  \item \href{https://arxiv.org/abs/2310.15111}{Matryoshka Diffusion Models}
  \item \href{https://arxiv.org/abs/2504.07963}{PixelFlow: Pixel-Space Generative Models with Flow}
\end{itemize}


% Evaluation and safety deep-reading notes.

\section*{评测、安全与红队}
\phantomsection
\addcontentsline{toc}{section}{评测、安全与红队}

\begin{conceptbox}
\textbf{收录范围：} benchmark、failure discovery、red teaming、provenance、synthetic evidence、VLM/VLA safety、reward/evaluator reliability。\\
\textbf{归档规则：} 如果评测论文主要服务某个具体方向，也可以放入该方向；这里保留横跨多个模型类型的评测和安全方法。
\end{conceptbox}

\note{当前暂无单篇精读。候选：ProEval、DIVERT、VLA Safety、Seeing Is No Longer Believing、RedVLA、ReVSI。}

% Systems and inference deep-reading notes.

\section*{系统与推理效率}
\phantomsection
\addcontentsline{toc}{section}{系统与推理效率}

\begin{conceptbox}
\textbf{收录范围：} 训练系统、推理加速、attention 优化、采样加速、长序列视频 DiT 工程、数据引擎和 benchmark pipeline。\\
\textbf{归档规则：} 如果系统技巧是某篇模型论文的附属贡献，保留在主方向；如果系统/效率是论文主贡献，放在这里。
\end{conceptbox}

\note{当前暂无单篇精读。候选：FreqFormer、video DiT 长序列 attention、Mean Flow / one-step generation、VLA data engine survey。}

